% !TEX program = tectonic
% Tokamak Plasma Simulator — Mathematical Models and Numerical Solution Guide
% 基于 simulator_solver_methodology.md 重组润色；排版风格参考 Atlas User Manual。
\documentclass[12pt,fontset=none]{ctexart}

\usepackage[a4paper,margin=2.4cm,top=2.8cm,bottom=2.8cm]{geometry}
\usepackage{amsmath,amssymb,bm}
\usepackage{booktabs,longtable,array,multirow}
\usepackage{enumitem}
\usepackage{seqsplit}
\usepackage{xcolor}
\usepackage{fancyhdr}
\usepackage{listings}
\usepackage{graphicx}
\graphicspath{{figures/}{docs/tokamak_control/figures/}}
\usepackage{tikz}
\usetikzlibrary{arrows.meta,positioning,shapes.geometric,fit,calc}
\usepackage[colorlinks=true,linkcolor=black,urlcolor=blue!60!black,citecolor=black]{hyperref}

% ---------- 字体 ----------
\IfFileExists{/usr/share/fonts/opentype/noto/NotoSerifCJK-Regular.ttc}{
\setCJKmainfont{Noto Serif CJK SC}[Path=/usr/share/fonts/opentype/noto/,Extension=.ttc,UprightFont=NotoSerifCJK-Regular,BoldFont=NotoSerifCJK-Bold]
\setCJKsansfont{Noto Sans CJK SC}[Path=/usr/share/fonts/opentype/noto/,Extension=.ttc,UprightFont=NotoSansCJK-Regular,BoldFont=NotoSansCJK-Bold]
\setCJKmonofont{Noto Sans Mono CJK SC}[Path=/usr/share/fonts/opentype/noto/,Extension=.ttc,UprightFont=NotoSansCJK-Regular,BoldFont=NotoSansCJK-Bold]
}{
  \setCJKmainfont{Songti SC}
  \setCJKsansfont{Heiti SC}
  \setCJKmonofont{Heiti SC}
}

% ---------- 版式 ----------
\numberwithin{equation}{section}
\numberwithin{table}{section}
\numberwithin{figure}{section}
\setlength{\parskip}{0.45em}
\setlength{\parindent}{2em}
\setlength{\headheight}{14pt}
\renewcommand{\arraystretch}{1.25}

\pagestyle{fancy}
\fancyhf{}
\fancyhead[L]{\footnotesize\sffamily Magnetic Fusion}
\fancyhead[R]{\footnotesize\sffamily\nouppercase{\leftmark}}
\fancyfoot[L]{\footnotesize\sffamily 磁约束聚变控制与设计学习报告}
\fancyfoot[R]{\footnotesize\sffamily \thepage}
\renewcommand{\headrulewidth}{0.4pt}
\renewcommand{\footrulewidth}{0.4pt}

% ---------- 代码与证据标签 ----------
\newcommand{\code}[1]{\texttt{\expandafter\seqsplit\expandafter{\detokenize{#1}}}}
\newcommand{\evimpl}{\textbf{[实现]}}
\newcommand{\evdata}{\textbf{[数据]}}
\newcommand{\evderiv}{\textbf{[推导]}}
\newcommand{\evnote}{\textbf{[说明]}}
\newcommand{\evbg}{\textbf{[背景]}}
\newcommand{\evunknown}{\textbf{[未知]}}

% ---------- 伪代码列表 ----------
\lstdefinestyle{pseudo}{
  basicstyle=\footnotesize\ttfamily,
  frame=single, rulecolor=\color{black!50},
  backgroundcolor=\color{black!4},
  breaklines=true, columns=fullflexible,
  keepspaces=true, showstringspaces=false,
  aboveskip=0.8em, belowskip=0.8em,
  morecomment=[l]{\#},
  mathescape=false
}

% ---------- 表格列类型 ----------
\newcolumntype{L}[1]{>{\raggedright\arraybackslash}p{#1}}

\title{}
\date{}

\begin{document}

% ================= 封面 =================
\hypersetup{pageanchor=false}
\begin{titlepage}
\centering
\vspace*{2.5cm}
{\sffamily\bfseries\Huge Magnetic Fusion\par}
\vspace{0.9cm}
{\sffamily\LARGE Tokamak Control and Stellarator Design\\Physics and Numerical Simulation\par}
\vspace{1.6cm}
{\sffamily\large 磁约束聚变：托卡马克控制与仿星器设计\\物理图景、放电过程与数值仿真指南\par}
\vspace{0.5cm}
{\sffamily\normalsize 通用控制物理导读 · HFM / EXL-50U 与 MEQ 实现案例\par}
\vfill
\begin{tabular}{rl}
\sffamily 文档性质： & 通用控制学习资料与可追溯仿真实现分析 \\[2pt]
\sffamily 更新日期： & 2026-10-24 \\[2pt]
\sffamily 目标读者： & 具备工科数学物理基础的初学者与控制研究人员 \\
\end{tabular}
\vspace{1.2cm}
\end{titlepage}

% ================= 关于本文档 =================
\pagenumbering{roman}
\hypersetup{pageanchor=true}
\section*{关于本文档}
\markboth{关于本文档}{关于本文档}
\addcontentsline{toc}{section}{关于本文档}

本文档面向磁约束聚变的物理学习、托卡马克控制与仿星器设计研究，从基本受力和守恒律解释气体导入、磁场建立、击穿烧穿、建流加热、稳态维持、位形转换与安全关停，补充仿星器线圈、三维磁场与装置性质的设计联系，并以 HFM 发行物中的 EXL-50U 强化学习环境及 MEQ 自由边界平衡求解器作为可追溯的实现案例。报告的通用部分覆盖微观碰撞、集体电磁场、典型结构与完整放电过程，并与具体实现分析衔接：前者明确物理假设并引用公开研究与机构资料，后者以指定版本的源码和保存数据为证据。通用物理范围宽于案例软件的求解范围，不能把导读中的过程推定为软件已有功能。

\textbf{实现案例的分析范围与版本约定。} 源码分析部分研究项目根目录下 \code{hfm-runtime-release-sealed-latest/hfm-runtime-release-sealed-latest/} 的 HFM 发行物，重点是 EXL-50U 强化学习环境、MEQ 自由边界平衡与回路演化求解器。源码审计日期为 2026-10-21；通用控制导读、粒子/集体结构与仿星器章节更新于 2026-10-24，未重新核验全部实现证据。源码分析过程只进行文件读取、Git 对象读取、MAT 文件读取及静态代数核对；未启动 MATLAB、仿真服务或参数扫描，未修改任何既有源码、配置和算法。

\textbf{可追溯性边界。} 发行工作树的高层 MATLAB 文件为 \code{.p}，但 \code{HFM_env/.git} 保留了原始 \code{.m}。本文通过只读 \code{git show} 分析该仓库 HEAD \code{0e2cb6012bf178da35e17e36506851446342d9f5} 的源文件，没有 checkout、解密或反编译。Socket 服务端只有 \code{socket_serve/hfm_socket_server.p}，未找到对应可读源码。Git 源码能够解释保存的算法实现，但当前材料\textbf{无法从现有源码确认 \code{.p}/\code{.mexa64} 产物与该 Git 快照逐字、逐版本对应}；本文不把源码分析冒充二进制运行验证。

引用前缀约定如表~\ref{tab:prefix} 所示，所有行号均从 1 开始，范围两端均包含。

\begin{table}[htbp]
\centering\small
\caption{引用前缀约定}
\label{tab:prefix}
\begin{tabular}{lL{5.8cm}L{5.8cm}}
\toprule
前缀 & 实际位置或读取方式 & 证据属性 \\
\midrule
\code{R/} & 项目内发行物根目录 & 发行工作树 \\
\code{Q/} & \code{R/HFM_env/meq/} & 工作树中的 C、头文件、构建文件及说明 \\
\code{G:} & \code{HFM_env} 仓库指定 HEAD 的路径 & Git 对象库中的原始源码（含已不存在的 \code{.m}） \\
\code{E/} & 项目根目录 \code{environment/} & 目标目录以外的 Python 调用与包装层 \\
\code{C/} & 项目根目录 \code{configs/} & 目标目录以外的场景与默认配置 \\
\code{T2/} & 项目根目录 \code{train_v2/} & 项目外 PPO 训练入口（Hydra 版）及其配置 \\
\code{TR/} & 项目根目录 \code{training/} & 项目外训练回调、奖励、参考轨迹与蒸馏模块 \\
\code{D/} & 发行物中的 MAT 数据 & 用完整字段路径、形状和值定位，无文本行号 \\
\code{V/} & \code{G:TOK_EXL50U/XL50U_fge/RL_ENV/} & 环境层 Git 源码的简写 \\
\bottomrule
\end{tabular}
\end{table}

例如，\code{G:meq/fgeF.m:460--482} 是 Git 源码的行号，\textbf{不是}同名 \code{.p} 的行号。可在项目根目录复核：
\begin{lstlisting}[style=pseudo]
git -C hfm-runtime-release-sealed-latest/hfm-runtime-release-sealed-latest/HFM_env show 0e2cb6012bf178da35e17e36506851446342d9f5:meq/fgeF.m
\end{lstlisting}

全文用以下标签区分结论：\evimpl{} 表示可见工作树或 Git 源码直接实现；\evdata{} 表示 MAT/CSV 中保存的配置或数值；\evderiv{} 表示从指定代码关系严格得到的数学结论；\evnote{} 表示随附 README 的声明，不能替代实现；\evbg{} 表示通用物理、教学模型与辅助理论，不是案例软件的额外求解模块；\evunknown{} 表示无法从现有源码确认。除明确写成``在线二进制''的地方外，所有 \code{.m} 实现结论均受上述 Git 快照边界限定。

本文的主要复现实例是 \code{D/ini_L_v2.2.mat} 与 \code{D/ini_13906_xpt_v2.2.mat}。它们也是项目外默认场景 \code{C/env_default.yaml} 和 \code{C/shots.yaml} 所选择的数据；这不证明不可见 Socket 服务端自身具有同一默认值。装置文件不同版本的区别见第~\ref{sec:mesh} 节。

\clearpage
\begingroup
\small
\tableofcontents
\endgroup
\clearpage
\pagenumbering{arabic}

% ================= 第 1 章：阅读入门 =================
\section{阅读入门：等离子体平衡与数值求解的物理直觉}\label{sec:primer}
\markboth{阅读入门}{阅读入门}

\textbf{本章定位。}本章面向没有等离子体仿真经验、但具备工科本科数学物理基础的读者，用最少的前置知识建立阅读后续章节所需的物理图像。本章采用\evbg{}标签，构成正文的基础理论部分：所有与案例实现相关的结论，仍以第 \ref{sec:intro}--\ref{sec:rl} 章给出源码证据的内容为准。已经熟悉托卡马克平衡计算的读者可以跳过本章。

\subsection{托卡马克与轴对称磁场}\label{sec:primer-tokamak}

托卡马克是一种环形磁约束聚变装置：等离子体被磁场约束在环形的``真空室''内。描述它需要使用柱坐标 $(R,\phi,Z)$，其中 $R$ 是到环中心轴的大半径，$\phi$ 是绕环的环向角，$Z$ 是竖直方向。磁场可以分解为两个部分：

\begin{itemize}[nosep]
  \item \textbf{环向磁场} $B_\phi$：沿绕环方向，主要由外部环向场线圈产生，近似与 $1/R$ 成正比；
  \item \textbf{极向磁场} $B_R,B_Z$：在 $R$--$Z$ 截面内，由等离子体电流和极向场（PF）线圈共同产生。
\end{itemize}

\textbf{轴对称}是托卡马克平衡计算中最重要的简化：装置关于环向角 $\phi$ 对称，因此所有物理量都不依赖 $\phi$。整个三维磁场问题于是可以化归为一个\textbf{二维 $R$--$Z$ 截面}上的问题——这正是本文档中所有平衡方程只出现在 $(R,Z)$ 平面内的原因。

环向场与极向场叠加后，磁力线呈螺旋形，扫出一族嵌套的\textbf{磁面}（magnetic surface）。在这里的轴对称静态、标量压强近似下，磁面也是总极向磁通与压力的等值面。最外一层闭合磁面称为 \textbf{LCFS}（Last Closed Flux Surface，最外闭合磁面）；磁面中心是\textbf{磁轴}（magnetic axis）；偏滤器位形的分离面可具有鞍形 \textbf{X 点}（X-point）：单零位形通常有一个相关 X 点，双零位形有上下两个，不能把所有截面都概括为统一的``8''字形。LCFS 是闭合磁面区的边界，并非所有等离子体物质的外边缘；限制器、偏滤器和 XPT 的物理意义见第~\ref{sec:boundary-exhaust} 节。这些几何对象不是额外假设的，而是由磁通场的拓扑后处理识别出来的（见第 \ref{sec:bcs} 节）。图~\ref{fig:crosssec} 给出这些概念的截面示意。

\begin{figure}[htbp]
\centering
\begin{tikzpicture}[font=\footnotesize,scale=1.08]
 \draw[black!30] (0,-2.2) rectangle (5.35,2.2);
 \draw[black!30] (6.15,-2.2) rectangle (11.5,2.2);
 \node at (2.7,1.95) {（a）限制器位形};
 \node at (8.8,1.95) {（b）单零偏滤器位形};
 \draw[thick] (2.7,0) ellipse (1.4 and 1.5);
 \draw[black!55] (2.7,0) ellipse (1.0 and 1.1);
 \draw[black!55] (2.7,0) ellipse (0.55 and 0.65);
 \fill (2.7,0) circle (1.5pt);
 \node[anchor=north] at (2.7,-0.08) {磁轴};
 \fill[black!30] (4.1,-0.22) rectangle (4.65,0.22);
 \draw[thick] (4.1,-0.22)--(4.1,0.22);
 \node[align=center] at (4.55,-0.75) {实体\\限制器};
 \node[anchor=east] at (1.25,0.85) {LCFS};
 \node at (2.7,-1.92) {闭合区边界与限制器接触};
 % A separate diverted equilibrium: the loop and legs share the X-point.
 \draw[thick] (8.8,-1.25)
 .. controls (9.3,-0.7) and (10.2,-0.7) .. (10.2,0.2)
 .. controls (10.2,1.0) and (9.6,1.5) .. (8.8,1.5)
 .. controls (8.0,1.5) and (7.4,1.0) .. (7.4,0.2)
 .. controls (7.4,-0.7) and (8.3,-0.7) .. (8.8,-1.25);
 \draw[black!55] (8.8,0.15) ellipse (0.95 and 0.95);
 \draw[black!55] (8.8,0.15) ellipse (0.5 and 0.5);
 \fill (8.8,0.15) circle (1.5pt);
 \node[anchor=north] at (8.8,0.07) {磁轴};
 \draw[thick] (8.8,-1.25).. controls (8.55,-1.55) and (8.1,-1.65)..(7.7,-1.9);
 \draw[thick] (8.8,-1.25).. controls (9.05,-1.55) and (9.5,-1.65)..(9.9,-1.9);
 \fill (8.8,-1.25) circle (1.6pt);
 \node[anchor=east] at (8.42,-1.1) {X 点};
 \node[anchor=west] at (10.1,0.9) {LCFS};
 \draw[line width=2.8pt,black!45] (7.2,-1.94)--(8.1,-1.94);
 \draw[line width=2.8pt,black!45] (9.5,-1.94)--(10.4,-1.94);
 \fill (7.7,-1.9) circle (1.4pt);
 \fill (9.9,-1.9) circle (1.4pt);
 \node[font=\scriptsize] at (10.55,-1.5) {靶板};
\end{tikzpicture}
\caption{两种不同平衡的极向截面示意（非求解结果）。左：限制器接触确定 LCFS；右：主 X 点定义分离拓扑，向下的分离面支路形成偏滤器腿，并与灰色靶板相交。内部细线为磁面；浅灰矩形仅表示数值覆盖范围，不能解释为真实第一壁或靶面。闭合区外仍可存在等离子体。}
\label{fig:crosssec}
\end{figure}

\subsection{极向磁通：用一个标量描述极向磁场}\label{sec:primer-flux}

在 $R$--$Z$ 截面内，极向磁场可以用一个标量函数——\textbf{极向磁通} $F(R,Z)$——完全描述。本文档采用的约定是
\[
  F=\int B_Z\,dS_Z \quad (\mathrm{Wb}),
\]
其几何构造是：把点 $P(R,Z)$ 绕 $Z$ 轴旋转一周，生成高度为 $Z$、半径为 $R$ 的水平圆盘，$F$ 等于极向磁场穿过该圆盘的磁通（构造与推导见附录 \ref{app:flux} 节及图 \ref{fig:fluxdisk}；韦伯，单位关系见式 \eqref{eq:wb}）。$F$ 好比流体力学中的\textbf{流函数}：它在截面上的等值线恰好是磁力线在截面上的投影。

\textbf{注意：磁通与磁感应强度不是一个物理量。}磁感应强度 $\bm{B}$ 是矢量场，单位 T（即 $\mathrm{Wb}/\mathrm{m}^2$，旧称磁通密度），描述局部磁场强弱；磁通是 $\bm{B}$ 对曲面的面积分（$\mathrm{Wb}$），是依赖于所选曲面的标量。$F$ 的本质是磁通——它能逐点取值只是因为每个点 $(R,Z)$ 锚定了各自的圆盘面；与 $\bm{B}=\mu_0\bm{H}$ 的区分另见第 \ref{sec:vars} 节变量表。一旦知道 $F$，极向磁场就由它的偏导数给出（式 \eqref{eq:bfield}）：
\[
  B_R=-\frac{1}{2\pi R}\frac{\partial F}{\partial Z},\qquad
  B_Z=\frac{1}{2\pi R}\frac{\partial F}{\partial R}.
\]
（式中的 $2\pi$ 来自``每弧度磁通''与``总磁通''的约定换算，见第 \ref{sec:wbu} 节：本文 $F$ 是 Wb 总磁通，有些文献用 $\psi=F/(2\pi)$，两者不能混用。）

因此，``求平衡''归根结底就是求一个满足平衡方程的标量场 $F(R,Z)$——这是一个二维椭圆偏微分方程问题，与求静电势的 Poisson 方程在数学上非常相似。

\subsection{力平衡与 Grad--Shafranov 方程}\label{sec:primer-gs}

等离子体要静止（不膨胀、不收缩），必须满足\textbf{力平衡}：压力梯度向外推，必须被磁场的洛伦兹力抵消：
\[
  \nabla p=\bm{j}\times\bm{B}.
\]
代入安培定律 $\mu_0\bm{j}=\nabla\times\bm{B}$ 并利用轴对称性，可以推导（这是标准教科书内容，属\evbg{}；完整推导见本节末``推导细节''一段，供复习用）：

\begin{enumerate}[nosep]
  \item 由 $\bm{j}\cdot\nabla F=0$ 可知电流必须\textbf{沿磁面流动}；
  \item 轴对称使电流只剩环向分量 $j_\phi$ 可自由分布，而 $j_\phi$ 必须由两个磁通函数决定：压力 $p(F)$ 和环向场函数 $T(F)\equiv RB_\phi$（注意 $T$ 不是温度，单位 T·m）；
  \item 代回安培定律，得到关于 $F$ 的椭圆方程（本文式~\eqref{eq:gsfull}）：
\end{enumerate}
\[
  \Delta^*F+4\pi^2\left(\mu_0R^2\frac{dP}{dF}+T\frac{dT}{dF}\right)=0,
  \qquad \Delta^*F\equiv R\,\partial_R\!\left(\tfrac1R\partial_RF\right)+\partial_Z^2F .
\]

\paragraph{推导细节（供复习，可跳读）。}以下推导基于理想 MHD 静态平衡（无流动、各向同性压强），只用到 0.1 的轴对称约定、0.2 的磁通构造结论与附录 \ref{app:ops}、\ref{app:ident} 的柱坐标公式；全部内容属\evbg{}标准数学。记 $\partial_\phi=0$，把磁场按 $F$ 与 $T$ 分解（磁通构造见附录 \ref{app:flux}）：
\[
B_R=-\frac{1}{2\pi R}\partial_ZF,\qquad B_Z=\frac{1}{2\pi R}\partial_RF,\qquad B_\phi=\frac{T}{R},
\]
其中 $B_R,B_Z$ 与 $F$ 的关系即式 \eqref{eq:bfield}；该表示自动满足 $\nabla\cdot\bm{B}=0$。

\textbf{第 1 步：力平衡迫使 $p$ 与 $T$ 都是磁通函数。}对 $\nabla p=\bm{j}\times\bm{B}$ 两边点乘 $\bm{B}$，得 $\bm{B}\cdot\nabla p=0$；磁力线躺在 $F$ 的等值线上（附录 \ref{app:flux} 末段），故 $p$ 沿等值线不变，即 $p=P(F)$。再用力平衡的另一种写法：由 $\nabla\cdot\bm{B}=0$ 与附录 \ref{app:ident} 的旋度恒等式，
\[
(\nabla\times\bm{B})\times\bm{B}=(\bm{B}\cdot\nabla)\bm{B}-\tfrac12\nabla(B^2),
\]
代入 $\mu_0\nabla p=(\nabla\times\bm{B})\times\bm{B}$ 并移项得 $(\bm{B}\cdot\nabla)\bm{B}=\nabla(\mu_0p+B^2/2)$。其 $\phi$ 分量左端为 $[\,(\bm{B}\cdot\nabla)\bm{B}\,]_\phi=\frac1R\,\bm{B}\cdot\nabla(RB_\phi)$——这里多出的 $B_\phi B_R/R$ 曲率项来自 $\bm{e}_\phi$ 随 $\phi$ 变化，正是附录 \ref{app:ident} 警告必须保留的项；右端梯度无 $\phi$ 分量，故 $\bm{B}\cdot\nabla(RB_\phi)=0$，即 $RB_\phi$ 沿磁力线不变：$T=RB_\phi=T(F)$。

\textbf{第 2 步：极向分量的力平衡决定极向电流。}把 $\bm{j}$ 分解为环向与极向两部分 $\bm{j}=j_\phi\bm{e}_\phi+\bm{j}_p$。力平衡的 $R$、$Z$ 分量为
\[
j_\phi B_Z-j_ZB_\phi=\partial_RP=P'(F)\,\partial_RF,\qquad
j_RB_\phi-j_\phi B_R=\partial_ZP=P'(F)\,\partial_ZF,
\]
解出极向电流（设 $B_\phi\neq0$）：
\[
j_Z=\frac{j_\phi B_Z-P'(F)\,\partial_RF}{B_\phi},\qquad
j_R=\frac{j_\phi B_R+P'(F)\,\partial_ZF}{B_\phi}.
\]
$\phi$ 分量 $j_ZB_R-j_RB_Z=0$ 由 $\bm{B}\cdot\nabla F=B_R\partial_RF+B_Z\partial_ZF=0$ 自动满足，不产生新条件：极向电流被 $j_\phi$ 唯一确定，没有额外自由度。

\textbf{第 3 步：安培定律的 $R$、$Z$ 分量联立，定出 $j_\phi$。}$\partial_\phi=0$ 时的旋度公式（附录 \ref{app:ops}）给出 $\mu_0j_R=-\partial_ZB_\phi$、$\mu_0j_Z=\frac1R\partial_R(RB_\phi)$；代入 $B_\phi=T/R$ 与 $T=T(F)$：
\[
\mu_0j_R=-\frac{T'(F)}{R}\partial_ZF,\qquad \mu_0j_Z=\frac{T'(F)}{R}\partial_RF.
\]
再代入第 2 步的解，两式都归结为同一关系
\[
j_\phi B_Z=\Bigl[P'(F)+\frac{T(F)\,T'(F)}{\mu_0R^2}\Bigr]\partial_RF
\]
（另一式只差一个由 $\bm{B}\cdot\nabla F=0$ 保证的符号因子，二者相容）。把 $B_Z=\frac{1}{2\pi R}\partial_RF$ 代入，得环向电流的显式表达：
\[
j_\phi=2\pi R\,P'(F)+\frac{2\pi}{\mu_0R}\,T(F)\,T'(F).
\]
即 $j_\phi$ 由\textbf{逆磁电流}（$P'$ 项，压强梯度驱动）与\textbf{维持电流}（$TT'$ 项，维持环向场分布）两部分组成——这就是 第~\ref{sec:primer-profiles} 节``为什么源项里是导数''的定量来源。

\textbf{第 4 步：环向安培分量给出 $F$ 的方程。}$\partial_\phi=0$ 时 $(\nabla\times\bm{B})_\phi=\partial_ZB_R-\partial_RB_Z$，代入 $B_R,B_Z$ 与 $F$ 的关系：
\[
\mu_0j_\phi=-\frac{1}{2\pi R}\Bigl[\partial_Z^2F+R\,\partial_R\Bigl(\frac1R\partial_RF\Bigr)\Bigr]=-\frac{1}{2\pi R}\Delta^*F.
\]
把第 3 步的 $j_\phi$ 代入并整理：
\[
\Delta^*F=-2\pi\mu_0R\,j_\phi=-4\pi^2\Bigl(\mu_0R^2\frac{dP}{dF}+T\frac{dT}{dF}\Bigr),
\]
与式~\eqref{eq:gsfull} 逐项一致。推导中的 $2\pi$ 幂次全部来自本文 Wb 总磁通约定 $F=2\pi RA_\phi$；若改用每弧度磁通 $\psi=F/2\pi$ 为未知量，算子形式不变、源项系数相应缩放，两者不可混用（约定见第 \ref{sec:gse} 节与附录 \ref{app:flux} 末段）。

阅读这个方程的要点：

\begin{itemize}[nosep]
  \item 左边第一项 $\Delta^*F$ 与线性椭圆算子类似（在 $R$ 很大时趋于普通 Laplace 算子），描述磁场自身的几何扩散；
  \item 右边（移项后）是\textbf{源项}：等离子体电流 $j_\phi$ 产生的磁通响应。式 \eqref{eq:gs} 给出与代码一致的形式 $\Delta^*F=-2\pi\mu_0R\,j_{\phi,\mathrm{source}}$；
  \item 电流 $j_\phi$ 本身又依赖于解 $F$（因为 $dP/dF$、$T\,dT/dF$ 都随 $F$ 变化），所以这是一个\textbf{非线性}方程；
  \item 给定 $P(F)$ 和 $T(F)$ 两个``剖面函数''后，平衡问题就封闭了——这两个函数是 GS 问题的``自由输入''。
\end{itemize}

\textbf{自由边界的含义：}等离子体的边界（LCFS）不是预先画好的，而是解出来的——边界条件由等离子体自身电流的场与外导体（线圈、室壁）的场叠加构成（式 \eqref{eq:fbdy}），所以称为\textbf{自由边界}问题。这与静电学中``导体边界形状待求''的问题类似，比固定边界问题更难。

\subsection{剖面函数与剖面系数：把``无限自由度''压缩成 3 个数}\label{sec:primer-profiles}

原则上 $P(F)$、$T(F)$ 可以是任意函数，即平衡问题有无穷多个自由输入。先明确这两个函数是什么。\textbf{$P(F)$ 是压强剖面}：由力平衡 $\nabla p=\bm{j}\times\bm{B}$ 两边点乘 $\bm{B}$ 得 $\bm{B}\cdot\nabla p=0$——压强沿磁力线不变，在嵌套磁面上即每个磁面取确定值，于是 $p$ 只能写成 $p(F)$，描述``压强如何沿磁面分布''。\textbf{$T(F)$ 是环向场函数剖面}：$T\equiv RB_\phi$（不是温度），描述环向磁场``绷得多紧''沿磁面的分布；轴对称力平衡同样强制 $RB_\phi$ 沿磁力线不变，故 $T=T(F)$。直观类比：环向场如同螺线管内的轴向场，维持它需要等离子体电流；真空的 $1/R$ 衰减不满足力平衡，$T(F)$ 规定了等离子体``重塑''环向场的方式。

为什么方程源项里是导数？GS 方程（式~\eqref{eq:gsfull}）中 $\mu_0R^2P'(F)+TT'(F)$ 恰是环向电流密度（式 \eqref{eq:gs}），两个导数对应两类电流来源：$P'(F)$ 是压强梯度驱动的逆磁电流（热压力与磁场的抗衡），$TT'(F)$ 是维持环向场分布所需的维持电流。因此``给定 $P(F)$、$T(F)$''的读法是：指定压强与环向场的分布，方程即给出所需电流及由电流产生的磁通 $F$——而电流又依赖 $F$，构成非线性闭环。

实际求解时需要把两个函数\textbf{参数化}。本文档仿真的做法是（第 \ref{sec:gse} 节）：

\begin{enumerate}[nosep]
  \item 选一族基函数 $g_\ell(F)$，把剖面导函数展开为有限项，例如当前配置用 3 个基：
  \[
    P'(F)=\frac{a_1\,g_1(F)}{2\pi\Delta S},\qquad
    TT'(F)=\frac{\mu_0\bigl[a_2\,g_2(F)+a_3\,g_3(F)\bigr]}{2\pi\Delta S};
  \]
  \item 展开系数 $\bm{a}_g=[a_1,a_2,a_3]^T$ 就是\textbf{剖面系数}——它们取代了``两个任意函数''，成为新的未知量；
  \item 3 个未知数需要 3 个独立条件才能定解：两个由控制目标给出（极向 beta $b_p$ 与轴上安全因子 $q_A$，式~\eqref{eq:rbp}、\eqref{eq:rqa}），第三个由电流演化方程（CDE，式 \eqref{eq:root} 第 4 行）给出。
\end{enumerate}

物理图像：$a_1$ 控制``压得多狠''（压力剖面形状与大小），$a_2,a_3$ 控制``环向场绷得多紧''（$T$ 的剖面）；$b_p$ 目标进一步钉住总压力水平，$q_A$ 目标钉住磁轴附近的磁场剪切，电流演化方程则决定总电流随时间如何变化。基函数在边界 $F=F_B$ 处自动取零，保证压力与 $T^2$ 在边界取参考值——这是一种\textbf{剖面参数化约束}，而不是把计算矩形边界设为零。

\subsection{电流基矩阵 $T_{yg}$：3 个系数如何变成整张电流分布图}\label{sec:primer-basis}

剖面系数本身不是电流。代码通过一个 $4032\times3$ 的矩阵 $T_{yg}$ 把剖面系数``翻译''成每个网格点上的电流（式 \eqref{eq:bfct3}）：
\[
  (T_{yg})_{(i,j),\ell}=\chi_{ij}\left[R_jf^P_\ell+R_j^{-1}f^T_\ell\right]g_\ell(F_{ij}-F_A,\,F_B-F_A),
  \qquad \bm{i}_y^{\mathrm{prof}}=T_{yg}\bm{a}_g .
\]
阅读要点：

\begin{itemize}[nosep]
  \item \textbf{矩阵的第 $\ell$ 列}是：令第 $\ell$ 个剖面系数为 1、其余为 0 时，整个 $63\times64$ 网格上出现的\textbf{电流分布图案}。乘 $R$ 还是乘 $1/R$ 取决于该基属于 $P'$ 还是 $TT'$（磁场关系式 (0.3) 中两项分别带 $R$ 与 $1/R$）；
  \item 域掩码 $\chi_{ij}$ 把等离子体域外的点置零——电流只存在于等离子体内部；
  \item $T_{yg}$ 随候选解 $F$ 更新（因为 $g_\ell$ 的自变量是磁通），所以它是\textbf{状态依赖}的，这也是平衡问题非线性的第二个来源。
\end{itemize}

于是平衡残差的第 1 行（式 \eqref{eq:root} 第 1 行）就非常直观了：
\[
  T_{yg}\bm{a}_g-\bm{i}_y=\bm{0}
  \quad\Longleftrightarrow\quad
  \text{``按剖面系数算出的电流''}=\text{``网格上实际的电流''}.
\]

\begin{figure}[htbp]
\centering
\begin{tikzpicture}[font=\footnotesize, node distance=6mm]
  \node[draw, rounded corners=2pt, inner sep=5pt, align=center] (ag) {$\bm{a}_g\in\mathbb R^3$\\剖面系数\\$[a_1,a_2,a_3]$};
  \node[draw, rounded corners=2pt, inner sep=6pt, right=14mm of ag, align=center] (tyg) {$T_{yg}$（$4032\times3$）\\电流基矩阵：第 $\ell$ 列 $=$\\单位剖面系数的电流分布图案};
  \node[draw, rounded corners=2pt, inner sep=5pt, right=14mm of tyg, align=center] (iy) {$\bm{i}_y\in\mathbb R^{4032}$\\网格环向电流\\每格一根电流丝 $I_{y,ij}$（A）};
  \draw[-{Stealth}, thick] (ag) -- (tyg);
  \draw[-{Stealth}, thick] (tyg) -- (iy);
  \node[below=9mm of tyg, align=center, text=black!75] {\small 每个内网格单元上安置一根等效环向电流丝，电流密度 $j_\phi=I_y/\Delta S$（A/m$^2$）。\\平衡要求：$\bm{i}_y^{\mathrm{prof}}=T_{yg}\bm{a}_g$ 与 $\bm{i}_y$ 一致（式 \eqref{eq:root} 第 1 行）};
\end{tikzpicture}
\caption{剖面系数、电流基矩阵与网格电流的关系。}
\label{fig:tyg}
\end{figure}

\subsection{导体回路与电路类比：等离子体也是一条``回路''}\label{sec:primer-circuits}

托卡马克的线圈和室壁导体可以看成集总\textbf{回路}：每个回路有电阻 $R$、与其他回路的互感 $M$，满足基尔霍夫电压定律——这正是式 \eqref{eq:circuit}：
\[
  M_{ee}\frac{\bm{i}_e^{n+1}-\bm{i}_e^n}{h}+M_{ey}\frac{\bm{i}_y^{n+1}-\bm{i}_y^n}{h}+D_R\bm{i}_e^{n+1}=\bm{v}_e^{n+1}.
\]
用电路语言读：$M\,di/dt$ 是感应电压（自感与互感），$RI$ 是电阻压降，右端是外加电源电压。等离子体本身也被当作一条``大回路''（式 \eqref{eq:cde}）：总电流 $I_p$ 在总电感 $L_p$ 上变化产生感应电压，电阻 $R_p$ 上有欧姆压降，还接受来自 537 条导体回路的互感耦合。也就是说：\textbf{平衡方程（磁场）与回路方程（电路）耦合在一起}，必须联立求解——这就是文档所说的耦合 DAE（微分--代数系统）。

\subsection{从连续方程到可计算的代数方程}\label{sec:primer-discrete}

计算机只能处理有限个数字，所以要经过两步离散（第 \ref{sec:discret} 章）：

\begin{enumerate}[nosep]
  \item \textbf{空间离散}：把连续的 $F(R,Z)$ 变成矩形网格上 $65\times66$ 个节点值；导数用相邻点的差分近似（五点格式，式 \eqref{eq:radial}）。网格是固定的矩形张量积网格，等离子体的弯曲边界由磁通值加几何掩码表示，不需要贴体网格；
  \item \textbf{时间离散}：连续的时间导数 $\dot{\bm{i}}$ 用后向差商 $(\bm{i}^{n+1}-\bm{i}^n)/h$ 近似（后向 Euler，式 \eqref{eq:resp}）。于是``推进一个时间步''就变成了：解一个 4572 维的非线性代数方程组。
\end{enumerate}

直觉上，后向 Euler 是``隐式''的：新时刻的所有未知量同时出现在方程里，解起来贵，但对快变过程稳定（式 \eqref{eq:amp} 的放大因子分析）。文档没有自适应步长控制，步长 $h$ 由初始化固定为 1 ms 或沿用载入数据的时间轴。

\subsection{非线性求解器在做什么}\label{sec:primer-solver}

每个时间步要求解 $\bm{\mathcal R}(\bm{x})=\bm{0}$（式 \eqref{eq:root} 的缩放版本）。求解器（第 \ref{sec:solver} 章）的行为可以这样理解：

\begin{itemize}[nosep]
  \item \textbf{Newton 迭代}：在当前猜测 $\bm{x}_k$ 附近把残差线性化，解线性方程组 $J_k\delta\bm{x}_k=-\bm{\mathcal R}(\bm{x}_k)$ 得到修正方向，再沿该方向走一步。反复迭代直到残差二范数小于容差 $\texttt{tolF}=10^{-4}$；
  \item \textbf{JFNK}：Jacobian 矩阵 $J_k$ 太大（$4572^2$），不显式构造，只在需要时用差分近似它与向量的乘积 $J_k\bm{v}$（式~\eqref{eq:jvec}）；
  \item \textbf{GMRES}：解 Newton 线性方程组时不求逆，而是在 Krylov 子空间里找最佳近似解（式 \eqref{eq:arnoldi}）；预条件器 $P^{-1}$（式 \eqref{eq:prec}）相当于给一个``近似逆''让迭代更快收敛；
  \item \textbf{线搜索与回退}：Newton 方向不一定直接走满，要沿线找一个使残差不增的步长（式 \eqref{eq:armijo}）；若彻底失败，保留上一时刻的物理状态，只推进时间（第 \ref{sec:linesearch} 节）；
  \item \textbf{容差的读法}：$10^{-4}$ 是\textbf{缩放后}残差的二范数，不是任何物理量的绝对误差（如不是 $10^{-4}$ Wb 或 $10^{-4}$ A），见第 \ref{sec:confirmable} 节。
\end{itemize}

\subsection{如何阅读本文档}\label{sec:primer-route}

本报告第~\ref{sec:primer} 章建立平衡与数值求解的物理直觉，第~\ref{sec:microphysics} 章补全粒子、碰撞、自洽场及典型结构，第~\ref{sec:lifecycle} 章解释托卡马克放电过程，第~\ref{sec:stellarator} 章解释仿星器线圈设计，第~\ref{sec:intro}--\ref{sec:rl} 章给出仿真实现案例。按读者目标给出路线图：

\begin{itemize}[nosep]\small
  \item \textbf{初次学习托卡马克控制}：读本章 $\to$ 第~\ref{sec:microphysics} 章的微观与集体图景 $\to$ 第~\ref{sec:lifecycle} 章的全过程导读 $\to$ 第 \ref{sec:model} 章（案例模型）与第 \ref{sec:interface} 章（控制接口）；
  \item \textbf{想理解仿星器如何设计}：第~\ref{sec:microphysics} 章（粒子与场）$\to$ 第~\ref{sec:stellarator} 章（线圈、性质与工程取舍）；
  \item \textbf{只想理解物理模型}：读本章 $\to$ 第 \ref{sec:model} 章（全部模型与方程）$\to$ 第~\ref{sec:approx} 节（近似与局限）；
  \item \textbf{想理解数值方法}：第 \ref{sec:discret} 章（离散）$\to$ 第 \ref{sec:solver} 章（求解器）$\to$ 第~\ref{sec:confirmable}--\ref{sec:complexity} 节（稳定性、复杂度）；柱坐标下的微分算子等数学工具见附录 B；
  \item \textbf{想运行或复现仿真}：第 \ref{sec:workflow} 章（工作流）$\to$ 第 \ref{sec:interface} 章（接口）$\to$ 附录 A（复现清单）；
  \item \textbf{想核对某个结论的出处}：第 \ref{sec:mapping} 章的方程--函数--数据索引表直接给出每个结论对应的源码位置；引用前缀（\code{G:}、\code{V/}、\code{D/} 等）的含义见``关于本文档''一节。
\end{itemize}

前四章采用 [背景] 标签；案例的证据标签按``关于本文档''阅读。建议先理解耦合方程组~\eqref{eq:root} 的物理任务，再逐章追踪其实现。

\clearpage

% ================= 微观过程、集体结构与模型层次 =================
\clearpage
\section{微观粒子、集体电磁场与等离子体结构}\label{sec:microphysics}
\markboth{粒子、场与结构}{粒子、场与结构}

\textbf{本章定位。}把等离子体想成``在磁笼里被加热的一团气体''可以入门，却不足以解释电流、碰撞和形状为何相互影响。本章沿着粒子运动、碰撞反应、场的自洽响应、空间分布和模型层次建立连续的物理图景，最后讨论各模型能解释的典型结构。所有内容属于\evbg{}通用理论；示意图不采用任何装置的实测数据，教学公式均说明近似条件。各过程在一次放电中的先后关系见第~\ref{sec:lifecycle} 章，案例求解器实际保留的方程见第~\ref{sec:fullsystem} 节。

\subsection{电场加速以后，什么叫电流，什么叫热}\label{sec:micro-heat}

\textbf{平均运动与随机运动是同一批粒子的两种统计量。}在一个足够小但包含许多粒子的体积中，第 $s$ 类粒子的瞬时速度可写成
\begin{equation}\label{eq:micro-velocity}
 \bm v=\bm u_s+\bm c_s,\qquad
 \bm u_s=\langle\bm v\rangle_s,\qquad \langle\bm c_s\rangle_s=\bm0.
\end{equation}
$\bm v$ 是单粒子速度，$\bm u_s$ 是该物种的局部平均流速，$\bm c_s$ 是相对平均流的随机速度，均以 m/s 计；尖括号表示按该物种的局部速度分布取平均。粒子即使热运动很快，平均速度也可能很小，因为不同方向相互抵消。对非相对论、近各向同性的分布，有
\begin{equation}\label{eq:micro-energy-split}
 \frac12m_sn_s\langle v^2\rangle_s
 =\frac12m_sn_su_s^2+\frac32n_sk_BT_s,
 \qquad k_BT_s=\frac{m_s}{3}\langle c_s^2\rangle_s.
\end{equation}
$m_s,n_s$ 是粒子质量（kg）和数密度（$\mathrm{m^{-3}}$）；左侧为总平动动能密度（$\mathrm{J/m^3}$），右侧第一项为整体流动能密度，第二项为随机热能密度。$k_B$ 为 Boltzmann 常数，$T_s$ 为以 K 表示的温度，$u_s=|\bm u_s|$。强各向异性时应使用压力张量或平行/垂直温度，单个标量 $T_s$ 不能完整描述速度分布。

电场可以同时改变平均流、分布宽度或非热尾部，具体结果取决于场的方向、时间变化与碰撞。均匀电场把一群完全无碰撞、初始同速的粒子一起加速时，可主要增加平均流动能，而不立即增加相对平均流的随机能；所以``粒子速度更大''不能无条件等同于``温度更高''。对热化中的真实分布，快粒子通过相互作用把能量交给更多粒子，随机速度分布通常变宽。电子与离子质量差别很大，加热也不必使 $T_e$ 与 $T_i$ 立即相等。

\textbf{电流取决于正负电荷的相对平均运动。}对于单电离氢同位素等离子体，在准中性 $n_i\simeq n_e=n$ 的近似下，
\begin{equation}\label{eq:micro-current}
 \bm j=en(\bm u_i-\bm u_e).
\end{equation}
$e>0$ 是元电荷，$\bm j$ 为电流密度（$\mathrm{A/m^2}$），$\bm u_i,\bm u_e$ 分别是离子、电子平均流速。电子朝一个方向漂移形成的常规电流朝相反方向；正负粒子若以同一速度整体运动，主要形成质量流，而净电流可很小。各向同性随机热运动本身并不自动产生宏观净电流。

为显示欧姆加热的能量来源，可暂取离子平均静止、电子惯性可忽略、无压强梯度的局部模型：
\begin{equation}\label{eq:micro-ohmheat}
 0=-e\bm E_\parallel-m_e\nu_{ei}\bm u_{e\parallel},\qquad
 \eta=\frac{m_e\nu_{ei}}{n_e e^2},\qquad
 p_{\rm Ohm}=\bm j_\parallel\cdot\bm E_\parallel
 =\eta j_\parallel^2.
\end{equation}
$\parallel$ 表示沿局部磁场方向，$\nu_{ei}$ 为电子相对离子的有效动量松弛频率（$\mathrm{s^{-1}}$）；第一式是电场力与平均碰撞阻力平衡。$m_e$ 为电子质量，$\eta$ 是这一简化模型的电阻率（$\Omega\,\mathrm m$），$p_{\rm Ohm}$ 是局域加热功率密度（$\mathrm{W/m^3}$）。电场不断向相对漂移输入能量，碰撞使其有序动量松弛，输入功最终进入随机能库。真实 Spitzer 电导率包含分布函数修正，本式用于说明机制，不能替代其完整系数。

孤立体系内的弹性碰撞守恒总能量，不能凭空造热；加热的能源是电源、射频波、注入束或聚变产物。电子--离子散射可以较快改变电子的方向与平均动量，却不意味着同样快地把电子的大部分能量交给重离子；电子内部热化、离子内部热化与电子--离子温度均衡有不同时间尺度。中性束先穿入、再电离形成快离子并减速，射频则可先使特定速度的粒子吸收波能，之后再发生碰撞传能；这些过程也可能保留非 Maxwell 尾部，不能都归结成一次机械撞击\cite{Pheat}。

\subsection{碰撞究竟碰到了什么，会不会产生新物质}\label{sec:micro-collisions}

\textbf{充分电离等离子体中的常见碰撞是电磁散射。}电子和离子不必像小球那样接触表面，便会因 Coulomb 力改变轨道。弱耦合条件下，大量远距离小角偏转的累计作用，比罕见的强偏转更重要。轨道相互作用可重新分配方向、动量和能量；只要没有原子或核反应，粒子种类与核组成就不改变。以经典两体弹性碰撞为例，在远离相互作用区比较入射和出射状态，
\begin{align}
 m_a\bm v_a+m_b\bm v_b&=m_a\bm v'_a+m_b\bm v'_b,\label{eq:micro-momentum}\\
 \frac12m_av_a^2+\frac12m_bv_b^2&=
 \frac12m_a{v'_a}^2+\frac12m_b{v'_b}^2.\label{eq:micro-elastic}
\end{align}
$a,b$ 标记两粒子，$m$ 为质量，$\bm v$ 和带撇的 $\bm v'$ 分别为碰撞前后速度；两式依次表达总动量与平动动能守恒。在相互作用过程内部还需计入势能；若有辐射、激发、电离或核反应，不能只用第二式的平动动能守恒来核算全部能量。Coulomb 散射的经典轨道可由两体力学推导，热化则涉及大量散射的统计作用\cite{Pcoulomb}。

\begin{table}[htbp]
\centering\small
\caption{不同``碰撞''改变的对象及其物理后果}\label{tab:micro-reactions}
\begin{tabular}{L{2.3cm}L{4.2cm}L{7.0cm}}
\toprule
过程 & 改变什么 & 是否改变原子核及控制相关后果 \\
\midrule
弹性散射 & 运动方向、各粒子的动能分配 & 核种不变；引起动量松弛、热化与碰撞输运 \\
激发/退激发 & 束缚电子的内部能级 & 核种不变；激发消耗电子能量，退激发可发光 \\
电离/复合 & 束缚电子数与电荷态 & 核种不变；改变自由电子数、辐射及导电性 \\
分子解离 & 分子键与组分 & 核种不变；例如分子变原子，需要解离能 \\
电荷交换 & 哪个粒子携带束缚电子 & 核种不变；可形成逃出磁约束的快中性粒子 \\
核聚变 & 核内质子/中子的组合 & 可生成不同核种；改变燃料组分并释放核能 \\
\bottomrule
\end{tabular}
\end{table}

以氘为例，激发、电离和辐射复合可分别表示为
\begin{align}
 e^-+\mathrm D^0&\longrightarrow e^-+\mathrm D^{0*},\label{eq:micro-excitation}\\
 e^-+\mathrm D^0&\longrightarrow \mathrm D^++2e^-,\label{eq:micro-ionization}\\
 e^-+\mathrm D^+&\longrightarrow \mathrm D^0+\gamma_{\rm ph}.\label{eq:micro-recombination}
\end{align}
$e^-$ 为自由电子，$\mathrm D^0$ 为中性氘原子，星号表示其束缚电子被激发，$\mathrm D^+$ 为失去束缚电子的氘离子；$\gamma_{\rm ph}$ 为带走能量的光子，不是增长率。电离式有两个出射电子，是因为入射电子仍然存在，同时释放出原先束缚的电子。电离与激发通常需要入射电子提供足够能量；辐射复合则通过光子释放结合能。三体复合、电离态级联及分子反应在相应条件下应另外计入。上述反应可生成新的电荷态或粒子组分，但氘核仍是氘核；不能把电离叫作生成氦。

氘离子与氘原子之间还可发生电荷交换：
\begin{equation}\label{eq:micro-cx}
 \mathrm D^+_{\rm fast}+\mathrm D^0_{\rm slow}
 \longrightarrow\mathrm D^0_{\rm fast}+\mathrm D^+_{\rm slow}.
\end{equation}
fast、slow 标记反应前后的快慢粒子群；在这一示意中，一个电子从慢原子转移给快离子，核的速度在短暂交换中近似保持。于是快粒子变成中性粒子，不能再直接被磁场约束，可能把能量带到壁面；慢粒子则变成离子参与等离子体运动。它与 D--T 聚变的区别是核组成没有改变。有关原子与分子反应过程及其反应系数可参考 EIRENE 文档\cite{Patomic}。核反应需要另用反应截面和核能收支，详见第~\ref{sec:cycle-fuels} 节。大量普通散射并不意味着每次散射都发生聚变。

\begin{figure}[!ht]
\centering
\begin{tikzpicture}[font=\small,>=Stealth]
 \node at (3,2.0) {（a）普通散射：偏转与能量交换};
 \draw[->,thick,blue!70!black] (0.3,0.6)--(2.2,0.6);
 \fill[blue!65] (1.7,0.6) circle (0.09);
 \draw[->,thick,blue!70!black] (2.4,0.6)..controls(3.1,0.65)and(3.2,1.2)..(4.6,1.35);
 \fill[orange!85!black] (2.8,0.1) circle (0.14);
 \draw[->,thick,orange!85!black] (2.95,-0.05)--(4.2,-0.8);
 \node[align=center] at (2.8,-1.55) {相互作用可在远距离发生\\核种类不变，不是每次碰撞都聚变};
 \node at (9.8,2.0) {（b）D--T 聚变：核组成改变};
 \fill[blue!65] (7.8,0.9) circle (0.18);
 \node[above] at (7.8,1.1) {D};
 \fill[orange!85!black] (7.8,-0.1) circle (0.21);
 \node[below] at (7.8,-0.35) {T};
 \draw[->,thick] (8.25,0.4)--(9.4,0.4);
 \fill[red!65] (10.1,0.95) circle (0.25);
 \node[above] at (10.1,1.2) {${}^4$He};
 \fill[black!55] (10.9,-0.2) circle (0.1);
 \node[right] at (11.1,-0.2) {$n$};
 \draw[->,red!70!black] (9.6,0.55)--(9.95,0.85);
 \draw[->,black!60] (9.6,0.25)--(10.7,-0.1);
 \node[align=center] at (10,-1.55) {满足核反应概率才生成新核\\释放能量由产物携带};
\end{tikzpicture}
\caption{把``碰撞''分成不同物理过程。左图用两粒子偏转示意普通散射；真实 Coulomb 松弛常由许多小角散射累积。右图以 D--T 为核反应例子，产物为氦核和中子；图中距离、角度及速度未按比例。电离、激发和电荷交换则属于前文单独列出的原子过程。本文原创教学图。}\label{fig:micro-collision-atlas}
\end{figure}

\subsection{密度怎样影响碰撞：概率、屏蔽和时间尺度}\label{sec:micro-rates}

\textbf{固定速度分布与靶种类时，密度增加会增加单粒子的相互作用机会。}对可以用有效二体截面描述的指定过程，
\begin{equation}\label{eq:micro-rate}
 \nu_{ab}=n_b\langle\sigma_{ab}v_{\rm rel}\rangle,\qquad
 P_{\rm event}(\Delta t)=1-e^{-\nu_{ab}\Delta t}.
\end{equation}
$\nu_{ab}$ 是一个 $a$ 粒子与 $b$ 靶粒子发生该类过程的频率（$\mathrm{s^{-1}}$），$n_b$ 为靶密度；$\sigma_{ab}$ 为该过程截面（$\mathrm{m^2}$），$v_{\rm rel}$ 为相对速度，尖括号表示对相关速度分布平均。若频率在 $\Delta t$ 内固定且采用独立 Poisson 事件模型，$P_{\rm event}$ 是至少发生一次事件的概率。$\nu_{ab}\Delta t\ll1$ 时，该概率近似为 $\nu_{ab}\Delta t$；较大时必须使用完整表达式，不能让线性近似超过 1。分布随加热变化时，截面平均也跟着改变，不能只看密度。

对完全相同的两类入射粒子，单位体积的二体事件率另有配对因子 $1/2$；对不同种类则为 $n_an_b\langle\sigma v\rangle$。这些因子解释了为什么聚变总反应率常与密度平方有关，但并不说明提高密度就足以解决能量平衡：更多粒子也要加热，某些辐射和原子反应损失也随密度增加，供气还会改变燃烧比例、边缘状态和稳定性。

\textbf{Coulomb 碰撞频率需指定``累计偏转多少''。}长程力使许多远处粒子同时施加微小作用，因而不能把每次极小偏转都当作一个有明确边界的小球碰撞。有效偏转频率通常表示累计到显著角度的松弛速率，动量松弛和能量交换的频率还不完全相同。一个常用的量级关系是
\begin{equation}\label{eq:micro-coulomb-scaling}
 \nu_{ei}\propto
 \frac{n_e Z_{\rm eff}\ln\Lambda}{T_e^{3/2}},\qquad
 \lambda_{\rm mfp}\sim\frac{v_{{\rm th},e}}{\nu_{ei}}.
\end{equation}
这里采用热电子、充分电离、弱耦合的经典碰撞近似；$Z_{\rm eff}$ 为有效离子电荷，$\ln\Lambda$ 为 Coulomb 对数，第一式仅给标度，不提供依赖温度单位的数值系数。$v_{{\rm th},e}$ 为代表性电子热速度，$\lambda_{\rm mfp}$ 为相应累计散射平均自由程（m）。固定 $T_e$ 时，密度更高通常使该频率更高；温度升高则可能使它下降。因此一个高温稠密的聚变等离子体仍可能在某些过程的时间尺度上弱碰撞\cite{Pcollisions}。

集体电荷响应还会屏蔽长程静电势。以电子提供屏蔽、离子背景近似固定的量级估计为例，
\begin{equation}\label{eq:micro-debye}
 \lambda_{De}=\sqrt{\frac{\epsilon_0k_BT_e}{n_ee^2}},\qquad
 N_D=\frac{4\pi}{3}n_e\lambda_{De}^3.
\end{equation}
$\lambda_{De}$ 是电子 Debye 长度（m），$\epsilon_0$ 为真空介电常数，$N_D$ 为半径 $\lambda_{De}$ 球体内的平均电子数。多物种能共同响应时应合并各物种的屏蔽贡献。尺度远大于 Debye 长度时，等离子体体区常近似准中性，但边界鞘层等区域不能强行如此处理。弱耦合通常要求一个屏蔽体积中有许多粒子，平均场和累计小角散射才有清楚的尺度划分。Coulomb 对数可理解为 $\ln(b_{\max}/b_{\min})$：$b$ 是碰撞参量，最大有效尺度与屏蔽有关，最小尺度则受强散射或量子效应限制，并非任意选取。

\subsection{粒子产生电磁场，电磁场又反过来改变粒子}\label{sec:micro-selfconsistent}

一个运动电荷产生电磁场，但宏观控制使用的不是把所有微观场逐项画出来的图。更有用的是由分布函数得到平均电荷和电流，再求平均场：
\begin{equation}\label{eq:micro-sources}
 \rho_q=\sum_s q_sn_s,\qquad
 \bm j=\sum_s q_sn_s\bm u_s.
\end{equation}
$\rho_q$ 是净电荷密度（$\mathrm{C/m^3}$），$q_s$ 是第 $s$ 类粒子的带符号电荷，$n_s,\bm u_s$ 是其密度和平均流速。每个符号在求和内按物种取值，正负电荷分别计入。准中性意味着这些电荷密度近乎抵消，不要求各物种的平均流相同，所以 $\rho_q\simeq0$ 与 $\bm j\neq\bm0$ 完全可以同时成立。

自洽场满足 Maxwell 方程：
\begin{equation}\label{eq:micro-maxwell}
 \begin{aligned}
 \nabla\cdot\bm E&=\rho_q/\epsilon_0,&\qquad \nabla\cdot\bm B&=0,\\
 \nabla\times\bm E&=-\partial_t\bm B,&
 \nabla\times\bm B&=\mu_0\bm j+\mu_0\epsilon_0\partial_t\bm E.
 \end{aligned}
\end{equation}
$\bm E,\bm B$ 分别为平均电场和磁场，$\partial_t$ 表示对时间偏导，$\nabla\cdot$ 和 $\nabla\times$ 分别为散度和旋度。四个关系依次说明电荷产生电场散度、磁场无孤立源、变磁场产生旋涡电场、电流和变电场产生磁场旋度；最后一项是位移电流。低频、非相对论 MHD 常忽略位移电流，但这不适用于所有高频波或 Debye 尺度过程。边界上外部线圈电流和电源条件也必须给定，不能只写方程而不指定装置。

因此有一条闭合反馈链：场改变粒子轨道 $\to$ 分布与电流改变 $\to$ 自场改变 $\to$ 轨道再次改变。等离子体电流不是一根永远保持形状的导线，压力和电流分布变化会使平衡与线圈响应一起改变。动理学中的平均 Lorentz 力描述集体场，碰撞算子描述粗粒化后微观相关性的效应；使用平均场模型时，不能又把已经计入平均场的作用重复当成碰撞阻力\cite{Poperator}。

\subsection{托卡马克里的总磁场：四类来源怎样拼成一幅图}\label{sec:micro-field}

\textbf{先把方向与来源分开。}环向是绕真空室大圆的一周，极向是围绕局部等离子体小截面的一周。TF 线圈的导体主要沿极向绕装置布置，却产生主要环向磁场；PF 线圈与 CS 的导体包含绕装置中轴的环向电流，产生位于 $R$--$Z$ 平面内的极向场。不能根据线圈名称，误把导体方向与其所产生的磁场方向混为一谈。实际磁体功能参见文献~\cite{C2,C3}。

在线性磁性近似下，可以对\emph{给定电流分布}的磁场作源项分解：
\begin{equation}\label{eq:micro-superposition}
 \bm B=\bm B_{\rm TF}+\bm B_{\rm CS}+\bm B_{\rm PF}
 +\bm B_{\rm plasma}+\bm B_{\rm wall}.
\end{equation}
右侧依次为环向场线圈、中心螺线管、极向场线圈、等离子体和被动导体电流产生的场，每项单位均为 T。叠加是 Maxwell 场方程对给定源的性质；等离子体电流和壁涡流本身受总场影响，故整个平衡演化问题并不因此成为线性。含铁磁饱和等非线性材料时，也不能原样采用该线性源项叠加。

\begin{table}[htbp]
\centering\small
\caption{从磁场来源理解分工和耦合}\label{tab:micro-field-sources}
\begin{tabular}{L{2.0cm}L{5.4cm}L{6.1cm}}
\toprule 来源 & 主要场图景 & 需要避免的误解 \\
\midrule
TF & 沿大圆方向的强环向场，近似随 $1/R$ 变化 & 不单独建立所需极向磁面和等离子体电流 \\
CS & 主要在自身内部储有轴向磁通，外部漏场参与极向场；变磁通产生环向感应电场 & 核心驱流作用来自磁通变化，不是靠其漏场取代全部约束场 \\
PF / VS & 极向场叠加，提供平衡、形状与快速位置控制所需外场 & 一只线圈一般同时影响多个几何量 \\
等离子体电流 & 主要环向电流围绕小截面产生极向场；极向电流也可修正环向场 & 其分布不是外部事先固定的，压力也参与自洽响应 \\
被动结构 & 感应涡流对场的变化产生屏蔽和延迟 & 无外加电源不意味着电流或磁力作用为零 \\
\bottomrule
\end{tabular}
\end{table}

理想长螺线管是理解 CS 的好起点：其内部可以有较强轴向场，外部局部场较小，却仍可因所链磁通随时间变化而在外围闭合回路中产生环电压。因此，不能只看等离子体位置的 CS 局部磁场大小，就判断它能否提供感应电场。真实有限长 CS 的漏场、PF 补偿及壁电流都要纳入启动真空场计算；这也是启动需协调多组线圈的原因，见第~\ref{sec:cycle-coils} 节。

\begin{figure}[htbp]
\centering
\begin{tikzpicture}[>=Stealth,font=\small]
 % Toroidal view.
 \begin{scope}[shift={(0,0)}]
 \node[font=\bfseries] at (2.5,3.0) {俯视：环向场};
 \fill[blue!5,even odd rule] (2.5,0.6) circle (1.9) (2.5,0.6) circle (0.95);
 \draw[black!65] (2.5,0.6) circle (1.9);
 \draw[black!65] (2.5,0.6) circle (0.95);
 \draw[blue!75!black,thick,->] (3.95,0.6) arc[start angle=0,end angle=95,radius=1.45];
 \draw[blue!75!black,thick,->] (1.05,0.6) arc[start angle=180,end angle=275,radius=1.45];
 \node[blue!75!black] at (4.6,1.8) {$B_\phi$};
 \draw[fill=black!12] (2.5,0.6) circle (0.45);
 \node[font=\footnotesize] at (2.5,0.6) {CS};
 \node[align=center,font=\footnotesize] at (2.5,-1.8) {TF 产生绕大圆的场\\图中不画实际 TF 导体};
 \end{scope}
 % Poloidal section.
 \begin{scope}[shift={(7.3,0)}]
 \node[font=\bfseries] at (2.0,3.35) {极向截面：磁面与外线圈};
 \draw[rounded corners=16pt,black!55] (0.45,-1.2) rectangle (3.8,2.45);
 \draw[fill=black!15] (-0.3,-1.25) rectangle (0,2.5);
 \node[rotate=90,font=\footnotesize] at (-0.15,0.6) {CS};
 \foreach \x/\y in {0.65/2.65,3.6/2.65,0.65/-1.45,3.6/-1.45}{
   \draw[fill=orange!25] (\x,\y) circle (0.15);
 }
 \node[font=\footnotesize] at (3.8,2.25) {PF};
 \draw[blue!70!black,thick] (2,0.8) ellipse (0.53 and 0.75);
 \draw[blue!70!black,thick] (2,0.8) ellipse (0.82 and 1.1);
 \draw[blue!70!black,thick] (2,-0.85) .. controls (0.7,-0.05) and (0.65,1.5) .. (1.5,2.0)
 .. controls (2.9,2.8) and (3.65,0.45) .. (2,-0.85);
 \draw[blue!70!black] (2,-0.85) -- (1.45,-1.2);
 \draw[blue!70!black] (2,-0.85) -- (2.65,-1.2);
 \fill (2,0.8) circle (0.035);
 \node[below,font=\footnotesize] at (2,0.73) {磁轴};
 \node[font=\footnotesize,fill=white,inner sep=1pt] at (2,-0.86) {X};
 \draw[red!75!black,thick,->] (2.82,0.8) -- (2.82,1.3);
 \node[red!75!black,right,font=\footnotesize] at (3.02,1.25) {$\bm B_p$};
 \node[blue!75!black] at (1.62,1.0) {$\odot$};
 \node[font=\footnotesize,align=center] at (2,-2.4) {$B_\phi$ 垂直截面，$\bm B_p$ 沿轮廓切向\\PF 符号只标位置，不指定电流正负};
 \end{scope}
\end{tikzpicture}
\caption{同一磁场的两种投影示意。左图只强调环向场方向；右图的蓝线表示磁面的截面轮廓，不是一幅在平面内闭合的真实三维总场线轨迹。$\odot$ 表示穿出纸面的分量，红箭头表示局部极向分量；画法不对应数值平衡或具体装置比例。}\label{fig:micro-fields}
\end{figure}

在等离子体电流已经建立的圆截面大长径比近似下，
\begin{equation}\label{eq:micro-pitch}
 \bm B\simeq B_\phi\bm e_\phi+B_\theta\bm e_\theta,\qquad
 \frac{d\theta}{d\phi}\simeq\frac{R_0B_\theta}{rB_\phi}=\frac1{q(r)}.
\end{equation}
$\phi$ 为环向角，$\theta$ 为小截面的极向角，$\bm e_\phi,\bm e_\theta$ 为相应单位向量；$R_0$ 为参考主半径，$r$ 为小半径，$q$ 为安全因子。第二式描述\emph{沿磁力线}前进时两个角度的变化比，不是粒子的回旋频率比。总场线一边绕大圆，一边绕小截面，通常近似留在同一磁面上；在非有理绕行比的理想嵌套磁面上，它可以不断绕行而不在有限圈后闭合。磁面是二维曲面，磁力线是一维曲线，两者不应互换。

在分离面 X 点，$\bm B_p=\bm0$，总场通常仍含 $B_\phi$。启动时若 $I_p$ 尚小，可能没有后续的完整闭合区；PF/CS 场主要帮助形成适合击穿的真空场和长连接长度。建流后等离子体自场与外部极向场共同形成磁面；位形转换时是这组场的自洽组合变化，不是把一团带固定自场的物质机械推到新位置。

\subsection{磁约束如何限制轨道，为什么仍有粒子逃逸}\label{sec:micro-orbits}

在尺度远大于回旋半径、变化慢于回旋周期的磁场中，可将粒子轨道理解成``快速回旋 + 导引中心沿场运动 + 缓慢漂移''。静磁场不做功，但约束轨道后可以减少直接横向逃逸。局部均匀场中的回旋频率与电漂移为
\begin{equation}\label{eq:micro-drifts}
 \Omega_s=\frac{|q_s|B}{m_s},\qquad
 \bm v_E=\frac{\bm E\times\bm B}{B^2}.
\end{equation}
$\Omega_s$ 为角频率（rad/s），$q_s,m_s$ 为电荷与质量；$\bm v_E$ 为导引中心的 $\bm E\times\bm B$ 漂移速度（m/s）。在该近似下电漂移对正负电荷相同；$\Omega_s$ 使用电荷绝对值只表示频率大小。此漂移不是沿电场无限加速，不能把所有外加电场的功都写成电漂移的加热\cite{Porbits}。

托卡马克中 $B$ 大致在内侧较强、外侧较弱，场线本身也有曲率，粒子因而有梯度与曲率漂移。若仅有纯环向磁场，这些漂移会使不同电荷产生不利的分离及自洽电场，不能获得所需的简单轴对称约束。增加极向场，使粒子沿场绕行时不断经过内外侧，有助于组织其整体轨道与平均漂移；这并不使跨磁面损失完全消失。有限轨道宽度、碰撞、波动、磁岛与边界都会影响约束。

沿非均匀磁场运动还可出现磁镜反射。在绝热、无碰撞且近似无平行电场时，
\begin{equation}\label{eq:micro-mirror}
 \mu_s=\frac{m_sv_\perp^2}{2B}\simeq\text{常数},\qquad
 \frac12m_sv_\parallel^2+\mu_s B\simeq\text{常数}.
\end{equation}
$\mu_s$ 为第一绝热不变量（J/T），不是磁导率；$v_\perp,v_\parallel$ 为相对局部磁场的垂直与平行速度。第一式说明进入强场区时垂直动能增大；第二式表示总动能近似守恒，故平行动能减少。若 $v_\parallel$ 降到零，导引中心反射；能走完整个环的称为通过粒子，在强弱场之间往返的称为俘获粒子。其在极向投影中的有限宽度轨道可呈香蕉形，相关碰撞输运称为新经典输运的一部分。边缘仍有沿开放场线到靶板的平行损失，详见第~\ref{sec:cycle-breakdown} 节。

\subsection{一团等离子体并不均匀：密度、温度、电流各有剖面}\label{sec:micro-profiles}

应把宏观量想成空间与时间的函数 $n_s(R,Z,t)$、$T_s(R,Z,t)$、$\bm j(R,Z,t)$，而不是每项只有一个数。同一时刻可以出现较热核心、较冷边缘、不同的电子与离子温度，以及非均匀电流密度；密度既可能中心峰化，也可能平坦或空心，取决于供给位置、粒子钳制、排气及输运。相同总粒子数和相同总电流，可以对应很不相同的局部状态。总电流为
\begin{equation}\label{eq:micro-profile-current}
 I_p=\int_{S_p}j_\phi(R,Z)\,dS.
\end{equation}
$S_p$ 是等离子体极向截面，$j_\phi$ 为环向电流密度（$\mathrm{A/m^2}$），$dS$ 为截面面积元；积分得到 A。固定此积分并不能确定每一点的 $j_\phi$，因此控制总电流不等于控制安全因子剖面。

在慢变、近轴对称且具有嵌套磁面的体区，常用磁面标签 $\varrho$ 表示径向位置：
\begin{equation}\label{eq:micro-fluxlabel}
 \varrho=\sqrt{\frac{F-F_A}{F_B-F_A}},\qquad
 n_s\simeq n_s(\varrho,t),\quad T_s\simeq T_s(\varrho,t).
\end{equation}
$F$ 为总极向磁通，$F_A,F_B$ 分别为磁轴及所取边界磁通；在同一闭合区内比值通常由 0 变化到 1，$\varrho$ 无量纲。用 $\varrho$ 而不用 $\rho$ 是为了与质量/电荷密度区别。后两式是假设同一磁面上的变化相对较小，不是普遍恒等式：强旋转、各向异性、边缘不对称和快速事件都可能破坏它；X 点及开放场线区域也需另外处理。

静态标量压强平衡能推出 $p$ 沿场线不变，从而在相应条件下写成 $p=P(F)$；它\textbf{不能独立推出密度和温度都各自完全均匀}，因为同一个压力可由不同 $n,T$ 组合产生。沿场热传导常比跨场热传导有效，因此等温磁面的近似有其动力学依据，但密度、流动和杂质仍需要各自守恒与输运方程。图~\ref{fig:micro-profiles} 只用归一化曲线展示可出现的剖面类型，不代表从平衡方程计算出的预测。

\begin{figure}[htbp]
\centering
\begin{tikzpicture}[>=Stealth,font=\small,x=5.7cm,y=2.6cm]
 \begin{scope}
 \draw[->] (0,0) -- (1.1,0) node[right] {$\varrho$};
 \draw[->] (0,0) -- (0,1.15) node[above] {$n/n_{\rm ref}$};
 \draw[blue!75!black,thick,domain=0:1,samples=60] plot (\x,{0.18+0.8*(1-\x*\x)});
 \draw[orange!85!black,thick,dashed,domain=0:1,samples=60] plot (\x,{0.24+0.55*exp(-((\x-0.55)/0.38)^2)});
 \node[blue!75!black,font=\footnotesize] at (0.35,1.04) {峰化};
 \node[orange!85!black,font=\footnotesize] at (0.68,0.85) {空心/偏轴峰};
 \node[below,font=\footnotesize] at (0,0) {磁轴};
 \node[below,font=\footnotesize] at (1,0) {边界};
 \end{scope}
 \begin{scope}[xshift=7.5cm]
 \draw[->] (0,0) -- (1.1,0) node[right] {$\varrho$};
 \draw[->] (0,0) -- (0,1.15) node[above] {$T/T_{\rm ref}$};
 \draw[blue!75!black,thick,domain=0:1,samples=60] plot (\x,{0.07+0.85*(1-\x*\x)});
 \draw[red!70!black,thick,dashed] (0,1.02) .. controls (0.4,1.0) and (0.65,0.79) .. (0.85,0.61)
 .. controls (0.9,0.56) and (0.92,0.12) .. (1,0.07);
 \node[red!70!black,font=\footnotesize] at (0.55,1.06) {带边缘陡梯度的例子};
 \node[below,font=\footnotesize] at (0,0) {磁轴};
 \node[below,font=\footnotesize] at (1,0) {边界};
 \end{scope}
\end{tikzpicture}
\caption{密度与温度剖面的教学示意。$n_{\rm ref},T_{\rm ref}$ 是各图使用的正归一化参考值；曲线不是数据，也不要求两图来自同一放电。边缘陡梯度对应输运与源汇共同作用的可能结果，不能据此认定所有位形都有 H 模台基。}\label{fig:micro-profiles}
\end{figure}

\subsection{从击穿到位形转换，哪些分布在怎样演化}\label{sec:micro-evolution}

\textbf{启动时的空间图景。}预充中性气体未必均匀，种子电子、电离率和连接长度也随位置改变。电离首先在有利区域增长，形成仍含中性粒子、且能量分布可能强非热的局域载流区。只有随着建流和自场发展形成足够好的闭合区后，才适合逐步采用磁面平均的剖面描述。因此，不能用一个已存在的 GS 平衡逐帧解释所有击穿细节。

\textbf{建流与加热时的空间图景。}电流从边界电场及内部扩散共同建立，局域温度改变电阻率，电阻率又影响电流穿透；电流剖面改变 $q$ 与平衡。气体常先在边缘电离，燃料未必立即补到核心；中性束、射频和颗粒注入则可有不同的沉积位置。压力梯度进一步影响自举电流和稳定性。可以用控制变量逐项作用理解因果关系，但真正的状态演化通常是耦合的。

一般连续性方程先保证局部粒子数核算：
\begin{equation}\label{eq:micro-continuity}
 \partial_t n_s+\nabla\cdot(n_s\bm u_s+\bm\Gamma_s)=S_s^{(n)}.
\end{equation}
$\partial_t n_s$ 为密度的局部变化率；$n_s\bm u_s$ 为平均对流粒子通量，$\bm\Gamma_s$ 表示当前粗粒化模型中额外保留的扩散或涨落通量，两者单位均为 $\mathrm{m^{-2}s^{-1}}$；$S_s^{(n)}$ 是注入与反应净源（$\mathrm{m^{-3}s^{-1}}$）。若 $\bm u_s$ 已包含全部平均粒子输运，就不应再把同一通量重复放入 $\bm\Gamma_s$。

为显示磁面间的积累，在\emph{磁面几何固定}的近似下，可写一个壳层平均教学模型：
\begin{equation}\label{eq:micro-shell}
 \frac{\partial\bar n}{\partial t}
 +\frac1{V'}\frac{\partial}{\partial\varrho}(V'\Gamma_{\varrho})
 =\bar S,
 \qquad \Gamma_{\varrho}=-D_{\varrho}\frac{\partial\bar n}{\partial\varrho}
 +U_{\varrho}\bar n.
\end{equation}
$V(\varrho)$ 为该磁面内体积，$V'=dV/d\varrho$，$\bar n,\bar S$ 分别为壳层平均密度和净源；$V'\Gamma_{\varrho}$ 定义为穿过该面向外的总粒子率。因为 $\varrho$ 无量纲，$\Gamma_{\varrho}$ 的单位为 $\mathrm{m^{-3}s^{-1}}$，$D_{\varrho}$ 和 $U_{\varrho}$ 均具有 $\mathrm{s^{-1}}$ 的单位，不能直接与用物理半径定义的 $\mathrm{m^2/s}$ 扩散系数混用。第二式仅为有效扩散与径向输运速度闭合，第一式表达守恒；要预测剖面必须再给定闭合系数及边界条件。平行快运动不消除跨磁面的源汇差异，所以不同壳层可以维持不同密度与温度。

\textbf{位形变化时的空间图景。}改变 PF 电流不仅移动边界，也会改变磁面体积、连接长度、曲率、X 点和加热/排热几何；与此同时电流和压力仍在演化。若某物质体积内无粒子源汇，粒子数 $N$ 近似固定而体积 $V$ 缩小，则
\begin{equation}\label{eq:micro-compression}
 \bar n=N/V,\qquad \frac{d\bar n}{dt}=-\frac{\bar n}{V}\frac{dV}{dt}.
\end{equation}
$\bar n$ 是该体积平均密度，$N$ 为粒子数，第二式是对第一式求导；它表明几何压缩本身也能提高密度。压缩还可能通过压力做功改变热能。因此，移动磁面的输运模型须计入几何雅可比与坐标移动，不能把式~\eqref{eq:micro-shell} 的固定几何形式直接用于任意快速位形转换。

\textbf{位形稳定后，剖面仍可能继续变化。}磁平衡响应、垂直运动、热输运和电流扩散的时间尺度不同；几何已跟踪到目标，不代表内部温度或 $q$ 剖面已稳态。后续正常关停同样是在同时降低热能、改变电阻和电流、调整几何及边界源汇。全过程需要不同阶段之间传递粒子数、热能、剖面和场的信息，而不只是传递一张轮廓图。

\subsection{动理学、双流体、磁流体分别看见怎样的等离子体}\label{sec:micro-models}

\subsubsection{动理学：在同一空间点保留完整速度分布}

动理学不是只给每个位置一个温度，而是追踪 $f_s(\bm r,\bm v,t)$：同一点可以同时有热电子、非热快电子和方向偏置的粒子群。这里用 $\bm r$ 表示空间坐标，以免与控制状态向量 $\bm x$ 混淆。其一般教学形式为
\begin{equation}\label{eq:micro-kinetic}
 \partial_t f_s+\bm v\cdot\nabla f_s+
 \frac{q_s}{m_s}(\bm E+\bm v\times\bm B)\cdot\nabla_v f_s
 =\sum_{s'}C_{ss'}[f_s,f_{s'}]+S_s^{\rm rxn}+S_s^{\rm ext}.
\end{equation}
$f_s$ 按 $n_s=\int f_s\,d^3v$ 归一化，$d^3v$ 为速度体积元，$\nabla_v$ 对速度求导。左侧依次为局部分布变化、空间运动和电磁加速引起的速度空间运动；右侧 $C_{ss'}$ 为物种间散射算子，$S_s^{\rm rxn}$ 为原子/核反应净源，$S_s^{\rm ext}$ 为外部注入源，均与 $\partial_t f_s$ 同量纲。忽略碰撞和源项便得到 Vlasov 方程；它仍与自洽场耦合，``无碰撞''不意味着粒子彼此毫无影响。

对累计小角 Coulomb 散射，可以把碰撞效应写成速度空间漂移与扩散的 Fokker--Planck 形式：
\begin{equation}\label{eq:micro-fp}
 C_s=-\sum_i\frac{\partial}{\partial v_i}(A_i f_s)
 +\frac12\sum_{i,j}\frac{\partial^2}{\partial v_i\partial v_j}(D_{ij}f_s).
\end{equation}
这里 $i,j$ 是速度的 Cartesian 分量索引，不是物种索引；$A_i$ 为单位时间的平均速度增量（$\mathrm{m/s^2}$），$D_{ij}$ 为速度增量协方差的增长率（$\mathrm{m^2/s^3}$）。第一项表示系统性减速或漂移，第二项表示随机散射使分布在速度空间扩散。自洽 Coulomb 系数依赖各物种分布；任意取常数阻力加随机噪声，不会自动保证多物种总动量和总能量守恒。Landau 算子给出这类弱耦合 Coulomb 碰撞的具体闭合\cite{Plandau}。

动理学方程可以在相空间网格上求解，也可用粒子方法采样。PIC（particle-in-cell）通常用带权重的计算粒子表示许多真实粒子，沉积电荷/电流后求网格场，再推进计算粒子；它不等于逐个真实粒子的原子级模拟。是否含碰撞、原子反应和辐射取决于具体模型，而不是由``粒子''二字自动保证。无碰撞波--粒共振还可在不依赖二体撞击的情况下转移场能与粒子能；要判断是否已变成热，应查看分布形状、相混合、粗粒化与后续松弛，而不能只看一个粒子的速度变化。

\subsubsection{由动理学取矩：双流体模型与闭合问题}

从分布得到流体量的桥梁是速度矩\cite{Pmoments}：
\begin{equation}\label{eq:micro-moments}
 \begin{aligned}
 n_s&=\int f_s\,d^3v,& n_s\bm u_s&=\int\bm v f_s\,d^3v,\\
 \mathsf P_s&=m_s\int\bm c_s\bm c_s f_s\,d^3v,&
 \bm q_s^{\rm heat}&=\frac{m_s}{2}\int c_s^2\bm c_s f_s\,d^3v.
 \end{aligned}
\end{equation}
$\bm c_s=\bm v-\bm u_s$ 为随机速度；$\bm c_s\bm c_s$ 表示二阶张量外积，$\mathsf P_s$ 为压力张量（Pa）；$\bm q_s^{\rm heat}$ 为相对平均流的热流密度（$\mathrm{W/m^2}$），上标用于区别电荷 $q_s$。各向同性时 $\mathsf P_s=p_s\mathsf I$，$p_s=n_sk_BT_s$，$\mathsf I$ 为单位张量。密度、平均流和温度是 $f_s$ 的压缩信息，不是与动理学毫无关系的另一套物理。

分别保留电子、离子的连续性、动量和热能方程便得到双流体图景。例如忽略反应动量源时，每个物种的动量方程为
\begin{equation}\label{eq:micro-twofluid}
 m_sn_s(\partial_t+\bm u_s\cdot\nabla)\bm u_s
 =q_sn_s(\bm E+\bm u_s\times\bm B)-\nabla\cdot\mathsf P_s+\bm R_s.
\end{equation}
左侧为惯性力密度；右侧依次为电磁力、压力张量梯度及碰撞动量交换力密度，单位均为 $\mathrm{N/m^3}$。$\bm R_s$ 是其他物种对该物种的碰撞力，在封闭的弹性碰撞体系中有 $\sum_s\bm R_s=\bm0$。两个物种可以有不同流速与温度，模型因而能保留相对漂移及电子压力效应。方程还需要给出 $\mathsf P_s$、热流及交换率的闭合；只对动理学取矩会产生更高阶矩，并不会自然得到有限个封闭方程。

在单电离、准中性、电子惯性可忽略且电子压力近各向同性的简化下，电子动量可导出
\begin{equation}\label{eq:micro-generalohm}
 \bm E+\bm u_i\times\bm B
 =\eta\bm j+\frac{\bm j\times\bm B}{en_e}
 -\frac{\nabla p_e}{en_e}.
\end{equation}
左侧为相对于离子流动的有效电场；右侧依次为电阻、Hall 和电子压力项，单位均为 V/m。$p_e$ 为电子压强，$en_e$ 为电子电荷密度大小；Hall 项来自电子流与离子流不同。是否可以忽略这些项，取决于关注尺度，而不是一看到有碰撞便都可以删掉。

\subsubsection{单流体 MHD：整体运动与磁场共同演化}

在长尺度、低频等假设下，可合并两物种的整体动量，把等离子体视作可导电、带压强的流体。一个忽略黏性和原子粒子源、采用标量压强的电阻 MHD 教学系统为\cite{Pmhd}
\begin{align}
 \partial_t\rho_m+\nabla\cdot(\rho_m\bm u)&=0,\label{eq:micro-mhdmass}\\
 \rho_m(\partial_t+\bm u\cdot\nabla)\bm u
 &=\bm j\times\bm B-\nabla p,\label{eq:micro-mhdmomentum}\\
 \partial_t\bm B&=\nabla\times(\bm u\times\bm B-\eta\bm j),
 \quad \mu_0\bm j=\nabla\times\bm B,\quad\nabla\cdot\bm B=0,\label{eq:micro-mhdfield}\\
 (\partial_t+\bm u\cdot\nabla)p+\gamma_{\rm ad}p\nabla\cdot\bm u
 &=(\gamma_{\rm ad}-1)\bigl(\eta|\bm j|^2-\nabla\cdot\bm q^{\rm heat}+Q_{\rm net}\bigr).
 \label{eq:micro-mhdenergy}
\end{align}
$\rho_m$ 为质量密度（$\mathrm{kg/m^3}$），不同于净电荷密度 $\rho_q$；$\bm u$ 是质量加权平均流速，$p$ 为电子与离子的总标量压强。式~\eqref{eq:micro-mhdmass} 表示质量守恒；式~\eqref{eq:micro-mhdmomentum} 将磁力、压力梯度与整体惯性相平衡，可研究加速、波动与不稳定性，而不仅是静态几何。

式~\eqref{eq:micro-mhdfield} 中 $\nabla\times(\bm u\times\bm B)$ 表示流动与场的耦合，$-\nabla\times(\eta\bm j)$ 表示电阻扩散，后两式分别是忽略位移电流的 Ampère 定律和磁场无散约束。式~\eqref{eq:micro-mhdenergy} 为本近似下的内部热能闭合：$\gamma_{\rm ad}$ 为绝热指数，三维非相对论理想气体常取 $5/3$；左侧第一项是随流压强变化，第二项为压缩/膨胀作用。右侧含欧姆热、热流的净流出及其他净体积热源 $Q_{\rm net}$，单位均为 $\mathrm{W/m^3}$。$Q_{\rm net}$ 可包含吸收加热减辐射损失，但不能重复计入已单列的欧姆热。这个单压力方程不能独立预测 $T_e,T_i$ 的分离。

理想 MHD 进一步取 $\eta=0$ 并忽略上述非理想项，在适当光滑条件下体现磁通随流冻结。它能描述许多整体扰动，却不能自洽处理所有需要打破磁场连接关系的局部过程。由完整 MHD 方程再令流速和惯性可忽略、采用轴对称标量压力，才回到 $\bm j\times\bm B=\nabla p$ 与 GS 平衡。\textbf{GS 是一类受限的瞬时平衡模型，不是完整 MHD 时间演化。}满足平衡只说明某种力学配置可能存在，不能自动判定它对扰动是否稳定。

\subsubsection{约化模型：保留目标过程，平均掉更快细节}

导引中心/漂移动理学平均掉快速回旋相位，回旋动理学进一步在规定的低频、小回旋半径等排序下保留特定微尺度涨落，常用于研究磁化等离子体湍流及输运\cite{Pgyro}。这并非任意条件下都比完整动力学``更正确''，而是在所需尺度上提高可计算性。磁面平均输运模型再把许多涨落效应折合为通量闭合，适合较慢的剖面演化；低维控制模型则进一步保留少数可估计、可调节的状态。

模型选择应比较回旋时间、碰撞松弛时间、波动时间以及粒子轨道和装置长度。即使沿场平均自由程很长，强磁化与低频排序下仍可构造某些流体闭合；但不能因此认为简单各向同性、局域热传导的闭合必然可靠。反过来，动理学也不是所有现象都必须采用的唯一方法：研究缓慢线圈响应和边界跟踪时，GS 加回路模型可能更合适。弱碰撞流体闭合的局限和可保留的信息见参考文献~\cite{Pclosure}。

\subsection{模型可以推导或发现哪些典型结构}\label{sec:micro-structures}

\textbf{典型结构来自方程、对称性、边界与驱动的共同作用。}模型既能解析推出某些结构，也能通过非线性仿真显示结构如何形成，但不是每种结构都由同一模型产生。给定图案的平衡构造、从小扰动长出的不稳定结构、从诊断重建出的实验结构，是三种不同证据。

\subsubsection{从屏蔽到鞘层：电荷分离具有空间层次}

体区准中性和壁面附近的电荷分离可以同时存在。线性小扰动、等温电子 Boltzmann 响应及固定离子背景下，静电势满足
\begin{equation}\label{eq:micro-shielding}
 \nabla^2\phi=\phi/\lambda_{De}^2\quad\text{（离开扰动源处）},
 \qquad \phi(r)\propto\frac{e^{-r/\lambda_{De}}}{r}.
\end{equation}
$\phi$ 为相对未扰动等离子体的静电势（V），$\lambda_{De}$ 为电子 Debye 长度，$r$ 为距局域点扰动的空间距离；假设 $|e\phi|\ll k_BT_e$。第一式由 Poisson 方程和线性电子密度响应联立得到，第二式为球对称、远处势趋于零的解，说明电荷的影响被屏蔽，而非像真空中一样只有 $1/r$ 衰减。相应理论说明了为什么宏观体区常可准中性处理\cite{Pdebye}。

壁面吸收电子和离子的响应不同，会建立鞘层电势并调节通向壁面的粒子流；鞘层是一个边界条件参与决定的非中性层。其详细尺度与磁场、碰撞、壁偏压等有关，不能把全部边缘区域都等同于一个 Debye 长度。准中性 MHD 可以把鞘层作用折合成边界条件，却不能在始终令 $\rho_q=0$ 的同时解析出其内部净电荷结构。要研究该层内部，应保留 Poisson 响应和适当的流体或动理学边界模型。

\subsubsection{磁面、分离面、磁镜与 Alfv\'en 波：有组织的场与运动}

轴对称 GS 方程在合适的源和边界下可构造嵌套磁面、磁轴与鞍点，配合拓扑追踪得到限制器/偏滤器分离面；它并不保证任意压力、电流和线圈约束都有所需解。单粒子模型通过式~\eqref{eq:micro-mirror} 可推出通过/俘获粒子分界，并解释香蕉轨道等相空间分层。这两种结构描述的对象不同：一个是磁场的几何，另一个是粒子在给定场中的轨道族。

线性化理想 MHD 则能推出传播结构。对静止、均匀背景场上的小幅剪切扰动，
\begin{equation}\label{eq:micro-alfven}
 \partial_t^2\delta\bm u_\perp
 =v_A^2\partial_\ell^2\delta\bm u_\perp,
 \qquad v_A=\frac{B_0}{\sqrt{\mu_0\rho_{m0}}}.
\end{equation}
$\delta\bm u_\perp$ 是相对背景场的垂直流速扰动，$\ell$ 为沿背景磁场的距离，$B_0,\rho_{m0}$ 是背景磁场大小和质量密度；$v_A$ 为 Alfv\'en 速度（m/s）。它来自磁张力与惯性的耦合，说明场线弯曲和流动可以作为波传播；长度 $L_\parallel$ 上的传播时间量级为 $L_\parallel/v_A$。这是解析模型给出速度和尺度、再用仿真或诊断检验的例子，并非从一张静态磁面图就能看出的结论。

\subsubsection{电流片、重联和磁岛：必须允许连接关系改变}

完整或约化的电阻 MHD、Hall/双流体与动理学模型，可以在相应条件下研究局部电流集中形成的电流片及磁重联；各模型的耗散和非理想机制不同，不能认为它们给出的层厚、增长率与重联速率天然一致。理想 MHD 可描述电流片趋于变薄等动力学，但若数值结果出现拓扑改变，必须检查是否来自未说明的数值耗散，不能直接当成理想冻结条件下的解析结论。

托卡马克中，螺旋扰动常在有理磁面附近形成磁岛。圆截面近似下，用极向/环向整数模数 $m,n$ 表示扰动相位 $\chi=m\theta-n\phi$，共振条件为 $q(r_s)=m/n$，其中 $r_s$ 为相应半径；这里 $m$ 不是质量，$n$ 不是密度。撕裂模可改变该处场线连接，形成 O 点、X 点与岛分离面，并影响热量沿场的输运\cite{Pislands}。

一个用于理解几何的无量纲单谐波局部模型是
\begin{equation}\label{eq:micro-island}
 \mathcal H=\frac12X^2-\varepsilon\cos\chi,\qquad
 \frac{dX}{d\zeta}=-\frac{\partial\mathcal H}{\partial\chi},\qquad
 \frac{d\chi}{d\zeta}=\frac{\partial\mathcal H}{\partial X}.
\end{equation}
$X=(r-r_s)/a$ 是归一化径向偏离，$a$ 为选定径向尺度；$\chi$ 为螺旋角，$\zeta$ 为沿磁力线前进的归一化参数，$\varepsilon>0$ 为模型中的归一化扰动强度，$\mathcal H$ 是相应守恒的场线 Hamilton 函数。这组方程描述磁力线相图，不是粒子的时间 Hamilton 方程；$\varepsilon$ 也不能不经转换就设为实测 $\delta B/B$。

保持 $\mathcal H$ 不变便得到闭合岛内轮廓与岛外绕行轮廓：$X=0,\chi=0$ 是 O 点，$X=0,\chi=\pi$ 是鞍形 X 点，分离面为 $\mathcal H=\varepsilon$。令 $\chi=0$ 可得径向半宽 $2\sqrt\varepsilon$，故此归一化模型的全宽为 $4a\sqrt\varepsilon$。这展示了\emph{由简单方程解析推出一种典型结构及宽度标度}；若要预测真实岛宽如何随时间增长，仍需撕裂稳定性、平衡电流及非线性演化方程，单靠这张相图不够。

\begin{figure}[htbp]
\centering
\begin{tikzpicture}[>=Stealth,font=\small,x=0.021cm,y=1.6cm]
 \draw[->] (-195,0) -- (200,0) node[right] {$\chi$};
 \draw[->] (0,-1.5) -- (0,1.55) node[above] {$X$};
 \draw[blue!70!black,thick,domain=-180:180,samples=121] plot (\x,{sqrt(max(0,0.5*(1+cos(\x))))});
 \draw[blue!70!black,thick,domain=-180:180,samples=121] plot (\x,{-sqrt(max(0,0.5*(1+cos(\x))))});
 \foreach \sg in {-1,1}{
  \draw[blue!50,domain=-61.31:61.31,samples=61] plot (\x,{\sg*sqrt(max(0,2*(-0.12+0.25*cos(\x))))});
  \draw[blue!50,domain=-113.57:113.57,samples=81] plot (\x,{\sg*sqrt(max(0,2*(0.10+0.25*cos(\x))))});
  \draw[black!50,dashed,domain=-180:180,samples=101] plot (\x,{\sg*sqrt(2*(0.45+0.25*cos(\x)))});
 }
 \fill (0,0) circle[x radius=2,y radius=0.03];
 \node[above right] at (4,0) {O};
 \node[above] at (-180,0) {X};\node[above] at (180,0) {X};
 \node[below] at (-180,-0.03) {$-\pi$};\node[below] at (180,-0.03) {$\pi$};
 \node[blue!70!black,right,font=\footnotesize] at (70,0.95) {岛分离面};
\end{tikzpicture}
\caption{式~\eqref{eq:micro-island} 在 $\varepsilon=1/4$ 时的解析等值线示意。左右边界表示同一周期角的接续，粗线是分离面，细闭合线是岛内磁力线相图；这里的闭合图形不是单粒子的回旋轨道。}\label{fig:micro-island}
\end{figure}

轴对称 GS 中的偏滤器 X 点，与三维螺旋扰动造成的磁岛 X 点，不是同一个模型结论。前者可在轴对称截面出现；后者涉及破坏原嵌套磁面的非轴对称结构。只允许轴对称的当前 HFM/MEQ 主链不能自行产生并解析后一类撕裂磁岛。

\subsubsection{湍流涡、带状流与输运屏障：多尺度层次}

电磁涨落使导引中心漂移，相关涨落可以形成净径向粒子和热通量。动理学或回旋动理学可研究漂移不稳定性、微尺度涡及其非线性耦合，并在适当条件下形成磁面上近轴对称的带状流；这些流的剪切又能改变湍流输运。磁面平均输运模型通常不解析每个涡，而使用其平均通量或有效系数。两种模型的结果可以通过通量、剖面和时间尺度相互传递，而不是必须在同一计算中解析所有尺度\cite{Pgyro}。

H 模边缘台基或内部输运屏障表现为较小有效输运区域中的陡密度/温度/压力梯度。实验诊断可确定其位置和演化；理论模型可研究形成与稳定性机制；多尺度计算可检验特定机制在所给条件下能否形成类似结构。若在输运代码中直接指定一个很小的扩散系数，得到陡剖面只是对该闭合的响应，不等于已经从第一性原理发现输运屏障。完整 L--H 转换与边缘台基还耦合中性粒子、边界条件及不同尺度，不应把任意一次湍流模拟都称为全过程预测\cite{Ppedestal}。

\begin{table}[htbp]
\centering\small
\caption{典型结构与能够支持它的模型信息}\label{tab:micro-hierarchy}
\begin{tabular}{L{2.1cm}L{5.4cm}L{6.0cm}}
\toprule 结构/层次 & 可采用的模型与典型结果 & 仅靠更简化模型缺少什么 \\
\midrule
Debye 屏蔽、鞘层 & Poisson 响应配合电子/离子动力学，得到电荷和电势层 & 全程准中性近似不含鞘层内部电荷分离 \\
回旋、俘获与香蕉轨道 & 单粒子/导引中心动力学，得到轨道族与相空间分界 & 静态磁面图不等于粒子轨道图 \\
磁面、磁轴、偏滤器 X 点 & 轴对称 GS 及边界条件，构造瞬时平衡几何 & 不包含完整稳定性和非轴对称演化 \\
Alfv\'en 波与整体位移 & 含惯性的 MHD，得到频率、传播与增长规律 & 准静态平衡序列不自动给出动力学谱 \\
电流片、磁岛与重联 & 含相应非理想项、维度和分辨率的场/流体或动理学模型 & 理想冻结条件与轴对称约束会排除部分过程 \\
湍流、带状流及屏障 & 动理学/回旋动理学与输运、边缘模型耦合 & 给定有效输运系数不能独立证明形成机制 \\
\bottomrule
\end{tabular}
\end{table}

\subsubsection{怎样辨别解析结论、仿真形成与实验发现}

解析推导可以确定一个结构在何种假设下存在、具有何种标度；仿真则能检验其形成路径、增长、饱和及与其他过程的耦合。若希望把仿真中出现的结构作为物理发现，应说明它是初始条件直接放进去的，还是由扰动演化得到；应检查允许的维度和对称性、边界、守恒收支以及空间/时间分辨率。结构位置、宽度、频率或增长率的收敛，比一张看起来相似的图更有判断力。

实验观测还要通过合适诊断检验模型可测量的输出，例如磁测量中的模数与频率、温度剖面的平台化、边缘梯度或粒子分布。某个模型成功解释一个结构，只支持其在该条件和尺度上的有效性；不同模型可能在共同适用区给出同一大尺度结构，却在薄层、快粒子与输运预测上不同。把目标结构所需的信息与模型实际保留的信息逐项对应，才能选择既能回答问题、又可计算和验证的模型。

\subsection{本章与后续控制报告的衔接}\label{sec:micro-reading}

从粒子到控制，可以沿一条因果链阅读：电场与注入向粒子供能，分布函数的平均流决定电流、随机部分决定热量；碰撞与反应重新分配能量和组分；电荷、电流产生自洽场，场又组织轨道、磁面和输运。密度、温度和电流的空间剖面共同影响平衡与稳定性，线圈和其他执行器则不断改变这个耦合系统的边界与驱动。

第~\ref{sec:lifecycle} 章按放电时序展开这条链，第~\ref{sec:model} 章进入已选定的 HFM/MEQ 磁平衡与回路近似。阅读案例时应先用表~\ref{tab:micro-hierarchy} 判断它能研究哪类结构，再分析求解器与控制接口；不能因为案例变量中有压力或磁通，就推定其同时求解了碰撞、鞘层、湍流和核反应。通用理论及相关研究统一列入文末参考文献，正文引用均可点击跳转。

% ================= 全过程导读 =================
\clearpage
\section{全过程导读：从中性气体到受控放电与安全关停}
\label{sec:lifecycle}
\markboth{全过程导读}{全过程导读}

\textbf{本章定位与适用范围。}本章是面向初学者的通用控制物理正文，沿用全文按章编号的公式体系。全文从粒子受力、电磁感应和守恒律出发，解释控制器为什么需要某类执行器与测量量；这里的``第一性原理视角''不表示只靠基本常数就能预测一次真实放电。碰撞截面、壁面释放、湍流输运、约束时间和不稳定性阈值仍需要实验数据或更高保真模型闭合。本章全部属于\evbg{}；其中的充气、击穿、热输运、聚变自加热及破裂过程，不代表 HFM/MEQ 已经实现这些模块。已有实现的精确边界仍见第~\ref{sec:intro}、\ref{sec:approx} 节。

\subsection{先分清三个目标：产生等离子体、维持放电、实现聚变点火}\label{sec:cycle-goals}

实验语境中常把``气体被打成等离子体''称为点火，但至少应区分三个阶段。\textbf{击穿（breakdown）}是自由电子在电场中获得能量并碰撞电离中性粒子，形成电子数快速增长的过程。\textbf{建流与加热}是在击穿之后建立可约束的等离子体电流、克服电离及辐射损失，并进入可维持的热等离子体状态。\textbf{聚变点火（ignition）}则要求聚变产物留在等离子体中的自加热足以补偿能量损失，不再依赖外部加热维持燃烧；它不是普通托卡马克放电建立的必要条件。该术语区分采用本章参考资料 \cite{C1}。

多数控制实验关注电流、位置、形状、密度和温度的可重复控制，不必使用氘氚燃料，也不必达到聚变点火。一个电流和边界保持不变的``平顶阶段''，仍可能持续消耗变压器磁通。这里\textbf{剩余变压器磁通}是指在当前线圈电流、磁体和电源约束下，沿预定感应驱动方向还能提供的外部磁通摆幅，也称剩余伏秒额度，单位 Wb 或 V$\cdot$s。中心螺线管像变压器原边，等离子体像次级：要持续提供有向环电压，就必须继续改变外部磁通；线圈到达可用摆幅终点后，这一方向的感应驱动便受到限制。它不是``还没有被等离子体吸收掉的磁场''，也不是磁能的另一种叫法。一个非感应稳态放电，也仍可能需要持续外部加热。\textbf{形状平稳、非感应电流稳态与聚变自持燃烧是三个不同条件。}

按因果关系，一次典型放电可读作：
\begin{center}
\small
气体准备与磁场准备 $\longrightarrow$ 击穿与烧穿 $\longrightarrow$ 电流爬升与加热\\[3pt]
$\longrightarrow$ 平顶维持 $\longleftrightarrow$ 位形转换 $\longrightarrow$ 受控降能与降流 $\longrightarrow$ 放电结束
\end{center}
这里是物理过程图，不是所有装置通用的设备时序。环向场建立、中心螺线管预充磁、抽气和气体预充可以交叠；偏滤器位形也常在电流爬升期间形成，并非一定等到平顶后才建立。异常发生时，控制系统还可能从任一阶段转入中止或保护路径。

\subsection{从粒子到控制状态：哪些量守恒，哪些量需要闭合}\label{sec:cycle-principles}

单个带电粒子的运动首先受 Lorentz 力支配：
\begin{equation}\label{eq:cycle-lorentz}
 m_s\frac{d\bm v}{dt}=q_s(\bm E+\bm v\times\bm B),
 \qquad \frac{d}{dt}\left(\frac12m_sv^2\right)=q_s\bm E\cdot\bm v.
\end{equation}
式中，$s$ 标记电子或某种离子，$m_s$ 为其质量（kg），$q_s$ 为带符号电荷（C，电子为负）；$t$ 为时间（s），$\bm v$ 为速度（m/s），$v=|\bm v|$。$\bm E$ 是电场（V/m），$\bm B$ 是磁感应强度（T）。左式的 $m_s d\bm v/dt$ 是动量变化率，$q_s\bm E$ 是电场力，$q_s\bm v\times\bm B$ 是磁场力；右式左端是单粒子动能变化率，$q_s\bm E\cdot\bm v$ 是电场对粒子做功的功率（W）。$\times$ 与 $\cdot$ 分别表示叉乘和点乘。
磁力与瞬时速度垂直，因而静磁场本身不直接给粒子增加动能；它改变轨道，帮助限制横向运动。电场能做功，而碰撞使有序运动能转化为随机热运动。回旋半径为
\begin{equation}\label{eq:cycle-gyroradius}
 \rho_s=\frac{m_sv_\perp}{|q_s|B}.
\end{equation}
这里 $\rho_s$ 为回旋半径（m），$v_\perp$ 是相对局部磁场方向的垂直速度大小，$B=|\bm B|$；分母使用 $|q_s|$，因为电荷符号只改变回旋方向，不改变半径。
它与装置尺度之比说明粒子是否充分磁化；沿磁场的运动仍需用闭合磁面和磁场螺旋性来约束。

对各物种分布函数 $f_s(\bm x,\bm v,t)$，更基础的描述为
\begin{equation}\label{eq:cycle-kinetic}
 \frac{\partial f_s}{\partial t}+\bm v\cdot\nabla f_s+
 \frac{q_s}{m_s}(\bm E+\bm v\times\bm B)\cdot\nabla_v f_s
 =C_s[f]+S_s.
\end{equation}
式中 $\bm x$ 是空间坐标，$\nabla$ 对空间求导，$\nabla_v$ 对速度求导；$f_s$ 按 $n_s=\int f_s\,d^3v$ 归一化。左侧三项依次是固定相空间位置的分布变化、粒子运动导致的空间输运、电磁力导致的速度空间输运。$C_s[f]$ 为依赖各物种分布的碰撞算子，方括号表示其函数依赖；$S_s$ 为外部源与反应源汇，二者与 $\partial f_s/\partial t$ 具有相同单位。
$C_s$ 表示碰撞，$S_s$ 表示注入、电离、复合等源汇。对速度积分可得到粒子数守恒，对动量和能量取矩可得到流体方程；但热流、应力及反应率不会自动消失，需要闭合。因此，控制模型通常保留少数与执行器直接相关的状态，例如
\begin{equation}\label{eq:cycle-state}
 \bm x=(N_0,N_e,W_e,W_i,I_p,\bm I_c,\bm I_w,\bm a_j,\bm a_p,\bm g),
 \qquad \bm y=h(\bm x)+\bm\epsilon.
\end{equation}
这里 $\bm x$ 改用于低维状态向量，不再表示上式中的空间坐标；$\bm y$ 是诊断输出，$h$ 是从状态到测量的映射，$\bm\epsilon$ 是测量误差。状态分量可以有不同单位，各误差项采用对应诊断的单位；用于数值优化前通常要归一化。
$N_0,N_e$ 是中性粒子和电子总数，$W_e,W_i$ 是热能，$I_p$ 是等离子体电流，$\bm I_c,\bm I_w$ 分别是主动线圈和被动结构电流，$\bm a_j,\bm a_p$ 是电流与压力剖面参数，$\bm g$ 是位置、边界和 X 点等几何量。测量 $\bm y$ 带有误差 $\bm\epsilon$；磁平衡约束使这些状态并不全部独立。这也解释了为什么不能仅凭一次边界图像确定内部温度和电流剖面。

\subsection{气体导入：先建立可电离的粒子库}\label{sec:cycle-gas}

抽气首先移除残余气体，壁面处理则影响后续杂质释放与粒子再循环。随后气阀向真空室导入工作气体。未电离且近似均匀、等温时，气体满足 $p_0V_v=N_0k_BT_0$。这里 $p_0$ 为中性气体压强（Pa），$V_v$ 是可用真空容积（$\mathrm{m^3}$），$N_0$ 为中性粒子数，$k_B$ 为 Boltzmann 常数（J/K），$T_0$ 是中性气体温度（K），不是未来等离子体的电子温度。以粒子每秒表示注入率 $\dot N_{\rm gas}$，简单粒子库模型为
\begin{equation}\label{eq:cycle-gas}
 \frac{dN_0}{dt}=\dot N_{\rm gas}+\dot N_{\rm wall}
 -\frac{S_{\rm eff}}{V_v}N_0-\dot N_{\rm ion}+\dot N_{\rm rec}.
\end{equation}
左端为中性粒子数的变化率。右侧 $\dot N_{\rm gas}$ 是气阀注入率，$\dot N_{\rm wall}$ 是壁面净释放率；$(S_{\rm eff}/V_v)N_0$ 是泵抽走的粒子率，$\dot N_{\rm ion}$ 是电离消耗率，$\dot N_{\rm rec}$ 是复合生成率。各项单位均为 $\mathrm{s^{-1}}$（粒子数每秒），点号表示对时间求导；式中正负号分别表示对中性粒子库的补充与消耗。
$S_{\rm eff}$ 为有效抽速，单位 $\mathrm{m^3/s}$；其余项分别描述壁面释放、电离消耗与复合补充。气体若主要是分子，必须另外计入解离、分子离子反应及每个分子所含的核数，不能把分子数直接当作电子数。后文为说明因果关系，使用单原子、单电离的简化模型。

预充过少时，种子电子产生新的电离碰撞的机会不足；预充过多时，电子从电场获得的能量更难克服碰撞、电离和辐射损失。因此，目标不是``气体越多越容易点燃''，而是在给定电场、磁连接长度及壁面条件下进入可击穿、可烧穿的区域。气阀开度也不是密度本身：阀门延迟、管路容积、抽气和再循环共同决定 $N_0(t)$。控制侧需要用真空计、阀门状态及后续密度诊断估计粒子库，而不是只执行一个固定阀门脉冲。

\subsection{线圈通电：约束磁场、感应电场与形状场各司其职}\label{sec:cycle-coils}

\textbf{环向场线圈（TF）}提供主要环向磁场，理想环形线圈近似给出
\begin{equation}\label{eq:cycle-tf}
 B_\phi(R)\simeq\frac{\mu_0N_{\rm TF}I_{\rm TF}}{2\pi R}.
\end{equation}
$B_\phi(R)$ 是大半径 $R$ 处的环向场（T），$R$ 以 m 计；$\mu_0$ 为真空磁导率（H/m），$N_{\rm TF}$ 为等效总匝数，$I_{\rm TF}$ 为每匝电流（A），$N_{\rm TF}I_{\rm TF}$ 是安匝数。分母 $2\pi R$ 为环向圆周长；近似忽略了离散线圈导致的场纹波。
\textbf{中心螺线管（CS）}通过改变穿过环向回路的磁通产生感应电场；等离子体类似变压器的次级。\textbf{极向场线圈（PF）}控制平衡位置和边界形状，也参与磁通预算；\textbf{快速垂直稳定线圈（VS）}用于抑制快速位置偏移。具体装置中这些功能会耦合，不能把每个物理目标机械地分配给唯一一组线圈。磁体分工及 CS 的感应作用见资料 \cite{C2,C3}。

对固定的环向闭合曲线 $\mathcal C$，Faraday 定律为
\begin{equation}\label{eq:cycle-faraday}
 V_{\rm loop}\equiv\oint_{\mathcal C}\bm E\cdot d\bm l
 =-\frac{d\Phi_{\mathcal C}}{dt},
 \qquad E_\phi\simeq\frac{V_{\rm loop}}{2\pi R_0}.
\end{equation}
$V_{\rm loop}$ 为沿固定闭合曲线 $\mathcal C$ 的环电压（V），$d\bm l$ 为沿回路定向的微小线元（m），积分累加切向电场。$\Phi_{\mathcal C}$ 为回路所围曲面的磁通（Wb）；负号表示感应方向服从 Lenz 定律。第二式取近似均匀环向电场：$R_0$ 为参考主半径（m），$E_\phi$ 为该半径处的环向电场，$2\pi R_0$ 是回路长度。
$\Phi_{\mathcal C}$ 是穿过该曲线所围曲面的总磁通，正负号取决于曲线与曲面定向。它不同于后文局部极向磁通函数的任意零点。恒定 CS 电流只维持磁场，并不持续产生变压器电压；真正驱动电流的是磁通随时间的变化。启动前的预充磁为后续磁通摆幅留出空间。

在线性磁性材料近似及给定几何下，主动与被动回路可统一写成
\begin{equation}\label{eq:cycle-coils}
 \begin{aligned}
 \bm V_c&=\bm R_c\bm I_c+
 \frac{d}{dt}(\bm L_{cc}\bm I_c+\bm M_{cw}\bm I_w+\bm M_{cp}I_p),\\
 \bm 0&=\bm R_w\bm I_w+\frac{d\bm\lambda_w}{dt}.
 \end{aligned}
\end{equation}
两个下标 $c,w$ 分别表示受电源驱动的主动线圈（coil）和无独立电源的被动导体（wall）。逐项含义为：
\begin{itemize}[nosep]
 \item $\bm V_c$ 为主动回路端电压向量（V），$\bm I_c$ 为主动线圈电流向量（A），$\bm R_c$ 为电阻矩阵（$\Omega$）；$\bm R_c\bm I_c$ 是电阻压降。
 \item $\bm L_{cc}$ 为主动回路的自感和相互之间的互感矩阵（H），$\bm L_{cc}\bm I_c$ 是这些电流产生的主动回路磁链（Wb，按匝数计链）。
 \item $\bm M_{cw}$ 为被动结构到主动回路的互感矩阵（H），$\bm I_w$ 为被动结构电流（A）；$\bm M_{cw}\bm I_w$ 表示被动电流给主动回路增加的磁链。
 \item $\bm M_{cp}$ 为单位等离子体总电流到主动回路的互感向量（H），$I_p$ 为等离子体总电流（A）；$\bm M_{cp}I_p$ 是这一集总近似下的等离子体磁链贡献。括号内总磁链的时间导数为感应压降（V），它与电阻压降相加平衡外加电压。
 \item 第二式左侧 $\bm0$ 表示被动回路没有外加端电压；$\bm R_w$ 为其电阻矩阵，$\bm R_w\bm I_w$ 为电阻压降，$d\bm\lambda_w/dt$ 为总磁链变化产生的感应压降。两者之和为零，并不要求每一项分别为零。
\end{itemize}
例如 $\bm\lambda_w=\bm L_{ww}\bm I_w+\bm M_{wc}\bm I_c+\bm M_{wp}I_p$：$\bm L_{ww}$ 包括被动结构的自感和互感，$\bm M_{wc},\bm M_{wp}$ 分别给出主动线圈和等离子体的耦合。几何和剖面变化时互感也可能随时间改变，不能总把所有矩阵移到时间导数之外。
$\bm\lambda_w$ 是被动结构的磁链。真空室中的涡流会延迟磁场变化，因此控制器发出的电压不能瞬间变成期望边界。启动时还需利用 PF 场抵消不利的杂散极向场，形成足够长的磁连接长度，使电子在碰壁前有机会电离气体。此时还没有成熟的嵌套磁面，不能从第一时刻就套用静态 GS 平衡。

\subsection{从种子电子到击穿：电离增长必须超过粒子损失}\label{sec:cycle-breakdown}

背景辐射或预电离产生少量自由电子；感应电场或射频波使其获得能量。电子与中性粒子碰撞后，可能发生激发、电离或弹性散射。单电离且均匀近似下，一个教学用电子数模型为
\begin{equation}\label{eq:cycle-avalanche}
 \frac{dn_e}{dt}=n_en_0K_{\rm ion}-n_en_iK_{\rm rec}
 -\frac{n_e}{\tau_{\rm loss}}+S_{\rm seed}.
\end{equation}
本式描述单位体积内电子数的变化，各项单位为 $\mathrm{m^{-3}s^{-1}}$：
\begin{itemize}[nosep]
 \item $n_e,n_0,n_i$ 分别为电子、中性原子及单电离离子的数密度（$\mathrm{m^{-3}}$）；左侧 $dn_e/dt$ 为电子密度的变化率。在均匀固定体积模型中不另写对流项。
 \item $n_en_0K_{\rm ion}$ 为碰撞电离产生电子的速率；$K_{\rm ion}$ 是电离速率系数。$n_en_iK_{\rm rec}$ 为有效两体复合造成的电子损失率；$K_{\rm rec}$ 是相应复合系数。若三体复合重要，需增加不同密度依赖的项，不能直接套用这一系数。
 \item $n_e/\tau_{\rm loss}$ 为电子离开有效体积的损失率，$\tau_{\rm loss}$ 为有效约束/损失时间（s）；这是一种集总闭合，并非完整空间输运方程。
 \item $S_{\rm seed}$ 是外部种子电子源率，例如预电离贡献。前三项随现有粒子密度变化，种子源则可在电子数很少时启动增长。
\end{itemize}
$K_{\rm ion}=\langle\sigma_{\rm ion}v\rangle$ 中，$\sigma_{\rm ion}$ 是电离截面（$\mathrm{m^2}$），$v$ 是相对速度，尖括号表示对速度分布取平均；$K_{\rm ion}$ 和 $K_{\rm rec}$ 的单位为 $\mathrm{m^3/s}$，由电子能量分布决定。早期分布未必是 Maxwell 分布，不能不加判断地只用一个 $T_e$ 计算反应率。若忽略复合且系数近似恒定，电子数按 $\exp(\gamma t)$ 增长，其中
\begin{equation}\label{eq:cycle-growth}
 \gamma=n_0K_{\rm ion}-\tau_{\rm loss}^{-1}>0
\end{equation}
$\gamma$ 是净电子增长率（$\mathrm{s^{-1}}$）：$n_0K_{\rm ion}$ 是每个电子引起电离增殖的平均频率，$\tau_{\rm loss}^{-1}$ 是逃逸损失频率。忽略种子源且系数固定时，$n_e(t)=n_e(0)e^{\gamma t}$，其中 $n_e(0)$ 为起始电子密度；$\gamma>0$ 才产生指数净增长。
是雪崩增长的必要条件之一。\textbf{平行粒子损失}是粒子沿局部磁场方向运动，沿着连接到壁面或靶板的开放磁力线离开等离子体有效体积的过程；``平行''指与磁场平行，不是指与装置中轴平行。磁力能弯曲垂直运动的轨道，却不能直接阻止沿场运动。因此，启动时即使粒子已磁化，若磁力线很快接触壁面，电子仍会迅速逃走。用沿场到达吸收边界的连接长度 $L_\parallel$（m）和代表性沿场速度大小 $|v_\parallel|$（m/s），可估计 $\tau_{\rm loss}\sim L_\parallel/|v_\parallel|$；该式只给自由流量级，实际还受碰撞、鞘层电场、漂移和壁面作用影响。在充分形成的等离子体中，电子和离子逃逸还会通过自洽电场相互耦合，不能简单分别按各自热速度计算净损失。降低不利杂散场、增加有效连接长度与提供适当预电离，都是在改善这一增长竞争。

\textbf{击穿成功还不等于烧穿成功。}烧穿（burn-through）指越过启动阶段中性粒子与低电离态杂质造成的能量损失障碍，使电离程度和电子温度能继续上升。资料 \cite{C4} 专门比较了雪崩条件与烧穿条件：只满足前者并不足以保证后者。用电子热能表示收支，可写为
\begin{equation}\label{eq:cycle-burnthrough}
 \frac{dW_e}{dt}=P_{\rm Ohm}+P_{\rm aux,e}
 -P_{\rm ion}-P_{\rm rad}-P_{e\to i}-P_{\rm tr,e},
 \qquad W_e=\frac32N_ek_BT_e.
\end{equation}
$W_e$ 为电子总热能（J），$N_e$ 为有效体积内电子总数；$T_e$ 是以 K 表示的平均电子温度，$k_BT_e$ 是相应热能尺度。$3/2$ 来自非相对论三维平动自由度。左侧为电子热能增长功率；右侧各 $P$ 均以 W 计：$P_{\rm Ohm}$ 为欧姆加热，$P_{\rm aux,e}$ 为辅助加热直接沉积到电子的部分，$P_{\rm ion}$ 为电离吸收的电子能量，$P_{\rm rad}$ 为向外辐射损失，$P_{e\to i}$ 为电子到离子的碰撞传热，$P_{\rm tr,e}$ 为其他输运损失。正号项补能，负号项耗能；若离子反而向电子传热，$P_{e\to i}$ 应按带符号能流处理。
这里以近似固定的初始有效体积说明问题；体积明显变化时还须计入压缩功及边界输运。$P_{e\to i}$ 是电子向离子传能，$P_{\rm tr,e}$ 是未另行计入的输运能流。简化损失估计为
\begin{equation}\label{eq:cycle-atomicloss}
 P_{\rm ion}\simeq V_p n_en_0K_{\rm ion}\mathcal E_{\rm ion},
 \qquad P_{\rm rad}=\int_{V_p}n_e\sum_z n_z L_z(T_e)\,dV.
\end{equation}
$V_p$ 为等离子体有效体积（$\mathrm{m^3}$），$V_pn_en_0K_{\rm ion}$ 是均匀近似下每秒的电离次数，乘单次电离能 $\mathcal E_{\rm ion}$（J）得到功率。辐射式中的 $z$ 标记计入的杂质种类或电荷态，$n_z$ 是其密度，$L_z(T_e)$ 是温度相关的辐射系数；先对 $z$ 求和，再对体积元 $dV$ 积分，将局域辐射功率密度（$\mathrm{W/m^3}$）变为总功率。
$\mathcal E_{\rm ion}$ 是一次电离所需能量；激发、解离等需单独补充。$L_z$ 为相应杂质状态的辐射损失系数，单位 $\mathrm{W\,m^3}$，本式将辐射作为离开等离子体的净能量处理。

即使 $dW_e/dt>0$，也不能直接断言 $dT_e/dt>0$，因为新电离出的电子也要分享能量：
\begin{equation}\label{eq:cycle-temperature}
 \frac{dT_e}{dt}=\frac{2}{3N_ek_B}\frac{dW_e}{dt}
 -\frac{T_e}{N_e}\frac{dN_e}{dt}.
\end{equation}
左侧是平均电子温度变化率（K/s）。第一项把热能变化率按现有 $N_e$ 换算成升温率；第二项来自电子数变化，$dN_e/dt>0$ 时表示有限热能分给更多电子造成的温度稀释，若电子数减少则符号相反。两项都由对 $W_e=(3/2)N_ek_BT_e$ 求导得到。
这给出一个重要直觉：启动同时要``造出足够多的载流粒子''和``让每个粒子的平均能量升高''。控制器通常联合观察电流增长、线辐射、环电压及密度；仅看到可见光变亮并不能判定已经进入可持续放电。

\subsection{建流与爬升：磁面形成、电阻加热和电流扩散}\label{sec:cycle-ramp}

当环向电流 $I_p$ 建立，它产生极向磁场；在大长径比、圆截面近似下，
\begin{equation}\label{eq:cycle-btheta}
 B_\theta(r)\simeq\frac{\mu_0I_p(r)}{2\pi r}.
\end{equation}
$r$ 为距磁轴的近似小半径（m），$I_p(r)$ 是该半径以内包围的环向电流（A），$B_\theta(r)$ 为极向场（T）。$\mu_0$ 为真空磁导率，$2\pi r$ 为近似圆形极向回路长度；等式来自沿该回路的 Ampère 定律。边缘 $r=a$ 时，$a$ 为小半径且 $I_p(a)=I_p$。
环向场与极向场叠加形成螺旋磁力线。在适当外场和压力条件下形成闭合磁面后，沿场运动可在闭合区内循环，不再仅因快速沿场飞行就直接撞壁，因此启动时的平行粒子损失显著减弱；刮削层中的开放磁力线仍存在向偏滤器靶板的平行粒子与热量损失。跨磁面输运和轨道损失也不会由此消失，才逐步进入可用磁平衡与输运模型描述的阶段。

若用一个集总电流表示等离子体，并暂忽略动生电动势及分布电流效应，可将外部施加的感应电压写成
\begin{equation}\label{eq:cycle-ip}
 V_{\rm ext}=R_p(I_p-I_{\rm NI})+\frac{d(L_pI_p)}{dt},
 \qquad I_{\rm NI}=I_{\rm CD}+I_{\rm BS}.
\end{equation}
$R_p$ 是等效等离子体电阻（$\Omega$），$L_p$ 是这一集总模型的等离子体电感（H）。$I_p-I_{\rm NI}$ 表示需要感应维持的电流部分，$R_p(I_p-I_{\rm NI})$ 为相应有效电阻电压；$d(L_pI_p)/dt=L_p\dot I_p+I_p\dot L_p$ 为自磁链变化电压，分别含建流及电感变化贡献。两项均以 V 计。$I_{\rm NI}$ 是非感应电流总量（A），由外加驱动 $I_{\rm CD}$ 和自举贡献 $I_{\rm BS}$ 构成。
$V_{\rm ext}=-d\Phi_{\rm ext}/dt$ 只包含外部回路产生的磁通变化；这里已把自磁通项移到右侧，所以不能再把总磁通的式~\eqref{eq:cycle-faraday} 原样叠加一次。$I_{\rm CD}$ 为外加非感应电流驱动，$I_{\rm BS}$ 为与压力梯度有关的自举电流。本式只是控制直觉模型，真实的非感应电流应作为局域 $j_{\rm NI}(r)$ 进入电流扩散方程；它一般不能精确折算为一个串联电压源。

电流方程中 $j_{\rm NI}(r)$ 表示非感应电流密度（$\mathrm{A/m^2}$）。在 $I_{\rm NI}=0$、电感不变时，$V_{\rm ext}=R_pI_p+L_p\dot I_p$，欧姆加热为 $P_{\rm Ohm}=R_pI_p^2$。更一般地，电阻耗散是 $\int\eta j_{\rm Ohm}^2dV$。充分电离、经典碰撞近似下，Spitzer 电阻率具有 $\eta\propto Z_{\rm eff}\ln\Lambda\,T_e^{-3/2}$ 的标度：固定电流时温度越高，欧姆加热越弱；固定电场时电流会变化，不能沿用``加热必然减少''的结论。这里 $\eta$ 为电阻率（$\Omega\,\mathrm m$），$j_{\rm Ohm}$ 为欧姆电流密度，$\int\eta j_{\rm Ohm}^2dV$ 是对整个体积的耗散积分；$Z_{\rm eff}=\sum_i n_iZ_i^2/n_e$ 为有效离子电荷，$n_i,Z_i$ 是第 $i$ 类离子的密度与电荷数，$\ln\Lambda$ 为描述碰撞尺度比的 Coulomb 对数。比例式只给温度依赖，完整数值系数依赖温度的单位约定。高温阶段通常需要中性束或射频加热，相关机制见资料 \cite{C5}。

从 $\bm E=\eta\bm j$、$\mu_0\bm j=\nabla\times\bm B$ 及 Faraday 定律出发，在常电阻率、无流动且忽略非感应源项时，有
\begin{equation}\label{eq:cycle-diffusion}
 \frac{\partial\bm B}{\partial t}=\frac{\eta}{\mu_0}\nabla^2\bm B,
 \qquad \tau_R\sim\frac{\mu_0a^2}{\eta}.
\end{equation}
$\bm j$ 为电流密度，$\nabla\times$ 表示旋度，推导同时使用 $\nabla\cdot\bm B=0$ 并忽略位移电流。式中 $\partial\bm B/\partial t$ 是局部磁场变化率，$\nabla^2$ 为逐分量空间 Laplace 算子；$\eta/\mu_0$ 是磁扩散系数（$\mathrm{m^2/s}$）。$\tau_R$ 是尺度为小半径 $a$ 的电阻扩散时间（s），$\sim$ 只表示量级，实际时间常数还取决于边界与剖面。
因此总电流达到设定值时，内部剖面不一定已达到稳态。较快的电流爬升、加热导致的电阻率变化以及局域电流驱动，都能改变电流穿透和安全因子剖面。圆截面大长径比近似下
\begin{equation}\label{eq:cycle-qbeta}
 q(r)\simeq\frac{rB_\phi}{R_0B_\theta(r)},
 \qquad \beta=\frac{2\mu_0\langle p\rangle}{B_0^2}.
\end{equation}
$q(r)$ 为安全因子，近似表示磁力线绕极向一周时绕环向的周数；$r,R_0$ 分别为小半径与参考主半径，$B_\phi,B_\theta$ 为相应环向和极向场。$\beta$ 是无量纲压力比，$\langle p\rangle$ 为体积平均总热压强（Pa），$B_0$ 为采用的参考磁场，$B_0^2/(2\mu_0)$ 是该场的磁压力；不能混用不同 $B_0$ 定义下的 $\beta$ 数值。
$q$ 描述磁力线绕行的螺距，$\beta$ 比较等离子体压力与磁压力。它们有助于组织稳定性问题，但不存在脱离形状、剖面和壁面条件的单一通用安全阈值。爬升控制必须同时考虑电源限幅、剩余磁通、内感变化及稳定性裕度。

\subsection{平顶与稳态：持续平衡粒子、热量、电流和排热}\label{sec:cycle-steady}

在已经形成的近似固定体积内，用总体收支理解密度和温度控制最直观：
\begin{align}
 \frac{dN_e}{dt}&=S_{\rm fuel}+S_{\rm recyc}-S_{\rm exhaust},\label{eq:cycle-particles}\\
 \frac{dW}{dt}&=P_{\rm Ohm}+P_{\rm aux}+P_{\alpha,\rm dep}
 -P_{\rm rad}-P_{\rm tr},\qquad W=W_e+W_i.\label{eq:cycle-energy}
\end{align}
式~\eqref{eq:cycle-particles} 中，$S_{\rm fuel}$ 是燃料电离后给电子库带来的净源，$S_{\rm recyc}$ 是壁面再循环后返回并电离的源，$S_{\rm exhaust}$ 是净排出的电子等效粒子率，各以粒子数每秒计；这些是已折合到电子数的有效源汇，多电荷离子时需按电荷态加权。式~\eqref{eq:cycle-energy} 中，$W_e,W_i$ 分别为电子与离子热能，$W$ 为总热能（J），$P_{\rm aux}$ 为总吸收辅助加热功率，$P_{\alpha,\rm dep}$ 是氘氚例子中快 $\alpha$ 粒子沉积的自加热功率，$P_{\rm rad},P_{\rm tr}$ 为辐射和输运损失；各功率均为 W。其他反应燃料应将 $P_{\alpha,\rm dep}$ 推广为所有带电聚变产物的沉积功率，见第~\ref{sec:cycle-fuels} 节。
这里将辐射与输运损失分别计入，所有功率均采用吸收或沉积后的功率；若定义 $\tau_{E,\rm tr}=W/P_{\rm tr}$，就不能再把辐射包含在同一个输运项中重复扣除。稳态要求平均源汇抵消，而不是所有源汇都为零。燃料注入、再循环和排气共同决定粒子数，辅助加热、辐射和跨磁面输运共同决定热能；排出功率还必须能被偏滤器和第一壁承受。

\textbf{为什么平顶仍会结束？}在纯感应、$\dot I_p\simeq0$ 的平台上仍有 $V_{\rm ext}\simeq R_pI_p$，剩余变压器磁通继续被消耗。其预算可概括为
\begin{equation}\label{eq:cycle-fluxbudget}
 \int_{t_0}^{t_1}V_{\rm ext}\,dt
 =\left[L_pI_p\right]_{t_0}^{t_1}
 +\int_{t_0}^{t_1}R_p(I_p-I_{\rm NI})\,dt.
\end{equation}
$t_0,t_1$ 是预算区间的起止时刻；左侧是该区间消耗的有向外部磁通摆幅。方括号 $[L_pI_p]_{t_0}^{t_1}$ 表示末时刻的自磁链减去初时刻的自磁链，是自磁链增量；最后的积分是为克服有效电阻而消耗的伏秒。三者单位均为 Wb。若已进入近似固定几何、纯感应平顶，且环电压方向不变，尚余摆幅 $\Delta\Phi_{\rm rem}>0$ 对应可维持时间量级为
\begin{equation}\label{eq:cycle-remainingflux}
 t_{\rm rem}\simeq\frac{\Delta\Phi_{\rm rem}}{|V_{\rm ext}|}.
\end{equation}
这里 $t_{\rm rem}$ 的单位是 s，$\Delta\Phi_{\rm rem}$ 是当前阶段可用于该方向感应驱动的剩余磁通额度（Wb），$|V_{\rm ext}|$ 是维持平顶所需的环电压大小（V），并假设两者定义在同一外部磁通约定下。该估计忽略电阻演化及为正常降流应保留的磁通，不能把剩余额度全部用于平顶。
真正长期的非感应稳态需要电流驱动和自举电流共同补偿所需电流剖面，同时保持热、粒子和部件热平衡；总量上 $I_{\rm NI}\simeq I_p$ 只是必要的整体条件。

\subsection{聚变反应方案：燃料、反应产物与自加热条件}\label{sec:cycle-fuels}

前述粒子与能量收支可以用于不同燃料，但反应率、产物能量分配和辐射损失必须随之改变。燃料方案与约束方案是不同的选择：D--T、D--D 等是核反应通道；托卡马克、仿星器和惯性约束则是使燃料达到所需状态的方法。本报告仍以托卡马克控制为主。普通氢或氦放电也能研究击穿、磁平衡与控制，并不意味着这些放电在进行有用的聚变燃烧。

\subsubsection{用质量亏损和反应率比较方案}

反应释放的核能来自反应前后静质量之差：
\begin{equation}\label{eq:cycle-massdefect}
 \mathcal E_{\rm fus}=\left(\sum_{\rm in}m-\sum_{\rm out}m\right)c^2.
\end{equation}
$\mathcal E_{\rm fus}$ 为一次反应的释放能量（J），两项求和分别为入射与出射粒子的静质量总和（kg），$c$ 为真空光速。核反应常以 MeV 标注能量，$1\,\mathrm{MeV}\simeq1.602\times10^{-13}\,\mathrm J$。正质量亏损说明反应放能，但不代表在任意温度下反应都足够快。带正电的核相互排斥，反应截面取决于相对能量和量子隧穿；提高温度可以增加有效反应率，同时也改变输运与辐射损失。

常见通道可写成下列近似低入射能的能量分配 \cite{C10,C12}：
\begin{align}
 \mathrm D+\mathrm T&\longrightarrow{}^4\mathrm{He}\,(3.5)+n\,(14.1),\label{eq:cycle-reaction-dt}\\
 \mathrm D+\mathrm D&\longrightarrow\mathrm T\,(1.01)+p\,(3.02),\label{eq:cycle-reaction-ddp}\\
 \mathrm D+\mathrm D&\longrightarrow{}^3\mathrm{He}\,(0.82)+n\,(2.45),\label{eq:cycle-reaction-ddn}\\
 \mathrm D+{}^3\mathrm{He}&\longrightarrow{}^4\mathrm{He}\,(3.67)+p\,(14.68),\label{eq:cycle-reaction-dhe}\\
 p+{}^{11}\mathrm B&\longrightarrow3\,{}^4\mathrm{He}+8.7\,\mathrm{MeV}.\label{eq:cycle-reaction-pb}
\end{align}
这里 $\mathrm D={}^2\mathrm H$ 为氘核，$\mathrm T={}^3\mathrm H$ 为氚核，$p$ 为质子，$n$ 在这些反应式中为中子符号（不是密度），$\mathrm{He}$ 为氦核，$\mathrm B$ 为硼核，左上角数字为质量数。前四式中每个产物后的括号表示其近似动能，单位均为 MeV，而不是化学计量系数；最后一式的 8.7 MeV 是三个 $\alpha$ 粒子共享的总释放能，不表示每个粒子各得 8.7 MeV。$\alpha$ 即 ${}^4\mathrm{He}$ 核。两体反应中轻产物通常获得较多动能；三体终态的能量分配具有分布，不能照搬两体固定分配。

\begin{itemize}[itemsep=0.35em]
 \item \textbf{D--T：较易获得显著热核反应率的常用基准。}总释放约 17.6 MeV，其中约五分之一进入带电 $\alpha$ 粒子，可通过减速碰撞加热等离子体；约五分之四由中子带走。控制侧除密度、温度和形状外，还要考虑燃料比例、氦灰排出及快粒子约束；氚燃料循环与中子对部件的负荷也是系统问题。相对于下述先进燃料，D--T 在较低离子温度下即可取得较高反应率 \cite{C11}。
 \item \textbf{D--D：两条主要分支与后续反应。}式~\eqref{eq:cycle-reaction-ddp}、\eqref{eq:cycle-reaction-ddn} 分别释放约 4.03 和 3.27 MeV；在常见热核能区两条主要分支具有相近概率，但精确分支比随能量变化。第一分支生成氚，第二分支直接产生中子；生成的氚和 ${}^3\mathrm{He}$ 若继续反应，会引入 D--T 与 D--${}^3$He 通道。因此 D--D 不是无中子或无氚方案，模型应跟踪这些次生产物及其反应，而不能把两个分支简单压成单一固定产物。
 \item \textbf{D--${}^3$He：主通道的产物都带电。}式~\eqref{eq:cycle-reaction-dhe} 释放约 18.35 MeV，主反应本身不产生中子；但燃料中仍有氘，D--D 副反应及其后续 D--T 反应会带来中子。相较 D--T，该方案要求更苛刻的高温与约束条件，还需考虑 ${}^3\mathrm{He}$ 供给、较快质子与 $\alpha$ 粒子的轨道及能量沉积，不能直接沿用 D--T 的自加热系数。
 \item \textbf{$p$--${}^{11}$B：带电产物占优，热平衡更具挑战。}主通道生成三个 $\alpha$ 粒子，不直接生成中子，但副反应与实际杂质意味着不宜把系统绝对称为零中子。硼核电荷高于氢同位素，反应需要更高的相对能量；高温电子的轫致辐射等损失也更突出。观测到该反应并不等于实现自持点火或净功率增益，须分别核对反应率、加热与辐射收支 \cite{C12}。
\end{itemize}

\subsubsection{不同燃料共用的功率收支，不能共用的数值门槛}

对热平衡分布下的二体反应通道 $r$，用同一个形式计算单位体积反应率：
\begin{equation}\label{eq:cycle-general-rate}
 \mathcal R_r=\frac{n_a n_b}{1+\delta_{ab}}\langle\sigma v\rangle_r,
 \qquad
 P_{{\rm fus},r}=\int_{V_p}\mathcal R_r\mathcal E_{{\rm fus},r}\,dV.
\end{equation}
$\mathcal R_r$ 是每单位体积、每秒的反应次数（$\mathrm{m^{-3}s^{-1}}$），$a,b$ 标记入射核种，$n_a,n_b$ 为其密度；$\delta_{ab}$ 在相同核种时为 1，否则为 0。因而 D--D 必须使用 $n_D^2/2$，避免把每一对核重复计数，而 D--T 使用 $n_Dn_T$。$\langle\sigma v\rangle_r$ 是该通道截面与相对速度的热平均（$\mathrm{m^3/s}$）；分别对 D--D 的两个分支使用对应的部分反应率系数。$\mathcal E_{{\rm fus},r}$ 是该通道单次总释能（J），积分得到该通道总聚变功率 $P_{{\rm fus},r}$（W）。若燃料有非热束或强非 Maxwell 分布，必须对实际速度分布求反应率，不能直接套用热平均值。

真正进入等离子体热库的是带电产物的沉积功率：
\begin{equation}\label{eq:cycle-general-selfheat}
 P_{\rm self,dep}=\sum_r\int_{V_p}\mathcal R_r
 \sum_{c\in\mathcal C_r}f_{r,c}\mathcal E_{r,c}\,dV,
 \qquad P_{\rm self,dep}\geq P_{\rm rad}+P_{\rm tr}.
\end{equation}
$\mathcal C_r$ 是反应 $r$ 的带电产物集合，$c$ 标记产物（此处不是光速），$\mathcal E_{r,c}$ 是单次反应赋予该产物的动能，$f_{r,c}$ 是最终沉积在等离子体中的比例；中子不在该求和内。重复出现的同类产物须按数量计入，非单能产物应作能谱平均。外层求和覆盖所有计入通道，$P_{\rm self,dep}$ 为自加热总功率。右侧不等式是忽略持续外部加热和欧姆功率后，瞬时自加热至少补偿辐射与输运损失的必要收支条件；真正稳态取等号，并须另满足燃料补充、灰分排出及稳定性条件。若持续自加热大于损失，热能会上升，不能称为已经达到稳态。

改变燃料还会改变电子数与热容量。准中性条件及非相对论、近 Maxwell 分布下的总热能为
\begin{equation}\label{eq:cycle-multifuel-energy}
 n_e=\sum_i Z_i n_i,
 \qquad W=\frac32\int_{V_p}k_B\left(n_eT_e+\sum_i n_iT_i\right)dV.
\end{equation}
$i$ 标记全部离子物种，包括燃料、灰分与杂质；$Z_i$ 是离子电荷数，$n_i,T_i$ 为其密度与温度，$n_e,T_e$ 为电子密度与温度。第一式表示正负电荷在宏观尺度近似抵消；第二式分别累加电子和各离子的随机热能，温度均以 K 表示。对于 D--${}^3$He 和 $p$--${}^{11}$B，不能再取 $n_e=n_D+n_T$ 或直接沿用 $W/V_p=3nk_BT$。高电荷离子同时改变辐射与碰撞，因此比较燃料时应比较完整功率收支，而不是只比较单次反应释放多少 MeV。

\subsubsection{D--T 等比例例子：由一般收支得到简化点火条件}

\textbf{以氘氚燃料作定量示例。}式~\eqref{eq:cycle-reaction-dt} 中的反应产生约 $3.5\,\mathrm{MeV}$ 的带电 $\alpha$ 粒子和约 $14.1\,\mathrm{MeV}$ 的中子。中子不受磁场约束，其能量不能算作等离子体自加热；$\alpha$ 能量也只有实际沉积的部分才进入热平衡。等温、均匀、等比例氘氚、无杂质及 $T_e=T_i=T$ 的教学近似下，取 $n=n_D+n_T=n_e$，有
\begin{equation}\label{eq:cycle-alpha}
 \frac{P_{\alpha,\rm dep}}{V_p}
 =f_\alpha\frac{n^2}{4}\langle\sigma v\rangle_{DT}\mathcal E_\alpha,
 \qquad \frac{W}{V_p}=3nk_BT,
\end{equation}
$n_D,n_T$ 是氘和氚核密度，$n$ 为燃料离子总密度（$\mathrm{m^{-3}}$）；等比例时 $n_Dn_T=n^2/4$。$\langle\sigma v\rangle_{DT}$ 为氘氚热平均反应率系数（$\mathrm{m^3/s}$），$\mathcal E_\alpha$ 为每次反应赋予 $\alpha$ 粒子的能量（J）；$f_\alpha$ 为其中最终沉积于等离子体的比例。左式是 $\alpha$ 自加热功率密度（$\mathrm{W/m^3}$）。右式中电子和离子各贡献 $(3/2)nk_BT$，相加得到热能密度 $W/V_p$（$\mathrm{J/m^3}$）。式中温度 $T$ 以 K 计；若用 keV 表示 $k_BT$，还须换算为 J 后再核对功率。
其中 $0\leq f_\alpha\leq1$ 是能量沉积分数。忽略欧姆功率和辐射，仅为显示量纲而取损失 $W/\tau_E$，点火收支给出
\begin{equation}\label{eq:cycle-lawson}
 n\tau_E\geq\frac{12k_BT}{f_\alpha\langle\sigma v\rangle_{DT}\mathcal E_\alpha}.
\end{equation}
$\tau_E$ 为本例中以 $W/\tau_E$ 表示能量损失的约束时间（s），$n\tau_E$ 为密度与约束时间之积（$\mathrm{m^{-3}s}$）。分子 $12k_BT$ 来自热能系数 3 与等比例反应率系数 $1/4$；分母是单位密度配对的沉积加热能力。等号表示功率恰好平衡，不等号表示在这些假设下自加热至少可补偿该损失。
这不是通用的数值点火门槛：反应率强烈依赖温度，真实剖面、杂质辐射、氦灰和快粒子损失都会改变结果。它说明密度、温度与约束时间必须共同满足能量条件，不能只追求高温。通常定义的等离子体聚变增益 $Q=P_{\rm fusion}/P_{\rm aux}$ 也不是电站净发电效率；有限的高 $Q$ 不等于零外部加热的点火极限 \cite{C1}。

\textbf{稳态还必须稳定。}令净加热为 $\mathcal P(W)=P_{\rm in}(W)-P_{\rm loss}(W)$，则平衡满足 $\mathcal P(W_*)=0$。其中 $P_{\rm in}$ 为进入热库的总功率，$P_{\rm loss}$ 为离开热库的总功率，$W_*$ 是平衡热能；星号标记平衡值。在其他状态冻结的标量近似下，$\delta\dot W=\mathcal P'(W_*)\delta W$，这里 $\delta W=W-W_*$ 为小热能偏差，$\delta\dot W$ 为其变化率，$\mathcal P'$ 是净功率对热能的导数（$\mathrm{s^{-1}}$），只有 $\mathcal P'(W_*)<0$ 才对小热扰动恢复；真实系统应检查密度、热量及电流等耦合状态的闭环动力学。这解释了燃烧控制为何仍需要反馈，而不是达到某个温度后停止调节。

\subsection{位形转换：改变磁面几何，同时保持动力学可达}\label{sec:cycle-shape}

所谓位形，是等离子体横截面、磁轴、最后闭合磁面及 X 点的几何和拓扑结构。它与 L/H 模式这样的约束状态不是同一概念：前者描述磁面形状，后者描述输运及边缘屏障性质，二者可以耦合但不能互相替代。常见过程包括限制器位形转为偏滤器位形、改变伸长比和三角形变，以及在装置能力允许时切换单零和双零位形。

为与后文一致，定义总极向磁通 $F=2\pi\psi$，于是
\begin{equation}\label{eq:cycle-flux}
 \bm B_p=\frac{1}{2\pi R}\nabla F\times\bm e_\phi,
 \qquad \bm j\times\bm B=\nabla p.
\end{equation}
$F$ 为总极向磁通（Wb），$\psi=F/(2\pi)$ 为每弧度极向磁通；$R$ 是大半径，$\bm e_\phi$ 为环向单位向量，$\nabla F$ 在 $R$--$Z$ 平面内取梯度。$\bm B_p$ 为极向磁场，$\bm j$ 为总电流密度（$\mathrm{A/m^2}$），$\bm B$ 为总磁场，$p$ 为标量总压强。$\bm j\times\bm B$ 是磁力密度，$\nabla p$ 是压力梯度，单位均为 $\mathrm{N/m^3}$；二者相等表示此近似下的力平衡。
当变化慢于力学平衡建立、惯性可以忽略且采用轴对称标量压强时，第二式导出 GS 平衡；线圈电流变化改变外部磁通，等离子体电流和压力剖面则改变内部响应。求一个目标平衡只能证明存在某个静态解，并不能证明有限电压、有限时间的线圈能到达它。

\textbf{限制器与偏滤器的区别。}限制器位形的闭合区边界由接触实体限制器确定；偏滤器位形则由含 X 点的分离面界定，分离面外的开放磁力线把粒子与热量引向靶板。局部 X 点满足
\begin{equation}\label{eq:cycle-xpoint}
 \nabla F(R_X,Z_X)=\bm0,
 \qquad \det\left[\nabla\nabla F(R_X,Z_X)\right]<0,
 \qquad F_{\rm sep}=F(R_X,Z_X).
\end{equation}
$R_X,Z_X$ 为 X 点的大半径和竖直坐标（m）。$\nabla F=\bm0$ 表示两个一阶空间导数同时为零；$\nabla\nabla F$ 为包含二阶偏导的 $2\times2$ Hessian 矩阵，$\det$ 取其行列式，负值对应两个主方向曲率异号。$F_{\rm sep}$ 为分离面磁通值，以该鞍点的磁通作为整条分离面的等值。
负 Hessian 行列式表示这是鞍点。这里 $B_p=0$，但 $B_\phi$ 通常不为零，所以 X 点不是总磁场消失的点。上述条件是局部条件，仍须追踪磁通等值线并检查它与壁面的连接，才能判定位形拓扑。资料 \cite{C6} 给出了闭合磁面、X 点与刮削层（SOL）的物理图像。

设目标边界控制点为 $(R_k,Z_k)$，可用 $e_{F,k}=F(R_k,Z_k)-F_{\rm sep}$ 构造边界误差，其中 $k$ 标记控制点，$e_{F,k}$ 是该点磁通偏离目标分离面值的误差（Wb），并联合控制磁轴位置、X 点坐标及靶板打击点。X 点附近 $|\nabla F|$ 很小，线性估计 $\delta\ell_n\simeq\delta F/|\nabla F|$ 会变得敏感，其中 $\delta\ell_n$ 是法向位移（m），$\delta F$ 是小磁通误差（Wb），$|\nabla F|$ 是法向磁通梯度大小（Wb/m）；采用 $\ell_n$ 而不用 $n$，以免与粒子密度混淆。在 X 点本身梯度为零，这一线性公式失效，因此不能仅靠磁通误差很小就认定几何位置误差也很小。

\textbf{转换过程的三个耦合限制。}首先，压力和电流剖面会随加热及几何变化演化，需要同时更新平衡；其次，式~\eqref{eq:cycle-coils} 的电感及壁涡流使形状响应滞后；最后，高伸长位形可能出现更强的垂直不稳定性。局部教学模型可写为
\begin{equation}\label{eq:cycle-vertical}
 \delta\dot Z=\gamma_Z\delta Z+b_Z\delta I_{\rm VS}+d_Z,
 \qquad L_{\rm VS}\delta\dot I_{\rm VS}+R_{\rm VS}\delta I_{\rm VS}=\delta V_{\rm VS}.
\end{equation}
第一式中 $\delta Z$ 为偏离参考平衡的竖直位移（m），$\delta\dot Z$ 为其速度；$\gamma_Z$ 为开环增长率（$\mathrm{s^{-1}}$），$\delta I_{\rm VS}$ 为垂直稳定线圈电流相对参考值的增量（A），$b_Z$ 为电流到位移变化率的局部响应系数，$d_Z$ 为按速度单位折合的扰动。第二式中 $L_{\rm VS},R_{\rm VS}$ 是等效电感和电阻，$\delta V_{\rm VS}$ 是增量电压（V）；电感压降 $L_{\rm VS}\delta\dot I_{\rm VS}$ 与电阻压降 $R_{\rm VS}\delta I_{\rm VS}$ 之和等于该电压。
当 $\gamma_Z>0$ 时，位置误差会自行增长；反馈必须在诊断延迟与电源限幅下提供足够快的响应。真实垂直稳定模型通常包含更多壁电流状态，不能由此两式决定某装置的控制增益。

实际轨迹应由一族可达平衡组成，并约束线圈电流、电压、壁间隙和热负荷。可用 $s(\xi)=3\xi^2-2\xi^3$、$0\leq\xi\leq1$ 平滑插值参考量，其中 $\xi=(t-t_a)/(t_b-t_a)$ 是转换起止时刻 $t_a,t_b$ 之间的归一化时间，$s$ 是由 0 变化到 1 的插值权重，可令任意参考量 $y_{\rm ref}=(1-s)y_a+sy_b$ 从初值 $y_a$ 过渡到终值 $y_b$，使起止速度为零；但平滑插值本身不保证中间存在平衡。单零转双零还须控制上下鞍点的磁通差 $\Delta F_X=F_{X,u}-F_{X,l}$，其中 $u,l$ 分别指上、下 X 点，$\Delta F_X$ 以 Wb 计；接近零时两侧排热分配可能对微小扰动敏感。转换完成的判据应包含几何误差、稳定性裕度及靶板负荷，而不是仅看边界轮廓是否像目标图形。

\subsection{边界位形与排热：限制器、偏滤器和 XPT}\label{sec:boundary-exhaust}

前一节说明怎样改变位形，本节解释为什么要改变它。对初学者而言，应把三个问题分开：\textbf{闭合磁面区在哪里终止，越过该边界的粒子和热量沿什么路径运动，最后由哪些材料部件接收。}限制器与偏滤器首先区别在第一个问题，先进偏滤器还进一步改变后两个问题。研究位形既是研究约束、稳定性与线圈可控性，也是研究聚变装置能否持续排热和排出粒子。

\subsubsection{``最外层边界''究竟是哪一种边界}

本文中的 LCFS 是\textbf{围绕磁轴的嵌套闭合磁面族的最外边界}，不是全部等离子体物质的外轮廓，也不是真空室内壁。闭合磁面可想成一层环形的洋葱皮：磁场处处与它相切，磁力线在理想轴对称平衡中留在该面上。``面闭合''不要求每一条磁力线有限绕行后都回到起点；在无理旋转变换的面上，单条场线可以不断绕行而不精确闭合。

由式~\eqref{eq:cycle-flux} 可得
\begin{equation}\label{eq:edge-flux-tangent}
 \bm B\cdot\nabla F=0,\qquad
 \bm B=\bm B_p+B_\phi\bm e_\phi .
\end{equation}
这里 $F$ 为总极向磁通，$\nabla F$ 是磁面的法向梯度，$\bm B_p$ 和 $B_\phi\bm e_\phi$ 分别为极向和环向磁场；点积为零表示磁场与 $F$ 等值面相切。等式使用轴对称条件 $\partial_\phi F=0$。但是粒子导引中心仍可因碰撞、漂移和湍流跨磁面输运，所以磁面不是不透物质的实体外壳。

\textbf{LCFS 外仍有等离子体。}紧邻闭合区、磁力线与材料表面相交的边缘区域称为刮削层（scrape-off layer，SOL）；它通过跨边界输运不断得到粒子和能量，并沿场向材料表面损失。更外侧还可能有较稀薄的等离子体及中性气体。因而密度、温度不会在 LCFS 上普遍地跳变为零。实验所说的光学发光边缘、密度阈值边缘和磁拓扑边界也不必重合。

还应区分\textbf{第一壁}与\textbf{数值边界}：第一壁是面向等离子体的真空室内部材料表面；计算矩形则只是求解磁通所用的网格覆盖范围。在偏滤器位形中，常把穿过主 X 点、分隔核心闭合磁面区与外部开场线区的磁通面称为\textbf{分离面}，其包围核心的部分构成 LCFS 的极限边界。限制器位形也有 LCFS，但通常没有用主 X 点来定义这一边界。这些区别见图~\ref{fig:crosssec} 和资料~\cite{C6}。

\subsubsection{限制器是什么，限制器位形怎样形成}

\textbf{限制器（limiter）是一个实体部件}，例如从第一壁向内突出的耐热块、环形构件或专门设计的壁面局部。它不是某一种线圈，也不是一条人为添加的磁力线。给定限制器几何后，线圈与等离子体电流共同确定磁面形状；围绕磁轴的磁面向外扩展，首先与限制器相切的那一层决定有效 LCFS。更外侧的场线会与材料相交，沿这些场线运动的粒子可被限制器接收。

因此，产生限制器位形一般不要求主动制造 X 点：PF 线圈调节磁轴、尺寸与形状，使闭合区与预定限制器接触即可。限制器的作用是\textbf{规定接触位置和闭合区尺寸，并用预定部件承受边缘粒子、热量与侵蚀}，避免接触位置完全失控；它不是把所有高能粒子反射回核心的机械挡板。

这种位形在启动、低功率调试和基础等离子体研究中很有用，几何与控制目标通常较简单。但热负荷可能集中在有限接触区，材料侵蚀产生的杂质又离热核心较近，持续高功率运行会遇到排热和污染困难。具体热负荷与约束表现依赖装置和工况，不能由``限制器''三个字唯一决定。限制器材料也依装置而异，可采用石墨、钨或其他面向等离子体材料。

\subsubsection{偏滤器、X 点、偏滤器腿和打击点}

\textbf{偏滤器（divertor）是一套磁拓扑与排热排粒子部件相配合的系统}，包括组织开放场线的磁场、接收热流的靶板、相应的挡板/腔室、冷却和抽气设施等。称某个平衡为``偏滤器位形''，重点是其闭合区由含主 X 点的分离面界定，核心 LCFS 可与材料壁保持间隙，外部场线则连接到预定排热区域。

式~\eqref{eq:cycle-xpoint} 已给出 X 点的数学判据。物理上，PF 线圈产生的极向场与等离子体电流等产生的极向场在某处抵消，便可能形成 $B_R=B_Z=0$ 的鞍点。不能只要求抵消：磁轴处也可有 $B_p=0$，但磁轴是磁通极值点，而普通 X 点是两主方向曲率异号的鞍点。普通 X 点附近可在局部正交坐标中展开为
\begin{equation}\label{eq:edge-xlocal}
 F-F_X\simeq\frac12(h_1\xi^2-h_2\eta^2),\qquad h_1,h_2>0 .
\end{equation}
$F_X$ 是 X 点磁通，$\xi,\eta$ 是以 X 点为原点、沿 Hessian 主方向的距离（m），$h_1,h_2$ 为两方向曲率绝对值（Wb/m$^2$）；忽略了三阶及更高阶项。令 $F=F_X$ 得到两条相交的局部直线，说明截面上为什么出现 ``X''。这是\textbf{极向投影的鞍形结构}：$B_\phi$ 通常仍很强，总磁场没有在这里消失，三维磁力线也不是彼此穿过或分叉的轨道。

沿分离面从 X 点向材料靶面延伸的支路称为\textbf{偏滤器腿}。分离面与靶面的交点在截面中称为\textbf{打击点}（strike point），轴对称三维装置中对应绕环的打击线；实际沉积热流有有限宽度，并非全部集中在一个数学点。上下只有一侧承担主要分离拓扑的常见位形称单零（single-null），上下各有一个相关主 X 点的称双零（double-null）。平衡双零要求上下相关 X 点的磁通匹配；稍有失配就可能改变两侧的连接与功率分配。

在单零截面中，X 点下方两腿之间还可有\textbf{私有磁通区}（private-flux region）：其中场线连接靶面，却不绕过核心外侧的上游 SOL。它不是第二个约束核心，也不应把所有 LCFS 外区域都当成同一种连接关系。具体连接须追踪场线确认。

\begin{figure}[!ht]
\centering
\begin{tikzpicture}[font=\small,>=Stealth,scale=0.93]
 \begin{scope}[shift={(2.5,0)}]
 \node at (0,2.2) {（a）磁轴：极值点};
 \foreach \r in {0.5,0.9,1.3,1.7}{\draw[blue!65!black] (0,0) ellipse ({\r} and {0.85*\r});}
 \fill[blue!65!black] (0,0) circle (2pt);
 \draw[->,black!40] (-2,0)--(2,0) node[right] {$\xi$};
 \draw[->,black!40] (0,-1.7)--(0,1.8) node[right] {$\eta$};
 \node[align=center] at (0,-2.3) {附近为嵌套闭合等值线\\局部模型：$F-F_A\propto\xi^2+\eta^2$};
 \end{scope}
 \begin{scope}[shift={(9.5,0)}]
 \node at (0,2.2) {（b）普通 X 点：鞍点};
 \draw[blue!70!black,very thick] (-1.55,-1.55)--(1.55,1.55);
 \draw[blue!70!black,very thick] (-1.55,1.55)--(1.55,-1.55);
 \foreach \c in {0.25,0.7}{
  \draw[orange!90!black,domain=-1.3:1.3,samples=70] plot (\x,{sqrt(\x*\x+\c)});
  \draw[orange!90!black,domain=-1.3:1.3,samples=70] plot (\x,{-sqrt(\x*\x+\c)});
  \draw[orange!90!black,domain=-1.3:1.3,samples=70] plot ({sqrt(\x*\x+\c)},\x);
  \draw[orange!90!black,domain=-1.3:1.3,samples=70] plot ({-sqrt(\x*\x+\c)},\x);
 }
 \fill (0,0) circle (2pt);
 \node[fill=white,inner sep=1pt] at (0,0.3) {$X$};
 \node[align=center] at (0,-2.3) {蓝线为 $F=F_X$，橙线为相邻磁通\\局部模型：$F-F_X\propto\xi^2-\eta^2$};
 \end{scope}
\end{tikzpicture}
\caption{先辨认磁轴与 X 点：两者均可有 $B_p=0$，但一个是极值点，一个是鞍点。本图按局部二次磁通模型绘制，坐标为任意一致单位；曲线是极向磁通等值线，即场线的极向投影，蓝线交叉不代表三维总磁力线相互穿过。只有检查全局连接后，才能确定哪一支属于核心分离面。本文原创教学图。}\label{fig:edge-o-x-atlas}
\end{figure}

\textbf{读图顺序。}先在图~\ref{fig:edge-o-x-atlas}(a) 找到一圈圈闭合等值线，再看 (b) 的两条蓝线怎样把邻域分成四区；不要把 (b) 当作整个真空室。回到图~\ref{fig:crosssec}(b)，才可辨认主 X 点上方的核心闭合区和下方通向靶板的腿。

\subsubsection{为什么开放场线能引导粒子和热流，为什么要引向靶板}

这里的``开放''是指场线在装置的等离子体可达区域内与材料表面相交，\textbf{不是磁力线在真空中凭空终止}；磁场仍满足 $\nabla\cdot\bm B=0$。在回旋半径远小于几何尺度的磁化近似下，带电粒子绕场回旋，同时沿场运动。磁力不能做功：$q_s(\bm v\times\bm B)\cdot\bm v=0$，其中 $q_s$ 是粒子带符号电荷，$\bm v$ 为速度；磁场的主要作用是约束横向运动并组织路径，沿场加速则可来自电场和集体压力驱动。

一条场线及其平行运动的简化表达为
\begin{equation}\label{eq:edge-follow}
 \frac{d\bm x_B}{d\ell}=\bm b,\qquad
 \dot{\bm X}\simeq v_\parallel\bm b+\bm v_d,\qquad
 L_\parallel=\int_{\rm upstream}^{\rm target}\frac{B}{B_p}\,d\ell_p .
\end{equation}
$\bm x_B$ 为场线位置，$\ell$ 为沿场弧长，$\bm b=\bm B/B$ 为单位场向量，$B=|\bm B|$；$\bm X$ 为导引中心位置，$v_\parallel$ 为平行速度，$\bm v_d$ 汇总代表漂移等垂直运动。$L_\parallel$ 是从所选上游位置到靶面的连接长度，$d\ell_p$ 是同一场线在极向截面的投影弧长，$B_p=|\bm B_p|$。最后一式适用于可按极向投影参数化的场线段，表明环向绕行会使实际路径比截面中的腿长长得多。理想分离面上趋近 X 点的连接长度可发散；实际热流须对有限宽 SOL、有限距离及输运进行计算，不能把该发散当成无限排热能力。

在 SOL 中，电子沿场的热传导和平行粒子流通常远强于横向的碰撞输运，因此越过 LCFS 的能量和粒子可沿开放场线到达靶板。电子较快到达材料表面所建立的鞘层电场，又使电子与离子通量相互约束，并影响离子的入射能量。前述图像不表示每个粒子严格贴着一条场线：有限轨道、漂移、湍流、电荷交换以及中性粒子运动都能改变沉积位置；中性粒子不受磁场直接引导。开放场线的两端都可能连接材料，实际流动方向和内外靶功率分配由源项、温度、压力、电势与漂移等决定，不能只凭一根箭头判定。

\textbf{为什么需要专门的靶板？}有限输运和持续加热意味着能量必然流出核心，持续供气及聚变燃烧也要求排出部分粒子和氦灰。若损失直接落在普通第一壁的局部，可能造成过热、溅射和杂质返流。偏滤器让主要带电粒子排热通道进入离核心较远、可冷却、可诊断、可更换并与抽气系统配合的区域。靶板接收粒子后会发生中和、反射、再循环与侵蚀；\textbf{到达靶板不等于已经被抽走}，中和后的气体还需通过抽气通道排出。磁场也不能把 DT 反应中不带电的中子导向靶板，中子负荷必须由包层、屏蔽等系统另行处理。

图~\ref{fig:edge-exhaust-chain} 概括了这条因果链。边缘输运与靶面相互作用的实验和模型仍需单独验证，不能只用磁平衡图代替\cite{Dsol,Ddetach}。

\begin{figure}[htbp]
\centering
\begin{tikzpicture}[font=\small,>=Stealth,
 box/.style={draw=black!55,rounded corners,align=center,text width=3.6cm,minimum height=1.0cm}]
 \node[box] (core) at (0,0) {核心闭合磁面区\\跨 LCFS 输运};
 \node[box] (sol) at (5.0,0) {SOL 开放场线\\沿场流动与热传导};
 \node[box] (div) at (5.0,-2.0) {偏滤器体积\\辐射、原子过程与展宽};
 \node[box] (target) at (0,-2.0) {靶板与冷却系统\\接收剩余热流和粒子};
 \draw[->,thick] (core)--(sol);
 \draw[->,thick] (sol)--(div);
 \draw[->,thick] (div)--(target);
 \node[align=center,font=\footnotesize] (pump) at (0,-3.55) {中和、再循环及抽气\\只有部分粒子最终被泵排出};
 \draw[->] (target)--(pump);
\end{tikzpicture}
\caption{从核心损失到工程排热的概念链。辐射在体积中分散能量，冷却系统带走材料接收的热量；图不表示所有能量只沿这一通道损失，也不表示辐射或抽气已由当前磁平衡仿真器求解。}\label{fig:edge-exhaust-chain}
\end{figure}

\begin{figure}[!ht]
\centering
\includegraphics[width=0.94\linewidth]{splend1d_geometry.pdf}
\caption{公开论文中的沿场几何：SPLEND1D 原文图 1\cite{Ddetach}。$\bm b$ 是单位磁场方向，$s$ 沿磁通管增加并指向 target（靶面）；$\theta$ 与 $\varphi$ 分别标记极向、环向坐标。\textbf{注意原图的 $\alpha$ 是磁场螺距角，不是式~\eqref{eq:edge-targetheat} 中磁场与靶面的入射角。}图像仅作 SVG 到 PDF 的格式转换，未改动内容；O. F\'evrier 等，2024，\href{https://arxiv.org/html/2402.04656v1\#S2.F1}{原图链接}，\href{https://creativecommons.org/licenses/by/4.0/}{CC BY 4.0}。}\label{fig:edge-published-connection}
\end{figure}

\textbf{怎样读论文图。}沿图中的磁通管从左向右追踪，看到的不是一条截面短线，而是一段带有环向位移的实际路径。因此式~\eqref{eq:edge-follow} 中要用 $B/B_p$ 把极向投影长度换成沿场长度。图右端的靶面给出这条开场线的材料边界条件；它不要求整台装置的磁场在该处终止。

\subsubsection{怎样减少靶板峰值热流，``脱靶''又是什么}

降低局部负荷有两类互补办法：让到达靶面的功率更少，以及让剩余功率分布在更大的有效面积上。在忽略横向热流和局部体积损失的投影近似下，
\begin{equation}\label{eq:edge-targetheat}
 q_{\rm tgt}\simeq q_\parallel|\bm b\cdot\bm n_{\rm tgt}|
 =q_\parallel\sin\alpha,\qquad
 \langle q_{\rm tgt}\rangle\sim\frac{P_{\rm tgt}}{A_{\rm wet}} .
\end{equation}
$q_\parallel$ 是靶前沿场热流密度，$q_{\rm tgt}$ 是单位靶面面积接收的热流密度（均为 W/m$^2$），$\bm n_{\rm tgt}$ 为靶面单位法向量，$\alpha$ 为磁场与靶面的夹角。$P_{\rm tgt}$ 是沉积在所考虑靶区的总功率（W），$A_{\rm wet}$ 为有效润湿面积（m$^2$），尖括号表示面积平均而非峰值。小入射角和磁通扩展可展宽足迹，但表面起伏、遮蔽、场线映射和漂移均会影响实际负荷，不能把入射角任意降至零来获得无穷大承载能力。

\textbf{磁通扩展}可由相邻磁面间距理解：在光滑位置，小磁通间隔 $\Delta F$ 对应的法向间距近似为 $\Delta\ell_n\simeq\Delta F/|\nabla F|=\Delta F/(2\pi R B_p)$。这里 $\Delta\ell_n$ 是间距（m），$\Delta F$ 是两面磁通之差（Wb）；其余量定义同前。减小局部 $B_p$ 可以增大极向间距，但靶面足迹还包含大半径和入射角的映射，且该线性近似在 X 点本身失效。增加连接长度则可能为热传导、横向输运及粒子--中性气体相互作用提供更多路径与体积；它并不在没有耗散的情况下自动消灭功率。

\textbf{脱靶（detachment）是边缘等离子体的运行状态，不是磁拓扑名称。}通过足够的密度、适量杂质辐射和中性粒子相互作用，可在到达靶板前显著降低等离子体温度、沿场动量以及靶面热流，较深脱靶时还可出现靶面离子通量下降。激发辐射、电离能量消耗、电荷交换和复合等过程的相对重要性随工况变化，不能都等同于``粒子消失''。辐射把能量以光子形式转移至更大面积的材料表面，能量守恒仍须满足。所谓脱靶也不意味着磁力线从靶板断开；偏滤器位形可处于附着或脱靶状态，XPT 本身不保证脱靶。需要防止辐射冷却区过度侵入核心，同时协调燃料、杂质与抽气\cite{Ddetach,Dadc}。

\begin{figure}[!ht]
\centering
\begin{tikzpicture}[font=\small,>=Stealth]
 \foreach \x/\title in {0/{（a）附着：较多功率到靶},7.2/{（b）脱靶：靶前体积耗散}}{
  \begin{scope}[shift={(\x,0)}]
   \node at (2.6,0.65) {\title};
   \draw[blue!65!black,thick] (0.1,-0.2)--(5.2,-0.2);
   \draw[blue!65!black,thick] (0.1,-1.3)--(5.2,-1.3);
   \fill[black!35] (5.2,0.05) rectangle (5.45,-1.55);
   \node[rotate=90] at (5.8,-0.75) {靶板};
   \node at (0.6,-1.75) {上游 SOL};
  \end{scope}
 }
 \draw[->,line width=3pt,red!75!black] (0.3,-0.75)--(5.05,-0.75);
 \node[align=center] at (2.8,-2.5) {靶前仍较热\\靶板承担较大热流};
 \draw[->,line width=3pt,red!75!black] (7.5,-0.75)--(9.5,-0.75);
 \fill[orange!12] (9.7,-1.25) rectangle (11.4,-0.25);
 \draw[->,line width=1.2pt,orange!90!black] (11.5,-0.75)--(12.25,-0.75);
 \foreach \dx in {-0.45,0,0.45}{
  \draw[->,orange!90!black] (10.5,-0.3)--({10.5+\dx},0.25);
  \draw[->,orange!90!black] (10.5,-1.2)--({10.5+\dx},-1.75);
 }
 \node[align=center] at (10.2,-2.5) {辐射与原子过程降低靶前能量\\磁场连接仍保留，剩余热流到靶};
\end{tikzpicture}
\caption{用同样的开场线连接区别附着与脱靶。蓝线表示磁通管的两侧，红/橙箭头粗细仅定性表示沿场剩余功率；向外箭头表示辐射等体积能量转移，不是带电粒子横穿磁场的轨道。脱靶依赖粒子、动量和能量收支，不能仅凭靶温低来判定。本文依据脱靶物理\cite{Ddetach}绘制的概念图，无数值标度。}\label{fig:edge-attachment-atlas}
\end{figure}

\textbf{读图提示。}比较图~\ref{fig:edge-attachment-atlas} 的左右两图，蓝线始终接到靶板；改变的是靶前的温度、动量与功率沉积。向外辐射的能量还会落到其他材料表面，并没有从系统中消失。

\subsubsection{靶板由什么组成，材料选择有哪些代价}

靶板不是一整块任意金属板，而通常是\textbf{面向等离子体的装甲、连接层、热沉/冷却管及支撑结构}组成的工程部件。以 ITER 型方案为例，可采用钨（W）装甲、铜连接层和铜铬锆（CuCrZr）合金冷却管，背后配合承力结构及主动水冷\cite{Dmaterial,Diter}。钨熔点高、在适当条件下侵蚀较小，适于高负荷区域；但仍会受热疲劳、开裂、再结晶、瞬态熔化和辐照损伤影响。钨还是高原子序数杂质，少量进入热核心就可能带来很强的辐射，因此必须同时控制侵蚀和杂质输运。

石墨和碳纤维复合材料也曾广泛用于限制器或偏滤器，具有耐热冲击等优点；但碳侵蚀及其与氢同位素的共沉积会带来燃料、尤其是氚的滞留问题\cite{Dcarbon}。不同装置的材料、冷却能力与可更换性并不相同。研究位形的目的之一，正是在材料和工程约束下把长期及瞬态负荷降低到可承受范围，而不是寻找一种可以无限承热的材料。

\begin{figure}[!ht]
\centering
\begin{tikzpicture}[font=\small,>=Stealth]
 \fill[gray!50] (1,0) rectangle (8.5,1.0);
 \fill[orange!35] (1,-0.35) rectangle (8.5,0);
 \fill[yellow!30!brown!25] (1,-2.1) rectangle (8.5,-0.35);
 \draw[thick] (1,-2.1) rectangle (8.5,1.0);
 \draw (1,0)--(8.5,0) (1,-0.35)--(8.5,-0.35);
 \fill[cyan!20] (4.7,-1.25) ellipse (1.4 and 0.55);
 \draw[blue!65!black,very thick] (4.7,-1.25) ellipse (1.4 and 0.55);
 \node at (4.7,0.5) {钨装甲：直接接收粒子与热流};
 \node at (4.7,-0.18) {铜连接层};
 \node at (4.7,-1.05) {冷却水通道};
 \node[anchor=west] at (9,-0.55) {CuCrZr 热沉/管壁};
 \draw[->] (9,-0.7)--(7.0,-0.9);
 \foreach \x in {2,3.5,5,6.5,8}{\draw[->,red!75!black,thick] (\x,2)--(\x,1.1);}
 \node at (5,2.4) {来自等离子体的热负荷};
 \draw[->,red!70!black,thick] (1.5,0.8)--(1.5,-1.25)--(3.15,-1.25);
 \node[align=center] at (1.3,-3.0) {热传导\\装甲 $\to$ 热沉 $\to$ 水};
 \draw[->,blue!70!black,thick] (3.8,-1.48)--(5.6,-1.48);
 \node[align=center] at (7.8,-3.0) {冷却水把热带出部件\\再由外部冷却系统移走};
\end{tikzpicture}
\caption{钨--铜--CuCrZr 型靶板的局部剖面概念图。材料堆叠和冷却通道按公开的 ITER 型 monoblock 资料\cite{Dmaterial}说明性绘制；不是制造图，厚度、通道截面和连接方式未按比例，真实 monoblock 中管道被材料包围，亦有独立支撑结构。红色箭头表示能量传递，蓝色箭头表示水流；本图不表示等离子体流进冷却水。本文原创教学图。}\label{fig:edge-target-section}
\end{figure}

\textbf{读图提示。}先找最上方与等离子体直接相对的钨表面，再沿热传导路径找到内部冷却通道。磁场负责把带电粒子损失组织到靶区，部件内部则靠导热和冷却介质排热，是两个不同环节。真实部件还要承受应力、循环载荷和辐照损伤。

\subsubsection{XPT 是什么，次 X 点与主 X 点有何区别}

\textbf{XPT 通常指 X-point target（X 点靶）先进偏滤器位形}：在核心主分离面所对应的主 X 点之外，进一步利用偏滤器腿/靶前区域的\textbf{次级 X 点}组织排热区磁场。``主''与``次''描述拓扑作用，不是两个 X 点具有不同的基本物理性质；它们都可满足局部鞍点判据。主 X 点界定核心闭合区的分离拓扑，次 X 点主要改变开放场线区的磁通扩展、连接长度、支路和耗散区位置。

次 X 点不必与主 X 点具有完全相同的磁通。可用
\begin{equation}\label{eq:edge-secondaryflux}
 \Delta F_{21}=F_{X,2}-F_{X,1},\qquad
 \epsilon_{X,2}=\frac{F_{X,2}-F_B}{F_B-F_A}
\end{equation}
描述其磁通接近程度。$F_{X,1}$ 和 $F_{X,2}$ 分别为所选主、次 X 点的磁通（Wb），$\Delta F_{21}$ 为二者之差；$F_B$ 为核心 LCFS 边界磁通，$F_A$ 为磁轴磁通，$\epsilon_{X,2}$ 为无量纲归一磁通偏差，假设 $F_B\ne F_A$。该偏差的符号依赖磁通约定，不能脱离约定直接解释为空间上的内外。次 X 点的位置与磁通差共同决定哪些有限宽 SOL 通道会感受到它；仅有 $n_X$ 增加不足以证明形成了有效的 XPT 排热通道。

XPT 研究希望在远离热核心的位置获得更有利的排热几何，并使辐射或脱靶前沿更易稳定。例如 TCV 的 X-point target radiator（XPTR）实验研究了在次 X 点附近局域辐射的状态\cite{Dxptr}；\textbf{XPTR 是特定辐射运行状态，XPT 是磁位形，二者不能直接画等号}。排热收益仍需结合中性粒子、杂质、靶板与挡板验证。

\textbf{本文项目中的特殊约定。}当前 F2 任务使用带上下主级和次级结构的目标，要求识别到四个 X 点，并将主级放在槽 1、2，次级放在槽 0、3；四个目标点按 $Z$ 从高到低排列时，外侧两槽对应这里的上下次级点。该槽位约定来自 \code{environment/xpt_utils.py}，相关特征和评分见第~\ref{sec:rl} 章。\textbf{``XPT 必须恰好四个 X 点''不是普遍定义}：单零 XPT 也可围绕一个主 X 点及相应次级结构构造。也不能要求所有次 X 点都与 $F_B$ 严格相等；项目的磁通偏差目标与容差是具体控制指标。

此外，本项目工具提取的``打击点''来自\textbf{计算矩形边界上的磁通等值交点}，是几何控制代理量；只有明确计算边界与真实材料靶面的对应关系后，才能把它解释为装置靶板上的物理打击位置。仅由 GS 平衡、X 点坐标或几何奖励变好，不能推出热流降低、脱靶实现或材料安全。

\clearpage
\noindent\textbf{XPT 读图页：先找到核心，再找到主、次 X 点。}
\begin{figure}[!ht]
\centering
\begin{minipage}[c]{0.41\linewidth}
 \centering
 \includegraphics[width=\linewidth,height=0.68\textheight,keepaspectratio]{project_xpt_13906.png}
\end{minipage}\hfill
\begin{minipage}[c]{0.53\linewidth}
 \small
 \textbf{从中间向外读这幅图}\par
 (1) 中间一圈圈细线：围绕磁轴的闭合磁通轮廓。横轴 $r$ 对应本文大半径 $R$，纵轴 $z$ 对应 $Z$，单位为 m。\par
 (2) 蓝色星号 $X_1,X_2$：本任务的上下主 X 点，位置约为 $Z=\pm1.0$ m，组织核心闭合区边界。\par
 (3) 橙色星号 $X_0,X_3$：上下次 X 点，位置约为 $Z=\pm1.4$ m，进一步改变腿/靶前区域的连接结构。\par
 (4) 红色三角：工具标出的计算边界交点。它们是打击点的几何代理，不能仅凭此图认定为真实靶板上的测量点。\par
 (5) 粗虚线强调边界及相关磁通支路；请同时比较主、次 X 点的磁通差，而非只看图形像不像字母 X。\par
 本图是项目已有的 \code{13906_500} 初态可视化，不是所有 XPT 都有四个 X 点，也不是新增控制算法的实验结果。
\end{minipage}
\caption{本项目的四 X 点磁通图及中文读图索引。原图为本仓库已有的 13906 炮磁通可视化，只按比例排版，未改动其点位或磁通线。它直接展示本任务中主/次 X 点槽位的空间含义；与前面的局部鞍点概念图一起阅读。}\label{fig:edge-project-xpt-atlas}
\end{figure}
\clearpage

\subsubsection{还有哪些位形，如何产生并选择它们}

常见方案并不构成单一优劣排序。表~\ref{tab:edge-configs} 按拓扑和排热设计列出几类，相关比较见资料~\cite{C6,Dgeometry,Dadc}。

\begin{table}[!ht]
\centering\small
\caption{边界与先进偏滤器位形：设计目的和主要代价}\label{tab:edge-configs}
\begin{tabular}{L{2.5cm}L{5.3cm}L{5.2cm}}
\toprule
方案 & 主要物理图景与研究目的 & 控制与工程方面的注意点 \\
\midrule
限制器 & 实体接触确定 LCFS；用于启动、调试和基础研究 & 接触区热负荷、侵蚀与核心杂质污染 \\
单零 / 双零 & 一侧或上下主 X 点组织腿和靶区；分离核心与主要材料接触区 & 打击点、上下平衡、漂移及功率分配 \\
雪花（snowflake） & 理想二阶极向场零点具有六支局部结构；实际常用两个邻近普通 X 点逼近 & 对线圈误差敏感；高阶零点不能用普通鞍点判据概括 \\
X-divertor & 强调腿向靶区的磁通扩展，以展宽排热足迹 & 须比较整段 SOL 映射；与雪花及 XPT 名称不等价 \\
长腿 / Super-X & 延长路径；Super-X 还把靶区布置到更大 $R$，利用几何扩展及耗散空间 & 真空室体积、挡板、线圈与中性气体条件 \\
XPT & 在靶前/腿区利用次 X 点组织扩展与耗散，研究脱靶可控性 & 次 X 点位置、磁通差、支路连接及稳定性 \\
\bottomrule
\end{tabular}
\end{table}

上述设计要素有时能够组合，并非互斥标签。圆形、D 形、伸长和正/负三角形变则主要描述截面形状，可与不同边界拓扑搭配；L/H 模式和附着/脱靶又分别属于约束状态与边缘排热状态。三维仿星器还可用岛偏滤器，见第~\ref{sec:stellarator} 章，不能把轴对称主/次 X 点分类原样套用到所有三维连接结构。

\textbf{这些位形怎样由线圈产生？}线圈和等离子体电流共同形成 $F(R,Z)$，改变 PF 电流及剖面便可改变梯度零点、鞍点磁通和材料交点。设计工作先用自由边界平衡求出满足边界/X 点/壁间隙目标的电流组合，再检查线圈电流、电压、磁体受力与可制造性；转换时用有限电压驱动电流变化，还需考虑电感、涡流、剖面演化及垂直稳定性。普通鞍点与磁轴等临界点可以随参数移动、成对出现或消失，不能把``添加一个 X 点''理解成独立安放一个物体。更复杂的目标若超出线圈布局与电源能力，就未必存在可达平衡。

要回答``哪种位形更好''，必须指定工况并比较：核心约束和稳定性是否保留，靶面峰值及瞬态负荷是否降低，辐射/脱靶前沿能否反馈稳定，抽气与杂质控制是否有效，以及整个转换路径是否可达。磁平衡提供几何起点；SOL 流体或动理学、原子/中性粒子和材料热响应模型才把这一起点连接到排热性能预测。

\subsection{从物理过程到控制系统：估计、反馈与阶段切换}\label{sec:cycle-control}

控制器并不直接读到所有物理状态。磁探针和磁通环约束外部磁场与平衡重建，电流诊断约束 $I_p$；干涉仪常给出沿视线积分的密度，Thomson 散射等诊断提供密度与温度信息；辐射测量、可见光和红外成像补充启动、杂质及壁面热负荷信息。空间采样、标定误差和延迟都应进入状态估计，尤其不能把线平均密度误当作局域中心密度。资料 \cite{C7} 给出实际形状控制、平衡重建和诊断协作的例子。各阶段的控制目标与判据概括见表~\ref{tab:cycle-stages}。

\begin{table}[htbp]
\centering\small
\caption{放电阶段、控制目标与可观测的阶段判据（概念表）}\label{tab:cycle-stages}
\begin{tabular}{L{2.1cm}L{4.0cm}L{3.8cm}L{3.6cm}}
\toprule
阶段 & 主要调节量 & 主要测量/估计 & 需要确认的状态 \\
\midrule
气体与磁场准备 & 气阀、抽气、TF/CS/PF 电流 & 压强、线圈电流、真空场 & 粒子库与磁场进入预定窗口 \\
击穿与烧穿 & 感应电场、预电离、辅助加热与供气 & 电流增长、光辐射、密度 & 电离与升温可持续，而非只有瞬时发光 \\
建流与爬升 & 环电压、PF、加热与密度参考 & 电流、磁平衡、剖面估计 & 闭合区维持，电源和磁通仍有裕度 \\
平顶维持 & 加热、电流驱动、供气、排气、形状控制 & 热能、密度、辐射、形状与稳定性指标 & 源汇平衡且扰动可恢复 \\
位形转换 & 时变边界/X 点参考及协调线圈电压 & 边界、壁间隙、打击点、垂直运动 & 路径可达且约束未越界 \\
受控结束 & 降热、降流、供气与排气协调 & 热能、电流、位置、部件负荷 & 低能、低流阶段仍受控 \\
\bottomrule
\end{tabular}
\end{table}

数学上，这是一个带阶段变量 $m$ 的混合动力系统：
\begin{equation}\label{eq:cycle-hybrid}
 \dot{\bm x}=f_m(\bm x,\bm u,\bm d),\qquad
 \bm y=h_m(\bm x)+\bm\epsilon,\qquad
 m^+=\mathcal G(m,\hat{\bm x},\text{事件}).
\end{equation}
$\bm x$ 为连续物理状态，$\dot{\bm x}$ 为其变化率，$m$ 是当前离散阶段，$f_m$ 为该阶段的动力学；$h_m$ 为测量映射，$\bm\epsilon$ 为测量误差。$\hat{\bm x}$ 是根据诊断得到的状态估计，$\mathcal G$ 是阶段切换规则，$m^+$ 为切换后的阶段；``事件''包含诊断阈值触发、设备状态或保护信号。
$\bm u$ 包含线圈电压、气体流量和加热/电流驱动指令，$\bm d$ 包含壁面释放等扰动。上层安排阶段和参考轨迹，中层协调密度、热量及形状，下层电源跟踪电流并处理快速稳定控制。阶段切换通常需要状态条件、持续时间及迟滞，避免在阈值附近来回切换；时钟到了但烧穿未完成时，不应机械地进入高性能阶段。

对约束控制，一种抽象写法是
\begin{equation}\label{eq:cycle-objective}
 \min_{\bm u}\int\bigl[\|\bm y-\bm y_{\rm ref}\|_Q^2+
 \|\Delta\bm u\|_R^2\bigr]dt,
 \quad |V_k|\leq V_{k,\max},\quad |I_k|\leq I_{k,\max},\quad g_j(\bm x)\leq0.
\end{equation}
最小化变量为执行器指令轨迹 $\bm u(t)$；积分在选定预测或控制时段内累计代价。$\bm y_{\rm ref}$ 为输出参考，$\bm y-\bm y_{\rm ref}$ 为跟踪误差；$\Delta\bm u$ 表示相邻采样指令之差，用于惩罚过快动作。定义 $\|\bm e\|_Q^2=\bm e^{\mathsf T}Q\bm e$，其中 $Q$ 为半正定跟踪权重矩阵，动作项的 $R$ 为相应权重矩阵，二者需结合变量量纲选取。约束中 $k$ 标记执行器通道，$V_k,I_k$ 为其电压和电流，带 max 下标者为幅值上限；$j$ 标记状态约束，$g_j(\bm x)\leq0$ 表示处在所定义的允许区域内。
$g_j$ 可表示壁间隙、热负荷或经验证的稳定性限制。这里 $Q,R$ 是权重矩阵，与安全因子 $q$ 及回路电阻不同。该表达式可以组织 MPC 或轨迹优化，也能帮助设计 RL 的状态、动作和评价指标；但给越界状态一个负奖励不等于硬约束得到保证。学习策略的输出仍需经过可行性检查，并保留独立于策略的装置保护功能。

\subsection{安全关停：在失去约束之前移走热能和磁能}\label{sec:cycle-shutdown}

\textbf{正常结束不是直接切断所有电源。}放电结束时既有热能 $W$，也有等离子体电流和磁体储能。集总模型中 $W_{\rm mag}\simeq L_pI_p^2/2$；完整耦合回路的储能则是 $\bm I^{\mathsf T}\bm L\bm I/2$。这里 $W_{\rm mag}$ 为磁能（J），$\bm I$ 收集全部计入的回路电流，$\bm L$ 为含自感和互感的对称电感矩阵，$\mathsf T$ 表示转置；该式在线性、准静态回路近似下成立。若在电流尚大时失去位置与形状控制，剩余能量可能快速沉积到局部部件。因此，正常关停应沿着一条预先验证的低能量路径逐步退出。

首先降低目标热能或聚变功率，并协调燃料注入、辅助加热和排气。减少加热会降低压力、改变自举电流和磁平衡；在有约束模式跃迁时，热能和边界响应还可能突变。因此不能把降热看成只改动一个功率设定值。随后或同时按装置方案降低电流，必要时降低伸长、调整打击点，转入更易控制的终止位形；是否回到限制器位形取决于装置及场景，不是普遍必须执行的一步。

在简化纯感应回路中，乘式~\eqref{eq:cycle-ip} 以 $I_p$ 可得
\begin{equation}\label{eq:cycle-shutdownenergy}
 V_{\rm ext}I_p=R_pI_p^2+
 \frac{d}{dt}\left(\frac12L_pI_p^2\right)+\frac12I_p^2\dot L_p.
\end{equation}
各项均为功率（W）。左侧 $V_{\rm ext}I_p$ 是外部感应驱动向该等效回路输送的功率，符号为负时表示反向能量交换；$R_pI_p^2$ 是欧姆耗散；$d(L_pI_p^2/2)/dt$ 是这部分磁储能的变化率；$(I_p^2/2)\dot L_p$ 是电感随形状或剖面变化带来的附加项。这里点号仍表示时间导数，式中已取 $I_{\rm NI}=0$。
最后一项提醒读者：位形改变时不能把电感视作常数；完整能量核算还要纳入互感和力学功。电流降速受电源电压、内部电流扩散和稳定性共同限制。温度下降会提高电阻，若冷却过快，电流演化可能偏离预定路径；降流过慢又会持续消耗磁通和部件运行裕度。控制问题是协调这些时间尺度，而不是追求最快归零。

只有当等离子体电流及热能进入装置规定的终止范围、后续位置演化可接受时，才结束相应等离子体反馈。TF、CS 和 PF 磁体本身仍可能储能，后续退磁、能量回收或耗散应按各自电源和磁体保护序列进行。粒子排出、壁面冷却与诊断确认也可能持续到等离子体消失之后。

\textbf{异常中止与破裂缓解是另一条路径。}破裂通常涉及快速热能损失（热猝灭）和电流衰减（电流猝灭），可能伴随垂直位移、晕电流、电磁力及逃逸电子。电流快速变化产生的感应电场可能使部分电子持续加速；低密度条件下碰撞阻力不足，会增加逃逸问题的复杂性。因此``关掉加热即可安全终止''和``电流掉得越快越安全''都不是可靠结论。相关风险与研究方向见资料 \cite{C8,C9}。

异常处理可概括为检测与预测、可恢复时的回避、受控提前结束，以及无法回避时的专门缓解。大量物质注入等缓解手段旨在重新分配能量沉积并控制电流衰减，但其杂质种类、注入量、时机及对逃逸电子的影响必须通过装置验证；本章不把它们写成普适操作指令。HFM/MEQ 中的越界终止或仿真异常标志仅表示软件停止条件，不能视为已经模拟或验证了真实装置的安全关停。

\subsection{把全过程与本报告的仿真案例连接起来}\label{sec:cycle-scope}

后文实现案例可精细研究的是\textbf{已有等离子体平衡附近，线圈电压如何通过主动/被动回路与自由边界平衡耦合，改变电流、位置和形状}。它适合解释本章的磁控制链路和部分位形演化；载入初始平衡的 \code{reset} 不等于从中性气体物理击穿，输入电阻轨迹不等于自洽求解电子温度，指定压力剖面也不等于完成加热和输运预测。

若要将该案例扩展为全过程模拟，至少需要相互耦合的中性粒子与原子反应、启动非平衡电离、电子/离子能量输运、电流剖面扩散、加热与电流驱动沉积、壁面再循环及排热模型；聚变燃烧还需燃料组分、氦灰和快粒子模型。破裂与逃逸电子通常需要另行采用适合快速非平衡过程的模型。不同阶段应有清楚的状态交接和守恒检查，不能把启动模型中的粒子总数或能量在切换到平衡模型时无故重置。

建议初学者围绕五个问题复习：为什么恒定 CS 电流不能持续感应电压；为什么电离发光不等于烧穿；为什么电流平顶不等于非感应稳态；为什么 X 点的极向场为零而总场未必为零；为什么关停时仍需维持控制。能用本章的粒子、热量、磁通和回路方程回答这些问题，再阅读后文数值残差和控制接口，会更容易理解每一个状态变量的物理含义。

\subsection{理论来源与适用边界}\label{sec:cycle-sources}

本章引用的公开资料统一列在文末参考文献中，可从正文编号跳转，在线资料访问日期为 2026-10-24。本章的简化收支、量纲分析与控制表达式是为教学组织的模型，假设已在对应位置说明；不是对资料中某一装置参数的复现，也不是对当前仿真器功能的新增声明。

\clearpage

% ================= 仿星器：从线圈设计到装置性质 =================
\clearpage
\section{仿星器：如何用线圈设计磁场与装置性质}\label{sec:stellarator}
\markboth{仿星器与线圈设计}{仿星器与线圈设计}

\textbf{本章定位。}本章把前文的粒子、场与结构图景扩展到仿星器（stellarator），重点解释设计变量怎样影响物理性质。它不是某一装置的建造参数说明，也不声称轴对称 HFM/MEQ 案例能计算三维仿星器。下文公式为\evbg{}理论与设计表达式，均说明近似；实际优化结果必须经过三维平衡、轨道、输运、工程与实验验证。

\subsection{从托卡马克到仿星器：把旋转变换更多地交给外线圈}\label{sec:stellarator-concept}

两类装置都利用螺旋绕行的磁场组织轨道、减少粒子和热量损失。常规托卡马克的主要环向场来自外线圈，而形成约束所需的极向场有重要部分来自等离子体环向电流；仿星器则通过具有三维空间结构的外部线圈，在很小净等离子体电流、乃至真空场条件下就能建立所需的磁力线旋转变换。这里的``真空场''指只计外部线圈及相关结构的场，不是装置运行时必须没有等离子体\cite{Sreview}。

这使仿星器可以避免把持续的大等离子体电流及变压器磁通摆幅作为建立基本约束结构的必要条件，但并不意味着它没有电流、没有不稳定性或不需要主动控制。有限压力会产生平衡电流与自举电流，辅助电流驱动也可用于调整位形；加热、燃料、杂质、边界排热、超导磁体及装置保护仍需协调。因此，仿星器的稳态潜力来自其磁场生成方式，真正长时间运行还受加热、排热和工程条件限制。

\textbf{复杂线圈不是装饰，而是在编码沿环向变化的磁场。}在不同环向角度处，截面可以发生旋转、伸长、偏移，磁场强弱区也可以按所需方式排列。外部线圈同时决定磁力线怎样绕行和粒子沿途看到怎样的 $B=|\bm B|$；后者直接影响磁镜俘获、漂移和碰撞输运。设计目标因而不是单纯把磁面画得弯曲，而是控制三维场对粒子和流体的作用。

\subsection{线圈的正问题：给定形状与电流，怎样得到磁场}\label{sec:stellarator-forward}

先采用磁静态、细丝线圈、无铁磁非线性的近似。第 $k$ 条闭合线圈中心线为 $\bm r_k(u)$，其电流为 $I_k$，则
\begin{equation}\label{eq:stellarator-biot}
 \bm B_c(\bm r)=\frac{\mu_0}{4\pi}\sum_k I_k
 \int_0^{2\pi}\frac{\bm r'_k(u)\times[\bm r-\bm r_k(u)]}
 {|\bm r-\bm r_k(u)|^3}\,du.
\end{equation}
$\bm r$ 为求场位置（m），$u$ 是无量纲线圈参数而非时间；撇号表示对 $u$ 求导，$\bm r'_kdu$ 是有向线元。$I_k$ 以 A 计，多匝绕组应逐匝求和或在适用时折合为等效安匝；$\mu_0$ 是真空磁导率，$\bm B_c$ 为线圈场（T）。分子描述电流元与观察方向的叉乘，分母体现距离依赖。式子说明改变线圈位置、局部切向、形状或电流，都能改变等离子体所在处的场。

一种可优化的几何参数化是
\begin{equation}\label{eq:stellarator-coil-fourier}
 \bm r_k(u)=\bm a_{k0}+\sum_{\ell=1}^{L_k}
 [\bm a_{k\ell}\cos(\ell u)+\bm b_{k\ell}\sin(\ell u)].
\end{equation}
$\bm a,\bm b$ 是具有长度单位的三分量 Fourier 系数，$\ell$ 为谐波序号，$L_k$ 是截断阶数；所有系数与电流共同构成设计向量。非平面中心线允许三维谐波自由度，平面圆线圈只是其中受限制的特例。增加阶数可表达更多形状，却也可能带来过大曲率、过近线圈和制造敏感性，故设计中不能只追求拟合精度。

在无电流的等离子体预定区域，线圈真空场满足 $\nabla\times\bm B_c=\bm0$ 与 $\nabla\cdot\bm B_c=0$，但这不要求整个环形区域的场线都像直线，也不禁止外部三维线圈产生旋转变换；装置的非单连通拓扑和外部电流边界必须一起考虑。有限压力运行时，需再与等离子体产生的场叠加，联立求解 $\bm j\times\bm B=\nabla p$ 等条件，不能把真空设计直接当作所有压力下不变的运行磁场。

\subsection{先看场线：旋转变换、磁剪切、磁面和磁岛}\label{sec:stellarator-topology}

采用适当的磁面角坐标，磁力线平均绕行可用旋转变换描述：
\begin{equation}\label{eq:stellarator-iota}
 \iota(s)=\lim_{\Delta\phi\to\infty}\frac{\Delta\theta}{\Delta\phi},
 \qquad q=\frac1\iota.
\end{equation}
$s$ 为磁面标签，常取归一化环向磁通；$\theta,\phi$ 为极向与环向角，$\Delta$ 表示沿同一磁力线累计的角增量。$\iota$ 为每绕环向一周平均绕极向的周数，无量纲，$q$ 采用与前文一致的倒数约定，方向和符号须同时约定。这里的角速度不是粒子回旋频率。$d\iota/ds$ 描述旋转变换随磁面的变化，是磁剪切的一种度量；换用其他径向变量时数值会改变，须说明定义。

线圈的倾斜、非平面绕组及其沿环向的排列，使不同位置的极向场分量和截面响应相互配合，产生真空旋转变换。磁轴扭转与截面形变也参与几何解释，但仅凭一幅旋转椭圆图不能算出 $\iota$：必须追踪由线圈实际产生的场。若 $\iota=n/m$，其中 $m,n$ 为全环极向/环向整数模数，相应有理磁面可能对某类误差场产生共振，出现磁岛；这与前文 $q=m/n$ 的约定一致。

设计时通常追踪许多初始位置的磁力线，并记录它们反复穿过同一环向截面的交点，形成 Poincar\'e 图。平滑闭合曲线对应嵌套磁面截面，岛链对应共振结构，散布的交点则可能提示混沌区域；有限追踪长度与数值误差也需检查。某些边缘磁岛可服务于偏滤器排热，而核心过大的岛或混沌区可能损害约束。因此，目标不是把所有磁岛一律消掉，而是把拓扑结构放在合适位置并控制其尺寸。

\subsection{再看场强：准对称为什么能影响粒子约束}\label{sec:stellarator-qs}

磁力线不离开良好磁面，并不保证导引中心也永远不跨面漂移。粒子沿场经历的强弱变化、磁镜俘获、曲率及梯度漂移共同决定其轨道。优化者因此会在专门的 Boozer 磁坐标中分析场强谱，而不只是看装置外形：
\begin{equation}\label{eq:stellarator-b-spectrum}
 B(s,\theta_B,\phi_B)=\sum_{m,n}B_{mn}(s)
 \cos(m\theta_B-nN_{\rm fp}\phi_B).
\end{equation}
$B$ 为磁场大小（T），$B_{mn}$ 为场强谱系数，$m,n$ 为整数谐波编号，$N_{\rm fp}$ 为全环重复的磁场周期数。$\theta_B,\phi_B$ 是 Boozer 角坐标，与普通几何极角不必相同；该坐标选择使场强对导引中心轨道的对称性有明确意义。此处为方便采用仿星对称且适当选取角原点的余弦展开；一般场还需相应正弦项。这里 $nN_{\rm fp}$ 才是对应的全环环向模数，不能与上节的全环 $n$ 混用。

准对称（quasisymmetry，QS）的理想目标是场强主要依赖单一螺旋角：
\begin{equation}\label{eq:stellarator-qs}
 B\simeq\mathcal B(s,M\theta_B-N\phi_B).
\end{equation}
$\mathcal B$ 为所需场强函数，$M,N$ 为指定整数，$N$ 按全环角度定义并应与周期性相容。准轴对称（QA）取 $N=0$，使场强近似不依赖 $\phi_B$；准螺旋对称（QH）取非零的 $M,N$，使场强主要依赖一个螺旋角。\textbf{这里对称的是特定坐标中的场强，不是要求线圈、磁面形状或整个磁场矢量具有相同的空间连续对称。}因此三维线圈可以产生 QA 场，而装置外形仍不是轴对称的。

当这一场强对称性足够好且导引中心近似适用时，粒子运动具有类似连续对称带来的约束，有助于减小某些轨道的长期径向漂移和新经典损失。实际有限体积、有限压力及可制造线圈只能以规定误差去实现目标，需同时检查快粒子轨道和输运。已有研究分别展示了高精度准对称目标场和能够近似实现该场的线圈设计；它们属于理论/计算设计证据，不能直接等同于建成装置的实验性能\cite{Sprecise,Scoils}。

\begin{figure}[htbp]
\centering
\begin{tikzpicture}[>=Stealth,font=\small]
 \begin{scope}
  \node[font=\bfseries] at (2,3.4) {QA：$B\simeq\mathcal B(s,\theta_B)$};
  \fill[blue!4] (0,0) rectangle (4,2.7);
  \foreach \y in {0.45,0.9,1.35,1.8,2.25}{\draw[blue!70!black,thick] (0,\y)--(4,\y);}
  \draw[->] (0,0)--(4.4,0) node[right] {$\phi_B$};
  \draw[->] (0,0)--(0,3.0) node[left] {$\theta_B$};
  \draw[dashed,black!35] (0,2.7)--(4,2.7)--(4,0);
  \node[below,font=\footnotesize] at (2,-0.3) {同一 $s$ 面的角坐标展开};
 \end{scope}
 \begin{scope}[xshift=7.3cm]
  \node[font=\bfseries] at (2,3.4) {QH：$B\simeq\mathcal B(s,\theta_B-\phi_B)$};
  \fill[blue!4] (0,0) rectangle (4,2.7);
  \begin{scope}
   \clip (0,0) rectangle (4,2.7);
   \foreach \b in {-2.25,-1.8,-1.35,-0.9,-0.45,0,0.45,0.9,1.35,1.8,2.25}{
    \draw[blue!70!black,thick] (0,\b)--(4,{\b+2.7});
   }
  \end{scope}
  \draw[->] (0,0)--(4.4,0) node[right] {$\phi_B$};
  \draw[->] (0,0)--(0,3.0) node[left] {$\theta_B$};
  \draw[dashed,black!35] (0,2.7)--(4,2.7)--(4,0);
  \node[below,font=\footnotesize] at (2,-0.3) {示意取 $M=N=1$，非装置参数};
 \end{scope}
\end{tikzpicture}
\caption{单个磁面上等场强轮廓的示意。蓝线表示 $B$ 的等值线，不是磁力线，也不是线圈。坐标边界按角度周期接续；示意只解释对称性形式，不能独立证明磁面、力平衡或线圈可实现性。}\label{fig:stellarator-qs}
\end{figure}

\subsection{准等动力与全向性：不必只有准对称一条路线}\label{sec:stellarator-qi}

更一般的目标是让俘获粒子的\emph{反弹平均径向漂移}足够小。全向性（omnigenity）可用第二绝热不变量对场线标签的独立性表达。在静态、嵌套磁面和导引中心近似下，设静电势近似为磁通函数，
\begin{equation}\label{eq:stellarator-jparallel}
 \begin{aligned}
 J_\parallel(s,\alpha,\mathcal E,\mu)
 &=2\int_{\ell_1}^{\ell_2}
 \sqrt{2m[\mathcal E-\mu B(s,\alpha,\ell)-q_s\Phi(s)]}\,d\ell,\\
 \left.\frac{\partial J_\parallel}{\partial\alpha}\right|_{s,\mathcal E,\mu}&=0.
 \end{aligned}
\end{equation}
$J_\parallel$ 为一次完整往返的平行动量积分（动量乘长度），$\ell$ 为沿场弧长；$\ell_1,\ell_2$ 为平行速度降至零的两个反射点，系数 2 表示往返。$m,q_s$ 为粒子质量和带符号电荷，$\mathcal E$ 是含静电势能的总能量，$\mu$ 是磁矩，$\Phi$ 为静电势（V）；$\mu B$ 为垂直动能。$\alpha$ 是同一磁面上区分磁力线的标签，局部可用 $\theta_B-\iota\phi_B$。第二式要求在固定 $s,\mathcal E,\mu$ 下，不同场线上的反弹积分相同；在这些近似下，它对应俘获轨道反弹平均不向外漂移的条件\cite{Sreview}。

准等动力（quasi-isodynamic，QI）是具有特定等场强轮廓拓扑的一类近全向性目标；其理想化等 $B$ 轮廓在 Boozer 角空间中沿极向闭合。它不等同于 QA 或 QH，也不等于要求磁场大小处处相同。设计关心的是强弱场井的沿场分布如何使不同俘获轨道的漂移在一个反弹周期中抵消。准对称是实现良好轨道的一条有力路线，却不是唯一的优化目标。

W7-X 的优化涉及接近 QI 的约束性质、较小的新经典输运、有限压力平衡和较小自举电流等目标；它不是一个要求 QA 的装置。自举电流小的价值之一是减小运行中 $\iota$ 与边缘岛位置随电流演化发生的偏移，使岛偏滤器的打击区更可控\cite{Sw7x,Sbootstrap}。这也说明``电流越小越好''不是独立口号，而须与目标拓扑、压力和排热方案一起评价。

\subsection{把设计指标翻译成可解释的装置性质}\label{sec:stellarator-properties}

线圈形状对性质的作用通常经过一串中间变量，而不是一个几何参数只对应一种性能。例如移动某段线圈可能同时改变法向场误差、$\iota$、场强谐波和部件间距；减小某一谱分量也未必同时改善湍流、快粒子和工程条件。设计需把目标拆成可计算、可验证的指标。

例如，某段场线上 $B_{\max}/B_{\min}$ 衡量强弱场变化，$B_{\max},B_{\min}$ 分别为所考察区段的最大与最小场强；该比值影响磁镜俘获，却不足以描述沿程井形、不同场线间的反弹积分及漂移。MHD 稳定性中所说的有利磁井则涉及磁面体积、曲率与压力扰动的整体性质，不能用这个局部镜比替代。所谓有效波纹也是对特定新经典输运效应的综合度量，而不是任意一个 $B_{mn}$ 的大小。

\begin{table}[htbp]
\centering\small
\caption{线圈设计目标、磁场结构与装置性质之间的联系}\label{tab:stellarator-properties}
\begin{tabular}{L{2.3cm}L{5.2cm}L{6.0cm}}
\toprule 关注性质 & 常用磁场/模型诊断 & 设计含义及需要同时检查的因素 \\
\midrule
良好磁面体积 & 法向场误差、场线追踪、岛宽与混沌区 & 实现目标边界并控制共振误差；边界拟合好不保证内部处处可积 \\
旋转变换与剪切 & $\iota(s)$、有理面位置、有限压力及电流响应 & 选择拓扑和稳定性工作区；不可只比较轴上一个 $\iota$ 值 \\
低新经典输运 & QS/QI 误差、有效波纹及动理学通量 & 改善俘获粒子反弹平均漂移；还依赖碰撞率、径向电场与剖面 \\
快粒子约束 & 导引中心/轨道跟踪、损失位置、能量沉积 & 关注有限轨道宽度与直接损失，不仅是热粒子输运指标 \\
有限压力与稳定性 & 三维力平衡、磁井/曲率及扰动增长 & 检查目标压力下是否有可用平衡；局部磁镜井不等同于 MHD 磁井判据 \\
排热与运行鲁棒性 & 边缘岛、连接长度、打击图、误差敏感度 & 在有意设置的边缘结构上组织排热；还要与靶板和辐射输运匹配 \\
\bottomrule
\end{tabular}
\end{table}

在一些低碰撞率区间，三维场的俘获粒子新经典输运可能很强，减小有效波纹或改善全向性因而有重要作用；但较低新经典损失不自动意味着总能量约束最好。梯度驱动湍流、径向电场、密度与温度剖面仍可决定实际损失。W7-X 已有针对降低新经典输运的实验研究支持其优化思路，但这一结论具有所分析的运行条件，不应推广成所有工况下所有损失都已消除\cite{Sneo}。

有限压力下还须同时考虑平衡位移和电流响应。用于构造三维嵌套磁面平衡的程序，可能把磁面的存在作为表示的一部分；若要检查磁岛或混沌，应追加不强制同一拓扑的场线/平衡分析。只观察一个成功收敛的平衡文件，既不能证明全部稳定性，也不能代替真实线圈、误差及压力扫描后的验证。

\subsection{逆问题：从目标性质求线圈，而不是只把线圈画弯}\label{sec:stellarator-inverse}

\subsubsection{两阶段设计：先优化等离子体，再拟合外部场}

一种传统路线先参数化等离子体边界、磁轴和压力/电流剖面，用三维平衡与轨道/输运计算寻找所需性质；再寻找能近似实现这一边界磁场的线圈。第一阶段的最优边界可能需要非常复杂或非常靠近等离子体的电流分布，所以物理阶段得到的好场，不一定有简单可造的线圈。

对给定目标磁面 $S$，磁场应沿表面切向，可用归一化法向场误差作为线圈拟合指标：
\begin{equation}\label{eq:stellarator-normal}
 \mathcal F_B=\frac1{A_SB_{\rm ref}^2}
 \int_S[(\bm B_c+\bm B_{\rm pl})\cdot\hat{\bm n}]^2\,dS.
\end{equation}
$A_S$ 为目标面面积，$B_{\rm ref}$ 为正的参考场强；$\hat{\bm n}$ 是单位法向量，$dS$ 为面积元。此处 $\bm B_c$ 为线圈场，$\bm B_{\rm pl}$ \emph{专指等离子体产生的场}，不是前文的极向分量；真空设计取 $\bm B_{\rm pl}=\bm0$。积分测量穿出目标磁面的场强误差，除以面积和场强平方后无量纲。用总场写式子，是为了在有限压力设计中避免遗漏等离子体场；优化过程中若边界或剖面变化，$\bm B_{\rm pl}$ 也应相应更新。

$\mathcal F_B$ 小只说明对所给表面的近似切向条件好，不能单独确定 QS、输运和稳定性。还须约束所需环向/极向磁通或场强尺度，避免零电流、零场之类的无意义解；并重新追踪实际线圈场中的内部磁面和边缘结构。

\subsubsection{绕组面与电流势：先解连续电流，再离散成线圈}

另一种常见表示是在等离子体外选一个绕组面，寻找面电流，再取电流流线形成离散绕组。对面内无散电流，可写
\begin{equation}\label{eq:stellarator-current-potential}
 \bm K=\hat{\bm n}_c\times\nabla_s\Phi_c.
\end{equation}
$\bm K$ 为面电流密度（A/m），$\hat{\bm n}_c$ 为绕组面单位法向量，$\nabla_s$ 为沿表面的梯度，$\Phi_c$ 为电流势（A），不是静电势。于是 $\bm K\cdot\nabla_s\Phi_c=0$，面电流沿 $\Phi_c$ 的等值线流动；在环面上若要表示净链接电流，还须包含合适的多值势部分。选取一定间隔的等势线可以得到候选离散绕组，但其匝数、电流、截面和连接方式仍需工程处理。

REGCOIL 类方法在拟合误差之外，惩罚绕组面上过强或过集中的电流，使连续逆问题得到正则化。这样可在场误差与绕组复杂性之间做可控取舍，再检查离散后线圈的间距和曲率；正则化不能保证任意目标都有可制造解，也不能取消离散化后的再验证\cite{Sregcoil}。

\subsubsection{为什么离等离子体越远，细节控制往往越困难}

局部平面、无电流区域可用 Laplace 方程理解外部场空间细节的衰减。一个表面波数为 $k_t$ 的谐波，离源面距离 $d$ 后有量级关系
\begin{equation}\label{eq:stellarator-distance}
 |\delta B_{k_t}(d)|\propto|\delta B_{k_t}(0)|e^{-|k_t|d}.
\end{equation}
$k_t$ 为沿源面的空间波数（$\mathrm{m^{-1}}$），$d$ 为距源面的法向距离（m），$\delta B_{k_t}$ 为该谐波的场强幅值。该式是局部线性真空场的说明，不是任意环面几何的精确传递函数。它表明较短波长的细节在远处衰减更快：若想用很远、很平滑的线圈实现很精细的边界场，常需更大的电流或更复杂的绕组，逆问题也更敏感。

这构成一个具体的物理--工程矛盾：线圈靠近可以改善高空间频率的控制能力，却会挤占真空室、屏蔽、包层、诊断和维护空间；把线圈移远有利于通道与维护，却可能降低实现同一目标场的能力。因此线圈--等离子体间距不是画图后才添加的外形限制，而应参与目标场与线圈的共同设计。

\subsection{把物理与工程一起优化：目标函数和不可替代的约束}\label{sec:stellarator-optimization}

可将候选边界参数、线圈 Fourier 系数及电流收集为 $\bm c$，写一个概念性的多目标问题：
\begin{equation}\label{eq:stellarator-objective}
 \min_{\bm c}\ \mathcal J
 =w_B\mathcal F_B+w_{\rm QS}\mathcal E_{\rm QS}
 +w_{\rm orb}\mathcal L_{\rm orb}+w_{\rm eng}\mathcal C_{\rm eng}.
\end{equation}
$w$ 为非负权重；$\mathcal F_B$ 为法向误差，$\mathcal E_{\rm QS}$ 为指定 QS 目标的归一化偏离，$\mathcal L_{\rm orb}$ 为某组粒子轨道的损失或漂移指标，$\mathcal C_{\rm eng}$ 为归一化工程复杂度。各项须先明确量纲、采样面和所代表工况；QI 设计可改用全向性或反弹积分相关误差，不能把 QS 项原封不动当作唯一目标。轨道指标通常在特定能量、俯仰角和出生位置采样，不能从一个采样集合的低损失推定所有粒子都无损失。

典型几何量可从线圈中心线直接计算：
\begin{equation}\label{eq:stellarator-geometry}
 \mathcal L_k=\int_0^{2\pi}|\bm r'_k|\,du,\qquad
 \kappa_k(u)=\frac{|\bm r'_k\times\bm r''_k|}{|\bm r'_k|^3}.
\end{equation}
$\mathcal L_k$ 为线圈长度（m），$\kappa_k$ 为曲率（$\mathrm{m^{-1}}$），双撇表示对几何参数 $u$ 求二阶导；最小弯曲半径为最大曲率的倒数。线圈长度影响材料和电阻/电感，过大曲率不利于绕组和导体；还需施加线圈--线圈最小间距、线圈--部件间隙、有限绕组截面以及维护开口等约束。单纯把这些项加入带有限权重的惩罚函数，不等于它们一定满足，应逐项检查最终可行性。

电磁受力也参与设计。在细丝互作用近似下，某线元所受来自其他电流源的力为 $d\bm F=I_k\,d\bm l\times\bm B_{\rm other}$，其中 $d\bm l$ 为有向线元，$\bm B_{\rm other}$ 为其他源在该处的磁场；力以 N 计。计算绕组自身的受力时，不能将奇异细丝自场直接代入，需使用与导体有限截面一致的模型。支撑结构、局部应力、最大导体场、超导电流裕度和冷却能力，都不能只用中心线曲率代表。

单阶段或联合优化把线圈与平衡变量一并调整，可避免两阶段流程中``目标场很好但难以造出线圈''的一部分问题；也增加了计算代价和多个目标之间的耦合。已有研究同时考虑真空和有限压力情形，说明这是一种设计方法，而不是保证所有目标可兼得的定理\cite{Ssingle}。

\begin{figure}[htbp]
\centering
\begin{tikzpicture}[>=Stealth,font=\small,
 st/.style={draw,rounded corners=3pt,align=center,text width=3.5cm,minimum height=1.1cm}]
 \node[st,fill=blue!5] (goal) at (0,0) {目标性质\\输运、稳定性、排热};
 \node[st,fill=blue!5] (field) at (5,0) {三维场与等离子体\\平衡、谱、$\iota$、轨道};
 \node[st,fill=orange!12] (coil) at (10,0) {线圈与电流\\形状、间距、受力};
 \draw[->,thick] (goal)--(field);
 \draw[->,thick] (field)--(coil);
 \node[draw,align=center,text width=9.0cm,fill=black!3] (test) at (5,-2.45)
 {制造/装配误差、有限压力和多运行场景扫描\\诊断与实验比对};
 \draw[->,thick] (coil.south)--(10,-2.45)--(test.east);
 \draw[->,thick] (test.north)--node[right,font=\footnotesize]{重算与反馈}(field.south);
\end{tikzpicture}
\caption{线圈设计的闭合验证链。两阶段路线可分别优化目标场和线圈；联合路线同时调整两者。只有用实际候选线圈重算之后仍满足的性质，才是该线圈方案所支持的性质。}\label{fig:stellarator-design}
\end{figure}

\subsection{制造误差与校正线圈：设计性能能否落到实物上}\label{sec:stellarator-robustness}

实际绕组与设计中心线不完全相同，整只线圈还可能有位移、倾斜、电流误差与受力变形。重要的不是仅给每个坐标一个统一容差，而是判断哪些误差模式会破坏关键磁面或产生不利的共振谐波。对小扰动，可用局部灵敏度表达
\begin{equation}\label{eq:stellarator-sensitivity}
 \delta\bm y\simeq G_c\,\delta\bm c_g+G_I\,\delta\bm I.
\end{equation}
$\bm y$ 为所选输出，例如法向谱系数、岛宽的局部代理量或打击位置；$\bm c_g$ 为几何参数，$\bm I$ 为各回路电流，$G_c,G_I$ 分别是相应导数矩阵，单位随输入和输出定义。该线性近似只适用于足够小且未跨越拓扑突变的扰动；直接岛宽等非光滑量附近尤其应谨慎使用。

鲁棒优化可在多个相关误差样本上控制平均性能和较差情形，而不只优化名义形状。样本应体现整只线圈的刚性偏移、装配相关性及局部绕组误差，而非把每个点当作完全独立噪声。较低的名义误差若伴随极高敏感度，未必是更好的实物设计。绕组面优化与灵敏度分析也可用于理解这些工程取舍\cite{Sadjoint}。

安装后可通过真空场测量、线圈校准和场线映射识别误差，再利用修正/调整线圈补偿部分分量。校正执行器只能产生其可达的场组合，不能无代价地恢复任意几何错误；严重挤占维护空间或超过导体受力的设计问题，更不能靠运行反馈弥补。

\subsection{三个例子：线圈架构与优化目标并不是同一个分类}\label{sec:stellarator-examples}

\begin{table}[htbp]
\centering\small
\caption{装置例子及其线圈--物理设计重点}\label{tab:stellarator-examples}
\begin{tabular}{L{1.7cm}L{5.4cm}L{6.4cm}}
\toprule 例子 & 线圈架构或优化重点 & 从设计中理解的性质 \\
\midrule
LHD & 两条连续螺旋线圈与三对圆形极向场线圈的超导磁体方案 & 用连续螺旋绕组建立三维约束场，极向场及分层绕组电流提供一定几何调节 \\
HSX & 模块化三维线圈，围绕准螺旋对称开展研究 & 使粒子所见场强接近单一螺旋角依赖，并可比较偏离该性质的配置 \\
W7-X & 五个场周期、50 个非平面与20 个平面超导线圈的方案 & 围绕优化的三维约束、较低新经典损失、低自举电流及岛偏滤器协调设计 \\
\bottomrule
\end{tabular}
\end{table}

LHD 说明连续螺旋绕组是一种磁场生成架构，官方资料也描述了利用不同线圈层电流改变有效线圈半径的设计灵活性\cite{Slhd,Slhdvariation}。这不意味着每次放电都能机械改变绕组形状，而是固定几何中存在不同电流组合。

HSX 说明准螺旋对称是一种\emph{场的优化性质}，可以由三维模块线圈实现；模块化并不自动意味着 QH，亦不保证所有模块化装置都有同样输运\cite{Shsx}。同理，QA 是场强目标而不是``把全部线圈做成平面圆环''的别名。线圈架构回答怎么制造与布置电流，QA/QH/QI 等目标回答希望这些电流产生什么样的粒子约束性质。

W7-X 则说明同一优化装置需要同时兼顾多项物理目标和工程约束。其不同线圈组的电流可以调整磁场配置；有限压力、自举电流和加热剖面还会改变运行中的实际结构\cite{Sw7x,Sw7xpower}。这些例子强调线圈架构与目标性质的对应关系。

\subsection{线圈造好以后，还能怎样改变装置性质}\label{sec:stellarator-operation}

应区分\textbf{设计自由度}和\textbf{运行自由度}。建造前可以改变线圈数目、中心线、间距及电源分组；建造后通常只能在硬件允许的范围内调整回路电流、加热与燃料等输入，几何设计限定了其可达配置集合。某些设计专门针对多个 $\iota$ 或磁面目标联合优化线圈与控制线圈，才能在同一套硬件上获得较好的多配置灵活性\cite{Smulticonfig}。

一个有用的解析例子是同比例提高所有线圈电流。在忽略等离子体和材料非线性、保持几何不变的真空场近似下，
\begin{equation}\label{eq:stellarator-current-scaling}
 I_k\longrightarrow aI_k\quad\Rightarrow\quad
 \bm B_c\longrightarrow a\bm B_c,\qquad a>0.
\end{equation}
$a$ 为无量纲倍率，所有回路采用同一倍率。因磁场方向场不变，真空磁力线的几何、$\iota$ 及归一化谱形在该近似下不变；但场强升高，对同一速度粒子的回旋半径减小，磁压力增大。若电流比改变，场方向和相对谐波通常也会改变，可能调整 $\iota$、磁井与边缘岛位置。因此，提高主场强度与改变位形是不同操作。

上述同比例结论不等于真实等离子体性质全部不变：固定压强而增大场强会降低 $\beta$，改变有限压力响应；最大导体场、线圈受力与电源限值也会变化。加热和供气又会改变密度、温度、碰撞率、自举电流及径向电场，进一步影响三维输运与排热。设计阶段应提前扫描这些场景，运行阶段再用状态估计与反馈管理其演化。

仿星器通常可利用微波等辅助方式建立与加热等离子体，无须沿用托卡马克的感应建流作为基本启动路径；但工作气体、电离、能量收支与稳定性仍要满足相应条件。将第~\ref{sec:lifecycle} 章的放电逻辑推广到仿星器时，应保留这些守恒与控制问题，同时重写线圈分工、启动方式和三维边界模型，而不是把名称替换后原样使用。

\subsection{从线圈方案到可信装置结论的验证顺序}\label{sec:stellarator-validation}

一套可复核的设计分析至少应形成以下闭环：
\begin{enumerate}[nosep]
 \item 明确目标工作区、场周期、磁通/场强尺度及要优化的输运、稳定性和排热指标；
 \item 用目标边界或直接线圈场构造三维磁面，检查 $\iota$、谱结构以及有理面附近拓扑；
 \item 在有限压力、多个剖面与电流场景下重算平衡，跟踪热粒子和快粒子，并计算所需输运；
 \item 对离散线圈、有限绕组截面、部件空间和误差样本逐项检查，补充受力与热/超导分析；
 \item 利用可测磁场、Poincar\'e 结构、剖面和热负荷等观测，把计算性质与实际配置核对。
\end{enumerate}

真空场追踪可揭示线圈本身产生的磁面和岛；三维平衡加入压力与电流响应；导引中心和动理学研究粒子约束；MHD 与回旋动理学分别研究合适尺度的稳定性、涨落与输运；工程模型再决定能否制造并持续运行。任何一层的优化成功都只覆盖其保留的物理。尤其不能把轴对称 GS 求解器不含的三维谐波或 QI 性质，从当前案例文件中间接``读出''。

本章最值得建立的设计直觉是：\textbf{线圈参数先改变磁场的几何与强弱结构，再通过轨道、输运和力平衡改变装置性质；工程约束与实际运行状态又反过来限定这条链。}同一项性质可以有不同线圈实现，同一套线圈也可以支持有限范围的多种配置，所以仿星器设计本质上是受物理与工程共同约束的多目标逆问题。

\clearpage
% ================= 仿真实现案例 =================
\section{引言与仿真器总览}\label{sec:intro}

\subsection{物理范围及模型层次}\label{sec:legacy-1-1}

\textbf{MEQ（Matlab Equilibrium Toolbox，MATLAB 平衡计算工具箱）}是本项目使用的托卡马克磁平衡计算软件库，包含静态正问题、平衡重建、逆平衡及平衡演化等功能。``MEQ''指整套工具箱，而不是某一个方程或单个时间步算法。其名称及功能范围由 \code{Q/README.md:1--9,205--215} 明确给出。

\textbf{FGE（Forward Grad--Shafranov Evolution，正向 Grad--Shafranov 演化求解器）}是 MEQ 中本项目使用的演化求解模块。这里按源码入口的自述解释名称：给定初始平衡、线圈电压及剖面约束等输入，联立计算后续磁平衡和主动/被动导体电流。它与只计算给定时刻静态平衡的 FGS（Forward Grad--Shafranov）有区别。证据：\code{G:meq/fget.m:1--3}、\code{Q/README.md:213--215}；具体耦合残差见 \code{G:meq/fgeF.m:362--482}。

\textbf{GS 平衡（Grad--Shafranov equilibrium）}指轴对称等离子体的磁流体力学平衡：在标量压力、忽略流动惯性等理想化条件下，压力梯度与电流受到的磁力平衡；利用轴对称性，可将其写成二维极向磁通的椭圆型方程。这里的``平衡''表示每一时刻磁场与给定剖面满足力平衡，并不表示整个放电过程不随时间改变。\evbg{} 这是 GS 概念的理论解释；本程序实际通过网格磁通、有限剖面基、约束及回路方程实现这一类模型，没有另行积分完整三维 MHD 力学方程。它的磁通约定和对应连续式见第~\ref{sec:gse}、\ref{sec:fullsystem} 节，实际代数残差见 \code{G:meq/fgeF.m:280--288,362--482}。

\evimpl{} MEQ 的 FGE 是``轴对称自由边界 Grad--Shafranov 平衡 + 主动/被动回路演化 + 电流演化闭合''的耦合求解器。\code{G:meq/fget.m:1--3} 明确其目标；\code{G:meq/fgeF.m:365--482} 实际拼接平衡、剖面约束、CDE 和导体回路残差。它不通过三维流体动量、质量和能量方程推进整个 MHD 状态。当前可见主链没有密度输运、电子/离子能量输运或湍流方程；不能把该模型描述为完整时变 MHD 仿真。

\evdata/\evimpl{} 五份 \code{ini_L_*.mat} 均保存 \code{L.P.algoNL='all-nl'}、\code{L.P.cde='OhmTor_rigid_0D'}、\code{L.P.lin=0}、\code{L.P.ivacuum=0}、\code{L.P.iLpext=1}。因此这些配置选择非线性自由边界平衡、导体耦合及\textbf{零维总等离子体电流}方程，后者使用输入等离子体电感 \code{LX.Lp}。库还包含一维弱式电流演化、线性化、Picard 和不同平衡未知量分支，不能混同为同一条默认流程。选择与方程见第~\ref{sec:model}、\ref{sec:solver} 节及 \code{G:meq/meqcdec.m:38--72}。

各物理过程在可见实现中的性质概括如表~\ref{tab:physproc}。

\begin{table}[htbp]
\centering\small
\caption{物理过程、可见实现的性质与关键证据}
\label{tab:physproc}
\begin{tabular}{L{2.6cm}L{4.6cm}L{6.2cm}}
\toprule
物理过程 & 可见实现的性质 & 关键证据与限制 \\
\midrule
极向磁平衡 & 二维轴对称非线性平衡；网格磁通与有限剖面系数耦合 & \code{G:meq/fgeF.m:198--288,365--443} \\
主动/被动导体 & 预计算互感与电阻的集总回路方程 & \code{G:meq/fgeF.m:460--482}；\code{D/L.Mee,Mey,Re} \\
等离子体电流演化 & 当前配置为刚性零维闭合；不是全径向电流扩散 & \code{G:meq/meqcdefun.m:59--163}；\code{D/L.P.cde,L.P.iLpext} \\
压力与环向场剖面 & 磁通函数基展开及约束，当前为三个 \code{bfab} 系数 & \code{Q/libmeq/bfab.c:25--39,169--185}；\code{D/L.bfp=[1,2]} \\
电导率/电阻 & 用户电导率幅值乘给定径向形状，经磁面积分得到输入电阻轨迹 & \code{V/my_fusion_environment.m:107--156}；\code{V/evalSQ_cp.m:11--32} \\
磁测量与形状输出 & 互感/场响应矩阵、磁轴/X 点/边界搜索及后处理 & \code{Q/documentation/CHEATSHEET.md:27--108}；\code{V/my_fusion_environment.m:249--428} \\
控制 & 外部 12 路电压；可选增量 PID 接管 VS & \code{V/FGE_state_env_std.m:222--277} \\
电源非理想性 & 内置饱和/FIFO 与项目外经验仿射模型；RL 普通步有旁路差异 & \code{G:meq/meqps.m:14--53}；\code{V/iterFGE.m:53--55}；\code{E/power_supply.py:109--161} \\
三维不稳定性、输运、湍流 & 当前主链未实现这些完整 PDE & 不应从``Tokamak/MHD''名称补写 \\
\bottomrule
\end{tabular}
\end{table}

\subsection{目录结构与软件分层}\label{sec:layertree}

\evimpl/文件清点 \code{HFM_env} 工作树含 695 个 \code{.p}、156 个 \code{.c}、8 个 \code{.h}、98 个 \code{.mexa64} 和 38 个 \code{.mat}；没有 \code{.m}，但有两份 \code{.asv} 备份。此清点不包含 \code{matlab/} 和 \code{sysroot/}，后两者是体积很大的第三方运行环境，不应当作 HFM 的物理源文件。\code{socket_serve/} 只有 shell 启动器和一个 \code{.p}。目录结构如下：
\begin{lstlisting}[style=pseudo, basicstyle=\scriptsize\ttfamily]
R/
|- run_hfm_server.sh                 单实例入口
|- start_hfm_pool.sh / stop_hfm_pool.sh
|- socket_serve/
|  |- hfm_server.sh                  MATLAB -batch 入口
|  `- hfm_socket_server.p            不可见的协议服务端
|- HFM_env/
|  |- .git/                          本报告读取的 MATLAB 源码快照
|  |- ADD_PATH.p
|  |- TOK_EXL50U/XL50U_fge/
|  |  |- RL_ENV/                     环境、初始化、迭代、PID、SQ
|  |  |- RL_test/                    离线动作与初始化扫描示例
|  |  `- Data/                       装置 L、初始 LX、实验 CSV
|  |- meq/
|  |  |- *.p / solver/*.p            高层 MATLAB 求解器发行产物
|  |  |- libmeq/                     平衡/几何/基函数 C 内核
|  |  |- mexc/ / include/            MATLAB C 接口与类型/BLAS 宏
|  |  |- csfc/                       Simulink S-function 接口
|  |  |- genlib/src/                 Green 核、插值及特殊函数
|  |  `- matmkl/ / libv/             线性代数适配层
|  `- utils/MDS2DATA/                数据获取/格式转换辅助工具
|- matlab/R2024a/                    MATLAB 运行安装
`- sysroot/                          Linux 动态库与运行资源
\end{lstlisting}

\evimpl{} 主调用关系如图~\ref{fig:callgraph} 所示；箭头跨入 Socket 服务端处只由启动器与客户端协议证明，服务端内部如何构造 MATLAB 环境仍不可见。

\begin{figure}[htbp]
\centering
\resizebox{\textwidth}{!}{%
\begin{tikzpicture}[
  font=\footnotesize,
  node distance=5mm and 8mm,
  box/.style={draw, rounded corners=1pt, align=left, inner sep=4pt},
  grp/.style={draw=black!50, dashed, rounded corners=3pt, inner sep=6pt},
  arr/.style={-{Stealth[length=2mm]}, thick},
  lbl/.style={font=\scriptsize\itshape, text=black!70}
]
% 外层：项目外 Python
\node[box] (sim) {项目外 HFMSimulator};
\node[box, right=of sim] (sock) {HFMSocketPredictor};
\node[box, right=of sock] (ps) {PowerSupplyModel};
\node[box, right=of ps] (dk) {DockerSocketPredictor};
% 服务端
\node[box, above=12mm of sock] (server) {hfm\_socket\_server.p（TCP 行协议，服务端黑箱）};
\node[box, left=of server] (sh1) {run\_hfm\_server.sh};
\node[box, right=of server] (sh2) {socket\_serve/hfm\_server.sh};
% MATLAB 主链
\node[box, below=14mm of sim, text width=88mm] (matlab) {%
\textbf{Git 可见的 MATLAB 环境主链}\quad
FGE\_state\_env\_std(config)\\
 $\to$ my\_fusion\_environment $\to$ init $\to$ [可选 init\_equi] $\to$ iterFGE\_init\\
 $\to$ get\_Rp $\to$ evalSQ\_cp\\[2pt]
step(u) $\to$ [可选 MATLAB PIDController] $\to$ my\_fusion\_environment.step\\
 $\to$ iterFGE $\to$ fgetk\_implicit $\to$ solveF $\to$ newton\_solver\\
 \hspace*{1em}$\to$ evaluate\_F $\to$ fgeF $\to$ meqFx / meqpdom / bfct / meqcde\\
 \hspace*{1em}$\to$ solve\_linear\_problem $\to$ reduce\_system $\to$ gmres\_jfnk\\
 \hspace*{1em}$\to$ linesearch $\to$ 接受、保留旧解或停止\\
 $\to$ meqLYdot / meqlim $\to$ H 更新 $\to$ get\_obs / get\_state};
\draw[arr] (sim) -- (sock);
\draw[arr] (sock) -- (ps);
\draw[arr] (ps) -- (dk);
\draw[arr] (sh1) -- (server);
\draw[arr] (sh2) -- (server);
\draw[arr, <->] (sock) -- node[lbl, right]{TCP 行协议} (server);
\draw[arr, dashed] (matlab.north -| sim) -- node[lbl, left, align=right]{服务端内部\\构造不可见} (sim.south);
\end{tikzpicture}}
\caption{软件分层与主调用关系。Socket 服务端为不可见二进制；MATLAB 主链为 Git 源码可复核路径。}
\label{fig:callgraph}
\end{figure}

MEQ 标准批量入口 \code{fget} 与这里的自定义 \code{iterFGE} 有重要差异，特别是电源输出和预条件器更新，见第~\ref{sec:onestep} 节。不能用标准 \code{fget} 的行为覆盖 RL 迭代器的实际代码。

% ================= 第 \ref{sec:model} 章 =================
\section{物理与数学模型}\label{sec:model}

本章从仿真器实际实现的视角给出物理系统的数学描述：先约定变量、单位与维度（第~\ref{sec:vars} 节），再依次建立轴对称平衡方程与剖面参数化（第~\ref{sec:gse} 节）、导体回路与零维等离子体电流演化（第~\ref{sec:circuits} 节）、电阻预计算（第~\ref{sec:resist} 节）、边界条件与域识别（第~\ref{sec:bcs} 节），随后列出库中真实存在但默认未选中的可选闭合分支（第~\ref{sec:altcde} 节）与剖面约束（第~\ref{sec:constraints} 节），最后在第~\ref{sec:fullsystem} 节集中给出默认配置下每个时间步实际求根的完整方程组。

\subsection{变量、单位及维度}\label{sec:vars}

表~\ref{tab:vars} 中的尺寸针对 \code{ini_L_v2.2.mat}，长度单位及磁通约定由 \code{Q/README.md:218--224}、\code{Q/documentation/CHEATSHEET.md:8--25,54--108} 与有关输入注释支持。矩阵单位为这些定义下的量纲推导；MAT 数据本身没有统一单位元数据。

\begin{table}[htbp]
\centering\footnotesize
\caption{主要数学量、定义、单位及对应代码变量（v2.2 网格）}
\label{tab:vars}
\begin{tabular}{lL{6.6cm}L{6.2cm}}
\toprule
数学量 & 定义、单位 & 代码变量与例示维度 \\
\midrule
$R,Z$ & 柱坐标径向/竖直位置，m；假设轴对称，场不依赖环向角 & \code{L.rx} 66、\code{L.zx} 65；内网格 \code{ry} 64、\code{zy} 63 \\
$F(R,Z)$ & 极向磁通 $F=\int B_Z\,dS_Z$，Wb；不是每弧度磁通，也不是穿过 $R$--$Z$ 极向截面的环向磁通 & \code{LY.Fx}，$65\times66$，共 $n_x=4290$ \\
$F_A,F_B$ & 磁轴/边界磁通，Wb & \code{FA,FB}；多域内部使用 \code{F0,F1} \\
$\xi$ & 归一化磁通，无量纲，$(F-F_0)/(F_1-F_0)$ & \code{FN}；\code{L.pQ.^2} 为磁面剖面采样点 \\
$\bm{i}_y$ & 内网格环向等离子体电流单元，A & \code{Iy(:)}，$n_y=4032$；\code{Iy} 为 $63\times64$ \\
$j_\phi$ & 电流密度，A/m$^2$，$I_{y,ij}/\Delta S$ & \code{Iy*L.idsx}，$\Delta S=\texttt{drx*dzx}$ \\
$\bm{i}_a,\bm{i}_u$ & 主动回路/广义被动回路电流，A & \code{Ia} 12、\code{Iu} 525 \\
$\bm{i}_e$ & $[\bm{i}_a;\bm{i}_u]$，A & \code{Ie}，$n_e=537$；物理 vessel 细丝 \code{Iv} 544 \\
$\bm{v}_a$ & 12 路主动回路电压，V & \code{Va}；主动标签 \code{G.dima=PF_001...PF_011,VS} \\
$\bm{a}_g$ & 压力导数/环向场函数导数的基系数；当前三列单位依次为 A/(m·Wb)、A·m/Wb、A·m/Wb$^2$（由 $T_{yg}a_g$ 的量纲推导） & \code{ag}，$n_g=3$；\code{bfab} 参数 \code{[1,2]} \\
$P,T$ & 压力 Pa；$T=RB_\phi$，T·m，\textbf{不是温度} & \code{PQ,TQ,PpQ,TTpQ}，\code{G:meq/meqprof.m:1--50} \\
$I_p$ & 等离子体总环向电流，A & \code{Ip}；单域为域内 \code{Iy} 求和 \\
$b_p,q_A$ & 极向 beta 与轴上安全因子，无量纲 & \code{bp,qA}；外部参数 \code{q0} 作为 \code{qA} 目标 \\
$\sigma$ & 给定等效电导率，S/m & \code{signeo}；默认参数名不表示由温度自洽计算 \\
$L_p,R_p,I_{\mathrm{ni}}$ & 总等离子体电感 H、电阻 $\Omega$、非感应电流 A & \code{LX.Lp,LX.Rp,LX.IniD} \\
$M_{ee},M_{ey}$ & 导体互感 H；导体与网格电流互感 H & \code{L.Mee} $537\times537$；\code{L.Mey} $537\times4032$ \\
$\bm{R}_e$ & 导体电阻 $\Omega$，代码以向量逐元乘 & \code{L.Re}，537；主动部分为 \code{G.Ra} \\
$h$ & 仿真时间步长，s & \code{dt=LXt.t-LYp.t}，不是几何网格间距 \\
$\bm{x}$ & 缩放后的非线性求解未知量，无量纲（按源码定义的尺度） & \code{xnl=xSI./L.xscal}，$n_N=4572$ \\
$\bm{\mathcal R}$ & 缩放/拼接后的代数残差；各块原始量纲不同 & \code{res}，4572；收敛用整体二范数 \\
\bottomrule
\end{tabular}
\end{table}

\subsubsection{磁通单位：韦伯}\label{sec:wbu}

\evbg{} Wb 是 weber（韦伯）的单位符号，用来度量磁通。若 $\bm{B}$ 是磁感应强度（单位 tesla，特斯拉，符号 T），$d\bm{S}$ 是带方向的面积元（m$^2$），则磁通为 $\Phi_B=\int_S\bm{B}\cdot d\bm{S}$，单位关系为
\begin{equation}\label{eq:wb}
1\ \mathrm{Wb}=1\ \mathrm{T\,m^2}=1\ \mathrm{V\,s}=1\ \mathrm{H\,A}=1\ \mathrm{kg\,m^2\,s^{-2}\,A^{-1}} .
\end{equation}
例如，1 T 的均匀磁场垂直穿过 $1\ \mathrm{m^2}$ 的平面时，磁通为 1 Wb；变化率 1 Wb/s 对应单匝回路感应电压的幅值 1 V。式~\eqref{eq:wb} 是单位关系，不表示仿真器另外求解了一个独立的感应电压 PDE。

\evimpl{} 本文 $F$、\code{Fx/FA/FB/FX} 是以 Wb 表示的极向总磁通；$F/(2\pi)$ 是每弧度磁通的另一种约定，不能不换算就代入同一 GS 方程。互感矩阵单位 H，所以 $MI$ 的单位为 Wb；相反，代码中的 $T=RB_\phi$（\code{TQ/rBt}）单位是 T·m，不是 Wb，也不是温度。磁通定义依据 \code{Q/README.md:218--224}，$T$ 的定义依据 \code{G:meq/meqprof.m:5--18}。这些区别解释了本文 GS 方程和剖面系数中为何出现 $2\pi$ 或 $4\pi^2$。

\subsection{轴对称平衡方程与剖面参数化}\label{sec:gse}

\evimpl{} \code{G:meq/dlstc.m:18--37} 组装五点稀疏算子，\code{G:meq/gszrc.m:17--38} 产生相同结构的快速求解系数；\code{G:meq/fgeF.m:198--219} 用该算子在磁通与电流之间换算。

\evderiv{} 定义 $h_R=\texttt{drx}$、$h_Z=\texttt{dzx}$（m）及内点 $R_j=\texttt{ry}(j)$。算子的连续形式为
\begin{equation}\label{eq:gs}
\Delta^*F\equiv R\frac{\partial}{\partial R}\left(\frac1R\frac{\partial F}{\partial R}\right)+\frac{\partial^2F}{\partial Z^2},\qquad \Delta^*F=-2\pi\mu_0R\,j_{\phi,\mathrm{source}} .
\end{equation}
这里 $\mu_0$ 为真空磁导率，源码取 $4\pi\times10^{-7}$ H/m；$j_{\phi,\mathrm{source}}$ 是所求磁通对应的总源密度（A/m$^2$），仅等离子体自场时才等于表~\ref{tab:vars} 的 $j_\phi$，总场的网格等效外源见第~\ref{sec:bcs} 节。式~\eqref{eq:gs} 中的 $2\pi$ 来自本实现采用 Wb 总磁通的定义。该方程是从实际 \code{dlst}、\code{rhsf} 和 $I_y=\Delta S\,j_\phi$ 关系还原，不是直接套用以 $F/(2\pi)$ 为未知量的其他 GS 约定。对应磁场关系为
\begin{equation}\label{eq:bfield}
B_R=-\frac{1}{2\pi R}\frac{\partial F}{\partial Z},\qquad B_Z=\frac{1}{2\pi R}\frac{\partial F}{\partial R};
\end{equation}
符号和 $2\pi$ 可由 \code{Q/genlib/src/greenem.c:65--71} 的磁通与场模式核对。

\evimpl/\evderiv{} 压力和环向场剖面用磁通函数的有限基展开。令 $g_\ell(F-F_A,F_B-F_A)$ 为第 $\ell$ 个物理基函数，$\chi_{ij}$ 为等离子体域掩码，$f^P_\ell,f^T_\ell$ 指示该基属于 $P'$ 还是 $TT'$，则 \code{bfct1} 实际组装
\begin{equation}\label{eq:tyg}
(T_{yg})_{(i,j),\ell}=\chi_{ij}\left[R_j f^P_\ell+R_j^{-1}f^T_\ell\right]g_\ell(F_{ij}-F_A,F_B-F_A),\qquad \bm{i}_y^{\mathrm{prof}}=T_{yg}\bm{a}_g .
\end{equation}
严格说，代码 \code{if ... else if} 对两个标志都为真时优先选择 $P'$，式~\eqref{eq:tyg} 针对发行数据使用的互斥分类。\code{bfct1} 的 \code{Tpg} 是各列离散和，\code{ITpg} 是相应原函数的加权离散和；此函数没有再乘单元面积，面积及物理常数已由系数转换处理。证据：\code{Q/libmeq/bfct.c:180--245}、\code{G:meq/fgeF.m:280--288}。

\code{bfct3} 将归一化基的系数转换为物理剖面系数：
\begin{equation}\label{eq:bfct3}
a_{P',\ell}=\texttt{idsx}\,\alpha'_{\ell}\,a_{g,\ell}\,f^P_\ell/(2\pi),\qquad
a_{TT',\ell}=\texttt{idsx}\,\mu_0\alpha'_{\ell}\,a_{g,\ell}\,f^T_\ell/(2\pi);
\end{equation}
压力及 $T^2/2$ 的系数将 $\alpha'_{\ell}$ 换为积分转换因子 $\alpha_\ell$。这里 \code{idsx} $=1/\Delta S$，$\alpha',\alpha$ 的幂次取决于基函数，不应删除这些因子后声称求解了另一种标准形式。证据：\code{Q/libmeq/bfct.c:281--310}。

\evdata/\evimpl{} 当前 \code{L.bfct} 指向 \code{bfabmex}，\code{L.bfp=[1,2]}，即一个 $P'$ 和两个 $TT'$ 基。\code{bfab\_dispatch} 跳过 \code{bfef} 常数基，将后续多项式分别分配给两类，故当前三个物理基为
\begin{equation}\label{eq:bases}
g_1=F-F_B,\qquad g_2=F-F_B,\qquad g_3=(F-F_A)(F-F_B),
\end{equation}
第一列乘 $R$，第二、三列乘 $1/R$。前两基相同但属于不同物理剖面，不能合并。其原函数分别为 $(F-F_B)^2/2$、$(F-F_B)^2/2$ 及
\begin{equation}\label{eq:q3}
Q_3(F)=\tfrac12(F-F_B)^2\left[\tfrac23(F-F_B)+(F_B-F_A)\right].
\end{equation}
这些函数在边界的导数基为零，是\textbf{剖面参数化约束}，不是把计算矩形外边界设为零。证据：\code{Q/libmeq/bfab.c:25--39,55--79,169--185}、\code{Q/libmeq/bfef.c:21--44}；优化实现 \code{bf3p} 在 \code{Q/libmeq/bf3p.c:27--64,192--239} 表达相同三基，但当前 MAT 指向 \code{bfabmex}，不能据此声称调用了 \code{bf3pmex}。

\evimpl{} \code{meqprof} 先由积分基得到 $H_T\equiv\texttt{hqTQ}$，再构造环向场函数。记 $T_0=\texttt{rBt}$（T·m），则
\begin{align}
T&=\operatorname{sign}(T_0)\sqrt{T_0^2+2H_T},\qquad iT=T^{-1} && (\texttt{smalldia=false}); \label{eq:tsqrt}\\
T&=T_0+H_T/T_0,\qquad\ \ \, iT=T_0^{-1}-H_T/T_0^3 && (\texttt{smalldia=true}). \label{eq:tsmall}
\end{align}
第二式~\eqref{eq:tsmall} 是代码的一阶小抗磁近似，不能将其当作第一式~\eqref{eq:tsqrt} 的精确计算。\code{evalSQ\_cp} 调用时没有传 \code{smalldia}，故其电阻积分使用函数默认的第一分支。证据：\code{G:meq/meqprof.m:25--50}、\code{V/evalSQ_cp.m:26}。

\subsection{导体回路与零维等离子体电流演化}\label{sec:circuits}

\evimpl{} 非线性回路残差直接使用当前及上一时间片的电流差。定义 $D_R=\operatorname{diag}(\bm{R}_e)$、$\bm{v}_e=[\bm{v}_a;\bm{0}_{525}]$，则实际时间步方程为
\begin{equation}\label{eq:circuit}
M_{ee}\frac{\bm{i}_e^{\,n+1}-\bm{i}_e^{\,n}}{h}+M_{ey}\frac{\bm{i}_y^{\,n+1}-\bm{i}_y^{\,n}}{h}+D_R\bm{i}_e^{\,n+1}-\bm{v}_e^{\,n+1}=\bm{0} .
\end{equation}
主动与被动回路都与等离子体网格电流互感耦合，被动回路没有外加电压项。它是预计算矩阵上的集总电路模型，没有逐步求解三维导体内部电磁扩散。原始量纲为 V，之后按 \code{L.resscal} 缩放。证据：\code{G:meq/fgeF.m:116--124,460--482}。

\evderiv{} 对足够光滑的离散状态轨迹，差商式~\eqref{eq:circuit} 对应半离散连续方程
\begin{equation}\label{eq:circuitc}
M_{ee}\dot{\bm{i}}_e+M_{ey}\dot{\bm{i}}_y+D_R\bm{i}_e=\bm{v}_e .
\end{equation}
代码直接求解的是后向差分式~\eqref{eq:circuit}，不能据此添加独立惯性或动量状态。\code{dt=0} 时 \code{idt=0}，退化为 $D_R\bm{i}_e=\bm{v}_e$；这是代码处理静态残差/线性化的特殊情形，不是零时间间隔的物理瞬态积分。

\evimpl{} 默认 CDE 分支 \code{OhmTor_rigid_0D} 对每个活动等离子体域 $d$ 定义域电流 $\bm{i}_{y,d}=\bm{i}_y\odot\bm{1}_{\texttt{Opy}=d}$、总电流 $I_{p,d}=\bm{1}^T\bm{i}_{y,d}$ 和权重 $\bm{w}_d=\bm{i}_{y,d}/I_{p,d}$。若 $I_{p,d}=0$，代码改为域内均匀权重。令 $M_{ye}$ 为导体电流到内网格磁通的互感（H），实际残差是
\begin{equation}\label{eq:cde}
\mathcal R_{p,d}=L_{p,d}^{\,n+1}\frac{I_{p,d}^{\,n+1}-I_{p,d}^{\,n}}{h}+R_{p,d}^{\,n+1}\bigl(I_{p,d}^{\,n+1}-I_{\mathrm{ni},d}^{\,n+1}\bigr)+\bm{w}_d^{T}M_{ye}\frac{\bm{i}_e^{\,n+1}-\bm{i}_e^{\,n}}{h}=0 .
\end{equation}
全部系数/权重按当前候选解或当前输入求值。电阻项包含非感应电流，不能漏写。当前 \code{L.P.iLpext=1}，所以 $L_p$ 来自 \code{LX.Lp}；若关闭该标志，单域由 $\bm{w}^T(F_y-M_{ye}\bm{i}_e)/I_p$ 求有效电感，多域走单位域电流的自磁通计算。证据：\code{G:meq/meqcdefun.m:59--114}。

式~\eqref{eq:cde} 是``刚性零维电流闭合''的具体含义：仍求解可变形平衡，但电流时间演化用总电流方程闭合，不是把等离子体整体几何固定成刚体。

\subsection{电导率剖面及电阻的预计算}\label{sec:resist}

\evimpl{} \code{my\_fusion\_environment.get\_Rp} 接受幅值轨迹 $\sigma_{0,n}$（S/m）：标量复制到全时间轴；短序列重复末值；长序列截断。设 $p_q=\texttt{L.pQ}(q)$ 为磁面径向参数，无量纲，则保存
\begin{equation}\label{eq:sigma}
\sigma_{q,n}=(1-0.9\,p_q)\,\sigma_{0,n} .
\end{equation}
这里乘的是 \code{pQ}，\textbf{不是} \code{pQ.\^{}2}。它是给定线性剖面，不是由温度、密度或 Spitzer 公式自洽计算。证据：\code{V/my_fusion_environment.m:107--136}。

\evimpl/\evderiv{} \code{evalSQ\_cp} 对 \code{LX} 各片利用保存/预构造的磁通和基系数计算磁面量；令 $\Phi_q=\texttt{FtQ}_q+\texttt{Ft0Q}_q$（环向磁通，Wb），$T_q=RB_\phi$，则代码积分为
\begin{equation}\label{eq:sq}
S_{Q,k}=\frac{T_0}{2\pi}\sum_{q=1}^{k-1}\frac12\left(\frac{\sigma_q}{T_q^2}+\frac{\sigma_{q+1}}{T_{q+1}^2}\right)(\Phi_{q+1}-\Phi_q),\qquad R_p=\frac{1}{S_{Q,\mathrm{edge}}} .
\end{equation}
这是 \code{cumtrapz} 的严格离散形式；$S_Q$ 的单位为 S，倒数为 $\Omega$。代码没有对 $\Phi$ 重新排序，积分符号沿实际数组顺序。\code{get\_Rp} 一次性改写 \code{LX.signeo} 和 \code{LX.Rp} 的全时序；普通 \code{step} 没有重新以更新后的 \code{LY} 调用这一积分，因此不能声称每步自洽更新电阻。证据：\code{V/evalSQ_cp.m:11--32}、\code{V/my_fusion_environment.m:138--156}。

\subsection{边界条件、外导体与域识别}\label{sec:bcs}

\evimpl{} 计算矩形边界使用自等离子体磁通与外导体磁通叠加：
\begin{equation}\label{eq:fbdy}
\bm{F}_b=\mathcal B(\bm{i}_y)+M_{be}\bm{i}_e .
\end{equation}
当前 \code{ilackner=2}，$\mathcal B(\bm{i}_y)=M_{by}\bm{i}_y$；\code{ilackner=1} 则先做零边界 Poisson 求解、\code{nfdbmex} 法向差分，再乘 \code{Tbc} 构造边界。故这是自由边界磁场问题中的计算边界处理，不能写为``LCFS 上统一规定 $F=0$''。证据：\code{G:meq/meqfbp.m:8--19}、\code{G:meq/meqFx.m:95--122}。

\evimpl{} 外导体在网格上的等效电流记为 $T_{ye}\bm{i}_e$（A）。\code{meqFx} 的默认 \code{gsxe=2} 分支求解
\begin{equation}\label{eq:gsxr}
\bm{F}_x=\operatorname{GSZR}\bigl(\mathcal B(\bm{i}_y)+M_{be}\bm{i}_e,\ \bm{i}_y+T_{ye}\bm{i}_e\bigr).
\end{equation}
因此总磁通的逆映射是
\begin{equation}\label{eq:invmap}
\bm{i}_y=-\operatorname{diag}(\bm{r}_f)^{-1}D_{*}\bm{F}_x-T_{ye}\bm{i}_e,
\end{equation}
其中 $D_{*}=\texttt{L.dlst}$、$\bm{r}_f=\texttt{L.rhsf}$。忽略 $T_{ye}\bm{i}_e$ 会把外导体在网格上的贡献误当作等离子体电流。证据：\code{G:meq/meqFx.m:51--63,115--129}、\code{G:meq/fgeF.m:188--219}。

\evimpl{} 磁轴、X 点、限制器接触和 LCFS 不是独立动态变量，主要由当前磁通的拓扑/几何后处理得到。\code{asxy} 通过邻点差符号变化找候选极值/鞍点，消除过近极值--鞍点对，再用六点二次插值定位；\code{fbnd} 按电流符号在限制器与 X 点候选磁通间选择边界；\code{pdom} 结合几何掩码与 $F\gtrless F_B$ 构造域。证据：\code{Q/libmeq/asxy.c:21--94,146--178,181--243}、\code{Q/libmeq/fbnd.c:13--40}、\code{Q/libmeq/pdom.c:4--15}。多域高层逻辑另由 \code{G:meq/meqpdom.m} 完成，不能仅用单域 \code{pdom.c} 概括所有拓扑。

\evimpl{} 若域识别等步骤失败，\code{fgeF} 的有关分支返回 NaN 残差与失败标记；这会影响有限差分、线搜索和步接受，而非自动保证恢复到有效平衡。证据：\code{G:meq/fgeF.m:221--277}、\code{G:meq/solver/evaluate_F.m:87--108}。

\subsection{可选物理闭合分支}\label{sec:altcde}

表~\ref{tab:altcde} 给出当前库真实存在但发行配置未默认选中的分支。触发由 \code{P.cde}、\code{meqcdec} 及活动域/自由度决定；$n_g=n_C+n_p$ 的自由度计数必须匹配。\code{G:meq/meqcde.m:63--110,146--153} 还处理无 CDE 和被关闭域的补充约束。

\begin{table}[htbp]
\centering\small
\caption{可选电流演化闭合分支（当前默认配置未选中）}
\label{tab:altcde}
\begin{tabular}{L{2.4cm}L{7.2cm}L{3.4cm}}
\toprule
分支及触发 & 实际残差/算法 & 源码 \\
\midrule
\code{OhmTor_0D} & $\bm{w}^T(F_y^{n+1}-F_y^n)/h+R_p(I_p-I_{\mathrm{ni}})=0$；直接差分总网格磁通，不作刚性电感分解 & \code{G:meq/meqcdefun.m:168--233} \\
\code{cde_ss_0D} & $10^{-6}\{I_p-I_{\mathrm{ni}}+S_Q[(F_0-F_0^n)+(F_1-F_1^n)]/(2h)\}=0$，$S_Q$ 由 \code{trapz} 得到；名字中的 \code{ss} 不代表源码没有时间差商 & 同文件 \code{238--330} \\
\code{static_Iy_1D} & $[G(F)^T\bm{i}_y-G(F^{\mathrm{ref}})^T\bm{i}_y^{\mathrm{ref}}]/I_{p,0}=0$；保持测试函数投影，不是瞬态扩散 & 同文件 \code{335--399} \\
\code{cde_OhmTor_1D} & $r_\ell=\Delta S\sum_y[(F_y-F_y^n)/(Rh)+2\pi(j_\phi-j_{\mathrm{ni}})/\sigma_y]g_\ell(\xi_y)$；有限基测试空间的电流演化；配置字符串按此函数名匹配 & 同文件 \code{404--517}，函数 \code{cde_OhmTor_1D}；\code{G:meq/meqcdec.m:43} \\
\bottomrule
\end{tabular}
\end{table}

一维弱式分支将 \code{LX.signeo} 对 \code{L.pQ.\^{}2} 按归一磁通做线性插值；标量 \code{Ini} 非零时只是警告，所用局部 $j_{\mathrm{ni}}$ 仍为零，不能把该接口描述成已实现任意非感应电流密度剖面。证据：\code{G:meq/meqcdefun.m:420--466}。各分支方程不应相互替代，也不能把库中的一维能力归到默认 \code{OhmTor_rigid_0D}。

\subsection{剖面约束：\texorpdfstring{$b_p$}{bp} 与 \texorpdfstring{$q_A$}{qA} 的实际意义}\label{sec:constraints}

\evdata/\evimpl{} 动态例子的三个剖面系数由两个代数约束 \code{bp,qA} 加一个 CDE 闭合，即 $n_g=3=n_C+n_p=2+1$。静态初始化把零维 CDE 自由度改成总电流约束，得到 \code{Ip,bp,qA} 三个代数约束；导体电流此时是给定输入。证据：\code{G:meq/fgec.m:15--22}、\code{V/init_equi.m:15--25,103--116}、\code{G:meq/fgsc.m:13--24}。

为明确 $b_p$ 不是逐网格压力赋值，记 $\bm{t}_{p,d}$ 为域内电流基积分行（\code{TpDg}）、$\bm{t}_{P,d}$ 为压力原函数积分行（\code{ITpDg} 中属于压力的列）、$I_{p,0}$ 为归一电流尺度、$R_0=\texttt{L.P.r0}$（m）。则 \code{meqagconfun} 实际使用
\begin{align}
I_{p,d}^{\mathrm{prof}}&=\bm{t}_{p,d}\bm{a}_g,\qquad W_N=10^{-7}\pi R_0\bigl(I_{p,d}^{\mathrm{prof}}\bigr)^2,\qquad W_{N,0}=10^{-7}\pi R_0I_{p,0}^2, \label{eq:wn}\\
r_{bp}&=\frac{\bm{t}_{P,d}\bm{a}_g-bp_{\mathrm{target}}W_N}{W_{N,0}},\qquad r_{Ip}=\frac{\bm{t}_{p,d}\bm{a}_g-I_{p,\mathrm{target}}}{I_{p,0}} . \label{eq:rbp}
\end{align}
压力积分的常数/单位采用源码内部定义；不能将这里的 $W_N$ 无条件替换成其他教材中定义的储能。证据：\code{G:meq/meqagconfun.m:87--98,145--164}。

\evimpl{} 磁轴 $q_A$ 约束明确标为小抗磁近似，虽 \code{L.smalldia=0}，这一个约束仍有自身近似。令 $F_{RR},F_{ZZ},F_{RZ}$ 是磁轴处磁通 Hessian（Wb/m$^2$），
\begin{equation}\label{eq:gq}
g_A=\frac{\sqrt{F_{RR}F_{ZZ}-F_{RZ}^2}}{|F_{RR}+F_{ZZ}|},\qquad
C_q=\frac{4\pi10^{-7}R_A^2T_0\,g_A}{\Delta S},
\end{equation}
$\bm{t}_A$ 是磁轴处 $g_\ell(R_A f^P_\ell+f^T_\ell/R_A)$，$\bm{J}_g=(2\times10^{-7}/\Delta S)\bm{Ig}_A$，其中 $\bm{Ig}_A$ 是基函数原函数在轴上的值。源码的约束为
\begin{equation}\label{eq:rqa}
r_{q_A}=\frac{\left[(\bm{J}_g\odot\bm{f}_T)^T-q_{A,\mathrm{target}}C_q\bm{t}_A\right]\bm{a}_g+T_0^2}{(R_0B_0)^2}=0,
\end{equation}
其中 $B_0=\texttt{L.P.b0}$（T）。原式各向量朝向在式~\eqref{eq:rqa} 中统一为行/列乘积。\code{q0} 通过这一局部 Hessian 与基系数约束影响平衡，没有单独的安全因子输运方程。证据：\code{G:meq/meqagconfun.m:216--248}。

\subsection{默认配置的完整求解方程组}\label{sec:fullsystem}

本节把前述分散在各模块的关系集中列出。\evimpl{} 以下针对本文实际读取的 \code{ini_L_v2.2.mat} 单域配置：\code{algoNL='all-nl'}、\code{cde='OhmTor_rigid_0D'}、\code{gsxe=2}、\code{ilackner=2}、\code{iLpext=1}、\code{bfab[1,2]}、\code{lin=false}；其他分支不能直接套用本节全部方程。连续式标为数学解释；最终决定代码行为的是第~\ref{sec:legacy-2-8-3} 节的离散残差。

\subsubsection{与代码一致的连续 GS 形式及剖面闭合}\label{sec:legacy-2-8-1}

\evderiv{} 定义 $P_F=dP/dF$（Pa/Wb）和 $TT_F=d(T^2/2)/dF=T\,dT/dF$（T$^2$·m$^2$/Wb），分别对应 \code{PpQ} 和 \code{TTpQ}。这里求导变量是 Wb 总磁通 $F$，不是 $F/(2\pi)$。由 \code{bfct1} 的电流基与 \code{bfct3} 的物理系数换算，等离子体域内的剖面电流密度满足
\begin{equation}\label{eq:jprof}
j_\phi=2\pi R\,P_F+\frac{2\pi}{\mu_0R}\,TT_F .
\end{equation}
因此在不含外导体物理电流的等离子体区域，连续 GS 对应式为
\begin{equation}\label{eq:gsfull}
\boxed{\;\Delta^*F+4\pi^2\left(\mu_0R^2\frac{dP}{dF}+T\frac{dT}{dF}\right)=0,\qquad \Delta^*F=R\partial_R(R^{-1}\partial_RF)+\partial_Z^2F\;}
\end{equation}
式~\eqref{eq:gsfull} 与第~\ref{sec:gse} 节的源项形式相容；若改用 $\psi=F/(2\pi)$，必须同时变换导数和算子，不能直接删去 $4\pi^2$。上述式子是从源码系数严格还原的连续解释，不代表代码在整个计算矩形中忽略外导体：外导体的网格等效源、自由边界和真实磁通求解仍完整保留在第~\ref{sec:legacy-2-8-2} 节。证据：\code{Q/libmeq/bfct.c:180--245,281--310}、\code{Q/libmeq/bfef.c:34--44,60--70}、\code{G:meq/meqprof.m:32--42}、\code{G:meq/dlstc.m:18--37}。

\evderiv{} 为把压力和环向场闭合也明确列出，令 $a_\ell=(a_g)_\ell$，$g_1=g_2=F-F_B$、$g_3=(F-F_A)(F-F_B)$，$Q_1=Q_2=(F-F_B)^2/2$、$Q_3$ 如式~\eqref{eq:q3}。其中 $Q_\ell$ 是 $g_\ell$ 关于 $F$ 的原函数，且在 $F=F_B$ 取零。则
\begin{align}
P_F(F)&=\frac{a_1g_1(F)}{2\pi\Delta S},& P(F)&=\frac{a_1Q_1(F)}{2\pi\Delta S}, \label{eq:pprof}\\
TT_F(F)&=\frac{\mu_0[a_2g_2(F)+a_3g_3(F)]}{2\pi\Delta S},& H_T(F)&=\frac{\mu_0[a_2Q_2(F)+a_3Q_3(F)]}{2\pi\Delta S}, \label{eq:ttprof}\\
T(F)&=\operatorname{sign}(T_0)\sqrt{T_0^2+2H_T(F)},& B_\phi(F,R)&=T(F)/R . \label{eq:tprof}
\end{align}
这里 $\Delta S=h_Rh_Z$ 是电流网格单元面积，$H_T=(T^2-T_0^2)/2$；$a_\ell$ 的单位见第~\ref{sec:vars} 节，$T_0=\texttt{rBt}$。由此得到 $P(F_B)=0$、$T(F_B)=T_0$，是此剖面构造的边界取值，并非矩形计算边界的磁通为零。式中使用当前 \code{L.smalldia=0} 的平方根分支；启用小抗磁分支时应按式~\eqref{eq:tsmall} 替换。证据：\code{Q/libmeq/bfab.c:25--39,55--79}、\code{Q/libmeq/bfef.c:21--44}、\code{Q/libmeq/bfct.c:281--310}、\code{G:meq/meqprof.m:32--51}。

\subsubsection{网格磁通、边界和几何辅助关系}\label{sec:legacy-2-8-2}

\evimpl/\evderiv{} 记 $\bm{F}\in\mathbb R^{4290}$ 为 \code{Fx(:)}（Wb），$\Pi_b\in\{0,1\}^{258\times4290}$ 为选取矩形边界节点的矩阵（与 \code{L.lxy} 的边界补集相对应），$D_{*}=\texttt{L.dlst}$、$\bm{r}_f=\texttt{L.rhsf}$。每个候选物理状态先满足如下线性磁通子问题：
\begin{equation}\label{eq:fluxsub}
\boxed{\;\begin{aligned}
D_{*}\bm{F}+\operatorname{diag}(\bm{r}_f)\bigl(\bm{i}_y+T_{ye}\bm{i}_e\bigr)&=\bm{0}_{4032},\\
\Pi_b\bm{F}-M_{by}\bm{i}_y-M_{be}\bm{i}_e&=\bm{0}_{258} .
\end{aligned}\;}
\end{equation}
合计 $4032+258=4290$ 个关系确定 4290 个网格磁通值，记该求解映射为 $\bm{F}=\mathcal F(\bm{i}_y,\bm{i}_e)$。实际默认路径用 GSZR 和预计算边界响应实现它，不是额外对这 4290 个量运行第二套非线性求根。证据：\code{G:meq/meqmesh.m:41--58}、\code{G:meq/meqFx.m:95--122}、\code{G:meq/meqfbp.m:13--19}；空间系数详见第~\ref{sec:fivepoint} 节。

对候选磁通，域掩码 $\bm\chi\in\{0,1\}^{4032}$、轴/边界磁通 $F_A,F_B$、轴位置及 Hessian 由 \code{meqpdom} 和局部几何搜索取得；$T_{yg}$ 与约束积分也随候选解更新。代码可概括为
\begin{equation}\label{eq:dommap}
(\bm\chi,F_A,F_B,R_A,Z_A,F_{RR},F_{ZZ},F_{RZ},\ldots)=\mathcal D(\bm{F}),\qquad I_p=\bm{1}^T(\bm\chi\odot\bm{i}_y),\qquad \bm{w}=\frac{\bm\chi\odot\bm{i}_y}{I_p} .
\end{equation}
$\mathcal D$ 是对实际几何/拓扑处理的记号，不是另外假定的可微物理方程。$I_p=0$ 时源码改用 $\bm{w}=\bm\chi/(\bm{1}^T\bm\chi)$；空域仍需按失败路径处理。证据：\code{G:meq/fgeF.m:221--288}、\code{G:meq/meqcdefun.m:81--95}。

\subsubsection{每个时间步实际求根的完整方程}\label{sec:legacy-2-8-3}

\evimpl/\evderiv{} 令 $n$ 为已经接受的前片，$h=t_{n+1}-t_n$，$\delta_t z=(z^{n+1}-z^n)/h$。以下未显式写上标的量全部取当前候选的 $n+1$ 端；输入 $\bm{v}_a,bp_\star,q_\star,L_p,R_p,I_{\mathrm{ni}},T_0$ 取该片 \code{LXt}。默认未知物理量为 $\bm{z}=[\bm{i}_y;\bm{a}_g;\bm{i}_e]\in\mathbb R^{4572}$，其中 $\bm{i}_e=[\bm{i}_a;\bm{i}_u]$。求解器使用 $\bm{x}=S_x^{-1}\bm{z}$，$S_x=\operatorname{diag}(\texttt{L.xscal})$。

压力约束的 $\bm{t}_P$、磁轴约束的 $\bm{t}_A,\bm{J}_g,\bm{f}_T,C_q,R_0,B_0$ 及归一尺度 $W_{N,0}$ 均沿用第~\ref{sec:constraints} 节的定义；$W_N=10^{-7}\pi R_0(\bm{t}_p\bm{a}_g)^2$，其中 $\bm{t}_p\bm{a}_g$ 是\textbf{剖面预测的总电流}。它在候选解尚未满足平衡时不能直接替换成 CDE 中从 \code{Iy} 求和得到的 $I_p$。完整根条件为
\begin{equation}\label{eq:root}
\boxed{\;\begin{aligned}
&T_{yg}\bigl(\mathcal F(\bm{i}_y,\bm{i}_e),F_A,F_B,\bm\chi\bigr)\bm{a}_g-\bm{i}_y=\bm{0}_{4032}, &&\text{网格平衡}\\
&\frac{\bm{t}_P\bm{a}_g-bp_\star W_N}{W_{N,0}}=0, && b_p\text{ 约束}\\
&\frac{\bigl[(\bm{J}_g\odot\bm{f}_T)^T-q_\star C_q\bm{t}_A\bigr]\bm{a}_g+T_0^2}{(R_0B_0)^2}=0, && q_A\text{ 约束}\\
&L_p\,\delta_t I_p+R_p\bigl(I_p-I_{\mathrm{ni}}\bigr)+\bm{w}^TM_{ye}\,\delta_t\bm{i}_e=0, && \text{零维等离子体电流演化}\\
&M_{ee}\,\delta_t\bm{i}_e+M_{ey}\,\delta_t\bm{i}_y+\operatorname{diag}(\bm{R}_e)\bm{i}_e-\begin{bmatrix}\bm{v}_a\\\bm{0}_{525}\end{bmatrix}=\bm{0}_{537} . && \text{导体电流演化}
\end{aligned}\;}
\end{equation}
式~\eqref{eq:root} 五行依次对应网格平衡、$b_p$ 约束、$q_A$ 约束、零维等离子体电流演化和导体电流演化；残差元素数为 $4032+1+1+1+537=4572$，与 $4032+3+537=4572$ 个未知量一致。这里列出的是未乘相应残差尺度的根条件，平衡块保留源码 $T_{yg}a_g-i_y$ 的符号；复现 Newton 残差及 Jacobian 时还须按第~\ref{sec:tdisc} 节使用实际缩放。证据：\code{G:meq/fgeF.m:280--288,362--482}、\code{G:meq/meqagconfun.m:145--164,216--248}、\code{G:meq/meqcdefun.m:81--114}。

从时间差商还原的半离散动态形式是
\begin{equation}\label{eq:dae}
L_p\dot I_p+R_p\bigl(I_p-I_{\mathrm{ni}}\bigr)+\bm{w}^TM_{ye}\dot{\bm{i}}_e=0,\qquad
M_{ee}\dot{\bm{i}}_e+M_{ey}\dot{\bm{i}}_y+\operatorname{diag}(\bm{R}_e)\bm{i}_e=\bm{v}_e,
\end{equation}
同时要求式~\eqref{eq:root} 的平衡/剖面约束始终成立。因此该配置可理解为\textbf{磁平衡代数约束与集总电流微分方程耦合的微分--代数系统（DAE）}，时间离散后解非线性代数方程。这里仅解释方程形式，不对其 DAE 指数作未经分析的断言。

\subsubsection{给定量、初值及方程组的闭合范围}\label{sec:legacy-2-8-4}

\evimpl{} 式~\eqref{eq:root} 每一步还依赖下列已给定关系：装置几何/互感/电阻及网格来自 \code{L}；前片电流与时间来自 \code{H.LYp/H.xnl}；$bp_\star,q_\star,I_{\mathrm{ni}},L_p$ 来自 \code{LX}；$R_p$ 通过第~\ref{sec:resist} 节的给定电导率剖面与参考磁面积分预生成；$\bm{v}_a$ 来自当前实际控制路径。首片通过静态 \code{Ip,bp,qA} 约束初始化，预推进与电阻覆写的先后顺序见第~\ref{sec:init}--\ref{sec:onestep} 节。证据：\code{V/init.m:3--11}、\code{V/init_equi.m:15--58}、\code{V/my_fusion_environment.m:53--59,107--156}、\code{V/iterFGE.m:45--91}。

代码没有在上述方程之外增加密度、流速或温度的输运闭合；$P(F)$ 是有限剖面基的代数量，$R_p$ 也不是由在线温度方程算得。默认电流演化保留的是 $L_p\dot I_p$ 与 $\bm{w}^TM_{ye}\dot{\bm{i}}_e$，不能额外加上 $I_p\dot L_p$ 或 $\dot{\bm{w}}$，也不能换成一维电流扩散 PDE。其他 \code{cde}、静态、Picard 和线性化配置的改动分别见第~\ref{sec:altcde}、\ref{sec:algonl}、\ref{sec:altalgo} 节。

% ================= 第 \ref{sec:discret} 章 =================
\section{数值离散}\label{sec:discret}

本章说明连续模型如何转化为可求解的代数系统：网格与预计算数据（第~\ref{sec:mesh} 节）、五点差分算子的构造与精度（第~\ref{sec:fivepoint} 节）、边界与外导体响应的离线构造（第~\ref{sec:offline} 节）、时间离散与全耦合代数方程（第~\ref{sec:tdisc} 节），以及不同平衡未知量与静态模式的选择（第~\ref{sec:algonl} 节）。

\subsection{网格、索引与保存的预计算数据}\label{sec:mesh}

\evimpl/\evdata{} 主 GS 求解采用 $R$--$Z$ 平面的\textbf{二维结构化、正交、张量积矩形网格}：
\begin{equation}\label{eq:grid}
\mathcal G_x=\{(R_j,Z_i):R_j=\texttt{rx}(j),\ Z_i=\texttt{zx}(i)\}.
\end{equation}
它由两个一维坐标数组的 \code{meshgrid} 构成，内点算子用上下左右五点模板；当前主求解不是非结构三角形网格有限元。证据：\code{G:meq/meqmesh.m:41--58}、\code{G:meq/dlstc.m:22--37}。局部磁轴搜索使用六点插值，并不改变主求解网格的这一性质。

\evdata/\evderiv{} 本文 v2.2 数据的 $R$ 方向、$Z$ 方向分别等距，使用常数 $h_R,h_Z$；但 $h_R\neq h_Z$，所以单元为长方形，不是正方形。``均匀''在此表示同一方向步长不随位置改变，并不表示两个方向步长相等。对保存的 \code{L.rx/L.zx} 逐差核对，$\max|\operatorname{diff}(\texttt{rx})-h_R|\approx2.78\times10^{-16}$ m、$\max|\operatorname{diff}(\texttt{zx})-h_Z|\approx4.44\times10^{-16}$ m，差别处于浮点舍入量级。\code{meqmesh} 以首个间隔设置 \code{drx/dzx}，其注释明确给出 \code{linspace} 网格，GSZR/五点算子据此使用单一间距；不能把接口能接收数组解释成已支持任意非均匀坐标。证据：\code{G:meq/meqmesh.m:1--13,38--47}、\code{G:meq/gszrc.m:13--20}、\code{D/ini_L_v2.2.mat:L.rx,L.zx,L.drx,L.dzx}。

计算矩形与真实等离子体截面不同：LCFS/限制器可以是弯曲边界，由网格磁通、几何掩码及轮廓后处理表示；矩形网格并不要求等离子体边界为矩形。\code{rq/zq} 的磁面采样点、限制器轮廓点及 \code{rz/zz} 扩展后处理网格也不应混称为主 GS 网格。证据：\code{G:meq/meqc.m:46--48}、\code{G:meq/meqmesh.m:76--82}、\code{V/my_fusion_environment.m:361--391}。

\evdata{} \code{D/ini_L_v2.2.mat:L} 的计算网格为 $n_Z=65,\ n_R=66$；内部电流网格去掉四周一层，尺寸 $63\times64$。坐标与间距为
\begin{equation}\label{eq:gridrange}
\begin{gathered}
R\in[0.2615546875,\ 1.3594453125]\ \mathrm{m},\qquad h_R=0.016890625\ \mathrm{m},\\
Z\in[-1.655873015873016,\ 1.6558730158730157]\ \mathrm{m},\qquad h_Z=0.05174603174603165\ \mathrm{m}.
\end{gathered}
\end{equation}
边界点数 $n_b=2n_Z+2(n_R-2)=258$，扩展后处理网格为 \code{rz} 31 点与 \code{zz} 41 点。\code{Fx} 的 MATLAB 存储以 $Z$ 为第一维，因此零基线性索引 $k=i_Z+n_Zj_R$；\code{Iy} 也按 $Z$ 优先。这由 \code{Q/libmeq/gszr.c:3--5} 和 \code{V/my_fusion_environment.m:381--391} 的 \code{(:)} 展平共同确认。

保存的主要矩阵与向量如表~\ref{tab:matrices} 所示。

\begin{table}[htbp]
\centering\small
\caption{保存矩阵/向量的尺寸（v2.2）与实际作用}
\label{tab:matrices}
\begin{tabular}{lL{4.4cm}L{6.2cm}}
\toprule
保存矩阵/向量 & 尺寸（v2.2） & 实际作用/来源 \\
\midrule
\code{L.dlst} & $4032\times4290$，稀疏，20160 非零 & 内点五点 GS 算子 \\
\code{L.rhsf} & $4032\times1$ & $2\pi\mu_0R_jh_Z/h_R$，将电流单元换成离散源 \\
\code{L.cx,cq,cr,cs} & $64\times1$；后三者 $64\times63$ & GSZR 源缩放及径向消元系数 \\
\code{L.Txy,L.Txb} & $4290\times4032$、$4290\times258$ & 内点/边界嵌入，分别 4032/258 非零 \\
\code{L.Mby,L.Mbe} & $258\times4032$、$258\times537$ & 边界自场、边界外导体场 \\
\code{L.Tye} & $4032\times537$ & 外导体场在离散算子上的等效源 \\
\code{L.Mxy,L.Mxe} & $4290\times4032$、$4290\times537$ & \textbf{离散磁通求解器的预计算响应}，不能都称为解析 Green 核 \\
\code{L.Mee,L.Mey,L.Re} & $537\times537$、$537\times4032$、537 & 导体方程 \\
\code{L.Jx,L.Ju,L.Jxdot} & $4572^2$、$4572\times15$、$538\times4572$ & 保存的线性化对象；存在这些数据不表示每次 JFNK 显式使用全部矩阵 \\
\code{G.Maa,Mau,Muu} & $12^2$、$12\times525$、$525^2$ & 主动/广义 vessel 互感 \\
\code{G.Tvu} & $544\times525$ & \code{Iv=Tvu*Iu}，从广义回路映射到 vessel 细丝 \\
\code{G.Bma,Bmu,Bmx} & $100\times12$、$100\times525$、$100\times4290$ & 磁探针场响应 \\
\code{G.Mfa,Mfu,Mfx} & $47\times12$、$47\times525$、$47\times4290$ & 磁通环响应 \\
\bottomrule
\end{tabular}
\end{table}

\evimpl{} \code{Mee=[[Maa,Mau];[Mau',Muu]]}，\code{Re=[Ra;Ru]}；\code{Mey=P.mey*[Mya';Myu']}，\code{Mya/Myu} 是 \code{G.Mxa/G.Mxu} 的内网格行。默认 \code{P.mey=1}。刚性 CDE 的 $M_{ye}$ 在 \code{P.anajac=true} 时取离散 \code{L.Mxe} 内网格行，否则取原几何响应内网格行；所以不能无条件写 $M_{ye}=M_{ey}^T$。证据：\code{G:meq/fgec.m:24--45}。

\evdata/\evnote{} 五份装置文件的主体尺寸相同，但不能混配版本。\code{v1,v1.1} 的 $R$ 范围为 $[0.2535,1.3585]$ m、$h_R=0.017$ m；\code{v2,v2.1,v2.2} 使用式~\eqref{eq:gridrange} 的新网格。\code{G.nl} 在 v1 为 549、v1.1 为 4999；\code{D/readme.md:1--7} 将 v2 说明为新 limiter、v2.1 为复杂案例、v2.2 为打开 PF3/4 的 XPT 配置。``更稳定/更耗时''仅是随附说明，本文没有运行实验验证。

\subsection{保守型五点差分的推导}\label{sec:fivepoint}

\evimpl/\evderiv{} 令 $\eta=h_Z/h_R$，对每一内网格径向列定义
\begin{equation}\label{eq:coefs}
a_j^-=\eta^2\frac{R_j}{R_j-h_R/2},\qquad a_j^+=\eta^2\frac{R_j}{R_j+h_R/2} .
\end{equation}
这两个系数来自在径向半网格上计算 $R^{-1}\partial_RF$。具体地，
\begin{equation}\label{eq:radial}
R_j\,\partial_R(R^{-1}\partial_RF)\approx\frac{R_j}{h_R^2}\left[\frac{F_{i,j+1}-F_{ij}}{R_j+h_R/2}-\frac{F_{ij}-F_{i,j-1}}{R_j-h_R/2}\right] .
\end{equation}
再对 $Z$ 二阶导数作中心差分，并乘 $h_Z^2$，得到实际稀疏行
\begin{equation}\label{eq:drow}
(D_{*}\bm{F})_{ij}=a_j^-F_{i,j-1}+F_{i-1,j}-(a_j^-+a_j^++2)F_{ij}+F_{i+1,j}+a_j^+F_{i,j+1} .
\end{equation}
结合 $I_{y,ij}=j_{\phi,ij}h_Rh_Z$，离散源为
\begin{equation}\label{eq:source}
D_{*}\bm{F}=-\bm{r}_f\odot\bm{I}_{\mathrm{source}},\qquad (r_f)_{ij}=\frac{2\pi\mu_0R_jh_Z}{h_R},\qquad \bm{I}_{\mathrm{source}}=\bm{i}_y+T_{ye}\bm{i}_e .
\end{equation}
$D_{*}$ 包含已知边界所在的列；边界值由第~\ref{sec:bcs} 节给出，移至右端后才是内点线性系统。证据：\code{G:meq/dlstc.m:18--37}、\code{G:meq/gszrc.m:17--20,38}、\code{G:meq/meqFx.m:115--122}。对光滑磁通、固定正半径与均匀网格，Taylor 展开给出内部空间截断误差 $O(h_R^2+h_Z^2)$；这是\textbf{局部算子推导}，不是带自由边界、拓扑切换和基展开的全仿真二阶收敛证明。

\evdata{} 静态数据核对：对 v2.2 保存的 20160 个非零元逐项比较，式~\eqref{eq:coefs}--\eqref{eq:drow} 的系数最大绝对差约 $7.11\times10^{-15}$；\code{rhsf} 与 \code{cx=rhsf/2} 的比较到浮点舍入量级。对 \code{D/ini_13906_xpt_v2.2.mat:LX.Iy} 的初始电流场，$\|D_{*}(M_{xy}\bm{i}_y)+\bm{r}_f\odot\bm{i}_y\|_2\approx2.16\times10^{-13}$。这是读取存储数据后的代数一致性核对，\textbf{不是一次动态仿真或精度试验}。

\subsection{边界及外导体响应的离线构造}\label{sec:offline}

\evimpl{} 当前 \code{P.ilackner=2} 的在线 $M_{by}\bm{i}_y$ 不意味着该矩阵由解析 Green 函数直接填充。\code{G:meq/meqc.m:243--256} 对每个电流单位向量先做零边界 GSZR、取 \code{nfdbmex} 法向差分、再经 \code{Tbc} 映射，把各结果存为 $M_{by}$。\code{P.ilackner=1} 每次在线执行同样的两阶段边界运算；\code{P.ilackner=0} 才直接取几何结构的 \code{G.Mby}。因此它是预计算 Lackner 离散边界算子。

\evimpl{} \code{P.gsxe=2} 以外导体的网格内场响应 $M_{x,ie}$ 计算
\begin{equation}\label{eq:tye}
T_{ye}=-\operatorname{diag}(\bm{r}_f)^{-1}D_{*}M_{x,ie} .
\end{equation}
\code{gsxetol} 可按每列相对阈值裁零，密度低时转为稀疏；当前 MAT \code{P.gsxetol=0} 不作该阈值裁剪。\code{gsxe=1} 则把内部导体细丝按最近内网格点落点，并检查不落入限制器域；\code{gsxe=3} 在 \code{meqFx} 直接矩阵乘 $M_{xy}\bm{i}_y+M_{xe}\bm{i}_e$。证据：\code{G:meq/meqc.m:204--233}、\code{G:meq/meqFx.m:115--129}。

\evimpl/\evderiv{} 若 $B$ 为单位边界到全网格的离散响应，$G_y$ 为零边界单位网格电流响应，$N_b$ 为法向差分，则预计算磁通映射满足
\begin{equation}\label{eq:resp}
M_{xy}=B\,T_{bc}N_bG_y+G_y,\qquad M_{xe}=B M_{be}+G_yT_{ye} .
\end{equation}
证据：\code{G:meq/meqFxjac.m:31--49}。这解释为何``JFNK 不显式构造当前 Jacobian''与``装置保存大量稠密响应矩阵''可以同时成立。

\subsection{时间离散、缩放与全耦合代数方程}\label{sec:tdisc}

\evimpl{} \code{fgetk_implicit} 使用后向 Euler：所有代数残差在 $t_{n+1}$ 求值，时间导数用 $(x^{n+1}-x^n)/h$。\code{fget} 与自定义 \code{iterFGE} 都把该函数作为非线性时间步内核。证据：\code{G:meq/fgetk_implicit.m:1--2,17--25,101--143}、\code{V/iterFGE.m:61--91}。

令 $S_x=\operatorname{diag}(\texttt{xscal})$，默认求解未知量是
\begin{equation}\label{eq:xscal}
\bm{x}=S_x^{-1}\begin{bmatrix}\bm{i}_y\\ \bm{a}_g\\ \bm{i}_a\\ \bm{i}_u\end{bmatrix}\in\mathbb R^{4572},\qquad
\bm{\mathcal R}(\bm{x};\bm{x}^n,\bm{u}^{n+1})=\begin{bmatrix}S_{GS}\bigl(T_{yg}(F)\bm{a}_g-\bm{i}_y\bigr)\\ \bm{r}_{bp,q_A}\\ r_p\\ S_e\bm{r}_e\end{bmatrix}=\bm{0} .
\end{equation}
这里 $S_{GS}$、$S_e$ 为相应残差缩放，$F=\mathcal F(\bm{i}_y,\bm{i}_e)$ 是式~\eqref{eq:fluxsub} 的离散磁通求解；$\bm{u}=[\bm{v}_a;\bm{C}\bm{o};I_{\mathrm{ni}}]\in\mathbb R^{15}$，$\bm{C}\bm{o}$ 为 2 个剖面约束目标，不是 2 个额外线圈。当前各块尺寸为 $4032,2,1,537$。证据：\code{G:meq/fgeF.m:151--219,280--288,365--482}；\code{D/L.ind}。各块的索引与保存缩放见表~\ref{tab:scaling}。

\begin{table}[htbp]
\centering\small
\caption{各未知/残差块的索引与 v2.2 保存缩放}
\label{tab:scaling}
\begin{tabular}{L{3.2cm}L{4.8cm}L{5.8cm}}
\toprule
块 & \code{L.ind} 的 MATLAB 1 基索引 & v2.2 保存的缩放 \\
\midrule
未知 \code{Iy} / 平衡残差 & \code{ixGS,irGS=1:4032} & \code{xscal=1391.7161817122487} A，\code{resscal=.0007185373089286678} A$^{-1}$ \\
\code{ag} & \code{ixg=4033:4035} & \code{[2229.9770780882873,1312.3991187777453,-4096.415637794624]}，按基列量纲 \\
剖面约束 & \code{irC=4033:4034} & \code{[1,1]} \\
CDE & \code{irp=4035} & 1；不应解释为其他自定义配置也恒为 1 \\
主动电流/回路残差 & \code{ixa,ira=4036:4047} & 各回路电流/电压尺度 \\
被动电流/回路残差 & \code{ixu,iru=4048:4572} & 电流尺度为约 515.754 或 5157.536 A；回路残差各有电压尺度 \\
\bottomrule
\end{tabular}
\end{table}

\code{ag} 第三个尺度为负值，不能把所有尺度都称为正对角归一化。\code{L.ind.irD=4035:4572} 含 1 个 CDE 和 537 个导体动态行，故保存的 \code{Jxdot} 为 $538\times4572$。\code{L.Jx} 是该状态/尺度约定下的数据，切换场未知量或静态模式必须重新组装，不能继续复用原尺寸。

\subsection{不同平衡未知量与静态模式}\label{sec:algonl}

\evimpl{} \code{algoNL} 分支由 \code{G:meq/fgeF.m:80--92,198--219,365--387} 选择，各模式的场未知量与平衡残差见表~\ref{tab:algonl}。

\begin{table}[htbp]
\centering\small
\caption{不同平衡未知量模式的场未知量、平衡残差与示例维数}
\label{tab:algonl}
\begin{tabular}{L{2.2cm}L{7.0cm}r}
\toprule
模式 & 场未知量与平衡残差 & 重组后的当前示例维数 \\
\midrule
\code{all-nl}（默认） & \code{Iy}；$S_{GS}(I_y^{\mathrm{prof}}-I_y)$ & 动态 4572，静态 4035 \\
\code{all-nl-Fx} & \code{Fx}；缩放后的 $\mathcal F(I_y^{\mathrm{prof}},I_e)-F_x$ & 动态 4830，静态 4293 \\
\code{Newton-GS} & \code{Fx}；内点 $\texttt{rhsf.*}(I_{y1}-I_y)$；边界 $F_b-\mathcal B(I_y^{\mathrm{prof}})-M_{be}I_e$ & 动态 4830，静态 4293 \\
\code{picard} & 使用 \code{fgeFpicard}，前一片 \code{Iy,Ie} 为初猜；不是上面三个残差的简单改名 & 须按该分支重新构造 \\
\bottomrule
\end{tabular}
\end{table}

静态 \code{fgst} 令当前/前片都为 \code{LXt}，固定导体电流且关闭动态导体未知量。非 Picard 模式仍调用同一个 \code{fgeF} 和 \code{solveF}，而非另一份名为 \code{fgsF} 的实现。证据：\code{G:meq/fgst.m:80--119}。静态初始化使用 \code{Ip,bp,qA} 三约束；动态输入维数 15 与静态输入维数 540（537 导体电流 + 3 约束）不能混用。

% ================= 第 \ref{sec:solver} 章 =================
\section{数值求解方法}\label{sec:solver}

本章按求解层次组织：内层磁通求解器 GSZR（第~\ref{sec:gszr} 节）、时间步初猜（第~\ref{sec:guess} 节）、外层 Newton--Krylov 迭代（第~\ref{sec:jfnk}--\ref{sec:linesearch} 节）、默认参数的有效作用域（第~\ref{sec:params} 节）、核心求解伪代码（第~\ref{sec:pseudo} 节）、库中其他数值模块的边界（第~\ref{sec:othermodules} 节）、可选求解分支（第~\ref{sec:altalgo} 节），以及 Green 核与磁面后处理（第~\ref{sec:green} 节）。

\subsection{GSZR 的分层消元与径向三对角求解}\label{sec:gszr}

\evimpl{} \code{gszrc} 要求内网格竖直点数 $N_Z$ 满足 $\log_2(N_Z+1)$ 为整数；当前 $N_Z=63=2^6-1$ 满足。它构造 $c_e=a^+$、$c_w=a^-$，以及与 $\lambda_\ell=2\cos(\ell\pi/(N_Z+1))$ 对应的谱参数 \code{two}。源码通过半角递推与反对称填充构造这些参数，\textbf{没有显式 FFT/DST 调用}。证据：\code{G:meq/gszrc.m:14--29}。

\evderiv{} 每个竖直谱参数对应径向三对角矩阵 $T_\ell$，其主对角 $d_{j\ell}=2+c_{e,j}+c_{w,j}-\lambda_\ell$，上/下对角为 $-c_{e,j},-c_{w,j}$。从末行向前消元，保存
\begin{equation}\label{eq:elim}
cr_{j\ell}=\frac{c_{e,j}}{d_{j+1,\ell}},\qquad d_{j\ell}\leftarrow d_{j\ell}-cr_{j\ell}\,c_{w,j+1},\qquad cq_{j\ell}=d_{j\ell}^{-1},\qquad cs_{j\ell}=c_{w,j}\,cq_{j\ell} .
\end{equation}
这正是 \code{G:meq/gszrc.m:30--37}。C 宏 \code{GSP} 对右端 $p$ 先作 $p_j\leftarrow p_j+cr_{j\ell}p_{j+1}$，再作 $q_0=cq_{0\ell}p_0$、$q_j=cq_{j\ell}p_j+cs_{j\ell}q_{j-1}$；\code{GSU} 把所得径向向量加回对应竖直行。证据：\code{Q/libmeq/gszr.c:7--17}。

GSZR 先以 $\texttt{cx}_j\cdot J_y$ 初始化内部源并加入边界项，再按竖直索引的 2 倍层级消元，反向重建偶/奇行；这是一套固定消元过程，没有残差容差迭代。$\texttt{cx}=\texttt{rhsf}/2$，内/外径向边界贡献系数分别 $\texttt{ci}=\texttt{cw\_first}/2$、$\texttt{co}=\texttt{ce\_last}/2$。证据：\code{Q/libmeq/gszr.c:28--121}、\code{G:meq/gszrc.m:18--20,38}。

\evimpl{} C 函数最后的 \code{dz} 参数是以网格索引差分表示的竖直位移修正，不是 $h_Z$。内点附加 $(F_{i+1}-F_{i-1})\,\texttt{dz}/2$，端点使用单边差分；\code{meqFx} 传入 \code{L.idzx*位移}，常规未要求位移时为零。证据：\code{Q/libmeq/gszr.c:124--139}、\code{G:meq/meqFx.m:115--122}。MEX 用 \code{Jy} 尺寸推导输出 $n_Z\times n_R\times ntot$，边界检查以实际 \code{CHECK\_NUMEL} 为准，不能使用其第 13 行含笔误的注释。证据：\code{Q/mexc/gszrmex.c:25--41,53--66}。

\subsection{时间步初猜及真空预测}\label{sec:guess}

\evimpl{} \code{fgetk_implicit} 支持 \code{vacuum}、\code{extrap}、\code{previous} 初猜；多候选时逐个计算残差平方和并选最小者。当前 MAT 为 \code{initguess=\{'vacuum'\}}。真空预测只更新导体电流，暂时忽略等离子体增量耦合：
\begin{equation}\label{eq:vacpred}
\bm{i}_e^{(0)}=(M_{ee}+hD_R)^{-1}\left(M_{ee}\bm{i}_e^{\,n}+h\begin{bmatrix}\bm{v}_a^{\,n+1}\\ \bm{0}\end{bmatrix}\right),
\end{equation}
其余未知量沿用上一解。外推候选为 $\bm{x}^{(0)}=\bm{x}^n+h\dot{\bm{x}}^n$，前片候选为 $\bm{x}^n$。这里的``真空''只是初猜构造，不表示后续求解忽略等离子体。证据：\code{G:meq/fgetk_implicit.m:13--52,81--99}。

\subsection{Newton--Krylov 与残差范数}\label{sec:jfnk}

\evbg{} Newton--Krylov 与无显式 Jacobian 方法的一般理论背景见文献~\cite{Ajfnk}；本节的参数、分支与状态回退仍以所列源码为准。

\evimpl{} \code{solveF} 对 \code{algoF='broyden'} 调用 Broyden，其余进入 \code{newton_solver}；当前为 \code{jfnk}。第 $k$ 次 Newton 迭代求
\begin{equation}\label{eq:newton}
J_k\,\delta\bm{x}_k=-\bm{\mathcal R}(\bm{x}_k),\qquad \bm{x}_{k+1}=\bm{x}_k+\alpha_k\,\delta\bm{x}_k .
\end{equation}
\code{evaluate_F} 的目标函数是 $v(\bm{x})=\tfrac12\sum_i\mathcal R_i^2$；实际主收敛判据为 $\|\bm{\mathcal R}\|_2\le\texttt{tolF}$，并检查迭代计数与解中无 NaN。虽然 \code{solveF} 的参数注释提到最大分量，实际 \code{newton_solver} 用二范数，本文以执行代码为准。证据：\code{G:meq/solver/solveF.m:47--75,105--109}、\code{G:meq/solver/evaluate_F.m:57--61}、\code{G:meq/solver/newton_solver.m:87--105,130--167,279--282}。

\evimpl{} 当前 JFNK 的 Jacobian 乘向量为单边差分，且扰动不按 $\|x\|/\|v\|$ 自适应归一：
\begin{equation}\label{eq:jvec}
J_k\bm{v}\approx\frac{\bm{\mathcal R}(\bm{x}_k+\epsilon_d\bm{v})-\bm{\mathcal R}(\bm{x}_k)}{\epsilon_d} .
\end{equation}
若计算得到非有限值，\code{robust_on} 使 $\epsilon_d$ 连续除以 10，直到有限或低于 \code{epsilon_d_tol=10*eps}。这些额外评估不完整计入 \code{nfeval}，所以统计值不应当作精确的残差调用总数。证据：\code{G:meq/solver/evaluate_F.m:48--50,64--108}。

\textbf{重要实现细节。} 虽然 MAT 保存 \code{L.P.epsilon_d=1e-8}，\code{fgetk_implicit} 构造传给 \code{solveF} 的参数列表时\textbf{没有传此字段}，也没有传 \code{epsilon_res}。因此该调用路径使用 \code{solveF} 默认 $\epsilon_d=\sqrt{\mathrm{eps}}$（双精度约 $1.49\times10^{-8}$）与内层绝对残差容差 $10^{-12}$。不能把 MAT 中每个参数都当作已经传到最底层。证据：\code{G:meq/fgetk_implicit.m:54--79}、\code{G:meq/solver/solveF.m:64,71,85}。

\subsection{右预条件 GMRES、Arnoldi 与降维}\label{sec:gmres}

\evimpl{} 自定义 \code{gmres_jfnk} 对 $JP^{-1}$ 构造 Krylov 空间，$\beta=\|\bm{b}\|_2$、$\bm{v}_1=\bm{b}/\beta$，每步先算 $\bm{w}=J(P^{-1}\bm{v}_j)$，然后用\textbf{修正 Gram--Schmidt} 逐列正交化：
\begin{equation}\label{eq:arnoldi}
h_{ij}=\bm{v}_i^T\bm{w},\quad \bm{w}\leftarrow\bm{w}-h_{ij}\bm{v}_i\quad(i=1,\ldots,j),\qquad h_{j+1,j}=\|\bm{w}\|_2 .
\end{equation}
以 $\bar H_j$ 为上 Hessenberg 矩阵，求最小二乘 $\min_{\bm{y}}\|\bar H_j\bm{y}-\beta\bm{e}_1\|_2$，最后 $\delta\bm{x}=P^{-1}V_j\bm{y}$。零右端立即返回零方向。证据：\code{G:meq/solver/gmres_jfnk.m:24--52,102--148}。

\code{sim} 每轮直接最小二乘；\code{giv} 逐次 Givens 旋转；当前 \code{aya} 递推逆 Hessenberg 子块 $R$ 和 $\bm{u}$，以 $c=(1+\bm{u}^T\bm{u})^{-1/2}$ 估计残差 $\beta c$，最后 $\bm{y}=\beta c^2R\bm{u}$。终止条件是内部\textbf{绝对}残差小于 \code{epsilon_res}，或达到 \code{min(mkryl,n_reduced)}；这一函数不是每步自动重启的 GMRES。\code{nrestarts=5} 指 Newton 方向失败后的外层重启，不能混称 GMRES restart。证据：\code{G:meq/solver/gmres_jfnk.m:54--114}、\code{G:meq/solver/invert_jacobian.m:24--27,47--74}。

\evimpl{} 有些残差行只含一个对角元；行掩码将系统排列成
\begin{equation}\label{eq:reduce}
J=\begin{bmatrix}J_1 & J_2\\ 0 & D\end{bmatrix},
\end{equation}
先解 $\delta\bm{x}_2=\bm{b}_2/D$，再解 $J_1\delta\bm{x}_1=\bm{b}_1-J_2\delta\bm{x}_2$。矩阵自由分支用一次 $J\delta\bm{x}_2$ 获得右端修正。证据：\code{G:meq/solver/reduce_system.m:5--10,35--71}。启用 \code{userowmask} 允许优化，不代表每个平衡都实际有可消除的行。

\evimpl/\evdata{} 当前 \code{prectype='inv'} 对线性化残差 Jacobian 显式取逆；动态行加入 $J_{\dot x}/h$，即
\begin{equation}\label{eq:prec}
J_{\mathrm{pre}}=J_x+E_DJ_{\dot x}/h,\qquad P^{-1}=\operatorname{inv}(J_{\mathrm{pre}}),
\end{equation}
这里 $E_D$ 将 538 动态行嵌回 4572 行。\code{ilu} 分支使用稀疏不完全 LU 三角求解；\code{user} 可使用函数句柄、矩阵或标量；直接求解分支无须预条件。证据：\code{G:meq/fgepre.m:6--18}、\code{G:meq/solver/preconditioner_init.m:30--80}。自定义 RL 初始化建立 \code{H.Prec}，普通 \code{iterFGE} 的重建块被注释，故复用这个预条件器；即使装置文件写 \code{precupdate='time'}，这条自定义普通步也不自动落实该时间更新策略。证据：\code{V/iterFGE.m:69--77}。

\subsection{线搜索、随机重启与接受/回退}\label{sec:linesearch}

\evimpl{} Newton 默认方向步长为 $1-\texttt{relax}$，当前 \code{relax=0}。以 $v_0=\tfrac12\|\mathcal R(\bm{x})\|^2$ 和近似方向导数 $v'_0=-2v_0$，采用
\begin{equation}\label{eq:armijo}
v(\bm{x}+\alpha\,\delta\bm{x})<(1+\tau_k)v_0+10^{-4}\alpha\,v'_0,\qquad \tau_k=(2k+1)^{-2},\quad k=0,1,\ldots .
\end{equation}
这里 $\tau_k$ 按代码在 Newton 第 1 次使用 $\tau_0$，允许早期残差增加；不能称其每次都保证单调下降。线搜索最多 10 次，NaN 时步长减半；后续用带指定原点导数的三次插值预测，太小或变化不足时改成前一试步的一半。\code{mina=.01} 是触发重新减半的阈值，\textbf{不是最终步长的硬下限}。证据：\code{G:meq/solver/newton_solver.m:72--84,138--160}、\code{G:meq/solver/linesearch.m:62--99,115--196,206--230}。

方向无法接受时，外层最多允许配置的 5 次重启，随机候选为 $(1+10^{-3}\mathcal N(0,1))\odot\bm{x}$；不随机时为 $(1+10^{-5})\bm{x}$。默认每次可尝试 10 个有限残差候选。接近迭代上限时会扩大局部 \code{kmax}，但最终 \code{isconverged} 仍比较原始 \code{P.kmax}，因此``循环执行超 50 步后得到小残差''与``被标记为收敛''不是同一判断。\code{lastwarn} 检出非空 warning ID 时还会回退初猜并置 \code{stopped_early=true}。证据：\code{G:meq/solver/newton_solver.m:193--282}。

\evimpl/\evdata{} \code{fgetk_implicit} 仅在 \code{solverinfo.isconverged}，或 \code{keepok \&\& res<tolok} 时接受候选并后处理；否则保留前片状态。当前保存 \code{keepok=0,tolok=.01}，故``足够好但未收敛''分支默认关闭。无论接受与否，输出仍写当前时间、电压、输入参数和求解统计；所以时间推进与物理状态成功更新必须分别检查。证据：\code{G:meq/fgetk_implicit.m:101--143}。

\subsection{默认参数的有效作用域}\label{sec:params}

默认求解参数及其有效作用域如表~\ref{tab:params} 所示。

\begin{table}[htbp]
\centering\small
\caption{默认求解参数：保存值、实际默认值与作用}
\label{tab:params}
\begin{tabular}{L{4.9cm}L{3.2cm}L{5.7cm}}
\toprule
参数 & 保存值或实际默认 & 作用与来源 \\
\midrule
\code{algoNL,algoF,algoGMRES} & \code{all-nl,jfnk,aya} & \code{D/L.P}，经 \code{fgetk_implicit:17--25,54--66} 分派 \\
\code{tolF,kmax,mkryl} & $10^{-4},50,200$ & 外层缩放残差二范数、Newton 原始上限、Krylov 上限 \\
\code{nrestarts,relax} & 5、0 & Newton 失败重启与初始步长 \\
\code{keepok,tolok,nnoc} & 0、.01、0 & 步接受宽松分支关闭；首个普通不收敛步即超出允许连续次数 \\
\code{usepreconditioner,prectype} & 1、\code{inv} & 初始化显式逆预条件器 \\
\code{precupdate} & 四版 \code{once}，v2.1 \code{time} & 标准 \code{fget} 会区分；自定义 \code{iterFGE} 重建块注释 \\
\code{anajac,userowmask} & 1、1 & 可用解析导数/掩码；不改变 \code{jfnk} 当前方向差分路径 \\
\code{L.P.epsilon_d} & $10^{-8}$，但此调用未传 & 有效值来自 \code{solveF} 的 $\sqrt{\mathrm{eps}}$ \\
\code{epsilon_res} & \code{solveF} 默认 $10^{-12}$ & 自定义 GMRES 的绝对残差阈值 \\
\code{tol,tolcs,tolq} & $10^{-5},10^{-4},10^{-3}$ & 平衡/曲面后处理等参数，不等同主 Newton 的 \code{tolF} \\
\code{implicit} & 1 & 保存配置；实际所读迭代器明确调用 \code{fgetk_implicit}，不靠此字段推断积分法 \\
\bottomrule
\end{tabular}
\end{table}

\evimpl{} 解析 Jacobian 分支也可复核：\code{fgeFJac} 先乘状态尺度，在线性场响应上传播 $\delta F$，再算
\begin{equation}\label{eq:jacprof}
\delta I_y^{\mathrm{prof}}=T_{yg}\,\delta\bm{a}_g+\frac{\partial I_y^{\mathrm{prof}}}{\partial F}\,\delta F;
\end{equation}
回路方向为 $(M_{ee}\delta I_e+M_{ey}\delta I_y)/h+D_R\delta I_e$，末尾乘残差尺度。\code{all-nl} 的平衡 Jacobian 块为
\begin{equation}\label{eq:jacblocks}
S_{GS}\Bigl[\tfrac{\partial I_y^{\mathrm{prof}}}{\partial F}M_{xy}-I\Bigr]S_{Iy},\qquad S_{GS}T_{yg}S_g,\qquad S_{GS}\,\tfrac{\partial I_y^{\mathrm{prof}}}{\partial F}M_{xe}S_{I_e} .
\end{equation}
这里 $S_{Iy},S_g,S_{I_e}$ 是 \code{xscal} 对应的状态尺度块，$S_{I_e}$ \textbf{不同于}第~\ref{sec:tdisc} 节的回路残差尺度 $S_e$。证据：\code{G:meq/fgeFJac.m:18--103}、\code{G:meq/fgeF.m:488--496}。这些导数用于可选 Newton、线性化和预条件等用途；默认 JFNK 的内层方向乘积仍按式~\eqref{eq:jvec} 计算。

\subsection{核心求解伪代码}\label{sec:pseudo}

下面伪代码表示当前 \code{all-nl + rigid_0D + jfnk + aya} 的 Git 源码主路径，省略绘图和日志；不表示不可见 Socket 服务端的内部代码。

\begin{lstlisting}[style=pseudo, caption={核心求解伪代码：残差评估与隐式时间步}, label=lst:pseudo]
residual(x, previous, current_input):
    [Iy, ag, Ia, Iu] <- unpack(x .* xscal)
    Ie <- [Ia; Iu]
    Fx <- meqFx(Iy, Ie)                    # 非齐次自由边界 GSZR
    domains, F0, F1, axis, Xpoints <- meqpdom(Fx)
    if geometry invalid: return NaN residual
    Tyg <- bfab_basis(Fx, F0, F1, domains)
    Iy_profile <- Tyg * ag
    rGS <- scaleGS .* (Iy_profile - Iy)
    rC  <- bp_constraint, qA_constraint
    rp  <- Lp*dIp/h + Rp*(Ip-Ini) + w' * Mye * dIe/h
    re  <- Mee*dIe/h + Mey*dIy/h + Re.*Ie - [Va; zeros(nu)]
    return [rGS; rC; rp; scaleE.*re]

implicit_step(previous, input):
    x0 <- vacuum circuit predictor, or best residual among configured guesses
    x <- x0; initialize/reuse preconditioner
    repeat Newton loop:
        r <- residual(x)
        if ||r||2 < tolF: stop
        Jv(v) <- [residual(x + epsilon*v) - r] / epsilon
        reduce diagonal-only rows when available
        dx <- right-preconditioned GMRES_aya(Jv, -r)
        alpha <- nonmonotone Armijo/cubic line search
        if accepted: x <- x + alpha*dx
        else: attempt configured random restarts or stop
        if warning ID set: restore x0; mark stopped_early; return
    if converged or (keepok and residual<tolok): postprocess accepted x
    else: restore previous physical state
    append current time/inputs/statistics
\end{lstlisting}

图~\ref{fig:flow} 以流程图形式给出单个隐式时间步内的求解控制流，可与清单~\ref{lst:pseudo} 对照阅读。

\begin{figure}[htbp]
\centering
\begin{tikzpicture}[
  font=\footnotesize,
  node distance=4.5mm,
  start/.style={draw, rounded corners=8pt, inner sep=5pt, fill=black!8},
  proc/.style={draw, rounded corners=1pt, align=center, inner sep=4.5pt, text width=62mm},
  dec/.style={draw, diamond, aspect=2.6, align=center, inner sep=1.5pt, text width=30mm},
  term/.style={draw, rounded corners=8pt, inner sep=5pt, fill=black!8},
  arr/.style={-{Stealth[length=2mm]}, thick},
  nlbl/.style={font=\scriptsize, text=black!70, right=1mm}
]
\node[start] (s0) {构造初猜 $\bm{x}^{(0)}$（真空电路预测等）};
\node[proc, below=of s0] (s1) {评估残差 $\bm{\mathcal R}(\bm{x}_k)$};
\node[dec, below=of s1] (d1) {$\|\bm{\mathcal R}\|_2<\texttt{tolF}$？};
\node[proc, below=of d1] (s2) {JFNK 方向：单边差分 $J\bm{v}$\\ 行掩码降维 + 右预条件 GMRES(\code{aya})};
\node[proc, below=of s2] (s3) {非单调 Armijo/三次插值线搜索};
\node[dec, below=of s3] (d2) {步接受？};
\node[proc, below=of d2] (s4) {更新 $\bm{x}_{k+1}=\bm{x}_k+\alpha_k\delta\bm{x}_k$};
\node[dec, below=of s4] (d3) {迭代超限/方向失败？};
\node[proc, right=14mm of d2, text width=34mm] (s5) {随机重启（至多 5 次，各 10 候选）};
\node[term, right=14mm of s4, text width=34mm] (t1) {回退：保留前片状态，置 \code{stopped\_early} 或重启};
\node[term, left=14mm of d1, text width=30mm] (t2) {接受候选，后处理并写统计};
\draw[arr] (s0) -- (s1);
\draw[arr] (s1) -- (d1);
\draw[arr] (d1) -- node[nlbl]{否} (s2);
\draw[arr] (s2) -- (s3);
\draw[arr] (s3) -- (d2);
\draw[arr] (d2) -- node[nlbl]{是} (s4);
\draw[arr] (d2) -- node[nlbl, above]{否} (s5);
\draw[arr] (s5) |- (s3);
\draw[arr] (s4) -- (d3);
\draw[arr] (d3) -- node[nlbl, above]{是} (t1);
\draw[arr] (d3.west) -- ++(-6mm,0) |- node[nlbl, pos=0.25, left]{否} (s1.west);
\draw[arr] (d1) -- node[nlbl, above]{是} (t2);
\end{tikzpicture}
\caption{单个隐式时间步的 Newton--Krylov 求解控制流。}
\label{fig:flow}
\end{figure}

\subsection{其他可见数值模块的边界}\label{sec:othermodules}

\evimpl{} \code{bslv} 做固定次数的块迭代 $a^{k}=a_1-Ab^{k-1}$、$b^{k}=b_1-Ba^{k}$，没有残差停止测试。误差传播为 $e_b^{k}=BAe_b^{k-1}$，因此 $\rho(BA)<1$ 是全初值渐近收敛条件的数学推导，不是代码检查。证据：\code{Q/libmeq/bslv.c:4--15}、\code{Q/include/meqh.h:68}。

\code{ipm2/ipm4} 是带预测/校正、packed Cholesky、\code{.99} 可行步长及容差/次数退出的内点内核；\code{minQ} 则是剖面极值的三点二次拟合，\textbf{不是二次规划求解器}；\code{locQ} 搜索目标剖面值的线性插值交点。证据：\code{Q/libmeq/ipm2.c:29--147}、\code{Q/libmeq/ipm4.c:21--101}、\code{Q/libmeq/minQ.c:7--68}、\code{Q/libmeq/locQ.c:18--54}。这些库函数的存在不能证明当前默认 FGE 每步调用了内点优化；所追踪的主时间步线性求解链为第~\ref{sec:gmres} 节的 GMRES。

\subsection{可选 Picard、显式 Newton、Broyden 和线性化分支}\label{sec:altalgo}

\evimpl{} Picard 要求单域、$n_p=0$ 且关闭预条件器，故当前 rigid CDE 配置不能只改 \code{algoNL} 就切换。未知量改为缩放后的 $[r_C,z_C,a_g]$；内层以磁轴 Hessian 修正补偿场 $F_{RR}^2+F_{ZZ}$，使轴接近目标 $(r_C,z_C)$。这里 $F_R,F_Z$ 是补偿场系数，不能与磁通的一阶偏导混淆。令 $\delta r=r_C-r_A,\ \delta z=z_C-z_A$，每次修正系数为
\begin{equation}\label{eq:picard}
\delta F_R=-\frac{F_{RR}\delta r+F_{RZ}\delta z}{2r_C},\qquad \delta F_Z=-\bigl(F_{ZZ}\delta z+F_{RZ}\delta r\bigr).
\end{equation}
域识别后以 $I_p^{\mathrm{target}}/(T_{pDg}a_g)$ 归一化剖面电流；演化情形直接解导体隐式电路，再重算磁通。证据：\code{G:meq/meqalgoNLc.m:27--29}、\code{G:meq/fgsp.m:103}、\code{G:meq/fgeFpicard.m:35--49,135--186}。

Picard 内层磁轴偏差阈值为 $r_0\max(\texttt{tolPicard}/100,10^{-14})$，最多 \code{kinner} 次；电流/磁通固定点最多 \code{kpic} 次，实际停止量是 $|\max(F-F^{\mathrm{old}})|$，\textbf{不是} $\max|F-F^{\mathrm{old}}|$，不能替源码修正成标准无穷范数。外层残差为缩放后的 \code{[FR;FZ;ag_constraints]}，请求解析 Jacobian 会报错。证据：\code{G:meq/fgeFpicard.m:78,135,189--217}。

\evimpl{} \code{algoF='newton'} 且 \code{anajac=true} 时请求解析 Jacobian 或作用句柄；否则 \code{evaluate_F} 对显式 Newton 分支调用 \code{jacfd} 构造完整差分矩阵。单右端可选 $J\backslash b$，多右端直接分支使用一次 LU；自带 GMRES/BlockGMRES 与 MATLAB GMRES 为其他分支。\code{direct_inversion} 不接受矩阵自由句柄。证据：\code{G:meq/solver/newton_solver.m:37--43}、\code{G:meq/solver/evaluate_F.m:38--50}、\code{G:meq/solver/invert_jacobian.m:29--74}。

\evimpl{} Broyden 要求 \code{anajac=true}，且禁止与 \code{Newton-GS} 组合。该实现会初始化显式 Jacobian/逆、拆成对角部分 $D$ 与余项 $H$，并对指定网格块作低秩近似；不是完全不计算 Jacobian。逆更新模式令接受步长列为 $S=[s_i]$、残差差为 $Y=[y_i]$、$\widetilde S=[s_i-D_iy_i]$，每组 \code{group_size} 步更新
\begin{equation}\label{eq:broyden}
E=\widetilde S-HY,\qquad T=\gamma\widetilde S^TH+(1-\gamma)Y^T,\qquad V=(TY)^\dagger T,\qquad H\leftarrow H+EV .
\end{equation}
这里 $\gamma=\texttt{ratio\_gb}$，$\dagger$ 为 QR/\code{pinv} 实现的伪逆；默认分组 4、秩参数 10、$\gamma=.5$。\code{large_system} 可延迟合并更新，失败会尝试 Newton/JFNK 回退。其无导数线搜索使用
\begin{equation}\label{eq:broyls}
v(\alpha)<\tfrac12\left[(.95+\tau_k)\|F\|_2-10^{-4}\alpha^2\|d\|_2^2\right]^2,\qquad \tau_k=\frac{20}{(2k+1)^2},
\end{equation}
必须保留平方形式；不能无条件等价成范数下降。证据：\code{G:meq/fgsp.m:75--80,105--107}、\code{G:meq/solver/broyden_solver.m:45--82,197--208,274--344,571--617,683--721}。其中后备线搜索有传空 \code{rho,tau} 的路径，可靠执行需额外验证，本文不将其描述成已验证恢复机制。

\evimpl{} \code{P.lin=true} 时，\code{fgess} 以 $S\dot{\bm{x}}=K\bm{x}+U\bm{u}$ 构造隐式 Euler 离散矩阵
\begin{equation}\label{eq:lindisc}
A_d=(S-hK)^{-1}S,\qquad B_d=(S-hK)^{-1}hU,\qquad \bm{x}_{0,d}=-(S-hK)^{-1}h\,\bm{r}_0 .
\end{equation}
$S$ 是包括互感及平衡位移修正的动态矩阵（H 量纲按状态定义），$K$ 是电阻/耦合矩阵，$U$ 连接 $[\texttt{Va},\texttt{Co},\texttt{Ini},d\texttt{Co}/dt]$；这些来自 \code{L.lin}，不是任意拟合矩阵。\code{fgetkl} 实际更新
\begin{equation}\label{eq:linstep}
\bm{x}_{n+1}=A_d(\bm{x}_n-\bm{x}_{\mathrm{ref}})+B_d(\bm{u}_{n+1}-\bm{u}_{\mathrm{ref}})+\bm{x}_{0,d},\qquad
\bm{y}=C\bm{x}_{n+1}+D(\bm{u}-\bm{u}_{\mathrm{ref}})+\bm{y}_{\mathrm{ref}},
\end{equation}
不计算的字段置 NaN/false。证据：\code{G:meq/fgess.m:9--18,47--84,152--180}、\code{G:meq/fgetkl.m:7--48}。局部线性化不等于默认非线性场求解；自定义 RL 初始化器的线性路径风险见第~\ref{sec:protect} 节。

\subsection{Green 核、插值和磁面后处理}\label{sec:green}

\evimpl{} 可见 \code{greenem} 使用坐标对 $(R_1,Z_1),(R_2,Z_2)$ 形成
\begin{equation}\label{eq:green}
s=(R_1+R_2)^2+(Z_1-Z_2)^2,\qquad m=\frac{4R_1R_2}{s},\qquad x=\frac{(R_1-R_2)^2+(Z_1-Z_2)^2}{s},
\end{equation}
磁通互感核为
\begin{equation}\label{eq:g0}
G_0=\mu_0\sqrt{s}\left[(1-m/2)K(m)-E(m)\right].
\end{equation}
这里 $G_0$ 单位 H，乘电流得到磁通 Wb；代码直接以多项式加对数近似计算椭圆积分，不是训练出的代理网络或逐次三维 Maxwell 求解。定义多项式系数按常数项到高阶排列，则实际使用
\begin{equation}\label{eq:elliptic}
K=A_K(x)-B_K(x)\log x,\qquad E=A_E(x)-B_E(x)\log x,
\end{equation}
多项式系数见表~\ref{tab:elliptic}。

\begin{table}[htbp]
\centering\small
\caption{椭圆积分近似的源码多项式系数（常数项至四次项）}
\label{tab:elliptic}
\begin{tabular}{lL{9.8cm}}
\toprule
多项式 & 系数 \\
\midrule
$A_K$ & \code{[1.38629436112, .09666344259, .03590092383, .03742563713, .01451196212]} \\
$B_K$ & \code{[.5, .12498593597, .06880248576, .03328355346, .00441787012]} \\
$A_E$ & \code{[1, .44325141463, .06260601220, .04757383546, .01736506451]} \\
$B_E$ & \code{[0, .24998368310, .09200180037, .04069697526, .00526449639]} \\
\bottomrule
\end{tabular}
\end{table}

模式 0 给互感，5/6 给径向/竖直磁场响应，其余模式给坐标导数与场梯度；MEX 对所有坐标对求值，输出 $n_1\times n_2$，并允许编译条件控制 OpenMP。证据：\code{Q/genlib/src/greenem.c:5--90}、\code{Q/genlib/mexc/greenemmex.c:22--67}。这说明几何响应库的计算能力，但当前 $M_{xy}/M_{xe}/M_{by}$ 的离散预计算来源仍以第~\ref{sec:offline} 节为准。$x\to0$ 时对数/倒数奇异，不能从此函数推断已实现所有同点自感正则化；其他几何构造是否作截面平均需逐调用另证。

\evimpl{} \code{bint} 用四个给定权重作网格值插值，严格为
\begin{equation}\label{eq:bint}
F_i=c_0F_k+c_1F_{k+n_Z}+c_2F_{k+n_Z+1}+c_3F_{k+1};
\end{equation}
权重由调用方给出，此函数本身不生成权重、不检查外推范围。\code{asxy1} 的局部二次多项式为
\begin{equation}\label{eq:asxy}
F=d_0+a\xi+b\eta+\tfrac{c}{2}\xi^2+\tfrac{d}{2}\eta^2+e\xi\eta,
\end{equation}
轴/鞍点解 $[c,e;e,d]\,[\xi,\eta]^T=-[a,b]^T$，解各分量截到 $[-1,1]$；零行列式时用零逆值。它是局部二次定位，不是完整非线性优化。证据：\code{Q/libmeq/bint.c:4--11}、\code{Q/libmeq/asxy.c:247--282}。

\evimpl{} \code{rtci} 沿给定射线作磁面交点 Newton 修正
\begin{equation}\label{eq:rtci}
\delta a=\frac{F_{\mathrm{target}}-F_{\mathrm{interp}}}{dF/da},
\end{equation}
双线性估计值/导数，按域检查纠偏并截在 $\pm\texttt{dap}$，返回最大步长 \code{dan}；该 C 调用只做一轮，外层迭代由调用方组织。证据：\code{Q/libmeq/rtci.c:18--61}。\code{meqc} 构造磁面积分的径向 cubic spline 微分算子（轴导数 0、外端 not-a-knot）与角向周期 cubic spline 算子；这属于磁面后处理网格，\textbf{不是}主 GS 空间离散改用有限元/谱元。证据：\code{G:meq/meqc.m:167--180}。

\evimpl{} \code{bf3i} 可选基是规则归一磁通网格上的三函数线性插值，并对线性段精确积分；\code{bfab} 当前例则为解析多项式。证据：\code{Q/libmeq/bf3i.c:24--73}。由此应区分``预计算响应矩阵''``剖面插值表''``显式基函数参数化''和``数据驱动模型''；不能因为出现 MAT 表就称整个模拟器是黑箱代理模型。

% ================= 第 \ref{sec:workflow} 章 =================
\section{仿真工作流}\label{sec:workflow}

本章按一次完整仿真的实际执行顺序组织：运行入口与服务协议边界（第~\ref{sec:entry} 节）、参数载入与静态初始平衡（第~\ref{sec:init} 节）、初始化器的预推进（第~\ref{sec:prestep} 节）、普通 RL 时间步（第~\ref{sec:onestep} 节）、保护/终止与异常传播（第~\ref{sec:protect} 节）、reset 路径（第~\ref{sec:reset} 节），以及项目外 Python 交互的完整顺序（第~\ref{sec:python} 节）。

\subsection{运行入口和服务协议边界}\label{sec:entry}

\evimpl{} \code{R/run_hfm_server.sh:4--20,22--56} 相对自身定位 \code{matlab/R2024a}、\code{sysroot} 和 \code{socket_serve}，允许 \code{MATLAB_ROOT} 覆盖，默认 \code{LISTEN_PORT=5558}，准备动态库/资源路径后 \code{exec hfm_server.sh}。后者验证入口 \code{.p} 或 \code{.m} 存在，执行 MATLAB \code{-batch} 调用 \code{hfm_socket_server}。其 \code{LOG_FILE} 变量只被打印，没有实际重定向，不能据此声称自动写出该文件。证据：\code{R/socket_serve/hfm_server.sh:4--24}。

\evimpl{} 池启动器顺序分配 \code{BASE_PORT+i}，为每个进程写日志/PID 文件；没有就绪等待或数值状态验证。停止器按保存 PID 结束子/父进程。这是多个独立实例，不能描述为单个 PDE 的域分解并行。证据：\code{R/start_hfm_pool.sh:8--9,18--49}、\code{R/stop_hfm_pool.sh:7--13}。

\evimpl{} 项目外客户端用 UTF-8 换行文本通过 TCP 发送 \code{INIT} 加 JSON、\code{RESET} 加可选 JSON、\code{STEP} 加空格分隔电压、\code{EXIT}。初始化期望 \code{INIT_OK}，\code{RESET/STEP} 解析 JSON 观测，\code{EXIT} 仅尝试读取响应行。两行消息在客户端打包为一次 \code{sendall}；TCP 仍是字节流，不能把一次发送说成服务端收到一个原子包。服务端命令解析、异常封装、字段序列化和 reset 参数支持的细节\textbf{无法从现有服务端源码确认}。证据：\code{E/docker_socket_predictor.py:55--70,127--188}；入口产物 \code{R/socket_serve/hfm_socket_server.p}。

\subsection{参数载入和静态初始平衡}\label{sec:init}

\evimpl{} 当前 Git 环境构造签名为 \code{FGE_state_env_std(config)}；不是 README 的多位置参数签名。\code{config.L_addr,LX_addr} 必须非空；默认 \code{signeo=1.2e7}、\code{bp=q0=ICS=NaN}、\code{debug_mode=false,use_pid_z=false,z_target=0}。\code{init} 只是 \code{load(...).L/LX}、添加 solver 路径和关闭 \code{debugplot}，不重新生成装置矩阵。证据：\code{V/FGE_state_env_std.m:52--152}、\code{V/init.m:3--11}。

载入后保存 \code{L0} 与首片 \code{LX0=meqxk(LX,1)}。只有 \code{bp,q0,ICS} 中至少一项不是 NaN，才触发 \code{init_equi}。\code{ICS} 的单位为 A，实际修改 \code{LX0.Ia(1)}；\code{bp} 修改 \code{LX0.bp}，\code{q0} 修改 \code{LX0.qA}。它们不是控制电压，也不是完整压力/安全因子剖面的直接替换。证据：\code{V/my_fusion_environment.m:20--47}、\code{V/init_equi.m:27--40}。

\code{init_equi} 将动态零维 CDE 改为静态 \code{Ip} 约束，再 \code{fgsc} 重组并 \code{fgst} 求平衡；然后构造时间轴
\begin{equation}\label{eq:tgrid}
t_n=t_0+n\times10^{-3}\ \mathrm{s},\qquad n=0,1,\ldots,500,
\end{equation}
共 501 片、总长度 0.5 s。接着 \code{fgex('xl',t,L,LX0)} 内部还会调用 \code{L.codeinit=fgeinit}；\code{fgeinit} 先 \code{meqinit} 生成初始电流/电压，再以 \code{kmax=200} 做\textbf{第二次静态平衡求解}，不收敛会抛异常。证据：\code{V/init_equi.m:15--25,40,49--58}、\code{G:meq/fgex.m:14--28}、\code{G:meq/fgeinit.m:9--34}。因此不能把 \code{fgex} 简化为复制初态 501 次。

\evimpl{} 对默认零维欧姆闭合，稳态初始化可生成维持电流的前馈导体斜率。记 \code{Iaind} 为按产生 1 Wb 归一的主动电流参考向量（A），
\begin{equation}\label{eq:ffwd}
\mathcal F_{pa}=\frac{\bm{i}_y^TM_{ya}\,\bm{Iaind}}{\sum_dI_{p,d}}\ \mathrm{(Wb)},\qquad R_{p,\mathrm{eff}}^{-1}=\sum_dR_{p,d}^{-1},
\end{equation}
则
\begin{equation}\label{eq:ffwd2}
\dot{\bm{i}}_{a,0}=-\frac{R_{p,\mathrm{eff}}\sum_dI_{p,d}}{\mathcal F_{pa}}\,\bm{Iaind},\qquad
\bm{v}_{a,0}=M_{aa}\dot{\bm{i}}_{a,0}+\bm{R}_a\odot\bm{i}_{a,0},\qquad
\dot{\bm{v}}_{a,0}=\bm{R}_a\odot\dot{\bm{i}}_{a,0} .
\end{equation}
\code{meqinit} 当前保留输入 $\bm{i}_u$，将旧的 vessel 电流修正方案注释。\code{fgex} 以 $\bm{v}_a(t)=\bm{v}_{a,0}+\dot{\bm{v}}_{a,0}(t-t_0)$、$\bm{i}_a(t)=\bm{i}_{a,0}+\dot{\bm{i}}_{a,0}(t-t_0)$ 生成参考时序，并展开剖面目标、\code{Ini,signeo,Lp,Rp}。证据：\code{G:meq/fgeIaind.m:9--12,69--83}、\code{G:meq/meqinit.m:12--23,82--120}、\code{G:meq/fgex.m:33--66}。这些前馈量在后面的 \code{get_Rp} 覆写电阻之前已生成；源码没有在覆写后重新计算同一组前馈斜率。

\evdata{} 若三个重初始化参数全为 NaN，则保留载入 \code{LX} 原时间轴。例如保存的 \code{ini_13906_xpt_v2.2.mat:LX} 有 1051 片，$t_0=0.5005000712699257$ s、$t_{\mathrm{end}}=1.5505000712699257$ s，间隔约 1 ms；保存初态 $I_p\approx401603.3114$ A、$R_A\approx0.9094973661$ m、$Z_A\approx-0.0049719888$ m，$L_p\approx1.012205789\times10^{-6}$ H、$R_p\approx1.521576885\times10^{-7}\ \Omega$、$T_0=0.7242$ T·m，保存电导率幅值为 $7\times10^7$ S/m。它们是文件中初始数据，不能当作本次运行结果。参数重初始化与第一次预推进可改变首个输出；电导率覆盖发生在预步之后，只改写后续 \code{LX} 输入，不会回溯重算已生成的 \code{LY(2)} 首观测。

\subsection{初始化器会先推进一片}\label{sec:prestep}

\evimpl{} \code{iterFGE_init} 从首片建立 \code{LYp}、\code{xnl=L.LX2x(LYp)}、零 \code{xnldot}、电源缓存、保护状态及预条件器；随后以 \code{it=2} 真正调用一次 \code{fgetk_implicit}。本次预步使用内置 \code{meqps} 返回的 \code{Vact}。求完存入 \code{LY(2)} 并返回 \code{H.it=3}。证据：\code{V/iterFGE_init.m:99--165,170--229,287--304}。

构造 \code{my_fusion_environment} 在此之后才调用 \code{get_Rp(signeo)}，更新 \code{LX.signeo/SQ/Rp}；外层缓存 \code{H0}。因此初始 \code{get_obs} 读的是 \code{LY(H.it-1)=LY(2)}，不是原始 \code{LX} 第一片；第一条用户动作推进的是第三片。501 片场景理论上最多剩 499 次普通推进，且可能更早终止。证据：\code{V/my_fusion_environment.m:53--59,207--227}、\code{V/FGE_state_env_std.m:152--177}；499 是由索引范围 \code{3:501} 得到的推导。

\subsection{一个普通 RL 时间步}\label{sec:onestep}

表~\ref{tab:step} 给出一个普通 RL 时间步的各阶段；以下顺序直接对应自定义环境，不用标准批量 \code{fget} 代替。

\begin{table}[htbp]
\centering\small
\caption{普通 RL 时间步的各阶段、操作与源码位置}
\label{tab:step}
\begin{tabular}{L{2.2cm}L{6.8cm}L{4.0cm}}
\toprule
阶段 & 数据与数学操作 & 源码 \\
\midrule
输入整形 & 将动作调整到主动回路数；可选 PID 仅替换末路 & \code{V/FGE_state_env_std.m:222--277} \\
取当前片 & \code{it=H.it}，从 \code{LX} 取 \code{LXt}，复用 \code{H} 各状态 & \code{V/iterFGE.m:1--45} \\
电源调用与旁路 & 调用 \code{meqps} 更新缓存；实际令 \code{LXt.Va=Vacmd} & 同文件 \code{53--55} \\
选模式 & \code{lin} 走 \code{fgetkl}；非线性走 \code{fgetk_implicit} & 同文件 \code{61--91} \\
隐式解 & 构造初猜、残差、JFNK/GMRES、线搜索、接受或保留前解 & \code{G:meq/fgetk_implicit.m:13--143} \\
导数更新 & \code{xnldot=(xnl-xnlp)/dt}；测量差商 & \code{V/iterFGE.m:91--100}；\code{G:meq/private_/meqLYdot.m:5--11} \\
保护与存储 & \code{meqlim}，写 \code{LY(it),LYp}，检查报警与非收敛 & \code{V/iterFGE.m:103--135} \\
结束判断 & 时间片耗尽也设 \code{dobreak=true}；写回 \code{H}，\code{H.it+=1} & 同文件 \code{161--162,180--结束} \\
返回 & \code{state=H.dobreak}；从 \code{LY(H.it-1)} 提取观测 & \code{V/my_fusion_environment.m:166--227} \\
\bottomrule
\end{tabular}
\end{table}

\evimpl{} 普通 \code{iterFGE} 虽调用电源函数，但第 54 行被注释的 \code{LXt.Va=Vact} 后，第 55 行实际写入 \code{Vacmd}。而初始化预步和标准 \code{fget} 使用 \code{Vact}。同样，普通迭代器的预条件器刷新块被注释，只复用 \code{H.Prec}，标准 \code{fget} 根据 \code{precupdate} 刷新。证据：\code{V/iterFGE.m:53--55,69--77}、\code{V/iterFGE_init.m:189--206}、\code{G:meq/fget.m:198--219}。这两处差异会直接影响电压限制和非线性收敛，必须在复现实验中保留。

\subsection{保护、终止与异常传播}\label{sec:protect}

\evimpl{} 内置电源 \code{meqps} 将命令截到 \code{G.Vamin/Vamax}，延迟步数 $d_i=\operatorname{round}(\tau_i/h)$，FIFO 保存历史，$d_i=0$ 用当前值。当前 MAT \code{Vadelay=0}，幅值为 $[1499.85,175.5\times10,81]$ V 的正负界。这些限制对采用 \code{Vact} 的路径有效，普通 RL 步的旁路不施加该输出。证据：\code{G:meq/meqps.m:14--53}；\code{D/L.G.Vadelay,Vamin,Vamax}。

\evimpl{} \code{meqlim} 的热保护状态 $\Theta_i$ 单位为 A$^2$·s，更新为
\begin{equation}\label{eq:it2t}
\Theta_i^{\,n+1}=\Theta_i^{\,n}+h\,(I_{a,i}^{\,n+1})^2 .
\end{equation}
它是 $I^2t$ 累计量，不含温度、散热或 $I^2R$ 热传导。电流限值只报警，不把 \code{Ia} 裁到边界；输出 \code{Iarel,Tarel,Parel} 为归一限值距离，\code{alarm} 汇总电流、累计量及可选等式约束报警。当前保存 \code{Iamin/Iamax=$\pm$Inf,Talim=Inf}，对应这些具体限制不触发有限阈值报警，不代表其他失败判据关闭。证据：\code{G:meq/meqlim.m:7--28,44--68}；\code{D/L.G}。

终止来源有：保护报警且 \code{stoponalarm} 开启、连续非收敛次数超过 \code{P.nnoc}、求解器 \code{stopped_early}、异常以及正常达到 \code{itend}。当前 \code{nnoc=0}，一个普通非收敛步即可超过阈值。\code{get_obs.is_failure} 直接等于 \code{H.dobreak}，所以\textbf{正常耗尽预定时间轴也会被标成 failure}；不能只凭此布尔值区分数值发散和正常结束。证据：\code{V/iterFGE.m:80--85,114--135,161--162}、\code{V/my_fusion_environment.m:207--227}。

\evimpl/需验证风险 \code{iterFGE} 的 catch 置 \code{dobreak} 后仍可能继续访问尚未赋值的 \code{LYt}，因此不保证所有异常都转换成结构化观测。\code{H.it} 无条件增加，终止后继续 \code{step} 也不保证安全。\code{iterFGE_init} 虽有线性分支，末尾仍无条件保存只在非线性分支赋值的 \code{solverinfo}，存在未定义局部量路径。本文未运行这些分支，不能宣称全部模式可直接成功执行。证据：\code{V/iterFGE.m:77--103,183}、\code{V/iterFGE_init.m:193--207,304}。

\subsection{reset 的三条路径}\label{sec:reset}

\evimpl{} 参数 reset 不简单等价于从当前平衡热启动。\code{bp,q0,ICS} 为有效新值且与缓存不同，才从保存的 \code{L0,LX0} 重做 \code{init_equi}、迭代器预步及电阻计算，并更新 \code{H0}；仅电导率变化只改 \code{LX} 的电导率/电阻再恢复 \code{H=H0}；无参数变化也只是恢复 \code{H0}。所有路径会清可选 PID 历史。证据：\code{V/FGE_state_env_std.m:289--352}、\code{V/my_fusion_environment.m:63--92}。

\code{init_equi} 在静态阶段结束后检查最近一次 \code{lastwarn}，关键词涵盖 LCFS、域数、workspace、Newton/FGS 不收敛；它不能保证捕捉阶段内部的每一次警告。函数实际返回的赋值只有 0/1，注释和扫描 README 中的 2 是外层异常类别，不应写成该函数一定正常返回 2。构造时初始化警告会回退载入 \code{LX} 并设置 \code{init_error}，reset 警告则保留原 \code{LX} 并设置 \code{reset_error}；外层异常仍可抛出。证据：\code{V/init_equi.m:3,38--62,77--95}、\code{V/my_fusion_environment.m:41--47,78--92}。

\subsection{项目外 Python 交互的完整顺序}\label{sec:python}

\evimpl{} 项目内训练入口调用的是 \code{HFMSimulator}，其 \code{HFMSocketPredictor.step} 先执行外部 \code{PowerSupplyModel}，再发送 12 路电压；这层经验模型不属于目标 MEQ 物理内核。完整接口时序如图~\ref{fig:sequence} 所示。

\begin{figure}[htbp]
\centering
\begin{tikzpicture}[
  font=\footnotesize,
  node distance=3mm,
  actor/.style={draw, rounded corners=2pt, inner sep=4pt, minimum width=27mm, align=center},
  msg/.style={-{Stealth[length=1.8mm]}},
  lbl/.style={font=\scriptsize, text=black!75}
]
\node[actor] (py) at (0,0) {Python 训练环境\\\scriptsize HFMSimulator};
\node[actor, right=30mm of py] (srv) {Socket 服务端\\\scriptsize hfm\_socket\_server.p（黑箱）};
\node[actor, right=30mm of srv] (mat) {MATLAB 环境\\\scriptsize FGE\_state\_env\_std 等};
\draw[black!60] (py.south) -- (py |- 0,-7.2);
\draw[black!60] (srv.south) -- (srv |- 0,-7.2);
\draw[black!60] (mat.south) -- (mat |- 0,-7.2);
\def\ya{-1.1}\def\yb{-1.9}\def\yc{-2.7}\def\yd{-3.5}\def\ye{-4.3}\def\yf{-5.1}\def\yg{-6.1}\def\yh{-6.9}
\draw[msg] (py |- 0,\ya) -- node[lbl, above]{TCP connect / INIT（需要时）} (srv |- 0,\ya);
\draw[msg] (py |- 0,\yb) -- node[lbl, above]{RESET(signeo,bp,q0)} (srv |- 0,\yb);
\draw[msg, dashed] (srv |- 0,\yc) -- node[lbl, above]{环境 reset / 初态观测} (mat |- 0,\yc);
\draw[msg] (mat |- 0,\yd) -- node[lbl, above]{观测 JSON} (srv |- 0,\yd);
\draw[msg] (srv |- 0,\ye) -- node[lbl, above]{原始 JSON $\to$ 固定 shape 转换} (py |- 0,\ye);
\node[lbl, anchor=west, align=left, text width=50mm] at ($(py |- 0,\yf)+(0.2,0)$) {每次 step：action $\to$ 外部延迟 $\to$ 变化率限制 $\to$ 仿射电源标定（PowerSupplyModel）};
\draw[msg] (py |- 0,\yg) -- node[lbl, above]{STEP（12 路电压文本）} (srv |- 0,\yg);
\draw[msg, dashed] (srv |- 0,\yh) -- node[lbl, above]{[可选 VS PID] $\to$ iterFGE 隐式一步} (mat |- 0,\yh);
\node[lbl, anchor=north west, align=left, text width=120mm] at (py |- 0,-7.3) {返回路径：LCFS/参考/failure 包装 $\to$ reward\_fn（若配置）$\to$ terminated/truncated/info};
\end{tikzpicture}
\caption{项目外 Python 训练环境与 MATLAB 仿真内核之间的完整交互时序。}
\label{fig:sequence}
\end{figure}

客户端 reset 公共签名只暴露 \code{signeo,bp,q0}，不能把 MATLAB 类的 \code{ICS,z_target} reset 能力当作该 Python API 已提供。\code{HFMSimulator} 默认奖励为 0；可选 \code{reward_fn} 能改奖励和 \code{terminated}，奖励不是 GS 方程的一部分。证据：\code{E/hfm_predictor.py:77--103}、\code{E/hfm_simulator.py:332--413}。

% ================= 第 \ref{sec:interface} 章 =================
\section{状态与控制接口}\label{sec:interface}

\subsection{三种``状态''不可混用}\label{sec:threestates}

\evimpl{} 数值求解状态是 4572 维缩放未知量 \code{H.xnl}；完整下一步所需的迭代记忆还包括 \code{xnldot,LYp,Prec,prev_H,psstate,Tstate,nnoc,it,dt} 以及输入时序 \code{LX} 和参数 \code{L}。\code{xnl} 本身不包括时钟、控制器累计量和保护历史，因此仅传一个平衡状态不足以完全恢复闭环推进。证据：\code{V/iterFGE_init.m:287--结束}、\code{V/iterFGE.m:180--结束}。

\evimpl{} \code{my_fusion_environment.get_state()} 实际只返回 \code{H.dobreak} 这一标量；\code{XL_state_1/XL_state_2} 的展平拼接代码整段已被注释。README 所称``完整状态向量''不符合这份 Git 实现。\code{get_info} 报告的是部分字段\textbf{数量}，也不是所有数组展开后的状态维数。证据：\code{V/my_fusion_environment.m:166--202}、\code{V/FGE_state_env_std.m:363--391}；对照 \code{R/HFM_env/TOK_EXL50U/XL50U_fge/RL_ENV/README.md:75--80}。

\evderiv{} 含控制与保护的离散系统可写为
\begin{equation}\label{eq:markov}
s_{n+1}=\mathcal T(s_n,\bm{u}_n;L,LX),\qquad o_n=\mathcal H(LY_n,L),
\end{equation}
其中 $s_n$ 包含上列求解器、时钟、电源、PID、保护记忆；观测为 $o_n$。式~\eqref{eq:markov} 只是对实际数据流的数学表示，不能把 \code{get_state()} 的单个布尔量视为这个 Markov 状态。

\subsection{真实观测结构及展平顺序}\label{sec:obs}

\evimpl{} \code{XL50Utok} 指定 27 个字段，展开共 4632 个数；\code{get_obs} 另附 \code{is_failure}，因此若再把标志计作 1 个元素，原始接口共 4633 个值。\code{get_info.observation_dim=length(XL_obs)=27} 是字段数。证据：\code{V/XL50Utok.m:124--190}、\code{V/my_fusion_environment.m:207--227}、\code{V/FGE_state_env_std.m:369--378}。各字段的定义顺序与提取方式见表~\ref{tab:obs}。

\begin{table}[htbp]
\centering\small
\caption{观测字段（保持定义顺序）、元素数与定义/单位}
\label{tab:obs}
\begin{tabular}{lL{1.6cm}L{7.6cm}}
\toprule
字段 & 元素数 & 定义/单位及实际提取 \\
\midrule
\code{Ip} & 1 & 总电流 A \\
\code{R,Z} & 各 1 & \code{rA,zA}，磁轴位置 m；不是 \code{rY,zY} 电流质心 \\
\code{I_PF} & 12 & \code{Ia}，A，包含末路 VS \\
\code{Rmax,Rmin} & 各 1 & \code{rgeom(end) +/- aminor(end)}，m；不是直接返回 \code{rmax/rmin} 字段 \\
\code{aminor} & 1 & 最外磁面小半径，m \\
\code{deltal,deltau,kappa} & 各 1 & 外磁面上下三角形变、拉长比，无量纲 \\
\code{rc,zc} & 各 1 & \code{rgeom,zgeom} 最后值，几何中心 m \\
\code{FA,FB} & 各 1 & 轴/边界磁通 Wb \\
\code{Bm} & 100 & 磁探针响应 T \\
\code{nrx,nzx} & 各 1 & 66、65，从 \code{L} 取，计数 \\
\code{Fx} & 4290 & $65\times66$ 磁通按 MATLAB 列优先展平，Wb \\
\code{rx,zx} & 66、65 & 网格坐标 m \\
\code{rX,zX} & 各 6 & 固定槽 X 点坐标 m；超长截断、过短补零 \\
\code{nX} & 1 & 检出的 X 点计数；应结合截断槽数解释 \\
\code{FX} & 6 & X 点磁通 Wb \\
\code{rB,zB} & 各 32 & \code{rq,zq} 最后一列，即外磁面参数化采样，m \\
\code{lX} & 1 & 边界是否由 X 点定义/是否 diverted，逻辑量 \\
\code{is_failure} & 1，额外标志 & \code{H.dobreak}，包含正常时间结束 \\
\bottomrule
\end{tabular}
\end{table}

具体别名、外磁面末值、边界和磁通的提取见 \code{V/my_fusion_environment.m:249--428}。默认向量路径按目标长度补零/截断，缺字段保持零；部分异常打印后仍返回零或部分结果。因此零值既可能是真实物理值，也可能是缺测/失败/填充。这里没有统一的逐字段有效性掩码，不能通过零值本身判定物理状态。

\evimpl{} Python \code{_parse_observation} 再按固定向量维度转换 float64，缺失补零；只有服务端缺少 \code{lX} 时才用 \code{nX>0.5} 推断。这个回退与``存在 X 点即 diverted''并非普遍等价，已有 \code{lX} 时优先用真实字段。证据：\code{E/hfm_predictor.py:68--75}。

\subsection{动作空间和电压通道}\label{sec:action}

\evimpl{} MATLAB 环境按 $n_c=\operatorname{size}(\texttt{LX.Ia},1)=12$ 整形动作，过长截断、过短补零；该补零分支期望列向量，不能承诺任意二维行输入都安全。动作单位 V，常用命名顺序为 \code{CS,PF1,...,PF10,VS}；MAT 主动标签则是 \code{PF_001,...,PF_011,VS}，首路是否称 CS 属接口命名差异。证据：\code{V/FGE_state_env_std.m:174--177,222--231}、目标 \code{RL_ENV/README.md:57--73}、\code{D/L.G.dima}。

\evimpl{} 可选 7 维动作包装按固定复制矩阵 $P_{7\to12}\in\{0,1\}^{12\times7}$ 生成
\begin{equation}\label{eq:dup}
\bm{u}_{12}=P_{7\to12}\,\bm{u}_7=[u_1,u_2,u_2,u_3,u_3,u_4,u_4,u_5,u_5,u_6,u_6,u_7]^T .
\end{equation}
归一化包装先将 $\bm{a}\in\mathbb R^7$ 截到 $[-1,1]$，乘 \code{[1499,230,172,172,348,348,270]} V 后再复制。这是项目外策略动作映射，目标环境仍消费 12 路。证据：\code{E/preprocessing.py:31--36}、\code{E/wrappers.py:54--89}；目标离线 \code{RL_test/README_actions.md:31--45} 仅提供相应示例。

\code{HFMSimulator.action_space} 声明的 12 路界为 $\pm$\code{[1499,230,172,172,348,348,270,270,348,348,270,270]} V，但普通 \code{step} 没有按 \code{Box} 自动裁剪，只尝试广播形状。\code{Box} 是接口声明，不是执行器幅值保护。它与 MATLAB 内置界不同，必须按实际调用路径分析。证据：\code{E/hfm_simulator.py:80--89,319--330,372--379}。

\subsection{在线可选 VS 增量 PID}\label{sec:pid}

\evimpl{} \code{use_pid_z=true} 时使用 MATLAB 值类 \code{PIDController}，\code{PID_step} 返回更新后的对象；外层接回该对象，保留误差历史和累计电压。采样间隔为实际 \code{LX.t(2)-LX.t(1)}。定义控制量 $y_j\in[R,Z,I_p,I_{a,1},\ldots,I_{a,12}]$、目标 $y_j^*$，误差采用\textbf{测量减目标} $e_{j,k}=y_{j,k}-y_j^*$。参数 $K_{p,j}$ 的量纲按对应控制量确定，$T_{i,j},T_{d,j}$ 为 s，实际增量为
\begin{equation}\label{eq:pid}
c_{j,k}=K_{p,j}(e_{j,k}-e_{j,k-1})+K_{p,j}\frac{h}{T_{i,j}}e_{j,k}+K_{p,j}\frac{T_{d,j}}{h}\bigl(e_{j,k}-2e_{j,k-1}+e_{j,k-2}\bigr).
\end{equation}
这来自三点系数 $[K_pT_d/h,\,-K_p(1+2T_d/h),\,K_p(1+T_d/h+h/T_i)]$ 与历史误差点积；不是另调用一个 ODE 积分器。证据：\code{V/PID/PIDController.m:93--112,130--166,223--248}；工作树可见的 Python 对照为 \code{R/HFM_env/TOK_EXL50U/XL50U_fge/RL_ENV/PID/PIDController.py:165--227,271--296}。

令 $M\in\mathbb R^{15\times12}$ 为 \code{ctrl2du.csv} 行映射，$\bm{w}\in\mathbb R^{15}$ 为无量纲权重，$\bm{s}_k\in\mathbb R^{12}$ 为 \code{sumdu}（V），则
\begin{equation}\label{eq:pidmap}
\Delta\bm{u}_k=M^T\operatorname{diag}(\bm{w})\,\bm{c}_k,\qquad \bm{s}_k=\bm{s}_{k-1}+\Delta\bm{u}_k,\qquad \bm{u}_k=\bm{u}_k^{\mathrm{ref}}+\bm{s}_k .
\end{equation}
首次调用复制三次同一状态，所以 $e_{-2}=e_{-1}=e_0$，首步仅 $c_{j,0}=K_{p,j}h\,e_{j,0}/T_{i,j}$，没有初始比例跳变。更改权重只影响当前增量，不能清除历史累计量。源码没有输出限幅、抗积分饱和或微分滤波。证据：同上 MATLAB 类 \code{93--112,150--166}。

当前在线权重只有 $w_Z=3$ 非零；读取最近完成片的 \code{rA,zA,Ip,Ia}，将 $R,I_p,I_a$ 目标设为当前值，$Z$ 目标设为 \code{z_target}；参考电压前 11 路来自动作，末路置零，PID 后只替换末路。\code{ctrl2du.csv} 的 Z 行只映射 VS，因此本模式前 11 路保持输入，VS 由 Z 环累计控制。证据：\code{V/FGE_state_env_std.m:195--201,234--262}；工作树 \code{RL_ENV/PID/ctrl2du.csv:2--16}。各控制环参数见表~\ref{tab:pid}。

\begin{table}[htbp]
\centering\small
\caption{PID 控制环参数与映射}
\label{tab:pid}
\begin{tabular}{lL{1.4cm}L{1.2cm}L{1.2cm}L{5.6cm}}
\toprule
环 & $K_p$ & $T_i$/s & $T_d$/s & 映射与证据 \\
\midrule
R & -4000 & .015 & .005 & 输出 10/11 路，CSV 参数第 11 行 \\
Z & -250 & .010 & .006 & 输出 12 路，第 12 行；当前在线仅此环启用 \\
$I_p$ & .015 & .020 & .005 & 输出 1--7 路系数 \code{[1,.09,.09,.045,.045,.065,.065]}，第 13 行 \\
各线圈电流 & 见 \code{PID_para.csv:14--25} & .005 或 .010 & .001 & 12 路单位阵映射；VS 电流环 $K_p=0$ \\
\bottomrule
\end{tabular}
\end{table}

参数 CSV 的 $I_p$ 注释称``纯 PI''，但 $T_d=0.005\neq0$，实际控制公式~\eqref{eq:pid} 仍有微分项，本文以数值/执行代码为准。Python 类对权重长度作补零/截断，MATLAB 类没有同一修复分支；不能把两个语言版本的所有输入处理视为一致。

\evimpl{}（独立未选分支）\code{V/PID/pid_Z.m} 是另一控制器，当前类不调用它：
\begin{equation}\label{eq:pidz}
J^*=J_{k-1}+he_k,\qquad v^*=K_pe_k+K_iJ^*-K_d(Z_k-Z_{k-1})/h,
\end{equation}
输出截到上下限；截断或 $|Z|$ 超阈值时回滚保存积分，但不重算当前输出。它有条件抗饱和，不能以主 \code{PIDController} 无限幅就声称项目所有 PID 都无限幅。缺省参数结构漏 \code{Z_aim} 却直接访问该字段，缺省调用有失败路径。证据：\code{V/PID/pid_Z.m:31--80}。

\subsection{项目外电源模型的离散递推}\label{sec:psmodel}

\evimpl{} \code{E/power_supply.py} 是经验模型，按延迟、变化率限制、仿射标定的顺序计算。对第 $i$ 路输入 $u_i[k]$，延迟 $\tau_i$（s）、采样间隔 $h_{ps}$（默认 .001 s）、每步变化上限 $c_i$（V/step）、标定系数 $a_i$（无量纲）、$b_i$（V），实际为
\begin{align}
d_i&=\operatorname{round}(\tau_i/h_{ps}),\qquad u_i^d[k]=u_i[\max(0,k-d_i)], \label{eq:psdelay}\\
u_i^r[0]&=u_i^d[0],\quad u_i^r[k]=u_i^r[k-1]+\operatorname{clip}\bigl(u_i^d[k]-u_i^r[k-1],-c_i,c_i\bigr),\quad v_i[k]=a_iu_i^r[k]+b_i . \label{eq:psrate}
\end{align}
首步不受变化率限制；延迟前史取首个输入。默认延迟区间均为 0，虽代码每步可以重采样，当前数值仍是零延迟。$\bm{c}=$ \code{[239.84,28,28,28,28,28,28,28,28,28,28,12.8]} V/step；完整斜率/截距见 \code{E/power_supply.py:19--37}，例如首路 $a_1=0.8580659942,\ b_1=245.0691566557$，末路 $a_{12}=1,\ b_{12}=0$。没有额外幅值裁剪。证据：同文件 \code{42--74,96--161}。

这一仿射截距使零命令不必得到零物理电压。变化率界限制标定前的 $u^r$，标定后每步界为 $|a_i|c_i$；这是由式~\eqref{eq:psrate} 得到的性质。外部 \code{dt} 不自动验证等于 MATLAB \code{LX} 的 \code{dt}，自定义时间网格时需要实验/配置一致性检查，不能声称所有采样周期自动同步。

\subsection{Gym 观测、参考轨迹与结束标志}\label{sec:gym}

\evimpl{} \code{HFMSimulator} 把 32 点 \code{rB,zB} 组成 \code{lcfs_points}（$32\times2$），添加 hold/trajectory 参考、X 点等参考字段和 \code{failure}。默认 \code{max_steps=100}、hold 初始参考；\code{terminated=raw.failure}，\code{truncated=(step_count>=max_steps)}，奖励函数还可覆盖前者。证据：\code{E/hfm_simulator.py:114--146,230--294,300--330,372--413}；\code{C/env_default.yaml:2--12}。

按当前固定 shape，27 个 HFM 字段 4632 元素加当前 LCFS 64、参考标量 4、参考 LCFS 64、XPT 参考 38、failure 1，共 4803 个数；这是对当前 Python \code{observation_space} 的展开计数，不是内部物理自由度。参考轨迹是任务输入，不参加 MEQ 平衡残差。正常核心时间耗尽也会映射成 Python \code{failure}，Python 的 100 步截断又是另一层停止机制。

% ================= 第 \ref{sec:mapping} 章 =================
\section{实现映射}\label{sec:mapping}

\subsection{方程—函数—数据的对应索引}\label{sec:index}

表~\ref{tab:index} 的定位足以从数学流程回到实现。\code{G:} 和 \code{V/} 均严格指``关于本文档''一节给定的 Git 快照；不应点击工作树中的 \code{.p} 并把它当作带这些行号的源码。

\begin{longtable}{L{2.4cm}L{5.4cm}L{5.4cm}}
\caption{方程/计算步骤与精确源码或数据位置的对应索引}\label{tab:index}\\
\toprule
本文内容 & 方程/计算步骤 & 精确源码或数据位置 \\
\midrule
\endfirsthead
\toprule
本文内容 & 方程/计算步骤 & 精确源码或数据位置 \\
\midrule
\endhead
\bottomrule
\endfoot
入口 & MATLAB batch 与服务进程 & \code{R/run_hfm_server.sh:4--56}；\code{R/socket_serve/hfm_server.sh:4--24} \\
通信 & INIT/RESET/STEP/EXIT 行协议 & \code{E/docker_socket_predictor.py:127--188}；服务端 \code{.p} 不可见 \\
环境构造 & config 校验、默认值、选择 PID & \code{V/FGE_state_env_std.m:52--201} \\
数据载入 & L/LX，保存首片 & \code{V/init.m:3--11}；\code{V/my_fusion_environment.m:20--23} \\
参数重初始化 & 改 \code{Ia(1)}、\code{bp}、\code{qA}；两次静态解 & \code{V/init_equi.m:15--58,103--116}；\code{G:meq/fgeinit.m:9--34} \\
静态求解 & 固定 $I_e$，\code{fgeF + solveF} & \code{G:meq/fgsc.m:7--32}；\code{G:meq/fgst.m:80--119} \\
维持电流前馈 & \code{Iadot0}、\code{Va0}、\code{Vadot0} & \code{G:meq/meqinit.m:82--120}；\code{G:meq/fgex.m:33--66} \\
预推进 & \code{it=2}，返回 \code{H.it=3} & \code{V/iterFGE_init.m:182--229,287--304} \\
电阻 & $\sigma$ 幅值/径向形状、$S_Q$ 梯形积分、$R_p$ 倒数 & \code{V/my_fusion_environment.m:107--156}；\code{V/evalSQ_cp.m:11--32} \\
时间步 & 输入、求解、导数、保护、H 更新 & \code{V/iterFGE.m:45--135,161--结束} \\
状态尺度/索引 & \code{Iy/ag/Ia/Iu} 与残差块 & \code{G:meq/meqalgoNLc.m}；\code{D/L.ind,xscal,resscal} \\
离散 GS & 半径半网格五点式 & \code{G:meq/dlstc.m:18--37}；\code{D/L.dlst,rhsf} \\
GSZR 系数 & 谱参数、三对角消元 & \code{G:meq/gszrc.m:14--38} \\
GSZR 执行 & 递归消元与重建 & \code{Q/libmeq/gszr.c:7--139}；\code{Q/mexc/gszrmex.c:25--66} \\
自由边界 & $M_{by}$ 构造及在线调用 & \code{G:meq/meqc.m:243--256}；\code{G:meq/meqfbp.m:8--19} \\
外部源 & $T_{ye}$/gsxe 与磁通逆映射 & \code{G:meq/meqc.m:204--233}；\code{G:meq/meqFx.m:51--63,115--129}；\code{G:meq/fgeF.m:198--219} \\
磁通响应 & 离散单位响应组合 & \code{G:meq/meqFxjac.m:31--49}；\code{D/L.Mxy,Mxe} \\
剖面基 & \code{bfab[1,2]}、$R$/$1/R$ 分配 & \code{Q/libmeq/bfab.c:25--39,169--185}；\code{Q/libmeq/bfct.c:180--310} \\
环向场 & 平方根/小抗磁近似 & \code{G:meq/meqprof.m:25--51} \\
$b_p$/$q_A$ 约束 & 积分压力与轴 Hessian 约束 & \code{G:meq/meqagconfun.m:145--164,216--248} \\
零维 CDE & $L_p\Delta I_p/h$、$R_p(I_p-I_{\mathrm{ni}})$、$\bm{w}^TM_{ye}\Delta I_e/h$ & \code{G:meq/meqcdefun.m:59--163} \\
其他 CDE & 磁通 0D、弱式 1D、静态投影 & \code{G:meq/meqcdefun.m:168--517}；\code{G:meq/meqcdec.m:38--72} \\
导体方程 & $M_{ee}\Delta I_e/h+M_{ey}\Delta I_y/h+R_eI_e-V_e$ & \code{G:meq/fgeF.m:460--482}；\code{G:meq/fgec.m:24--45} \\
Jacobian & 链式作用、物理/数值尺度 & \code{G:meq/fgeF.m:488--496}；\code{G:meq/fgeFJac.m:18--103} \\
隐式推进 & 初猜、solveF、接受/回退 & \code{G:meq/fgetk_implicit.m:13--143} \\
非线性迭代 & JFNK/Newton/Broyden 分派 & \code{G:meq/solver/solveF.m:45--109}；\code{G:meq/solver/newton_solver.m:87--282} \\
差分 $J\bm{v}$ & 单边扰动、非有限时缩步 & \code{G:meq/solver/evaluate_F.m:48--108} \\
Krylov & Arnoldi/MGS、aya/giv/sim、右预条件 & \code{G:meq/solver/gmres_jfnk.m:24--148} \\
预条件 & $J_x+J_{\dot x}/h$、inv/ILU/user & \code{G:meq/fgepre.m:6--18}；\code{G:meq/solver/preconditioner_init.m:30--80} \\
线搜索 & 非单调 Armijo、插值与折半 & \code{G:meq/solver/linesearch.m:62--230} \\
Picard & 轴补偿、定电流归一、固定点 & \code{G:meq/fgeFpicard.m:35--217} \\
Broyden & 分组多割线、低秩压缩、回退 & \code{G:meq/solver/broyden_solver.m:45--82,274--373,485--721} \\
线性化 & $S,K,U\to A_d,B_d$；输出偏置 & \code{G:meq/fgess.m:47--84,152--180}；\code{G:meq/fgetkl.m:7--48} \\
轴/X/LCFS & 局部二次定位、边界候选选择 & \code{Q/libmeq/asxy.c:21--282}；\code{Q/libmeq/fbnd.c:13--40}；\code{G:meq/meqpdom.m} \\
保护 & 电流判界、$I^2t$ 累计、报警 & \code{G:meq/meqlim.m:7--68} \\
观测 & 27 字段、固定 shape、缺失补零 & \code{V/XL50Utok.m:124--190}；\code{V/my_fusion_environment.m:207--428} \\
控制 & 增量 PID、映射、VS 替换 & \code{V/PID/PIDController.m:93--166,223--248}；\code{V/FGE_state_env_std.m:234--262} \\
外部电源 & 延迟、变化率、仿射映射 & \code{E/power_supply.py:19--74,109--161} \\
Gym 输出 & 参考、LCFS 拼接、reward、终止 & \code{E/hfm_simulator.py:230--294,372--413} \\
PPO 训练入口 & Hydra 配置、向量环境、PPO、LR 调度 & \code{T2/train_ppo.py:417--619}；\code{T2/utils.py:7--29} \\
观测特征 & 146/188 维每步特征、历史窗口、动作历史 & \code{T2/wrappers.py:199--348} \\
动作映射 & 7D 归一化、absolute/delta、7$\to$12 复制 & \code{T2/wrappers.py:417--453}；\code{E/preprocessing.py:28--37} \\
稠密奖励 & 评分核、分量、掩码、早停 & \code{T2/rewards.py:37--359} \\
reset 扰动与参考 & 参考采样、signeo/bp/q0 扰动 & \code{T2/wrappers.py:32--183} \\
训练回调 & 奖励分量日志、checkpoint & \code{TR/callbacks_x.py:90--238} \\
去噪与蒸馏 & 残差去噪器、教师--学生蒸馏 & \code{TR/denoising.py:1--60}；\code{TR/distill_noise_student.py:1--50} \\
\end{longtable}

\subsection{文件版本、说明与实现冲突的处理}\label{sec:conflicts}

容易误读的材料与实际证据优先级的对照见表~\ref{tab:conflict}。

\begin{table}[htbp]
\centering\small
\caption{容易误读的材料与实际实现/证据优先级}
\label{tab:conflict}
\begin{tabular}{L{4.6cm}L{8.2cm}}
\toprule
容易误读的材料 & 实际实现/证据优先级 \\
\midrule
发行 README 称 socket server ``sources'' & 仅 \code{.p}；服务端内部仍未知，不能补写其类构造与 JSON 生成 \\
RL README 多位置参数构造 & Git 类实际接受 config struct；以 \code{FGE_state_env_std.m:52} 为准 \\
RL README 完整 state 展平 & 实际 \code{get_state=dobreak}，展平实现被注释 \\
\code{get_info.observation_dim=27} & 27 是字段数；4632 是字段数值元素总量 \\
README 指向 v1 & 项目外默认 registry 选择 v2.2，目标包包含多个版本 \\
\code{PID_para.csv} 的``纯 PI''注释 & $I_p$ 行非零 $T_d$ 仍进入微分公式 \\
\code{implicit=1}、\code{epsilon_d=1e-8} 等保存字段 & 先确认调用链是否读取/传递，再认定有效参数 \\
内置 \code{Vamin/Vamax} & 普通 RL 步使用 \code{Vacmd} 旁路，幅值数据不等于执行裁剪 \\
\code{is_failure=true} & 也可能是正常时间耗尽；需联合求解统计/时钟分析 \\
测试 README 性能/稳定性措辞 & 本任务未运行相应实验，不能当作测得性质 \\
\bottomrule
\end{tabular}
\end{table}

工作树的 \code{RL_test/convert_npy_to_mat.py} 仅转换离线动作数组；MDS 数据脚本、实验 CSV 和画图/视频辅助功能不在当前隐式求解调用主链。它们能提供初始化/参考数据，不证明主求解器是神经网络或经验状态转移模型。证据：该转换脚本 \code{9--25}、\code{V/init.m:3--11} 及第~\ref{sec:workflow} 节追踪的调用链。

% ================= 第 \ref{sec:properties} 章 =================
\section{数值性质与局限}\label{sec:properties}

\subsection{可由代码与代数确认的性质}\label{sec:confirmable}

\evderiv{}（条件性稳定性）对线性化系统 $S\dot{\bm{x}}=K\bm{x}$，若 $S$ 可逆，特征模态 $\dot z=\lambda z$ 的后向 Euler 放大因子为
\begin{equation}\label{eq:amp}
g(h\lambda)=\frac{1}{1-h\lambda};\qquad \operatorname{Re}\lambda\le0\ \Rightarrow\ |g|\le1 .
\end{equation}
式~\eqref{eq:amp} 说明该积分公式对线性衰减模态的 A 稳定性；不证明自由边界非线性耦合、随机重启、PID、非光滑域切换或异常回退的整个仿真稳定。若物理线性化本身有增长模态，也不应因``隐式''就声称消除其增长。离散矩阵证据：\code{G:meq/fgess.m:152--164}。

\evderiv{}（阶数）在光滑解、固定拓扑、求解足够准确条件下，内点五点算子为二阶空间局部近似，后向 Euler 为一阶时间近似。主要结论来自第~\ref{sec:fivepoint} 和 \ref{sec:tdisc} 节公式的 Taylor 展开，不是全程序的收敛阶实测。边界离散、基剖面截断、曲面插值和失败步保留旧解可改变整体误差，故不可将其统称为``全模拟器二阶精度''。

\evimpl/\evderiv{} 没有在主路径看到按 CFL 或局部误差自动选择 $h$ 的控制器。重新初始化固定 $h=0.001$ s；保留载入场景则使用其等距时间轴。批量 \code{fget} 校验多片时间等距，自定义迭代器依赖初始化得到的 \code{dt}。隐式方法不意味着任意步长都能得到可信平衡；非线性初猜、磁轴/LCFS 搜索、控制采样和条件数仍会限制可用步长。证据：\code{V/init_equi.m:49--52}、\code{G:meq/fget.m:79--100}、\code{V/iterFGE.m:45--91}。

\evimpl{} 主残差按块缩放，容差 $10^{-4}$ 没有一个通用的 SI 物理单位；特别不能解释为磁通绝对误差 $10^{-4}$ Wb 或电流误差 $10^{-4}$ A。判断物理误差应反缩放各块并分别检查平衡、电路和 CDE。\code{fgeF} 中 \code{resy=norm(Iy1-Iy)}、\code{resFx=norm(Fx1-Fx)} 的输入仍为二维矩阵，MATLAB 默认是矩阵二范数（最大奇异值），不能混同总残差向量的欧氏范数或场的 Frobenius 范数。证据：\code{G:meq/fgeF.m:723--736}、\code{G:meq/solver/evaluate_F.m:61}。

\evimpl/\evderiv{} NaN、warning ID、线搜索失败与 \code{isconverged} 都会影响状态接受；仅记录 \code{res} 小不足以解释全部退出。\code{keepok=0} 时未收敛步物理状态沿用前片，时间与输入仍更新，故会出现``时间增加但平衡状态重复''的合法代码路径。相对 Newton 步长只被打印，不是当前终止条件；内层 GMRES 达到次数上限也未独立保证方程已精确解出。

\subsection{计算复杂度与瓶颈}\label{sec:complexity}

令 $N_R=64$、$N_Z=63$、$n_y=N_RN_Z=4032$、$n_e=537$、$n=4572$、$m\le200$。表~\ref{tab:complexity} 为循环/矩阵尺寸推导的渐近代价，\textbf{不是本机计时}。

\begin{table}[htbp]
\centering\small
\caption{主要计算的渐近代价与源码依据}
\label{tab:complexity}
\begin{tabular}{L{3.6cm}L{5.8cm}L{3.6cm}}
\toprule
计算 & 时间/存储估计 & 源码依据 \\
\midrule
一次 GSZR & 约 $O(N_RN_Z\log N_Z)$；系数 $O(N_RN_Z)$，单实例临时径向向量 $O(N_R)$ & \code{Q/libmeq/gszr.c:56--121}；\code{Q/mexc/gszrmex.c:64--66} \\
显式边界 $M_{by}\bm{i}_y$ & $O(n_bn_y)$，当前 $n_b=258$ & \code{G:meq/meqfbp.m:13--19} \\
导体残差 $M_{ey}\Delta\bm{i}_y$ & $O(n_en_y)$，$M_{ee}\Delta\bm{i}_e$ 为 $O(n_e^2)$ & \code{G:meq/fgeF.m:460--475} \\
Krylov 正交化 & 一次内层最多 $O(nm^2)$；方向/基存储 $O(nm)$ & \code{G:meq/solver/gmres_jfnk.m:27--39,137--148} \\
JFNK 残差调用 & 每个方向至少一次新 F 评估；线搜索/稳健缩步另加 & \code{G:meq/solver/evaluate_F.m:87--108} \\
稠密逆预条件 & 通常 $O(n^3)$ 初始化、$O(n^2)$ 存储及应用；具体 BLAS/LAPACK 由依赖实现 & \code{G:meq/solver/preconditioner_init.m:61--74} \\
\code{gsxe=3} 磁通映射 & $O(n_xn_y+n_xn_e)$ 稠密乘法 & \code{G:meq/meqFx.m:123--129} \\
Green 坐标对构造 & 每模式 $O(n_1n_2)$，矩阵存储同阶 & \code{Q/genlib/mexc/greenemmex.c:56--67} \\
电源/PID 历史 & 保存全历史，随步数线性增长 & \code{E/power_supply.py:122--134}；可见 PID 类状态追加 \\
\bottomrule
\end{tabular}
\end{table}

\evdata/\evderiv{} v2.2 保存的稠密 \code{Jx} 单独占 167225472 字节（约 159.48 MiB），\code{Mxy} 占 138378240 字节（约 131.97 MiB）；$4572\times201$ 个双精度 Krylov 基值约 7.01 MiB。这些量还没有包括 MAT 解压对象、复制、预条件器、\code{LY} 全历史、其他互感矩阵和 MATLAB 运行开销。因而``方向差分不组装当前完整 $J$''不等于低内存程序；初始化与预计算的实际瓶颈需要运行采样确认。

\evimpl{} C 中存在单/双精度宏及 BLAS、MKL、MATLAB 线性代数适配，MEX 默认调用 \code{dgszr} 并分配 double。不能仅凭目录含 MKL/OpenMP 就承诺当前发行二进制的线程数、加速比或实时性。证据：\code{Q/include/meqh.h:13--35,53--84}、\code{Q/Makefile:30--32,62--99}、\code{Q/mexc/gszrmex.c:53--66}。GSZR 的 \code{ntot} 外循环本身是串行代码；源码注释说可并行的若干行需要独立 workspace，不是已经执行的并行实现。

\subsection{物理近似与误差来源}\label{sec:approx}

物理近似及其可确认的限制见表~\ref{tab:approx}。

\begin{table}[htbp]
\centering\small
\caption{物理近似、从代码可确认的限制及其对结果解释的影响}
\label{tab:approx}
\begin{tabular}{L{2.8cm}L{5.4cm}L{5.0cm}}
\toprule
来源 & 从代码可确认的限制 & 对解释结果的影响 \\
\midrule
轴对称平衡 & 单一 $R$--$Z$ 磁通场，有限剖面参数 & 不提供三维 MHD 模态、波传播或完整流体动量/能量演化 \\
三基剖面 & 默认一个 $P'$、两个 $TT'$ 函数 & 可表示剖面族有限，满足 $b_p$/$q_A$ 不等于任意真实剖面已复现 \\
刚性零维 CDE & 只闭合总 $I_p$；$L_p$ 输入，$R_p$ 参考场积分 & 不能解释径向电流扩散细节或温度自洽输运 \\
$q_A$ 约束 & 采用局部轴 Hessian 及小抗磁近似 & 与完整安全因子剖面求解有区别 \\
导体模型 & 525 广义被动回路对应 544 细丝；互感/电阻预计算 & 不是三维有限电导率固体的全场瞬态解 \\
空间/时间网格 & 固定网格与后向 Euler & 有空间截断、时间数值耗散；未见局部自适应细化 \\
自由边界后处理 & 候选拓扑识别、局部二次/双线性/样条 & X 点出现/消失可导致非光滑残差与插值误差 \\
Green 与几何响应 & 椭圆积分近似、离散响应矩阵 & 近奇异几何与预计算误差需专门测量 \\
经验电源模型 & 固定仿射标定与逐步限速 & 是执行器近似，不是线圈电源电子电路的完整物理模型 \\
PID & 三点差分、无主类输出抗饱和/滤波 & 误差噪声、时步和累计量影响闭环；尚无整体稳定性证明 \\
观测固定槽 & 截断、补零、边界仅 32 点 & 观测损失信息，零与缺失有歧义 \\
\bottomrule
\end{tabular}
\end{table}

表~\ref{tab:approx} 各项依据分别见第~\ref{sec:model}--\ref{sec:interface} 节对应方程及源码，不推定这些近似在所有场景下的误差大小。实验初值、保存的响应矩阵和拟合电源参数是数据依赖，但当前主状态转移仍由平衡/回路/CDE 方程构成，不应误称纯数据驱动代理模型。

\subsection{源码可见但尚未实验验证的风险}\label{sec:risks}

\evimpl/需验证 当前 JFNK 单边差分使用固定绝对扰动且有非光滑域识别；浮点相消、尺度不一致与拓扑切换可能降低有效 Jacobian 方向精度。显式逆预条件器在普通 RL 步被反复复用，状态远离初始化平衡时其有效性无法由代码保证。它们是由实现结构识别出的风险，本文没有测量失效区间。

\evimpl/需验证 $I_p=0$ 或空域的权重、$S_Q=0$ 的倒数、$T_0=0$、环向场平方根负被开方量、Hessian 退化、局部插值越界以及奇异三对角/Jacobian 均有潜在非有限值来源；应按实际调用分支检查，不可声称全局输入合法性已自动保证。源头分别为 \code{meqcdefun:81--114}、\code{evalSQ_cp:9--32}、\code{meqprof:46--51}、\code{asxy.c:275--278} 与相关消元代码。当前 MAT 初态有效不证明所有控制动作或 reset 参数都有效。

\evimpl{}（实现差异）\code{fgeF} 在 CDE 是否请求 Jacobian 的两条分支上有残差缩放写法差异；当前五版 \code{resscal(irp)=1}，对这些案例无数值区别，自定义非单位尺度时需另核。证据：\code{G:meq/fgeF.m:445--458} 与 \code{D/L.resscal}。Broyden 后备空线搜索参数、线性初始化未赋值量、异常后继续访问输出及部分 reset 参数同步路径，也不能仅凭意图断言可靠；这里记录源码路径，不修改算法或配置。

\evunknown{} 未开展网格加密、时间步减半、非线性容差扫描、能量/磁通守恒预算、与实验/其他平衡码交叉验证、闭环扰动测试或性能基准，因此无法报告全局误差、可信参数域、最大稳定步长、平均迭代次数或实时吞吐。保存 \code{LX.isconverged=1} 和文件中的小残差只是历史数据，不能充当本次验证。

PID 测试也不能替代物理验证：\code{RL_ENV/PID/test/test_PIDController.py:302--304} 的所谓集成测试让 R/Z 按目标误差自行衰减，状态更新不使用控制输出；误差下降不能证明 Tokamak 闭环稳定。\code{PID/PID_test.py:108--124} 的植物更新被注释，README 的计时为示例。本任务没有执行这些会写临时 CSV/PNG/JSON 的测试。

% ================= 第 \ref{sec:rl} 章 =================
\section{强化学习训练算法}\label{sec:rl}

\textbf{本章定位。}第 \ref{sec:intro}--\ref{sec:properties} 章描述仿真内核与状态/控制接口；本章描述项目外训练栈：基于 Stable-Baselines3（SB3）PPO 的 F1/F2 任务统一训练入口（\code{T2/}）、观测特征与动作整形、与比赛评分对齐的稠密奖励、参考轨迹与 reset 扰动构成的隐式课程、学习率与归一化调度、监控与 checkpoint，以及观测去噪和教师--学生蒸馏两个辅助算法。与第 \ref{sec:workflow}--\ref{sec:interface} 章重叠的接口行为（如 \code{is_failure} 同时表示正常时间结束）不再重复，仅回引。

本章全部结论为 \evimpl{} 项目根目录 Python 源码的直接实现，行号可直接复核，不受第 \ref{sec:intro} 章 Git 快照边界限定；除明确标注 \evbg{} 的 PPO/GAE 标准定义外，不引入外部算法描述。本章未执行任何训练或评估运行，所有维度与更新次数均为静态推导值（\evderiv{}）。

\subsection{训练入口与运行架构}\label{sec:rlentry}

\evimpl{} 统一入口 \code{T2/train_ppo.py} 由 Hydra 驱动：默认配置 \code{T2/configs/base.yaml} 与任务覆盖 \code{T2/configs/task/f1.yaml}、\code{T2/configs/task/f2.yaml} 合并，任意嵌套键可用点号在命令行覆盖（如 \code{ppo.learning\_rate=1e-4}）；\code{smoke\_rollout=true} 走单环境零动作调试路径，否则进入 \code{run\_training}。证据：\code{T2/train_ppo.py:629--634}、\code{T2/configs/base.yaml:7--13}。调试路径内嵌两处 \code{breakpoint()}（\code{T2/train_ppo.py:396,405}），直接调用会进入交互式调试器，它不是无头验证脚本。

环境由工厂函数 \code{make\_env} 逐环境构建后装入向量环境：多于一个环境用 \code{SubprocVecEnv(start\_method="spawn")}，否则退化为 \code{DummyVecEnv}。仿真服务端口按 \code{BASE\_PORT+i} 顺序分配，shot 从 \code{shot\_groups} 展开后按下标轮转指派，每环境种子为 \code{base\_seed+i}。证据：\code{T2/train_ppo.py:75--126,351--360,464--492}。F1 默认 64 环境、\code{max\_steps=300}，F2 默认 32 环境、\code{max\_steps=500}。证据：\code{T2/configs/task/f1.yaml:7--8}、\code{T2/configs/task/f2.yaml:7--8}。

单个训练环境的包装链（自内向外）为：
\begin{lstlisting}[style=pseudo, caption={单个训练环境的包装链}, label=lst:envchain]
HFMSimulator(reward_fn=SemifinalDenseReward)
  -> EpisodeResetWrapper / HoldInitialResetWrapper   # 参考轨迹与 reset 参数采样
  -> NormalizedAction7DTo12DWrapper                  # 7D 归一动作 -> 12D 电压
  -> RewardInfoWrapper                               # 奖励分量注入 info
  -> [ObservationDelayNoiseWrapper]                    # 可选：观测延迟 + 高斯噪声
  -> HistoricalOBservationWrapper                      # dict 观测 -> 定长特征向量
\end{lstlisting}
证据：\code{T2/train_ppo.py:167--212}。包装顺序有明确语义：污染包装位于动作映射之后、特征整形之前，奖励与评分始终基于真值观测，被延迟/加噪的观测只进入策略输入。证据：\code{T2/train_ppo.py:196--211}、\code{T2/wrappers.py:351--357}。

\subsection{PPO 算法与超参数}\label{sec:ppo}

\evbg{} 训练器为 SB3 的 PPO（\code{MlpPolicy}）：策略网络与价值网络各为 \code{net\_arch} 指定的 MLP（默认 $[256,256]$、tanh 激活），对角高斯策略的初始对数标准差为 \code{log\_std\_init=-1.0}。PPO 的代理目标为裁剪形式
\begin{equation}\label{eq:ppoclip}
L^{\mathrm{CLIP}}(\theta)=\hat{\mathbb E}_t\Bigl[\min\bigl(\rho_t\hat A_t,\ \operatorname{clip}(\rho_t,1-\varepsilon,1+\varepsilon)\,\hat A_t\bigr)\Bigr],\qquad \rho_t=\frac{\pi_\theta(a_t\mid s_t)}{\pi_{\theta_{\mathrm{old}}}(a_t\mid s_t)},
\end{equation}
优势估计采用 GAE($\lambda$)：
\begin{equation}\label{eq:gae}
\hat A_t=\sum_{l=0}^{\infty}(\gamma\lambda)^l\,\delta_{t+l},\qquad \delta_t=r_t+\gamma V(s_{t+1})-V(s_t).
\end{equation}
式~\eqref{eq:ppoclip}--\eqref{eq:gae} 是 PPO/GAE 的标准定义\cite{Appo,Agae}，仅用于说明配置键的含义；具体实现位于 SB3 库内，不属于本项目源码。

\evimpl{} 配置键经 \code{build\_ppo\_init\_kwargs} 逐项转为 \code{PPO.\_\_init\_\_} 关键字参数，\code{net\_arch/activation\_fn/log\_std\_init} 由 \code{build\_policy\_kwargs} 单独组装进 \code{policy\_kwargs}；\code{total\_timesteps、lr\_anneal、warmup\_updates、progress\_bar} 属训练流程参数，不进入构造器。证据：\code{T2/train_ppo.py:221--295}。主要超参数的默认值与任务覆盖见表~\ref{tab:ppo}。

\begin{table}[htbp]
\centering\small
\caption{PPO 训练默认超参数（base 与任务覆盖；未列出的任务列继承 base）}
\label{tab:ppo}
\begin{tabular}{L{2.9cm}L{2.2cm}L{2.2cm}L{2.2cm}L{3.9cm}}
\toprule
参数 & base.yaml & F1 (task/f1) & F2 (task/f2) & 说明 \\
\midrule
\code{num\_envs} & 64 & 64 & 32 & 向量环境数 \\
\code{max\_steps} & 300 & 300 & 500 & 单回合环境步上限 \\
\code{total\_timesteps} & 10000 & 32768000 & 32768000 & 总采样步数 \\
\code{learning\_rate} & 3.0e-4 & 继承 & 1.0e-4 & 初始学习率 \\
\code{n\_steps} & 512 & 256 & 256 & 单环境每次更新采样步数 \\
\code{batch\_size} & 1024 & 继承 & 继承 & 每次 epoch 的 minibatch \\
\code{n\_epochs} & 6 & 继承 & 继承 & 每批数据复用轮数 \\
\code{gamma} & 0.995 & 继承 & 继承 & 折扣因子（式~\eqref{eq:gae}）\\
\code{gae\_lambda} & 0.95 & 继承 & 继承 & GAE 衰减（式~\eqref{eq:gae}）\\
\code{clip\_range} & 0.2 & 继承 & 继承 & 式~\eqref{eq:ppoclip} 中 $\varepsilon$ \\
\code{clip\_range\_vf} & null & 继承 & 继承 & 未设值则不裁剪价值 \\
\code{normalize\_advantage} & true & 继承 & 继承 & minibatch 内归一优势 \\
\code{ent\_coef / vf\_coef} & 0.0 / 0.5 & 继承 & 继承 & 熵与价值损失系数 \\
\code{max\_grad\_norm} & 0.5 & 继承 & 继承 & 梯度裁剪 \\
\code{target\_kl} & 0.03 & 继承 & 继承 & KL 超限时提前结束 epoch \\
\code{net\_arch / activation} & [256,256] / tanh & 继承 & 继承 & 双网同形、参数独立 \\
\code{log\_std\_init} & -1.0 & 继承 & 继承 & 对角高斯初始对数标准差 \\
\code{actions.mode} & delta & absolute & absolute & 动作模式（第~\ref{sec:rlaction} 节）\\
\code{features.history\_steps} & 4 & 4 & 4 & 观测历史窗口 $H$ \\
\code{lr\_anneal / warmup\_updates} & true / 10 & 继承 & 继承 & 学习率调度（式~\eqref{eq:lrschedule}）\\
\bottomrule
\end{tabular}
\end{table}

证据：\code{T2/configs/base.yaml:16--19,36--43,56--82}、\code{T2/configs/task/f1.yaml:7--8,24,27,60--61}、\code{T2/configs/task/f2.yaml:7--8,27,30,54--56}。

\evderiv{} 每次参数更新消耗的环境步数为
\begin{equation}\label{eq:tpu}
N_{\mathrm{upd}}=n_{\mathrm{steps}}\times n_{\mathrm{envs}}
\end{equation}
（代码中的 \code{timesteps\_per\_update}）。代入任务默认：F1 为 $256\times64=16384$，F2 为 $256\times32=8192$；配 \code{total\_timesteps=32768000} 得 F1 约 2000 次、F2 约 4000 次更新。这是配置值的直接算术结果，不是实测训练曲线。

\evimpl{} 学习率调度为 warmup 加 cosine annealing 的分段函数。记 $\eta_0$ 为初始学习率、$u_w$ 为 warmup 更新数、$p\in[0,1]$ 为 SB3 传入的剩余进度（$p=1$ 起点、$p=0$ 终点）、$p_w$ 为 warmup 结束时的剩余进度（由 $u_wN_{\mathrm{upd}}/T$ 换算），则
\begin{equation}\label{eq:lrschedule}
\eta(p)=\begin{cases}
\dfrac{\eta_0}{10}+\dfrac{9\eta_0}{10}\cdot\dfrac{1-p}{1-p_w}, & p\ge p_w,\\[8pt]
\dfrac{\eta_0}{2}\Bigl(1+\cos\bigl(\pi\,\tfrac{p_w-p}{p_w}\bigr)\Bigr), & p<p_w\ \text{（仅当 \code{lr\_anneal} 开启）}.
\end{cases}
\end{equation}
即前 $u_w$ 次更新内线性升到 $\eta_0$，随后按半余弦衰减到 0；\code{lr\_anneal=false} 时 warmup 后保持 $\eta_0$。证据：\code{T2/utils.py:7--29}、\code{T2/train_ppo.py:534--555}。

\evimpl{} 向量环境外再套 SB3 \code{VecMonitor} 与 \code{VecNormalize}（默认 \code{norm\_obs=true}、\code{norm\_reward=true}、\code{clip\_obs=clip\_reward=10}、\code{epsilon=1e-8}；方差折扣率缺省取 \code{ppo.gamma}）。观测归一化按运行均值/方差 $\hat\mu,\hat\sigma^2$ 作 $x\mapsto\operatorname{clip}\bigl((x-\hat\mu)/\sqrt{\hat\sigma^2+\epsilon},-10,10\bigr)$，奖励归一化基于折扣回报的二阶矩统计并同样裁剪。它对第~\ref{sec:rlfeatures} 节的整个特征向量整体统计，不做逐物理字段分组。证据：\code{T2/train_ppo.py:493--521}、\code{T2/configs/base.yaml:85--92}。

\subsection{观测特征工程}\label{sec:rlfeatures}

\evimpl{} \code{HistoricalOBservationWrapper} 把 dict 观测转成定长向量。每步基础特征 146 维，依次为
\[
\underbrace{1+1+1+1+1}_{I_p,\ \Delta I_p,\ I_p^{\mathrm{ref}},\ lX,\ \Delta lX}
\;+\;\underbrace{64}_{\mathrm{LCFS}}\;+\;\underbrace{64}_{\Delta\mathrm{LCFS}}
\;+\;\underbrace{12}_{\bm i_{PF}\!/\,\bm i_{\max}}\;+\;\underbrace{1}_{\mathrm{coil\_margin}}\;=\;146,
\]
F2 任务再追加 42 维 XPT 特征：4 个 X 点槽各 $[x_{\mathrm{valid}},\Delta x_r,\Delta x_z,\Delta x_{fx}]$、计数 $n_X$ 与 $n_{\mathrm{strike}}$、8 个打击点槽各 $[\mathrm{valid},\Delta r,\Delta z]$，合计 $16+2+24=42$，每步共 188 维。证据：\code{T2/wrappers.py:199--212,242--312}。

策略输入为最近 $H$ 步观测特征拼接最近 5 步动作历史：
\begin{equation}\label{eq:obsdim}
d=d_{\mathrm{step}}\cdot H+d_a\cdot 5,
\end{equation}
默认 $H=4$。动作历史维度 $d_a$ 由 \code{action\_mode} 决定：\code{absolute} 存策略输出的 7 维归一化动作（与旧 checkpoint 兼容），\code{delta} 存映射后的 12 维物理电压（经 \code{info["physical\_action"]} 注入）。代入默认配置：F1 为 $146\times4+7\times5=619$（absolute）或 $146\times4+12\times5=644$（delta）；F2 为 $188\times4+35=787$ 或 $188\times4+60=812$。这些是式~\eqref{eq:obsdim} 的算术结果；\code{HFMSimulator} 原始 4803 维 dict 展开（第~\ref{sec:gym} 节）中只有上述被选字段进入策略输入。证据：\code{T2/wrappers.py:229,314--348}、\code{T2/configs/base.yaml:36--43}。

\evimpl{} reset 时观测历史以零向量左填充、动作历史以首步动作右填充，episode 起点策略会收到``一步真实特征 + 其余全零''的历史窗口。XPT 特征与奖励的 XPT 分量同构：X 点磁通偏差按 $(F_X-F_B)/(F_B-F_A)$ 归一，即相对 LCFS--轴磁通差的 $\psi_X$ 偏差；无效槽乘 0 掩码；$n_X$ 作为显式特征输入，因为比赛评分对 X 点计数有硬掩码（第~\ref{sec:rlreward} 节），策略需要感知拓扑计数。证据：\code{T2/wrappers.py:242--250,317--328}。

\subsection{动作空间与电压映射}\label{sec:rlaction}

\evimpl{} 策略输出 $\bm a\in[-1,1]^7$，\code{NormalizedAction7DTo12DWrapper} 提供两种模式。\code{absolute}（绝对电压）模式为
\begin{equation}\label{eq:absmode}
\bm u_7=\operatorname{clip}(\bm a,-1,1)\odot\bm s,\qquad
\bm s=[1499,172,348,270,348,270,230]\ \mathrm V ;
\end{equation}
\code{delta}（增量电压）模式维护绝对电压累加器，reset 时清零：
\begin{align}
\bm u_7^{(k)}&=\bm u_7^{(k-1)}+\operatorname{clip}\bigl(\bm a^{(k)},-1,1\bigr)\odot\Delta\bm u_{\max},\label{eq:deltamode}\\
\Delta\bm u_{\max}&=0.16\times[1499,175,175,175,175,175,80]\ \mathrm V .\notag
\end{align}
两模式随后统一经固定索引 \code{[0,1,1,2,2,3,3,4,4,5,5,6]} 复制为 12 路（即第 \ref{sec:action} 节式~\eqref{eq:dup} 的 $P_{7\to12}$，与 \code{ACTION\_7D\_TO\_12D\_INDEX} 一致），转换后的 12 维物理电压既发给底层环境，也注入 \code{info["physical\_action"]} 供动作历史记录。证据：\code{T2/wrappers.py:19,417--453}、\code{E/preprocessing.py:28,31--37}。

两点边界需要保留。第一，\code{T2} 版的标定向量 $\bm s$ 与 \code{E/preprocessing.py} 文档值 \code{[1499,230,172,172,348,348,270]} 不同（通道顺序有差异），属实现版本差异，复现实验时不可混用两套标定。第二，$\Delta\bm u_{\max}$ 的数值 $[239.84,28,28,28,28,28,28,12.8]$ V 恰与第~\ref{sec:psmodel} 节经验电源模型的逐步变化率上界一致，但两者作用位置不同：前者限制策略\textbf{命令}的每步增量，后者限制\textbf{电源模型内部}的电压变化；普通 RL 时间步旁路内置 \code{Vamin/Vamax}（第 \ref{sec:onestep} 节），动作映射是训练栈对执行器幅值的唯一主动整形。

\subsection{与比赛评分对齐的稠密奖励}\label{sec:rlreward}

\evimpl{} 奖励类 \code{SemifinalDenseReward}（有状态，每个环境实例持有一个）按公开半决赛评分脚本设计。评分核为分段函数：记误差 $\varepsilon\ge0$、阈值 $\varepsilon_{\max}$，下界模式为
\begin{equation}\label{eq:scorelb}
s(\varepsilon)=\max\bigl(0,\ 1-\varepsilon/\varepsilon_{\max}\bigr),
\end{equation}
非下界模式为
\begin{equation}\label{eq:scoreulb}
s(\varepsilon)=\begin{cases}
2-\varepsilon/\varepsilon_{\max}, & 0\le\varepsilon<\varepsilon_{\max},\\[2pt]
\exp\bigl(1-\varepsilon/\varepsilon_{\max}\bigr), & \varepsilon\ge\varepsilon_{\max}.
\end{cases}
\end{equation}
非有限误差一律记 0。证据：\code{T2/rewards.py:37--65}。各误差分量、阈值与任务权重见表~\ref{tab:reward}。

\begin{table}[htbp]
\centering\small
\caption{稠密奖励的误差分量、阈值与任务权重}
\label{tab:reward}
\begin{tabular}{lL{5.2cm}L{1.6cm}L{1.2cm}L{1.2cm}}
\toprule
分量 & 误差定义 & $\varepsilon_{\max}$ & F1 权重 & F2 权重 \\
\midrule
\code{ip\_score} & $\lvert I_p-I_p^{\mathrm{ref}}\rvert/\lvert I_p^{\mathrm{ref}}\rvert$ & 0.05 & 0.3 & 6 \\
\code{lcfs\_score} & LCFS 与参考的逐点距离 RMS，cm & 3.0 & 0.7 & 6 \\
\code{main\_x\_dist\_score} & 主 X 点（槽 1、2）与参考距离 RMS，cm & 3 & --- & 4 \\
\code{main\_x\_psi\_score} & $(F_X-F_B)/(F_B-F_A)$ 的 RMS & 0.003 & --- & 4 \\
\code{second\_x\_dist\_score} & 次级 X 点（槽 0、3）距离 RMS，cm & 5 & --- & 3 \\
\code{second\_x\_psi\_score} & 同上归一磁通偏差的 RMS & 0.005 & --- & 4 \\
\code{strike\_dist\_score} & 打击点与参考距离 RMS，cm & 10 & --- & 3 \\
\bottomrule
\end{tabular}
\end{table}

证据：\code{T2/rewards.py:137--230}；权重 \code{T2/configs/task/f1.yaml:41--42}、\code{T2/configs/task/f2.yaml} 权重段。

\evimpl{} 奖励计算前先处理两类终止。观测 \code{failure}$>0.5$ 时直接返回 \code{failure\_penalty}（默认 $-5$）；线圈电流越限判据为
\begin{equation}\label{eq:coillimit}
\max_i\frac{\lvert I_{PF,i}\rvert}{I_{\max,i}}\ge1,\qquad \bm I_{\max}=[45000,\ \underbrace{14000,\ldots,14000}_{10},\ 4000]\ \mathrm A,
\end{equation}
触发时置线圈掩码为 0、返回 \code{coil\_current\_exceed\_limit\_penalty}（默认 $-5$），两种情况都把 \code{terminated} 置真并提前返回。证据：\code{T2/rewards.py:12,106--131}。

未触发早停时，训练默认模式的单步奖励为
\begin{equation}\label{eq:densereward}
r_t=\frac{\sum_j w_j\,s_j}{\sum_j w_j}\cdot m^{\mathrm{soft}}_{ip}\cdot m_{\mathrm{topology}},
\end{equation}
其中 $w_j$ 来自 \code{reward\_weights} 配置。两个掩码：$I_p$ 绝对误差 $\le5\times10^4$ A 时 $m^{\mathrm{soft}}_{ip}=1$，在 $5\times10^4$--$10^5$ A 区间线性降到 0，超出为 0；位形类型 $lX$ 与目标不一致时 $m_{\mathrm{topology}}=0.5$。证据：\code{T2/rewards.py:232--248,277--309}。

\evimpl{} \code{use\_competition\_score=true} 时改用比赛单步得分直接作奖励：
\begin{align}
S_{F1}&=0.3\,s_{ip}+0.7\,s_{lcfs},\label{eq:compf1}\\
S_{F2}&=\bigl(6s_{ip}+6s_{lcfs}+4s_{mXd}+4s_{mX\psi}+3s_{sXd}+4s_{sX\psi}+3s_{str}\bigr)/30,\label{eq:compf2}
\end{align}
再乘硬掩码 $m_{ip}^{\mathrm{hard}}\cdot m_{\mathrm{coil}}\cdot m_{\mathrm{topology}}$；$m_{ip}^{\mathrm{hard}}$ 在 $\lvert\Delta I_p\rvert>5\times10^4$ A 时为 0（软衰减版只用于训练模式），F2 的 X 点分量还乘 $n_X=4$ 硬计数掩码和对应槽位的有效性比例。两套计分使用不同的评分核模式：比赛得分固定用式~\eqref{eq:scorelb}（下界），训练分量按 \code{score\_lower\_bound}（默认 false）选式~\eqref{eq:scoreulb}，后者在超阈值后指数衰减但不为零。训练奖励、比赛得分与各掩码的具体耦合不可凭直觉互换。证据：\code{T2/rewards.py:144--152,159--230,250--263,277--278}、\code{T2/configs/task/f1.yaml:38,57}。

\evimpl{} 各分量经 \code{RewardInfoWrapper} 注入 \code{info["reward\_components"]} 供监控回调使用；奖励函数可覆盖 \code{terminated}，但它不参与 MEQ 平衡残差，也与第 \ref{sec:python} 节描述的接口一致。证据：\code{T2/wrappers.py:186--196}、\code{T2/rewards.py:314--359}。

\subsection{参考轨迹与 reset 扰动}\label{sec:rlreset}

\evimpl{} \code{reference\_mode=trajectory} 时，每个 episode 由 \code{EpisodeResetWrapper} 从参考库随机抽一条参考（仅一条时固定取 0），并向底层环境注入 \code{reference} 与 \code{reset\_params}；\code{hold\_initial} 时由 \code{HoldInitialResetWrapper} 强制初值保持参考。扰动采样器 \code{ResetPerturbationSampler} 按 shot 的 \code{reset\_defaults} 对 \code{signeo/bp/q0} 做缩放型扰动（默认值乘 $[\mathrm{lo},\mathrm{hi}]$ 内均匀采样）或增量型扰动（加 $[\mathrm{lo},\mathrm{hi}]$ 内均匀增量）；\code{apply\_mode=all} 时三个参数同时扰动，\code{one} 时每 episode 随机选其一；\code{fix\_reset\_params=true} 时同进程内 episode 间复用首次采样。两个采样器种子分别偏移 $+50000$（参考）与 $+10000$。F1/F2 任务配置均为缩放型 $[0.95,1.05]$、\code{apply\_mode=all}、\code{fix\_reset\_params=true}。证据：\code{T2/wrappers.py:32--183}、\code{T2/configs/task/f1.yaml:31--35}、\code{T2/configs/task/f2.yaml:31--35}。

由此构成隐式课程：目标位形在多组参考间切换，初始电导率、压力与安全因子参数在 $\pm5\%$ 内扰动，shot 在多进程间轮转（第~\ref{sec:rlentry} 节）。这些扰动走第 \ref{sec:reset} 节的参数 reset 路径，参数变化会触发重建静态初始平衡与初始化预步，数值开销集中在 reset 时刻而非普通步。

\subsection{观测延迟与噪声污染}\label{sec:rldelay}

\evimpl{} 可选 \code{ObservationDelayNoiseWrapper} 对策略可见观测施加 $k$ 步延迟（第 $t$ 步看到第 $\max(0,t-k)$ 步的观测快照，前史取 reset 观测）与逐键高斯噪声（\code{delay\_noise\_std} 把观测键映射到标准差，如 \code{\{Ip: 1000.0, R: 0.005\}}），模拟赛题推理侧的 \code{obs\_perturb}。默认 \code{delay\_step=0}、\code{noise=false}，即训练默认无观测污染，包装仅在显式开启时构造。污染只影响策略输入；奖励与比赛评分仍用真值观测。证据：\code{T2/wrappers.py:351--414}、\code{T2/configs/base.yaml:38--40}。

\subsection{监控、checkpoint 与恢复}\label{sec:rlops}

\evimpl{} 训练挂两个回调。\code{RewardComponentsCallback} 从每个 \code{infos} 里聚合 \code{reward\_components} 与超时计数，按 \code{num\_envs}$\times$\code{n\_steps} 的间隔追加写入 \code{reward\_components.csv}，并在 wandb 上注册同名指标。证据：\code{TR/callbacks\_x.py:24--51,90--160}、\code{T2/train\_ppo.py:576--584}。\code{PlasmaCheckpointCallback} 在每次 rollout 结束时判断 \code{n\_calls \% save\_freq == 0}（默认 \code{save\_freq=10}），保存模型及配套 \code{VecNormalize} 统计；注意其计数单位是 rollout 次数而非环境步。证据：\code{TR/callbacks\_x.py:193--238}、\code{T2/configs/base.yaml:97--100}。

\evimpl{} 训练结束后保存 \code{ppo\_<task>\_latest.zip}、\code{vecnormalize.pkl} 与均值/方差统计 \code{vecnormalize\_stats.npz}，运行配置存为 \code{configs.pkl} 供复现与 resume。恢复路径用 \code{PPO.load(..., custom\_objects=ppo\_kwargs)}：以当前训练配置覆盖已存模型的超参数（网络结构与权重随模型加载，不由配置重建），并重置时间步计数。wandb 通过 \code{sync\_tensorboard} 同步 SB3 的 TensorBoard 指标。证据：\code{T2/train\_ppo.py:455--462,562--619}、\code{T2/configs/base.yaml:103--116}。

\subsection{辅助算法：观测去噪与教师--学生蒸馏}\label{sec:rldenoise}

\evimpl{} \code{TR/denoising.py} 提供 SB3 特征提取器 \code{ResidualDenoisingFeatureExtractor}：残差 MLP（默认 2 个 256 宽 tanh 隐层）把观测映射回同维空间，
\begin{equation}\label{eq:denoise}
\tilde{\bm o}=\bm o+\gamma\, f_\phi(\bm o),
\end{equation}
输出层以零权重零偏置初始化，初始恰为恒等映射 $\tilde{\bm o}=\bm o$；$\gamma=\code{residual\_scale}$ 控制修正幅度。它保持下游 actor/critic MLP 的结构与维度不变，因此可从既有 PPO 策略复制权重后再按需微调，不必重训整个策略。证据：\code{TR/denoising.py:1--60}。

\evimpl{} \code{TR/distill\_noise\_student.py} 在其上构建监督式噪声蒸馏：教师策略读取干净归一化特征并保持冻结，学生策略在 actor/critic 之前插入可训练去噪层、读取带噪归一化特征，训练目标是对齐教师冻结的 actor/critic 输出；噪声样本由第~\ref{sec:rldelay} 节的观测污染管线生成，蒸馏本身不引入新的环境交互，也不要求仿真内核提供梯度。证据：\code{TR/distill\_noise\_student.py:1--50}。

\subsection{本章边界}\label{sec:rlboundary}

本章描述的 PPO 训练栈只通过第 \ref{sec:interface} 章的 Gym 接口与仿真内核交互；奖励、掩码与评分函数都在环境外层计算，不构成 GS 平衡残差的一部分。训练栈内部还有 \code{TR/features.py}（旧版特征向量）、\code{TR/reward\_f2.py}（旧版 F2 奖励）、\code{TR/train\_ppo.py} 与 \code{TR/train\_ppo\_x.py}（argparse 版训练入口）等历史实现，\code{T2/} 是其 Hydra 重构版；两者训练逻辑保持一致，但特征与标定细节以实际选用的入口为准，不能跨版本混用超参数与 checkpoint。本文未运行任何训练、评估或蒸馏，所有性能结论均不在本章作出。

% ================= 时间截断的理论分析 =================
\clearpage
\section{时间截断、GAE 与五百步评估：一项理论分析}\label{sec:gae-time-analysis}

\subsection{问题不是有没有用到数据，而是这些数据在估计什么}

考虑这样一个疑问：训练回合在第 500 步结束，每次只收集 256 步；前面的动作能看到较多后续结果，靠近第 500 步的动作却看不到第 501 步以后。这是否使优势估计具有位置依赖？把训练回合延长，是否能让评估前 500 步得到更一致的学习信号？

答案需要拆开。每个已收集的动作都有自己的 GAE，通常也都参与小批量训练；但不同动作的优势估计使用的实际后续长度确实不同。\textbf{长度不同不自动意味着有偏，更不自动意味着后面的动作学得更差。}关键在于：任务是否真的在第 500 步终止，截断时有没有正确自举，critic 是否准确，以及训练主要访问哪些状态。

本章不沿用前面的实现审计叙述，而以一个固定策略下的估计问题为起点，再讨论策略更新与当前实验。以下图表均为解析计算或说明性例子，\textbf{不是 HFM 训练收益或消融实验结果}。全文引用的 GAE、时间限制理论与 SB3 实现资料见文末\cite{Agae,Atimelimits,Asb3rollout,Asb3buffer}。本章只分析现有设置和实验方案，不改变训练配置或仿真器。

\subsection{先固定符号：三个时钟与两个目标}

设 $H_{\rm ep}$ 为单回合时间上限，$R_{\rm roll}$ 为每个环境每次更新前采集的步数，$T_{\rm eval}$ 为评估长度。当前 F2 未经命令行覆盖时为
\begin{equation}\label{eq:gaet-threeclocks}
 H_{\rm ep}=500,\qquad R_{\rm roll}=256,\qquad T_{\rm eval}=500.
\end{equation}
三者都是步数；这里采用不同下标，以免与历史观测窗口或大半径混淆。$g=0,1,\ldots$ 表示某个环境的累计交互索引，$\tau$ 表示该回合中\emph{已经执行}的步数，因此 $\tau=499$ 对应即将执行第 500 次动作。一次动作产生 $(s_\tau,a_\tau,r_\tau,s_{\tau+1})$。理论中的 $s$ 暂指充分的 Markov 状态；实际网络读取的是四步历史特征 $o$，不预设该特征已完全消除部分可观测性。

若控制任务在第 $T$ 步真正结束，价值是有剩余时间的量；若只是人为停止采样，价值描述的是假设不重置时的继续运行：
\begin{align}
 V^{\pi,T}_\tau(s)&=\mathbb E_\pi\!\left[\sum_{j=0}^{T-\tau-1}\gamma^j r_{\tau+j}\mid s_\tau=s\right],
 &V^{\pi,T}_T(s)&=0,\label{eq:gaet-finitevalue}\\
 V^{\pi,\infty}(s)&=\mathbb E_\pi\!\left[\sum_{j=0}^{\infty}\gamma^j r_{\tau+j}\mid s_\tau=s\right].&&\label{eq:gaet-contvalue}
\end{align}
$\pi$ 是固定策略，$\gamma\in[0,1)$ 是奖励折扣，$r$ 是该目标所使用的奖励，期望对环境随机性与策略采样取平均。第二式还假定参考输入等已纳入状态，使继续过程有明确含义。\textbf{不能一边把 500 步当成真实任务终点，一边又不加说明地使用无限时域价值。}同一物理状态在剩余 5 步和剩余 300 步时，有限时域价值可以不同；此时应向价值估计提供时间信息，或采用相应的时间相关价值函数\cite{Atimelimits}。

当前评估代码用 \texttt{model.predict} 执行动作并记录评分，到上限或结束信号即停止，并不计算 GAE。因此``让 eval 的 GAE 同分布''应改写成一个可检验目标：\textbf{让训练中前 500 步的状态--动作覆盖、价值目标和优势估计误差，更适合固定策略的 500 步评估。}若需要研究 eval 上的 GAE，应额外离线定义估计器和 critic，而不是假定评估器本来就在计算它。

\subsection{分类：真正终止、时间上限和采样窗口不能共用一个掩码}

令 $z_t$ 表示转移 $t$ 是否真正终止，$u_t$ 表示是否因时间上限被重置，$w_t$ 表示是否位于本次 rollout 的最后一条转移。定义两种不同的掩码：
\begin{equation}\label{eq:gaet-masks}
 b_t=1-z_t,\qquad
 c_t=(1-z_t)(1-u_t)(1-w_t).
\end{equation}
$b_t$ 决定是否保留下一状态价值，$c_t$ 决定是否还能在\emph{当前已收集片段}中累积下一项优势。它们取值为 0 或 1；时间上限只有在被解释为外加截断时才使用 $b_t=1$。令 $\bar V$ 为本次采样时固定的 critic，则统一写法为
\begin{equation}\label{eq:gaet-maskedrecursion}
 \delta_t=r_t+\gamma b_t\bar V(s_{t+1}^{\rm final})-\bar V(s_t),
 \qquad
 \hat A_t=\delta_t+\gamma\lambda c_t\hat A_{t+1}.
\end{equation}
$\lambda\in[0,1]$ 控制优势迹长；$s_{t+1}^{\rm final}$ 是这条转移实际到达的最终状态，\textbf{不是自动 reset 后的新初始状态}。窗口最后一项的递推在本片段停止，但其 TD 误差仍可使用下一状态价值。

\begin{table}[!ht]
\centering\small
\caption{同样是``停止''，估计语义并不相同}\label{tab:gaet-boundaries}
\begin{tabular}{L{3.2cm}L{2.3cm}L{2.3cm}L{5.0cm}}
\toprule
情形 & 自举 $b_t$ & 跨边界递推 $c_t$ & 估计含义 \\
\midrule
失败等真正终止 & 0 & 0 & 终止后价值为零；需按真实任务定义终止代价 \\
继续任务的时间上限 & 1 & 0 & 用最终状态价值补未来；不拼接新回合奖励 \\
rollout 窗口结束 & 1 & 0 & 环境保留状态，更新后继续；本次 GAE 在此收尾 \\
500 步为任务本身终点 & 0 & 0 & 对有限时域目标正确；价值需区分剩余时间 \\
普通内部转移 & 1 & 1 & 使用下一状态价值并继续累积 TD 误差 \\
\bottomrule
\end{tabular}
\end{table}

当前环境把达到步数上限标为 \texttt{truncated}；失败或奖励逻辑可令 \texttt{terminated} 为真。SB3 的向量环境保存最终观测，并在时间截断且非真正终止时标记 timeout；采样器把 $\gamma\bar V(s^{\rm final})$ 加到该转移的训练奖励，再让 GAE 在回合处断开。这与表中第二行等价，不意味着把下一回合的奖励接到旧回合上\cite{Asb3rollout,Asb3buffer}。实现核对采用公开的 SB3 2.7.0 源码；项目依赖只要求 SB3 $\geq2.0$，正式实验仍应记录实际安装版本。

\subsection{一个五百步回合，怎样穿过二百五十六步窗口}

暂时排除提前失败，并假定每个回合恰好 $H_{\rm ep}$ 步、首次采样与首次 reset 同步。转移 $g$ 剩余的可用 TD 误差数为
\begin{equation}\label{eq:gaet-support}
 K_g=\min\!\left\{
 R_{\rm roll}-(g\bmod R_{\rm roll}),\;
 H_{\rm ep}-(g\bmod H_{\rm ep})
 \right\}.
\end{equation}
$\bmod$ 表示整数余数，$K_g\geq1$ 包含当前 TD 误差。出现真正终止时，还须把到该终止的剩余长度纳入最小值。该式统计的是已收集的误差项数量，不是网络输入的历史长度，也不是每个动作后额外运行的仿真长度。

\begin{table}[!ht]
\centering\small
\caption{当前 F2 的前两次采样：每条样本都有优势，但支持长度不同}\label{tab:gaet-example}
\begin{tabular}{L{3.5cm}L{3.8cm}L{2.0cm}L{3.7cm}}
\toprule
样本位置（动作序号） & 所在片段 & $K$ & 片段结束的原因 \\
\midrule
回合 1，第 1 步 & 第一批第 1 条 & 256 & rollout 到 256 \\
回合 1，第 100 步 & 第一批第 100 条 & 157 & rollout 到 256 \\
回合 1，第 256 步 & 第一批第 256 条 & 1 & rollout 到 256 \\
回合 1，第 257 步 & 第二批第 1 条 & 244 & 回合到 500 \\
回合 1，第 499 步 & 第二批第 243 条 & 2 & 回合到 500 \\
回合 1，第 500 步 & 第二批第 244 条 & 1 & 回合到 500 \\
回合 2，第 1 步 & 第二批第 245 条 & 12 & 累计步数到 512 \\
回合 2，第 12 步 & 第二批第 256 条 & 1 & 累计步数到 512 \\
\bottomrule
\end{tabular}
\end{table}

\begin{figure}[!ht]
\centering
\includegraphics[width=\linewidth]{gae_sampling_timeline.pdf}
\caption{上：回合时钟与采样窗口时钟；竖直红虚线为 reset，灰点线为 rollout 边界。下：式~\eqref{eq:gaet-support} 对每条转移给出的 $K_g$。横轴是从零计数的累计转移索引。第二批中的 $244+12$ 并不是人为补数据，而是窗口跨过了 reset。该图按固定长度、无提前失败的假设直接计算。}\label{fig:gaet-timeline}
\end{figure}

后续窗口不会永远在同一个回合年龄结束。因为 $\gcd(500,256)=4$，两个时钟在 $\operatorname{lcm}(500,256)=32000$ 步后重新对齐，即 64 个回合或 125 个 rollout。这能改变某个年龄在不同回合遇到的窗口位置，却不能消除第 500 步始终发生的时间上限。若 32 个环境同时启动且从不提前结束，它们的这类边界还会同步；不同随机种子本身不保证时间相位不同。提前失败可打乱相位，但也引入生存选择。

\subsection{有限 GAE 的展开：边界以什么方式进入优势}

先看一个没有内部终止、长度为 $K$ 的片段。记 $q=\gamma\lambda$，有限 GAE 为
\begin{equation}\label{eq:gaet-finite}
 \hat A_t^{(K)}=\sum_{l=0}^{K-1}q^l
 \left[r_{t+l}+\gamma\bar V(s_{t+l+1})-\bar V(s_{t+l})\right].
\end{equation}
这里 $q$ 是无量纲迹衰减系数，$K$ 是当前样本到片段末端的剩余转移数，所有奖励和价值必须使用同一尺度。整理价值项，得到
\begin{align}\label{eq:gaet-expand}
 \hat A_t^{(K)}={}&-\bar V(s_t)
 +\sum_{l=0}^{K-1}q^l r_{t+l}\\
 &+\gamma(1-\lambda)\sum_{j=1}^{K-1}q^{j-1}\bar V(s_{t+j})
 +\underbrace{\gamma q^{K-1}}_{\beta_K}\bar V(s_{t+K}).\nonumber
\end{align}
$\beta_K$ 是\emph{只改变片段端点价值、固定其他项时}，该优势对端点价值的线性敏感度；中间状态的价值也以非零系数参与。$K=1$ 时，优势就是一步 TD 误差；$\lambda=1$ 时，中间价值项消去，成为 $K$ 步折扣回报加端点自举、再减当前价值。$\lambda=0$ 时所有样本都只使用当前 TD 误差，因此单纯增加可用 $K$ 不会延长优势迹。

如果只比较同一条可继续轨迹、同一策略和同一 critic 上的长短估计，则
\begin{equation}\label{eq:gaet-tail}
 \hat A_t^{(\infty)}-\hat A_t^{(K)}
 =q^K\hat A_{t+K}^{(\infty)}.
\end{equation}
这是在级数收敛时的恒等式。若 $|\delta_j|\leq D$，则差值不超过 $Dq^K/(1-q)$，其中 $D$ 是 TD 误差幅值界。\textbf{$q^K<1\%$ 只是迹的相对衰减标准，既不是相对优势误差小于 1\%，也不是实际控制误差小于 1\%。}当真实优势接近零时，谈相对误差尤其不稳健。

\subsection{偏差与方差：为什么不能简单说短估计一定更差}

把 critic 写成 $\bar V(s)=V^\pi(s)+e(s)$，$e$ 为对所选继续任务真实价值的误差。用相同片段分别代入两种价值，优势误差恰为
\begin{equation}\label{eq:gaet-valueerror}
 \Delta\hat A_t^{(K)}=-e(s_t)
 +\gamma(1-\lambda)\sum_{j=1}^{K-1}q^{j-1}e(s_{t+j})
 +\beta_K e(s_{t+K}).
\end{equation}
因此，较短 $K$ 通常具有较大的\emph{端点}误差系数，但实际总误差还取决于各处 critic 误差的大小、符号和相关性。例如所有状态都有同一个常数误差 $e(s)=c$ 时，误差反而合并为
\begin{equation}\label{eq:gaet-constant-error}
 \Delta\hat A_t^{(K)}
 =-(1-\gamma)c\,\frac{1-q^K}{1-q}.
\end{equation}
$c$ 为价值单位的常数偏移。这个反例说明，不能拿 $\beta_K$ 单独充当 critic 全部误差的测量，更不能据此断言增加 $K$ 总会降低误差。

在更理想的比较中，若状态 Markov、策略固定且 on-policy、$\bar V=V^\pi$ 精确，并采用确定的支持长度 $K$，则
\begin{equation}\label{eq:gaet-unbiased}
 \mathbb E[\delta_t\mid s_t,a_t]=A^\pi(s_t,a_t),\qquad
 \mathbb E[\hat A_t^{(K)}\mid s_t,a_t]=A^\pi(s_t,a_t).
\end{equation}
$A^\pi=Q^\pi-V^\pi$ 是该折扣目标的真实优势。后一结论来自后续策略动作平均后 TD 残差的条件期望为零；此处不把结论不加证明地推广到任意依赖未来结果的选择窗口。不同确定 $K$ 可以具有相同条件均值，但不同方差，所以\textbf{同均值与同分布不是同一要求}。精确 critic 下，增加后续项可能只增加噪声。

为看清这种可能性，构造一个只用于演示的模型：给定起始状态和动作，残差的噪声相互独立、方差均为 $\sigma_\delta^2$。则
\begin{equation}\label{eq:gaet-noisevar}
 \operatorname{Var}(\hat A_t^{(K)}\mid s_t,a_t)
 =\sigma_\delta^2\sum_{l=0}^{K-1}q^{2l}
 =\sigma_\delta^2\frac{1-q^{2K}}{1-q^2}.
\end{equation}
真实等离子体控制残差有时间相关性，还应加入 $2\sum_{i<j}q^{i+j}\operatorname{Cov}(\delta_{t+i},\delta_{t+j})$，不能从这个独立噪声例子推断实验方差。实际需要权衡的是 critic 误差、截断尾项和残差噪声，而不是把``看见的步数越多''作为唯一指标。

\subsection{代入当前参数：优势迹短，不代表控制目标只看几十步}

当前 $\gamma=0.995$、$\lambda=0.95$，所以 $q=0.94525$。迹权重半衰期 $\log(1/2)/\log q$ 约为 12.3 步，且 $q^{82}<0.01$。但奖励折扣半衰期 $\log(1/2)/\log\gamma$ 约为 138.3 步；$\gamma^{500}\simeq0.0816$。\textbf{GAE 如何混合残差，与价值函数应预测多远，是两件不同的事。}

\begin{table}[!ht]
\centering\small
\caption{同一长度 $K$ 下的三种不同权重，数值由当前参数直接计算}\label{tab:gaet-weights}
\begin{tabular}{rrrr}
\toprule
$K$ & 尾部迹因子 $q^K$ & 端点系数 $\beta_K$ & 奖励折扣 $\gamma^K$ \\
\midrule
1 & 0.945250 & 0.995000 & 0.995000 \\
12 & 0.508815 & 0.535595 & 0.941623 \\
13 & 0.480958 & 0.506271 & 0.936915 \\
20 & 0.324290 & 0.341358 & 0.904610 \\
50 & 0.059887 & 0.063039 & 0.778313 \\
82 & 0.009882 & 0.010402 & 0.662968 \\
100 & 0.003586 & 0.003775 & 0.605770 \\
244 & $1.080\!\times\!10^{-6}$ & $1.137\!\times\!10^{-6}$ & 0.294328 \\
256 & $5.495\!\times\!10^{-7}$ & $5.784\!\times\!10^{-7}$ & 0.277146 \\
\bottomrule
\end{tabular}
\end{table}

例如第 500 步的 $K=1$ 对端点误差几乎不衰减，但第 257 步的 $K=244$ 对同一端点误差的直接系数约为百万分之一。另一方面，处于第 500 步附近的状态，其价值可能包含相当重要的后续回报；不能因为优势迹在 82 步后变小，就把更远未来的价值当成已知为零。

\begin{figure}[!ht]
\centering
\includegraphics[width=\linewidth]{gae_weights_variance.pdf}
\caption{左：$q^K$、$\beta_K$ 与 $\gamma^K$ 的衰减不同，前两条接近但含义不同。右：独立同方差残差噪声例子的解析方差，极限约为 $9.39\sigma_\delta^2$；它说明更长片段可以增加噪声，并不测量 HFM 的真实噪声。坐标与曲线均由本章公式生成。}\label{fig:gaet-weights}
\end{figure}

\subsection{延长回合能改善什么：端点覆盖，而非自动消除边界}

对当前的继续任务解释，训练只到 $s_{500}$ 就 reset，会出现一种具体的价值学习困难：$s_{500}$ 可被用于自举，却没有在这个回合中产生后续动作和奖励。它的价值只能依靠相似状态的训练数据与函数泛化。若这些边缘状态与早期状态差别很大，或控制风险在 500 步后才显现，端点价值就可能不可靠。延长训练放电可以直接访问这些后续状态，改善这种覆盖，亦可惩罚``刚好撑到评估终点但已经走向失败''的动作。

然而延长 $H_{\rm ep}$ 不改变 $R_{\rm roll}=256$。以首次回合的第 500 次动作为例，$g=499$：
\begin{equation}\label{eq:gaet-lastaction}
 K_{499}(H_{\rm ep}=500)=1,\qquad
 K_{499}(H_{\rm ep}\geq512)=13.
\end{equation}
此时窗口结束于累计第 512 次转移，故从第 500 次动作到该窗口末端仅有 13 条转移。$H_{\rm ep}=600$ 和 1000 对这个样本的\emph{当次}支持长度相同，端点系数均约为 0.506，而不是趋近于零。它们仍可能通过后续训练改变 critic 的误差；那是学习效果，不能从 $K$ 的公式直接算出。

\begin{figure}[!ht]
\centering
\includegraphics[width=\linewidth]{gae_extend_episode.pdf}
\caption{首次回合、固定窗口相位下延长时间上限的几何效果。$H_{\rm ep}=600$ 和 1000 的曲线在前 500 步完全重合，因为都先遇到第 512 步的 rollout 边界。上图是实际残差支持长度，下图是\emph{仅对端点价值误差}的敏感度。图不意味着两个设置训练后的 critic 或策略相同。}\label{fig:gaet-extension}
\end{figure}

若只想让 $T_{\rm eval}$ 之前的样本离\emph{时间上限}至少有 $L$ 步，可选择 $H_{\rm ep}\geq T_{\rm eval}+L$，这是保守的长度条件。例如取 $L=82$ 时，600 步提供了这一量级的缓冲。但要使每个评价前缀样本实际具有至少 $L$ 条残差，还必须使它离 rollout 边界足够远；仅改 $H_{\rm ep}$ 无法做到。增加 $R_{\rm roll}$ 也只是减少窗口边界频率，窗口末尾仍会有短支持样本。

\subsection{一个不直观的代价：更长回合会减少启动阶段的训练比例}

在固定策略、固定初始分布、无提前失败且每步均匀纳入训练的理想情形下，训练采样的年龄分布为
\begin{equation}\label{eq:gaet-occupancy}
 d_H^\pi(s,a)=\frac1H\sum_{\tau=0}^{H-1}d_\tau^\pi(s,a),
 \qquad
 d_{\rm eval}^\pi(s,a)=\frac1{500}\sum_{\tau=0}^{499}d_\tau^\pi(s,a).
\end{equation}
这里简记 $H=H_{\rm ep}$，$d_\tau^\pi$ 是同一固定策略从初始分布运行到年龄 $\tau$ 时的状态--动作分布。训练中的真实策略不断变化、可能提前失败，且评估可采用确定性动作，所以实际偏移更多；此式只隔离回合长度本身的影响。

还可把某个回合年龄上的估计分布写成混合：
\begin{equation}\label{eq:gaet-mixture}
 \mathcal L_{\rm train}(\hat A\mid\tau)
 =\sum_k\Pr_{\rm train}(K=k\mid\tau)\,
 \mathcal L_{\rm train}(\hat A^{(k)}\mid\tau,K=k).
\end{equation}
$\mathcal L$ 表示概率分布，求和遍历可能的支持长度。延长回合可以改变混合权重及被访问的状态，却不会自动令不同年龄的这些分布相同。即使固定 $K$，启动与保持阶段的状态、动作、真实优势以及 critic 误差也可能不同。研究时更应比较\emph{同年龄、同状态--动作条件下的估计误差}，而不是要求所有年龄的原始优势具有相同直方图。

当 $H\geq500$，前 500 步占训练样本比例为 $500/H$。例如从 500 延到 1000，固定总采样预算下，启动与前缀样本约减少一半；前 100 步的比例由 $20\%$ 降至 $10\%$。如果主要难点是最初的位形转换，而后面大多是容易的保持状态，延长回合反而可能稀释关键样本。对 32768000 条总交互预算，完整 500 步回合约为 65536 个，1000 步则约为 32768 个；该粗略计数忽略初末不完整回合和提前失败。

\begin{table}[!ht]
\centering\small
\caption{改变时间上限的确定性算例；固定 $R_{\rm roll}=256$，无提前失败}\label{tab:gaet-horizons}
\begin{tabular}{rrrrr}
\toprule
$H_{\rm ep}$ & 前 500 步比例 & 距 timeout 不超过 82 步 & $K\leq82$ 比例 & 说明 \\
\midrule
500 & 100.00\% & 16.40\% & 43.175\% & 当前设置 \\
600 & 83.33\% & 13.67\% & 41.313\% & 增加前缀后的缓冲 \\
768 & 65.10\% & 10.68\% & 32.031\% & reset 与窗口边界重合 \\
1000 & 50.00\% & 8.20\% & 37.600\% & 后半回合也参与训练 \\
2000 & 25.00\% & 4.10\% & 34.813\% & 启动样本更稀疏 \\
\bottomrule
\end{tabular}
\end{table}

第三列是 $82/H$；第四列通过式~\eqref{eq:gaet-support} 在 $\operatorname{lcm}(H,256)$ 个转移的完整周期上枚举，取到两个边界中较早者，\textbf{不能把第三列再与 $82/256$ 直接相加}。768 的窗口对齐使两个边界的近邻完全重叠，所以数值小于 1000；这不是 768 具有更好控制性能的证明。整除对齐还会使同一年龄总遇到同样的窗口位置，可能保留条件估计差异。表中比例是短支持的计数，不是有偏样本或无效样本的比例。

\subsection{训练怎样受到影响：从优势误差到 actor 梯度}

在 PPO 更新起点附近，忽略裁剪和小批量优势标准化时，优势误差 $\epsilon_{A,t}=\hat A_t-A^\pi(s_t,a_t)$ 引入的期望策略梯度差为
\begin{equation}\label{eq:gaet-gradienterror}
 \Delta g=\mathbb E\!\left[
 \nabla_\theta\log\pi_\theta(a_t\mid s_t)\,\epsilon_{A,t}
 \right].
\end{equation}
$\theta$ 为 actor 参数，$g$ 为欲最大化回报的梯度，期望在当前采样分布下取值。零均值但高方差的误差会增加梯度噪声；与状态、动作或回合年龄系统相关的误差可能改变梯度方向。纯粹只依赖状态、与动作无关的基线误差在精确 on-policy 期望下可抵消，但式~\eqref{eq:gaet-valueerror} 中未来状态误差通常会通过动作影响而相关，不能全部当作可抵消常数。

当前 F2 每批有 $32\times256=8192$ 条样本，小批量 1024，最多复用 6 轮，即最多 48 次优化器更新。GAE 在这 6 轮中保持固定，actor 的概率比却持续改变；SB3 还按小批量标准化优势，并可能触发 KL 提前停止。因此一次边界估计偏差可以随同一批数据被多次使用，而不是每轮都用新 critic 重新修正。标准化不能补回缺少的未来信息，也不能保证纠正优势符号。不同年龄的样本每次被等权平均损失，不代表它们的估计噪声、训练覆盖和梯度贡献完全相同。

实际训练启用奖励归一化及裁剪；本章的 $r$ 与 $\bar V$ 应理解为\emph{同一个归一化尺度}上的量。若改变回合长度，回报累计及 reset 频率会改变奖励归一化统计，随后改变价值拟合和 PPO 的训练尺度。对比实验应记录原始奖励、训练奖励及归一化统计，不能把所有差异都归因于 GAE 的 $K$。

\subsection{如果五百步是任务终点，延长训练可能在优化另一个问题}

对同一策略、同一继续动力学和同一奖励，定义 $J_{500}=\mathbb E[\sum_{t=0}^{499}\gamma^t r_t]$，继续任务满足
\begin{equation}\label{eq:gaet-objectivegap}
 J_\infty-J_{500}
 =\gamma^{500}\mathbb E[V^{\pi,\infty}(s_{500})].
\end{equation}
$J$ 是从相同初始分布出发的期望折扣回报。如果评分是无折扣的 500 步平均/总分，它与两者又都不完全相同。对于每步奖励恒为 1 的说明性例子，$V^\pi=1/(1-\gamma)=200$，第 500 步以后的折扣尾部仍约为 $0.995^{500}\times200\simeq16.3$；这不是本项目实际奖励的数值。

因此可分成三种研究目标。其一，500 步只是报告窗口，希望控制器之后仍能安全保持：把时间上限作为外加截断并研究长回合是合理的，但仍需看前 500 步的性能。其二，任务严格只评价这 500 步且第 501 步没有任何效用：应研究有限时域价值和剩余时间信息，盲目加入后续回报会改变优化目标。其三，希望优先提高前 500 步得分，同时避免终点后立即失稳：可把后续安全性作为显式次级指标或约束，说明权重，而不是把两种目标混称为估计器同分布。

若把一个继续任务的时间上限错误地当成真正终止，固定其余量，片段开头的优势差为
\begin{equation}\label{eq:gaet-wrongzero}
 \hat A_t^{\rm zero}-\hat A_t^{\rm bootstrap}
 =-\beta_K\bar V(s_{t+K}).
\end{equation}
这指出了缺失自举的具体影响；当前标准 SB3 timeout 路径并非如此。反过来，若确实是有限任务终点，强行加上继续价值也会产生相反的目标错配。必须先定义终点语义，再判断哪种处理正确。

\subsection{当前代码的实际约束：不能只把 YAML 中的 500 改大}

一个已经确认的接口约束是：\texttt{training/references.py} 的 \texttt{build\_reference} 在 F2 分支固定生成 500 个 $I_p$ 参考点；\texttt{train\_v2/train\_ppo.py} 调用它时没有传入可变长度。环境的 \texttt{\_coerce\_scalar\_series} 要求非标量序列长度与 \texttt{max\_steps} 一致。因此在当前默认 trajectory 路径下，\textbf{只把 F2 的 max\_steps 改成 600 或 1000，会造成参考长度不匹配，而不是自动获得延长后的训练轨迹。}

真正进行长度消融时，应让参考生成器接受目标长度，保留前 500 步的逐步参考和物理时间尺度；对于其后的部分，明确选择保持末端目标还是继续变化。不能用整体重采样把原来的 100 步电流过渡拉长，否则同时改变了控制任务。还应验证服务端核心时间范围、初始炮段可用时域、参考索引、失败条件和保护机制确实支持继续推进。Python 计数允许更长并不能证明核心可运行更久。本章未启动 HFM 长放电，故不把这些必要验证写成已完成的事实。

另一个尺度是观测：当前四步历史特征不显式输入剩余时间。若把 500 步定义为真正有限任务，要检查同观测不同剩余时间的状态混淆；若定义为继续任务，也要检查目标参考及阶段是否给出了足够的动态信息。延长回合不会自动修复这些可观测性问题。

\subsection{怎样把假说变成可判定的实验}

建议先把问题表述为：\emph{在不改变前 500 步参考的前提下，额外训练后续安全保持，能否降低末段价值误差，并提高 500 步评估的稳定性或得分？}这比要求所有位置的 GAE 具有同一分布更具体，因为不同时刻的等离子体状态、可行动作和真实优势本来就可能不同。

\begin{table}[!ht]
\centering\small
\caption{区分三个原因的对照设计；均为待实施方案}\label{tab:gaet-design}
\begin{tabular}{L{1.4cm}L{2.6cm}L{2.4cm}L{6.4cm}}
\toprule
组别 & $H_{\rm ep}$ & $R_{\rm roll}$ & 主要检验的问题 \\
\midrule
A & 500 & 256 & 当前基线，核对 timeout 自举与最终观测 \\
B & 600 或 1000 & 256 & 保留窗口，只增加终点后的数据与生存训练 \\
C & 500 & 512 & 保留 reset，只改变窗口长度及更新频率 \\
D & 与 B 相同 & 512 & 两种长度是否有交互作用 \\
E & 与 B 相同 & 自定义前缀加缓冲 & 只对前 500 步目标样本更新，额外段用于估计；需另写采样器，非当前 SB3 配置已有功能 \\
\bottomrule
\end{tabular}
\end{table}

比较 A--D 时应同时报告相同总交互预算和前缀样本数量；仅匹配总步数并不能匹配启动样本。增加 rollout 长度还会改变每次更新的样本数、策略更新频率和学习率调度进度，应记录这些变化。各组使用多个训练种子和配对的评估初态；首次消融先保持奖励、$\gamma,\lambda$ 与观测设置不变。前缀加缓冲的方案也必须把缓冲计入采样成本，并在采样时固定策略，不能偷偷把更新后策略的未来轨迹当成旧策略的 on-policy 数据。

区分\textbf{估计器实验}与\textbf{训练收益实验}尤其有用。前者固定一个策略、critic、归一化统计和一组长轨迹，再离线按 500、600、1000 的 reset 假设与不同窗口相位重新分段，在相同前缀样本上比较 GAE。这样可以隔离分段效应，不混入策略学习的差异。离线长参考估计也不是绝对真值，须报告其长度、末端自举和估计误差；若要评估继续价值，应在相应状态下进行配对的后续仿真，而非把第 500 步后的值强行设零。后者才是从头训练各组，再在固定的 500 步评估协议上比较结果。

可按年龄段 $1$--$100$、$101$--$256$、$257$--$400$、$401$--$450$、$451$--$500$ 统计，同时保留到 rollout 边界的距离。应记录 $K$、终止类型、最终状态的价值、原始优势及其标准化版本，并比较以下量：
\begin{equation}\label{eq:gaet-diagnostics}
 \operatorname{Bias}_{\mathcal B}=\mathbb E_{\mathcal B}[\hat A-\hat A_{\rm ref}],\quad
 \operatorname{RMSE}_{\mathcal B}=\sqrt{\mathbb E_{\mathcal B}[(\hat A-\hat A_{\rm ref})^2]},\quad
 p_{\rm flip,\mathcal B}=\Pr_{\mathcal B}(\hat A\hat A_{\rm ref}<0).
\end{equation}
$\mathcal B$ 是同一批配对样本的年龄/边界距离分组，$\hat A_{\rm ref}$ 是明确报告构造方法的参考估计；这里的 Bias 是相对参考估计的差，未必是对真实优势的偏差。接近零的参考优势应单独报告，避免把符号翻转率夸大。任务指标则包括 500 步得分、存活率、末段误差、约束越界，以及在允许继续时的终点后失稳率。提前失败的样本不能仅因没有活到末段就被删掉；末段统计须注明幸存样本条件。

\subsection{对延长 max\_steps 这一想法的判断}

延长回合是一个合理的\emph{待验证假说}：它能让训练访问第 500 步以后的状态，可能改善用于末段自举的 critic，并让终点后的风险进入学习。但它不能保证 GAE 条件分布一致，不能移除 256 步窗口边界，也可能降低启动样本比例和改变训练目标。

更精确的设计标准是：先确认 500 步是评估窗口还是任务真正终点；若属于前者，可从保留前缀的 600 步缓冲消融开始，同时比较较长窗口，并用上述配对估计器实验判断改善来自哪里。600 只是满足当前迹衰减量级的一种候选，不是计算出的最优训练长度。只有末段价值误差、前 500 步性能和继续安全性得到实测支持后，才能把延长回合称为有效改进。\textbf{目标应是估计语义一致、误差可控和评估有效，而不是让本来处于不同状态的所有时刻拥有相同数值分布。}


% ================= 附录 =================
\clearpage
\appendix
\section{复现与审计说明}\label{app:repro}

\subsection{从当前材料复现主要算法所需的信息}\label{app:fixinfo}

复现当前数值流程应固定``关于本文档''一节给定的 Git 提交、成对的 L/LX 版本、\code{xscal/resscal}、所有求解参数和动作预处理；先按第~\ref{sec:workflow} 节区分``保留载入初态''与``参数重求初态''。随后保留初始化预步、电阻覆写顺序、普通步电源旁路与预条件器复用，按第~\ref{sec:pseudo} 节组装并解隐式残差。换用同一物理概念但不同边界、缩放或初始顺序，不构成对此实现的逐步复现。

对于当前已封装的运行时，可以调用随附 \code{.p}/\code{.mexa64} 来作运行对比；若要从源码重新构建，则还需要核实 MATLAB/编译器/BLAS 版本、构建选项和 Git 快照与二进制的对应关系。本文提供了这些因素的位置与数学机制，\textbf{没有声称已完成二进制重建或数值轨迹复现}。Socket server 内部源码缺失仍是从 TCP 消息到具体类实例之间的剩余证据边界。

\subsection{不运行仿真器的只读核对示例}\label{app:readonly}

下面示例只读取 MAT 和 Git 对象，适合复核尺寸、默认模式和源行号；\code{scipy.io} 的函数句柄对象无需执行或导入 MATLAB。
\begin{lstlisting}[style=pseudo, caption={只读核对示例}, label=lst:readonly]
from pathlib import Path
import subprocess
from scipy.io import loadmat

root = Path("hfm-runtime-release-sealed-latest/hfm-runtime-release-sealed-latest")
repo = root / "HFM_env"
data = repo / "TOK_EXL50U/XL50U_fge/Data/FGE_INI"
L = loadmat(data / "ini_L_v2.2.mat", simplify_cells=True)["L"]
LX = loadmat(data / "ini_13906_xpt_v2.2.mat", simplify_cells=True)["LX"]
print(L["nN"], L["ny"], L["G"]["na"], L["G"]["nu"])
print(L["P"]["algoNL"], L["P"]["algoF"], L["P"]["cde"])
print(L["dlst"].shape, L["dlst"].nnz, L["Mee"].shape, LX["t"].shape)
commit = "0e2cb6012bf178da35e17e36506851446342d9f5"
source = subprocess.check_output(
    ["git", "-C", str(repo), "show", f"{commit}:meq/fgetk_implicit.m"]
).decode("utf-8", errors="replace")
for number, line in enumerate(source.splitlines(), start=1):
    print(f"{number}: {line}")
\end{lstlisting}

复核源文件可能同时含 UTF-8 与旧中文编码，数学关系应优先核对 ASCII 执行语句；本文未把解码乱码作为变量名或物理证据。MAT 文件没有行号，因此已将数据结论定位到文件名、完整字段与 shape，而非伪造文本行号。

\subsection{本文核查范围}\label{app:scope}

已对入口、调用链、主要方程、网格/索引、矩阵形状、实际参数传递、分支触发、接口计数和文档中的符号作交叉核查，并区分工作树、Git 源码、保存数据、项目外包装和不可见服务端。GS 离散矩阵与保存响应完成了静态代数一致性检查；未运行任何物理仿真或重写原代码。所有无法由当前材料建立的二进制构建对应性、运行表现与整体精度/稳定性结论，均保留为待验证项。

\section{数学工具备查}\label{app:math}

\textbf{本附录定位。}本附录汇总阅读第 \ref{sec:model}--\ref{sec:solver} 章所需的数学工具：柱坐标系下的微分算子、轴对称简化、极向磁通的构造依据，以及 Grad--Shafranov 算子 $\Delta^*$ 的性质。全部内容属\evbg{}标准数学，可在任何矢量分析教材中查到；这里采用的记号与本文正文一致（$R$ 大半径、$\phi$ 环向角、$Z$ 竖直坐标）。

\subsection{柱坐标系与轴对称约定}\label{app:cyl}

柱坐标 $(R,\phi,Z)$ 与直角坐标的关系为
\[
  x=R\cos\phi,\qquad y=R\sin\phi,\qquad z=Z,
\]
其中 $R\ge 0$ 是到对称轴（$Z$ 轴）的距离，$\phi\in[0,2\pi)$ 是绕轴的环向角。对应的单位矢量 $(\bm{e}_R,\bm{e}_\phi,\bm{e}_Z)$ 构成右手系（$\bm{e}_R\times\bm{e}_\phi=\bm{e}_Z$），但不同于直角坐标的是：$\bm{e}_R$ 与 $\bm{e}_\phi$ 的方向\textbf{随 $\phi$ 变化}，这是柱坐标中一切``额外项''（曲率项）的来源：
\[
  \frac{\partial\bm{e}_R}{\partial\phi}=\bm{e}_\phi,\qquad
  \frac{\partial\bm{e}_\phi}{\partial\phi}=-\bm{e}_R,\qquad
  \frac{\partial\bm{e}_Z}{\partial\phi}=\bm{0}.
\]
对任意标量场 $f$ 与矢量场 $\bm{A}$ 求导时必须对单位矢量同步求导，由此产生 \ref{app:ops} 节各微分算子中的 $1/R$ 型修正项。

\textbf{轴对称约定}（见第~\ref{sec:primer-tokamak} 节）：托卡马克位形关于 $\phi$ 不变，即
\[
  \frac{\partial}{\partial\phi}=0
\]
对所有物理量成立。这使三维问题化归为 $(R,Z)$ 平面上的二维问题，也消除了 \ref{app:ops} 节中所有含 $\partial_\phi$ 的项。

\subsection{柱坐标下的梯度、散度、旋度与 Laplacian}\label{app:ops}

柱坐标下各微分算子的一般形式如下（$\bm{A}=A_R\bm{e}_R+A_\phi\bm{e}_\phi+A_Z\bm{e}_Z$），每个算子同时给出轴对称（$\partial_\phi=0$）化简结果。

\paragraph{梯度}
\[
  \nabla f=\bm{e}_R\,\partial_R f+\bm{e}_\phi\,\tfrac{1}{R}\partial_\phi f+\bm{e}_Z\,\partial_Z f
  \qquad\xrightarrow{\ \partial_\phi=0\ }\qquad
  \nabla f=\bm{e}_R\,\partial_R f+\bm{e}_Z\,\partial_Z f .
\]

\paragraph{散度}
\[
  \nabla\cdot\bm{A}=\tfrac{1}{R}\partial_R(RA_R)+\tfrac{1}{R}\partial_\phi A_\phi+\partial_ZA_Z
  \qquad\xrightarrow{\ \partial_\phi=0\ }\qquad
  \nabla\cdot\bm{A}=\tfrac{1}{R}\partial_R(RA_R)+\partial_ZA_Z .
\]

\paragraph{旋度}
\[
  \nabla\times\bm{A}=\bm{e}_R\!\left(\tfrac{1}{R}\partial_\phi A_Z-\partial_ZA_\phi\right)
  +\bm{e}_\phi\left(\partial_ZA_R-\partial_RA_Z\right)
  +\bm{e}_Z\,\tfrac{1}{R}\bigl(\partial_R(RA_\phi)-\partial_\phi A_R\bigr),
\]
\[
  \nabla\times\bm{A}=-\bm{e}_R\,\partial_ZA_\phi
  +\bm{e}_\phi\left(\partial_ZA_R-\partial_RA_Z\right)
  +\bm{e}_Z\,\tfrac{1}{R}\partial_R(RA_\phi)
  \qquad(\partial_\phi=0).
\]

\paragraph{标量 Laplacian}
\[
  \nabla^2f=\tfrac{1}{R}\partial_R\!\left(R\,\partial_Rf\right)+\tfrac{1}{R^2}\partial_\phi^2f+\partial_Z^2f
  \qquad\xrightarrow{\ \partial_\phi=0\ }\qquad
  \nabla^2f=\partial_R^2f+\tfrac{1}{R}\partial_Rf+\partial_Z^2f .
\]

\textbf{记忆要点}：散度中 $RA_R$ 的组合来自``流过一个圆柱侧面的通量随 $R$ 增长''这一几何事实——高斯定理要求面积元本身含因子 $R$。同理，标量 Laplacian 中比平面情形多出的 $\frac{1}{R}\partial_R f$ 项也是几何（曲率）修正，在 $R\to\infty$ 时消失，退化为平面 Laplacian $\partial_R^2+\partial_Z^2$。

\subsection{\texorpdfstring{由 $\nabla\cdot\bm{B}=0$ 构造极向磁通 $F$}{由 div B = 0 构造极向磁通 F}}\label{app:flux}

这是正文中``一个标量场描述极向磁场''（第~\ref{sec:primer-flux} 节）的推导依据。麦克斯韦方程要求 $\nabla\cdot\bm{B}=0$。对轴对称磁场 $\bm{B}=B_R\bm{e}_R+B_\phi\bm{e}_\phi+B_Z\bm{e}_Z$ 应用 \ref{app:ops} 节的轴对称散度公式：
\[
  \nabla\cdot\bm{B}=\frac{1}{R}\frac{\partial(RB_R)}{\partial R}+\frac{\partial B_Z}{\partial Z}=0
  \quad\Longleftrightarrow\quad
  \frac{\partial(RB_R)}{\partial R}+\frac{\partial(RB_Z)}{\partial Z}=0 .
\]
$(R,Z)$ 平面上的矢量场 $(RB_R,\,RB_Z)$ 无散，因此存在流函数 $F(R,Z)$（差一个常数唯一），使
\[
  RB_R=-\frac{\partial}{\partial Z}\!\left(\frac{F}{2\pi}\right),\qquad
  RB_Z=\frac{\partial}{\partial R}\!\left(\frac{F}{2\pi}\right),
\]
即正文式 \eqref{eq:bfield}。除以 $R$ 即得磁场分量。

\begin{figure}[htbp]
\centering
\begin{tikzpicture}[font=\footnotesize]
  % Z 轴
  \draw[-{Stealth}] (0,-0.7) -- (0,3.6) node[right] {$Z$};
  \node[below left] at (0,-0.65) {$O$};
  % 水平圆盘（椭圆透视）
  \draw[fill=black!8] (0,2.1) ellipse (3.0 and 0.6);
  \draw[black!50] (0,2.1) -- (3.0,2.1);
  \draw[-{Stealth}] (0.2,2.34) -- (2.8,2.34);
  \node[above] at (1.5,2.42) {$R$};
  \node at (0.1,3.0) {水平圆盘 $S_Z$（高度 $Z$，半径 $R$）};
  % 点 P 及其垂线
  \fill (3.0,2.1) circle (1.6pt);
  \node[right=1pt] at (3.0,2.1) {$P(R,Z)$};
  \draw[dashed, black!50] (3.0,2.1) -- (3.0,-0.05);
  \draw[dashed, black!50] (0,2.1) -- (0,2.72);
  \node[left, black!60] at (0,2.4) {$Z$};
  % B_Z 穿盘
  \draw[-{Stealth}, thick] (1.7,0.7) -- (1.7,1.8);
  \node[right] at (1.75,1.25) {$B_Z$ 穿过盘面};
  % B_R 穿侧壁
  \draw[-{Stealth}, thick] (3.32,1.45) -- (3.95,1.45);
  \node[right] at (3.98,1.45) {$B_R$ 穿过侧壁};
  \draw[black!50] (3.0,0.55) -- (3.32,0.55);
  \draw[black!50] (3.0,2.75) -- (3.32,2.75);
\end{tikzpicture}
\caption{极向磁通的几何构造：点 $P(R,Z)$ 绕 $Z$ 轴旋转一周生成水平圆盘 $S_Z$；$F(R,Z)$ 是 $\bm{B}$ 穿过 $S_Z$ 的通量。三个磁场分量中只有 $B_Z$ 对盘面通量有贡献；取相邻高度的两只共轴圆盘，由 $\nabla\cdot\bm{B}=0$，盘面通量差等于从圆柱侧壁流出的通量，由此得 $B_R$ 与 $\partial_ZF$ 的关系。}
\label{fig:fluxdisk}
\end{figure}

用圆盘语言重述上述推导。由定义
\[
  F(R,Z)=\int_{S_Z}\bm{B}\cdot d\bm{S}=2\pi\int_0^R B_Z(R',Z)\,R'\,dR' ,
\]
盘面上 $\bm{e}_\phi$ 沿盘缘切向、$\bm{e}_R$ 躺在盘面内，均不产生通量，故只剩 $B_Z$。对 $R$ 求导立即给出 $B_Z=\frac{1}{2\pi R}\partial_RF$。再取高度 $Z$ 与 $Z+dZ$ 的两只共轴圆盘，它们与圆柱侧壁围成封闭体积；由 $\nabla\cdot\bm{B}=0$，净通量为零，即
\[
  \underbrace{F(Z+dZ)-F(Z)}_{\text{两盘面通量差}}+\underbrace{B_R\,(2\pi R\,dZ)}_{\text{侧壁流出}}=0
  \quad\Rightarrow\quad
  B_R=-\frac{1}{2\pi R}\frac{\partial F}{\partial Z}.
\]
负号是约定使然：侧壁外法向沿 $+\bm{e}_R$，当 $B_R>0$（场向外）时，下游（更高处）盘面捕获的通量反而减少。这正是二维不可压缩流体力学中流函数构造的翻版。

等价地，由磁矢势 $\bm{B}=\nabla\times\bm{A}$ 出发：轴对称下旋度公式只剩 $\bm{e}_\phi$ 分量可作``势''，即 $\bm{A}=A_\phi(R,Z)\,\bm{e}_\phi$，且
\[
  B_R=-\partial_ZA_\phi,\qquad B_Z=\frac{1}{R}\partial_R(RA_\phi),\qquad F=2\pi R\,A_\phi .
\]
（本文 $F$ 取 Wb 总磁通约定，故含 $2\pi$；若采用每弧度磁通 $\psi=F/2\pi$，则 $\psi=RA_\phi$。）$F$ 的等值线是磁力线在截面上的投影：在等值线上 $dF=\partial_RF\,dR+\partial_ZF\,dZ=2\pi R(-B_Z\,dR+B_R\,dZ)=0$，即 $(dR,dZ)$ 与 $(B_R,B_Z)$ 平行。

\subsection{\texorpdfstring{Grad--Shafranov 算子 $\Delta^*$}{Grad--Shafranov 算子 Delta*}}\label{app:gsop}

正文式~\eqref{eq:gsfull} 出现算子
\[
  \Delta^*F\;\equiv\;R\,\partial_R\!\left(\frac{1}{R}\partial_RF\right)+\partial_Z^2F
  \;=\;\partial_R^2F-\frac{1}{R}\partial_RF+\partial_Z^2F .
\]
它的来源是把 Ampère 定律的环向分量写回 $F$：由 $B_Z=\frac{1}{2\pi R}\partial_RF$ 及正文式 \eqref{eq:gs} 的 $j_{\phi,\mathrm{source}}$ 关系，环向电流对应的磁场项组合后恰为 $\frac{1}{2\pi}\Delta^*F$。三个等价写法有助于理解其性质：
\[
  \Delta^*F=R^2\,\nabla\cdot\!\left(\frac{1}{R^2}\nabla F\right)
  \qquad\text{（散度型，体现守恒结构）},
\]
\[
  \Delta^*F=\nabla^2F-\frac{2}{R}\partial_RF
  \qquad\text{（对标量 Laplacian 的显式修正）}.
\]

\textbf{性质}：
\begin{itemize}[nosep]
  \item \textbf{椭圆性}：含一阶导数的 $-\frac{1}{R}\partial_RF$ 是低阶项，主部与平面 Laplacian 相同（$\partial_R^2+\partial_Z^2$），故 $\Delta^*$ 是椭圆算子——这是第 \ref{sec:discret} 章五点差分、第 \ref{sec:solver} 章快速椭圆求解（GSZR）适用的前提；
  \item \textbf{近轴行为}：修正项带 $1/R$，在大 $R$ 处趋于平面 Laplacian；靠近对称轴时需与正则边界条件配合（物理上 $B_Z$ 在轴上应有限）；
  \item \textbf{守恒离散}：第 \ref{sec:fivepoint} 节五点格式中半网格系数 $a_{j\pm}$ 的设计（式 \eqref{eq:radial} 一带），正是为了在离散层面保留下散度型 $\nabla\cdot(R^{-2}\nabla F)$ 的通量守恒结构，而不是直接差分展开式 $\partial_R^2-\frac1R\partial_R$——后者在粗网格上会丢失通量守恒。
\end{itemize}

\subsection{常用矢量恒等式与曲线坐标下的注意事项}\label{app:ident}

阅读推导时经常用到：
\[
  \nabla\times(\nabla f)=\bm{0},\qquad
  \nabla\cdot(\nabla\times\bm{A})=0,\qquad
  \nabla\times(\nabla\times\bm{A})=\nabla(\nabla\cdot\bm{A})-\nabla^2\bm{A}.
\]
\textbf{注意事项}：最后一个恒等式中 $\nabla^2\bm{A}$ 是``对每个直角坐标分量取 Laplacian 后再换回柱坐标分量''的组合，\textbf{不等于}``对柱坐标分量 $A_R,A_\phi,A_Z$ 分别取标量 Laplacian''。曲线坐标下标量 Laplacian 作用于矢量分量会产生额外的曲率耦合项（例如 $\nabla^2(A_\phi\bm{e}_\phi)$ 中含 $-A_\phi/R^2$ 型项）。从 Ampère 定律环向分量推导 GS 方程时若按分量展开，必须保留这些项，最终才会正确归并出 $\Delta^*$；这也是文献中更常从磁通/矢势形式直接推导的原因。类似地，$\bm{j}\cdot\nabla F=0$、$\nabla p=\bm{j}\times\bm{B}$ 等恒等式的分量展开都要记得单位矢量随 $\phi$ 求导的贡献（见 \ref{app:cyl} 节）。

\subsection{符号对照}\label{app:symbols}

\begin{center}
\begin{tabular}{@{}l l l@{}}
\toprule
记号 & 含义 & 在柱坐标/轴对称下的角色 \\
\midrule
$F$（Wb） & 极向磁通（总磁通约定） & 标量场，$F=2\pi RA_\phi$ \\
$T$（T·m） & $RB_\phi$（环向场函数） & 标量场；注意 $T$ 不是温度 \\
$\bm{e}_R,\bm{e}_\phi,\bm{e}_Z$ & 柱坐标单位矢量 & 随 $\phi$ 变化，右手系 \\
$\nabla\cdot$ / $\nabla\times$ & 散度 / 旋度 & 见 \ref{app:ops} 节 \\
$\nabla^2$ & 标量 Laplacian & 含曲率修正 $\frac1R\partial_R$ \\
$\Delta^*$ & GS 算子 & $\partial_R^2-\frac1R\partial_R+\partial_Z^2$，椭圆 \\
\bottomrule
\end{tabular}
\end{center}


\clearpage
\phantomsection
\addcontentsline{toc}{section}{参考文献}
\renewcommand{\refname}{参考文献}
\begin{thebibliography}{99}
\small
\bibitem{Pheat}
Max-Planck-Institut f\"ur Plasmaphysik，\href{https://www.ipp.mpg.de/15108/plasmaheizung}{Plasma heating}：电阻、波与快粒子加热的物理机制。

\bibitem{Pcoulomb}
Richard Fitzpatrick，University of Texas at Austin，\href{https://farside.ph.utexas.edu/teaching/plasma/lectures1/node35.html}{Two-Body Coulomb Collisions}：两体散射与守恒关系。

\bibitem{Patomic}
EIRENE Team，\href{https://www.eirene.de/Documentation/eirene.pdf}{EIRENE Documentation}：原子/分子反应、反应产物和 HYDHEL/AMJUEL 反应率资料的使用。

\bibitem{Pcollisions}
Richard Fitzpatrick，University of Texas at Austin，\href{https://farside.ph.utexas.edu/teaching/plasma/Plasma/node8.html}{Collisions}：累计散射频率与平均自由程。

\bibitem{Poperator}
Richard Fitzpatrick，University of Texas at Austin，\href{https://farside.ph.utexas.edu/teaching/plasma/Plasma/node31.html}{Collision Operator}：平均场、电荷电流与碰撞算子的分工。

\bibitem{C2}
ITER Organization，\href{https://www.iter.org/machine/magnets}{Magnets}：环向场、极向场与磁体系统分工。

\bibitem{C3}
US ITER / Oak Ridge National Laboratory，\href{https://usiter.ornl.gov/hardware/central-solenoid/}{Central Solenoid}：通过磁通变化启动和维持等离子体电流。

\bibitem{Porbits}
Richard Fitzpatrick，University of Texas at Austin，\href{https://farside.ph.utexas.edu/teaching/plasma/lectures1/node16.html}{Motion in Uniform Fields}：回旋、导引中心及电漂移。

\bibitem{Plandau}
Richard Fitzpatrick，University of Texas at Austin，\href{https://farside.ph.utexas.edu/teaching/plasma/lectures1/node37.html}{Landau Collision Operator}：小角 Coulomb 碰撞的动理学闭合。

\bibitem{Pmoments}
Richard Fitzpatrick，University of Texas at Austin，\href{https://farside.ph.utexas.edu/teaching/plasma/lectures1/node47.html}{Moments of Distribution Function}：由分布函数到密度、流速、压力与热流。

\bibitem{Pmhd}
Richard Fitzpatrick，University of Texas at Austin，\href{https://farside.ph.utexas.edu/teaching/plasma/lectures1/node59.html}{MHD Equations}：单流体近似及其排序。

\bibitem{Pgyro}
I. G. Abel 等，\href{https://arxiv.org/abs/1209.4782}{Multiscale Gyrokinetics for Rotating Tokamak Plasmas: Fluctuations, Transport and Energy Flows}，\textit{Reports on Progress in Physics} 76 (2013) 116201，\href{https://doi.org/10.1088/0034-4885/76/11/116201}{DOI}。

\bibitem{Pclosure}
Richard Fitzpatrick，University of Texas at Austin，\href{https://farside.ph.utexas.edu/teaching/plasma/lectures1/node61.html}{Closure in Collisionless Magnetized Plasmas}：弱碰撞磁化条件下的闭合。

\bibitem{Pdebye}
Richard Fitzpatrick，University of Texas at Austin，\href{https://farside.ph.utexas.edu/teaching/plasma/lectures1/node7.html}{Debye Shielding}：线性屏蔽与准中性尺度。

\bibitem{Pislands}
Richard Fitzpatrick，University of Texas at Austin，\href{https://farside.ph.utexas.edu/teaching/plasma1/Fusionhtml/node14.html}{Tearing Modes}：有理磁面、非理想效应和磁岛。

\bibitem{Ppedestal}
Ian Abel、Axel Hallenbert，\href{https://arxiv.org/abs/1711.10403}{Multiscale Modelling for Tokamak Pedestals}，2018 年修订版，\href{https://doi.org/10.1017/S0022377818000326}{DOI}：边缘台基与多尺度模型耦合。

\bibitem{C1}
ITER Organization，\href{https://www.iter.org/fusion-glossary}{Fusion Glossary}：点火、聚变增益及放电阶段术语。

\bibitem{C4}
H.-T. Kim、A. C. C. Sips 与 EFDA-JET contributors，\href{https://arxiv.org/abs/1309.4890}{Physics of plasma burn-through and DYON simulations for the JET ITER-like wall}，\textit{Nuclear Fusion} 53 (2013) 083024：击穿、辐射障碍、壁面杂质与烧穿。

\bibitem{C5}
ITER Organization，\href{https://www.iter.org/machine/supporting-systems/external-heating-systems}{External heating systems}：欧姆、中性束与射频加热。

\bibitem{C10}
MIT Plasma Science and Fusion Center，\href{https://www-internal.psfc.mit.edu/research/MFEFormulary/files/MFEFormulary.pdf}{Magnetic Fusion Energy Formulary}，第 9 章：反应通道、分支与产物能量。

\bibitem{C12}
R. M. Magee 等，\href{https://doi.org/10.1038/s41467-023-36655-1}{First measurements of $p{}^{11}$B fusion in a magnetically confined plasma}，\textit{Nature Communications} 14 (2023) 955：质子--硼反应观测及功率平衡挑战。

\bibitem{C11}
U.S. Department of Energy，\href{https://www.energy.gov/science/doe-explainsdeuterium-tritium-fusion-fuel}{DOE Explains: Deuterium-Tritium Fusion Fuel}：燃料选择及 D--${}^3$He、质子--硼方案的较高温度要求。

\bibitem{C6}
EPFL Swiss Plasma Center，\href{https://www.epfl.ch/research/domains/swiss-plasma-center/tcv-plasma-shapes/}{TCV Plasma Shapes}：磁面、X 点及位形。

\bibitem{Dsol}
N. Walkden 等，\href{https://doi.org/10.1038/s42005-022-00906-2}{The physics of turbulence localised to the tokamak divertor volume}，\textit{Communications Physics} (2022)：跨场输运、SOL、私有磁通区与靶区热流的联系。

\bibitem{Ddetach}
O. F\'evrier 等，\href{https://arxiv.org/abs/2402.04656}{SPLEND1D, a reduced one-dimensional model to investigate the physics of plasma detachment} (2024)：沿场输运、原子过程、能量/动量损失与脱靶建模。

\bibitem{Dadc}
H. Raj 等，\href{https://arxiv.org/abs/2207.13305}{Improved Heat and Particle Flux Mitigation in High Core Confinement, Baffled, Alternate Divertor Configurations in the TCV tokamak}，\textit{Nuclear Fusion} (2022)，\href{https://doi.org/10.1088/1741-4326/ac94e5}{DOI}：几何、挡板及杂质注入共同影响排热的实验。

\bibitem{Dmaterial}
Idaho National Laboratory，\href{https://tmap8.inl.gov/examples/divertor_monoblock/}{TMAP8: Divertor Monoblock During Pulsed Operation}：ITER 型钨--铜--CuCrZr 多层部件与热/氚输运模型。

\bibitem{Diter}
ITER Organization，\href{https://www.iter.org/machine/divertor}{Divertor}：靶板、钨装甲、主动冷却及偏滤器系统。

\bibitem{Dcarbon}
ITER Organization，\href{https://www.iter.org/node/20687/design-review-tungsten-divertor-shows-way-ahead}{Design review for tungsten divertor shows way ahead} (2013)：碳侵蚀、氚滞留与钨装甲选择的背景。

\bibitem{Dxptr}
K. Lee 等，\href{https://doi.org/10.1103/PhysRevLett.134.185102}{X-Point Target Radiator Regime in Tokamak Divertor Plasmas}，\textit{Physical Review Letters} 134 (2025) 185102：远离核心的次 X 点辐射与脱靶控制。

\bibitem{Dgeometry}
M. Kotschenreuther 等，\href{https://arxiv.org/abs/1309.5289}{Magnetic Geometry and Physics of Advanced Divertors: The X-Divertor and the Snowflake}，\textit{Physics of Plasmas} 20 (2013)，\href{https://doi.org/10.1063/1.4824735}{DOI}：先进偏滤器几何和 SOL 磁通扩展的区分。

\bibitem{C7}
EPFL Swiss Plasma Center，\href{https://www.epfl.ch/research/domains/swiss-plasma-center/tcv-plasma-control/}{TCV Plasma Control}：形状控制、MEQ 与诊断协同。

\bibitem{C8}
EPFL Swiss Plasma Center，\href{https://www.epfl.ch/research/domains/swiss-plasma-center/tcv-disruptions-and-runaway-electrons/}{TCV Disruptions and Runaway Electrons}：破裂与逃逸电子问题。

\bibitem{C9}
MIT Plasma Science and Fusion Center，\href{https://disruptions.mit.edu/}{Disruptions}：破裂预测、稳定性观察与缓解研究。

\bibitem{Sreview}
Per Helander，\href{https://doi.org/10.1088/0034-4885/77/8/087001}{Theory of plasma confinement in non-axisymmetric magnetic fields}，\textit{Reports on Progress in Physics} 77 (2014) 087001；期刊页面另列有勘误。

\bibitem{Sprecise}
Matt Landreman、Elizabeth Paul，\href{https://arxiv.org/abs/2108.03711}{Magnetic fields with precise quasisymmetry for plasma confinement}，\textit{Physical Review Letters} 128 (2022) 035001。

\bibitem{Scoils}
Florian Wechsung 等，\href{https://doi.org/10.1073/pnas.2202084119}{Precise stellarator quasi-symmetry can be achieved with electromagnetic coils}，\textit{PNAS} 119 (2022) e2202084119。

\bibitem{Sw7x}
Max-Planck-Institut f\"ur Plasmaphysik，\href{https://www.ipp.mpg.de/3766317/W7X_en_2023.pdf}{Wendelstein 7-X: Research objectives and device overview}，2023 年英文介绍资料：磁体结构与优化目标。

\bibitem{Sbootstrap}
Max-Planck-Institut f\"ur Plasmaphysik，\href{https://www.ipp.mpg.de/4631577/Bootstrap-current}{Bootstrap current}：W7-X 自举电流、旋转变换与岛偏滤器的联系。

\bibitem{Sneo}
C. D. Beidler 等，\href{https://doi.org/10.1038/s41586-021-03687-w}{Demonstration of reduced neoclassical energy transport in Wendelstein 7-X}，\textit{Nature}，2021。

\bibitem{Sregcoil}
Matt Landreman，\href{https://arxiv.org/abs/1609.04378}{An improved current potential method for fast computation of stellarator coil shapes}，\textit{Nuclear Fusion} 57 (2017) 046003。

\bibitem{Ssingle}
R. Jorge 等，\href{https://arxiv.org/abs/2302.10622}{Single-Stage Stellarator Optimization: Combining Coils with Fixed Boundary Equilibria}，2023，\href{https://doi.org/10.1088/1361-6587/acd957}{DOI}。

\bibitem{Sadjoint}
E. J. Paul、M. Landreman、A. Bader、W. Dorland，\href{https://arxiv.org/abs/1801.04317}{An adjoint method for gradient-based optimization of stellarator coil shapes}，2018：绕组面优化与灵敏度。

\bibitem{Slhd}
National Institute for Fusion Science，\href{https://soken.nifs.ac.jp/archives/about/program_eng}{LHD device overview}：连续螺旋与极向线圈结构。

\bibitem{Slhdvariation}
National Institute for Fusion Science，\href{https://www.nifs.ac.jp/report/nifs383.html}{NIFS-383}：LHD 分层螺旋线圈电流与配置灵活性。

\bibitem{Shsx}
University of Wisconsin--Madison，\href{https://www.hsx.wisc.edu/}{Helically Symmetric eXperiment (HSX)}：模块线圈与准螺旋对称研究。

\bibitem{Sw7xpower}
\href{https://pure.mpg.de/pubman/item/item_2137690}{Power supplies for the WENDELSTEIN 7-X stellarator}，MPG.PuRe 文献记录：不同线圈组的独立供电与配置调节。

\bibitem{Smulticonfig}
Brandon F. Lee、Elizabeth J. Paul、Georg Stadler、Matt Landreman，\href{https://arxiv.org/abs/2208.01096}{Stellarator coil optimization supporting multiple magnetic configurations}，2022。

\bibitem{Ajfnk}
D. A. Knoll、D. E. Keyes，\href{https://doi.org/10.1016/j.jcp.2003.08.010}{Jacobian-free Newton--Krylov methods: a survey of approaches and applications}，\textit{Journal of Computational Physics} 193 (2004) 357--397。

\bibitem{Appo}
John Schulman 等，\href{https://arxiv.org/abs/1707.06347}{Proximal Policy Optimization Algorithms}，arXiv:1707.06347，2017。

\bibitem{Agae}
John Schulman 等，\href{https://arxiv.org/abs/1506.02438}{High-Dimensional Continuous Control Using Generalized Advantage Estimation}，arXiv:1506.02438。

\bibitem{Atimelimits}
F. Pardo、A. Tavakoli、V. Levdik、P. Kormushev，\href{https://arxiv.org/abs/1712.00378}{Time Limits in Reinforcement Learning}，\textit{PMLR} 80 (2018) 4042--4051：有限任务终点、外加时间截断与部分回合自举的区别。

\bibitem{Asb3rollout}
Stable-Baselines3，\href{https://stable-baselines3.readthedocs.io/en/v2.7.0/_modules/stable_baselines3/common/on_policy_algorithm.html}{OnPolicyAlgorithm source, v2.7.0}：窗口采样、最终观测的 timeout 自举及训练循环；实际安装版本应在实验中记录。

\bibitem{Asb3buffer}
Stable-Baselines3，\href{https://github.com/DLR-RM/stable-baselines3/blob/v2.7.0/stable_baselines3/common/buffers.py}{RolloutBuffer source, v2.7.0}：GAE 倒序递推、回合边界掩码及小批量遍历。

\end{thebibliography}

\end{document}
